\documentclass[11pt]{article}
\usepackage[a4paper,margin=26mm]{geometry}
\usepackage{amsmath,amssymb,amsthm,booktabs}
\usepackage[colorlinks=true,linkcolor=blue,urlcolor=blue,citecolor=blue]{hyperref}
\newtheorem{theorem}{Theorem}
\newtheorem{lemma}[theorem]{Lemma}
\newtheorem{corollary}[theorem]{Corollary}
\newtheorem{proposition}[theorem]{Proposition}
\newcommand{\PP}{\mathbb P}
\newcommand{\one}{\mathbf 1}
\newcommand{\eps}{\varepsilon}
\newcommand{\CN}{\mathcal C(N)}
\newcommand{\HN}{\mathcal H_N}
\newcommand{\LL}{\log N}
\title{Transferring quantitative Chen bounds to imbalanced semiprime representations}
\author{Research working draft}
\date{8 October 2026}
\begin{document}
\maketitle
\begin{abstract}
We derive an upper sieve for shifted finite prime-product sequences
with coefficient $K=6.7725$. An extended first-prime stop and variable
third-prefix levels are followed by a complete fourth-prefix refinement
and a fifth-prefix gain, with legal odd first exits. Exact rational
certificates bound these two sequential gains. A restricted lower level,
optimized Buchstab subtraction and two disjoint finite corrections give
a separate Chen coefficient $c=1.8321$ and representations $N=p+rq$,
with $r=1$ or prime and $r\le q^{0.7631}$, for every sufficiently large
even $N$. Accepting the published Chen constant $1.9728$ as an external
input instead gives $r\le q^{0.7473}$. Four certified finite feedback
rounds strengthen the second seed on its admitted stable states; its
use in the lower count proceeds through licensed descendants.
The prefix and seed certificates have separate numerical replays and
targeted analytic audits. The complete floating recursion and enlarged
threshold searches remain excluded from these constants.
This remains a research proof draft for independent review; neither
formal verification nor priority over all existing work is asserted.
\end{abstract}

\paragraph{Scope and status.}
All statements concern every sufficiently large even integer, with the
threshold allowed to depend on fixed parameters. No effective threshold or
verification for small integers is claimed. The argument invokes published
theorems as inputs; the conservative calculation does not use the numerical Chen bound in
\cite{Li2026}. The local OpenAI manuscript collection is not used as a
mathematical premise. Literature novelty remains unestablished.
\textbf{Input status:} applications explicitly using $c=1.9728$,
including the later replacement-tree consequence $b=0.7473$, retain
the external-input dependency discussed in Section~\ref{sec:input-audit}.
The separately derived calculation in Section~\ref{sec:conservative}
avoids that input and supplies its own weight proof and numerical
certificate. The further adaptations in
Sections~\ref{sec:primeproducts55}--\ref{sec:restricted-lower-level},
the sharper support argument, and the two-prime prefix improvement
first give $c=1.643$ and $b=0.796$ in
Theorem~\ref{thm:double-prefix-chen}. The further construction in
Section~\ref{sec:pruned-upper-sieve} gives $c=1.678$ and $b=0.790$
in Theorem~\ref{thm:pruned-chen}. The retained-base construction in
Section~\ref{sec:retained-base-upper} and its certified integral
give $c=1.731$ and $b=0.780$ in
Theorem~\ref{thm:retained-base-chen}. The variable-prefix construction
in Section~\ref{sec:variable-prefix-upper} now gives $c=1.783$,
$b=0.770$, and the external-input consequence $b=0.74865$ in
Theorem~\ref{thm:variable-prefix-chen}. Extending the first stop with
a single-variable fallback in Section~\ref{sec:extended-prefix-upper}
further gives $c=1.785$, $b=0.7695$, and the external-input implication
$b=0.7485$ in Theorem~\ref{thm:extended-prefix-chen}. Its numerical
coefficient is strictly below the converted printed $G_1$ bound.
The restricted-lower and Buchstab derivation in
Section~\ref{sec:buchstab-restricted-lower} subsequently improves
the separate calculation to $c=1.817$ and $b=0.766$ in
Theorem~\ref{thm:buchstab-lower-chen}, without double-sieve corrections. The alternative
limiting Buchstab derivation is given in
Section~\ref{sec:g1-localized-interface}. Every use of the L23
factorization is supplemented by the explicit repairs in
Section~\ref{sec:two-factor-greedy-repair}. A further selected-subproduct
construction in Section~\ref{sec:selected-subproduct-factorization}
justifies specific strict-gate branches beyond the largest-prime
restriction; it is not used to change either numerical consequence.
The new four-prefix and threshold formulas below remain unevaluated
by rigorous interval bounds; their numerical probes are not new
theorems. A direct application of Wu's first weighted inequality,
with a newly certified saving rather than an imported correction table,
is proved in Section~\ref{sec:same-new-level-double-sieve-route}.
It improves the separate consequence to $c=1.818$ and $b=0.7659$ in
Theorem~\ref{thm:same-level-one-step-chen}. The later fourth-prefix
descendant refinement and repaired second seed give $K=6.792$,
$c=1.8249$ and $b=0.7645$ in
Theorem~\ref{thm:fourth-prefix-second-seed-chen}; the two certified
correction blocks are disjoint. The generalized and monotone
double-sieve formulas are proved as additional analytic options,
but their floating total gains are not included in that theorem.
The requested advance beyond every existing result is not certified as achieved: these
are research proof drafts, and bibliographic priority has not been
established.

\paragraph{Current local update.}
The replacement fourth-prefix certificate, the retained fifth-prefix
gain and four certified finite seed updates give $K=6.7725$,
$c=1.8321$, and $b=0.7631$ in
Theorem~\ref{thm:stronger-fourth-finite-seed-chen}.
The earlier $K=6.786$, $c=1.8271$, $b=0.76407$ consequence in
Theorem~\ref{thm:fifth-prefix-omega-chen} is retained as an intermediate
result. The replacement fourth gain is never added to the old fourth
and tail gains. The full floating recursion still requires outward
certification and is not used in either theorem.

\section{Definitions and the deduction}
Write $\Omega(m)$ for the number of prime factors with multiplicity, and put
\[
 \CN=\prod_{\ell>2}\left(1-\frac1{(\ell-1)^2}\right)
       \prod_{\substack{\ell\mid N\\\ell>2}}\frac{\ell-1}{\ell-2},
 \qquad \HN=\CN\frac{N}{(\log N)^2}.
\]
The letters $p,r,q,\ell$ denote primes except that $r=1$ is expressly
allowed in the following definition. For $1/2<b<1$, define
\[
 D_b(N)=\#\{p<N:p\in\PP,\ N-p=q\in\PP\text{ or }
             N-p=rq,\ r,q\in\PP,\ r\le q^b\}.
\]
Each prime $p$ is counted once. As $b<1$, the two factors in a composite
representation are ordered uniquely. Set
\[
 D_{1,2}(N)=\#\{p<N:p\in\PP,\ \Omega(N-p)\le2\}.
\]
The convention $\Omega(1)=0$ can change a comparison by at most one.

\begin{theorem}[Transfer statement]\label{thm:transfer}
Suppose a known quantitative Chen theorem gives
$D_{1,2}(N)\ge(c-o(1))\HN$ for every sufficiently large even $N$.
For every fixed $b\in(1/2,1)$,
\[
 \liminf_{\substack{N\to\infty\\2\mid N}}\frac{D_b(N)}{\HN}
 \ge c-8\log(1/b).
\]
In particular $D_b(N)>0$ eventually if $8\log(1/b)<c$.
\end{theorem}

The input in \cite[Theorem 1.1]{Li2026} has exactly the above normalization
and states $D_{1,2}(N)\ge1.9728\HN$. It follows from
Theorem~\ref{thm:transfer} that
every fixed $b>e^{-1.9728/8}$ is admissible. Numerically,
\[
 e^{-1.9728/8}=0.781453212308\ldots .
\]
Thus $b=0.8$ (the notation ``$1+1.8$'') and $b=0.782$ are examples.
The defining inequality is exact; the displayed decimal is explanatory.

\section{Counting the half-products to be removed}
Let $K=10$, $T=N/(\log N)^K$, and define
\[
 \mathcal A=\{m=rq:T\le rq\le N,\ q^b<r\le q,\ r,q\in\PP\}.
\]
Every $m$ has one ordered factorization in this set, including squares.
Both factors satisfy
\begin{equation}\label{eq:rough}
 r,q\ge T^{b/(1+b)}.
\end{equation}
Indeed $q<r^{1/b}$ and $T\le rq<r^{1+1/b}$.

\begin{lemma}\label{lem:mass}
For fixed $b\in(1/2,1)$,
\[
 |\mathcal A|=\left(\log(1/b)+o(1)\right)\frac{N}{\log N}.
\]
\end{lemma}
\begin{proof}
First omit the lower cutoff, which removes at most $T$ integers. Put
$\alpha=b/(1+b)$. Summation over the smaller prime $r$ gives, up to
$O(\sqrt N)$ endpoint terms,
\[
 \sum_{r\le\sqrt N}
 \left\{\pi\bigl(\min(N/r,r^{1/b})\bigr)-\pi(r)\right\}.
\]
The range $r\le N^\alpha$ contributes $O_b(N/\log^2N)$ by the prime
number theorem and partial summation, as does $\sum_{r\le\sqrt N}\pi(r)$.
For $N^\alpha<r\le\sqrt N$, the prime number theorem applied uniformly
to $N/r$, followed by partial summation over $r$, gives the main term
\[
 \frac{N}{\log N}\int_{\alpha}^{1/2}\frac{du}{u(1-u)}
 =\log(1/b)\frac{N}{\log N}.
\]
All integration ranges are fixed compact subintervals of $(0,1)$.
The errors are $o_b(N/\log N)$, and $T=o(N/\log N)$.
\end{proof}

\section{The distribution input and curved boundaries}
For a finitely supported arithmetic function $f$, write
\[
 \Delta(f;a(d))=\sum_{n\equiv a\ (d)}f(n)
       -\frac1{\varphi(d)}\sum_{(n,d)=1}f(n).
\]
We use the generalized Bombieri--Vinogradov theorem in precisely its
unconditional range: \cite[Claim 12 and Theorem 14]{Polymath2014}.
For each fixed $0<\vartheta<1/2$ and $\eps>0$, divisor-bounded sequences
$u,v$ supported on $[U,2U]$ and $[V,2V]$, respectively, with
$x^\eps\ll U,V\ll x^{1-\eps}$ and $UV\asymp x$, have
\begin{equation}\label{eq:geh}
 \sum_{d\le x^\vartheta}\max_{(a,d)=1}
 |\Delta(u*v;a(d))|\ll_A x(\log x)^{-A},
\end{equation}
provided the stated Siegel--Walfisz condition holds for $v$.
The maximum over primitive residue classes is part of the theorem.
Here is the required uniform verification. If $J\subseteq[V,2V]$ is
any interval and $v=\one_{\PP\cap J}$, then for every fixed $A>0$
\[
 |\Delta(v\one_{(\cdot,t)=1};a(d))|
 \ll_{A,\eps}\tau(dt)^{O(1)}V(\log x)^{-A}
 \qquad ((a,d)=1,\ d,t\ge1).
\]
For $d\le(\log x)^B$, with fixed $B$, subtracting the endpoint forms
of the classical Siegel--Walfisz theorem and using partial summation
gives the assertion for $v$, uniformly in $J$. Replacing the usual
integral main term by the exact coprime mass uses the ordinary prime
number theorem with an arbitrarily large logarithmic saving.
For $d>(\log x)^B$, the trivial estimates for an interval and
$\varphi(d)\gg d^{1/2}$ give
$|\Delta(v;a(d))|\ll Vd^{-1/2}+1$; choosing $B>2A+2$
suffices since $V$ is a fixed positive power of $x$.
Finally, removing primes dividing $t$ changes the discrepancy by at
most $2\omega(t)\le2\tau(t)$, which is absorbed by the allowed
divisor factor and the polynomial size of $V$. This proves exactly
the auxiliary-coprimality condition of Claim 12. All constants are
independent of the interval endpoints.

\begin{lemma}[Distribution of the curved semiprime region]\label{lem:bv}
For each fixed $0<\theta<1/2$ and $A>0$, with $D=N^\theta$,
\[
 \sum_{d\le D}\max_{(a,d)=1}
 \left|\#\{m\in\mathcal A:m\equiv a\pmod d\}
       -\frac{|\mathcal A^{(d)}|}{\varphi(d)}\right|
 \ll_{A,b,\theta}N(\log N)^{-A},
 \qquad
 \mathcal A^{(d)}=\{m\in\mathcal A:(m,d)=1\}.
\]
\end{lemma}
\begin{proof}
Choose a fixed $\vartheta$ with $\theta<\vartheta<1/2$.
Partition the possible $r$ and $q$ into half-open grid intervals of ratio
$1+\delta$, $\delta=(\log N)^{-L}$. There are
$O((\log N)^{2L+2})$ pairs of grid intervals. For any rectangle meeting
the region, its product scale $x$ satisfies
$x\gg N/(\log N)^K$, so $N^\theta\le x^\vartheta$ eventually.
Both factors are between fixed positive powers of $x$ and their
complementary powers by \eqref{eq:rough}. On each rectangle the count
is a convolution of two interval-restricted prime indicators; thus
\eqref{eq:geh} applies. Apply it with a logarithmic saving larger than
$A+2L+10$, and sum over the rectangles.

It remains to justify replacing a curved region by rectangles without
multiplying a trivial boundary error by $D$. Let $f^-$ and $f^+$ be the
convolution sums for the inside and intersecting rectangles respectively;
they count ordered prime pairs, including any multiplicity in the
boundary rectangles. Then $0\le f^-\le\one_{\mathcal A}\le f^+$.
Put $u=f^+-f^-\ge0$. The union of boundary rectangles lies in
multiplicative $1+O_b(\delta)$ neighborhoods of
\[
 rq=T,\qquad rq=N,\qquad r=q,\qquad r=q^b.
\]
Even counting integer pairs rather than prime pairs, its total mass is
\begin{equation}\label{eq:boundary}
 \|u\|_1\ll_b\delta N\log N+N^{1/(1+b)}(\log N)^{O(1)}.
\end{equation}
For product boundaries, sum $O(\delta N/r+1)$ over the smaller
coordinate; for the power boundary, sum $O_b(\delta q^b+1)$ for
$q\ll N^{1/(1+b)}$; the diagonal contributes
$O(\delta N+\sqrt N)$. These also cover grid-endpoint roundings.

Let $h=\one_{\mathcal A}-f^-$. Since $0\le h\le u$, for any primitive
class $a(d)$,
\[
 |\Delta(h;a(d))|\le |\Delta(u;a(d))|+
                  2\|u\|_1/\varphi(d).
\]
The rectangle theorem already controls the summed discrepancy of $u$.
Since $\sum_{d\le D}1/\varphi(d)\ll\log N$, \eqref{eq:boundary}
is harmless after choosing $L>A+4$. This proves the assertion. The
choice of $L$ precedes the larger saving in \eqref{eq:geh}.
\end{proof}

\section{An upper sieve for bad representations}
Fix $2/5<\theta<1/2$, and take
\[
 D=N^\theta,\qquad z=N^{1/5},\qquad s=\frac{\log D}{\log z}=5\theta.
\]
Then $2<s<5/2$, entirely within the range where the dimension-one
upper sieve has $F(s)=2e^\gamma/s$.
Let $\mathscr P_N$ be the primes not dividing $N$ and
$P_N(z)=\prod_{\ell<z,\,\ell\nmid N}\ell$.
We sieve the sequence $\{N-m:m\in\mathcal A\}$. For sufficiently large
$N$, all $m\in\mathcal A$ are odd, so these are positive integers.

For a squarefree $d\mid P_N(z)$ all factors of $d$ are smaller than $z$.
By \eqref{eq:rough} and $b>1/2$, every $m\in\mathcal A$ is coprime to
$d$ for sufficiently large $N$. Consequently Lemma~\ref{lem:bv} gives
the exact main term
\[
 \#\{m\in\mathcal A:d\mid N-m\}
       =\frac{|\mathcal A|}{\varphi(d)}+R_d,
 \quad
 \sum_{\substack{d\le D\\d\mid P_N(z)}}|R_d|
       \ll_A N(\log N)^{-A}.
\]
Here $(d,N)=1$, so $a=N\pmod d$ is primitive; its variation with $N$
is explicitly allowed by the maximum in Lemma~\ref{lem:bv}.

The local density is $g(\ell)=1/(\ell-1)$ for $\ell\nmid N$ and zero
for excluded primes. Since $N$ is even, the dimension-one product is
\begin{equation}\label{eq:v}
 V_N(z)=\prod_{\substack{\ell<z\\\ell\nmid N}}
             \left(1-\frac1{\ell-1}\right)
       =\left(1+o(1)\right)\frac{2e^{-\gamma}\CN}{\log z}.
\end{equation}
There are at most five prime divisors $\ell\mid N$ with $\ell\ge z$,
so the omitted tail factor is $1+O(1/z)$, uniformly in $N$.
For an absolute constant $C_0$, the required one-sided dimension bound is
\[
 \frac{V_N(w)}{V_N(v)}\le\frac{\log v}{\log w}
             \left(1+\frac{C_0}{\log w}\right)\quad(2\le w\le v).
\]
It follows from the product over all odd primes, since excluding prime
divisors of $N$ only decreases this ratio. This avoids assuming a
two-sided density estimate uniform in the prime divisors of $N$.

For precision, apply the standard linear upper sieve as stated in
\cite[Theorem 2.10 and Remark 2.9]{Lichtman2025}, with level $D$ and
$g(d)=1/\varphi(d)$ on the allowed squarefree moduli. Remark 2.9
explicitly permits the one-sided product bound for this upper sieve.
For fixed sufficiently small $\eta>0$, and $D$ sufficiently large in
terms of $\eta$, it supplies at most $L_\eta=\exp(\eta^{-3})$
well-factorable coefficient sequences of level $D$, each of absolute
value at most one. Thus the combined error is at most
$L_\eta\sum_{d\mid P_N(z),\,d\le D}|R_d|=o_\eta(\HN)$.
The main-term error is $O_{C_0}(\eta)$.
First let $N\to\infty$ for each fixed $\eta$, and then let
$\eta\to0$. This gives the asymptotic upper bound
\begin{align*}
 S(\mathcal A;N,z)
 &:=\#\{m\in\mathcal A:(N-m,P_N(z))=1\}\\
 &\le |\mathcal A|V_N(z)(F(s)+o(1))+o(\HN)\\
 &\le\left(\frac{4\log(1/b)}{\theta}+o(1)\right)\HN.
\end{align*}
For a source of the linear-sieve normalization see
\cite[Section 1, equations (1.4) and (3.2)]{Wu2008}; the same
switching coefficient $8\log(1/b)$ at the square-root limit appears
at $b=0.9$ in \cite[equations (5.36)--(5.43)]{LiLiu2026}.
No inference from that single value of $b$ is used here.

Let $B_b(N)$ count the bad prime representations $N=p+rq$ with
$q^b<r\le q$. The terms with $rq<T$ number at most $T$.
The terms with $p<z$ number at most $z$; each remaining prime $p=N-m$
survives the sieve. Thus
\begin{equation}\label{eq:bad}
 B_b(N)\le
 \left(\frac{4\log(1/b)}{\theta}+o(1)\right)\HN.
\end{equation}
Square semiprimes are already counted in $B_b$ and require no omission.
There is no assumption that $m$ is coprime to $N$; primes dividing $N$
were correctly excluded from the sifting set.

\section{Completion of the transfer and safe constants}
Every Chen representation, except possibly $N-p=1$, is either counted
by $D_b$ or by $B_b$. Hence
\[
 D_b(N)\ge D_{1,2}(N)-B_b(N)-1.
\]
For every fixed $\theta\in(2/5,1/2)$, \eqref{eq:bad} proves
\[
 \liminf_{\substack{N\to\infty\\2\mid N}}\frac{D_b(N)}{\HN}
 \ge c-\frac{4\log(1/b)}{\theta}.
\]
Taking the supremum over these fixed values of $\theta$ gives
Theorem~\ref{thm:transfer}. This is an order of limits, not an invocation
of generalized Elliott--Halberstam at its unproved endpoint $1/2$.

\begin{corollary}[A numerical deduction assuming the published input]
Assuming \cite[Theorem 1.1]{Li2026} as the published input, for all
sufficiently large even $N$ the above deduction gives
\[
 D_{0.8}(N)\ge0.14\HN>0.
\]
\end{corollary}
\begin{proof}
Use $\theta=49/100$ throughout, and $\vartheta=99/200$ in the rectangle
distribution theorem. The limiting coefficient is
\[
 1.9728-\frac{400}{49}\log(5/4)=0.151219989\ldots .
\]
Floating-point values are not needed for positivity. The alternating
series for the logarithm gives $\log(5/4)<43/192$, so the coefficient is
strictly larger than
\[
 \frac{1233}{625}-\frac{1075}{588}
 =\frac{53129}{367500}>\frac7{50}=0.14.
\]
The positive strict gap absorbs the $o(1)$ errors for sufficiently large
$N$.
\end{proof}

\begin{corollary}[A smaller exponent, with the same input dependency]
Assuming the same published input, for all sufficiently large even $N$,
\[
 D_{0.782}(N)\ge\frac1{250}\HN>0.
\]
In particular the proposition called ``$1+1.782$'' holds by this deduction.
\end{corollary}
\begin{proof}
Set $b=391/500$, $\theta=4999/10000$ and
$\vartheta=9999/20000$. Write $u=109/891$, so that
$(1+u)/(1-u)=500/391$. The positive power series for the logarithm gives
\[
 \log\frac{500}{391}
 =2\sum_{j\ge0}\frac{u^{2j+1}}{2j+1}
 \le2\left(u+\frac{u^3}{3}+\frac{u^5}{5(1-u^2)}\right)
 =\frac{1020142152278647}{4148595849915000}
 <\frac{123}{500}.
\]
Consequently the limiting coefficient in the lower bound exceeds
\[
 \frac{1233}{625}-\frac{40000}{4999}\frac{123}{500}
 =\frac{13767}{3124375}>\frac1{250}.
\]
Again the strict gap absorbs all asymptotic errors. Both distribution
exponents are fixed and strictly smaller than $1/2$.
\end{proof}

\section{A conservative Chen bound with an explicit quotient weight}\label{sec:conservative}
This section gives a separate deduction avoiding the numerical input
in \cite{Li2026}. It uses a directly proved weight with the full
fourth-prime range, a derived prime-prefix distribution lemma, and
exact integer numerical enclosures.
\begin{theorem}[Conservative deduction]\label{thm:conservative}
For every sufficiently large even $N$,
\[
 D_{1,2}(N)\ge\frac{61}{50}\HN,
 \qquad D_{0.86}(N)\ge\frac1{100}\HN.
\]
In particular, $N=p+rq$ with $p,q$ prime, $r=1$ or prime,
and $r\le q^{0.86}$.
\end{theorem}

Put
\[
 a=\frac{100}{1149},\qquad h=\frac{50}{309},\qquad
 m=\frac{25}{128},\qquad d=\frac12-3a,\qquad e=\frac{57}{224}.
\]
Here $h$ is a sieve cutoff, unrelated to the double-sieve function; no
double-sieve correction is used in this section. Define
\[
 v(t)=\begin{cases}(1+t)/2,&t\le m,\\
                    57/64-3t/2,&t\ge m.
       \end{cases}
\]
This is the minimum of the three single-variable levels in
\cite[Proposition 6.1]{Lichtman2023}, with the large moving residue
parameter equal to one and spectral parameter $7/32$.
For example, the differences from the omitted middle branch are
$(25-107t)/78$ before the junction and $(1568t-175)/2496$ afterwards;
they are nonnegative on the respective ranges.

% Integration fragment. Requires amsmath, amssymb and amsthm.
% This supplies a combinatorial proof, independently of Wu's final range
% contraction. It does not establish the prime-prefix distribution interface.

\subsection{An explicit quotient weight with the full fourth-prime range}
\label{sec:repaired-chen-weight}

Put $a=100/1149$, $h=50/309$, and $d=1/2-3a$.
Thus $a<h<d<1/3$ and $h+d+2/3>1$.
For a positive integer $n$ set $P^-(1)=+\infty$ and define
\[
 Q_b(y;n)=\mathbf 1_{b\mid n}\,
          \mathbf 1_{P^-(n/b)\ge y}.
\]
The quotient in this definition is essential. In every sum below the
letters $p_i$ denote distinct primes and the prime producing
$n=N-P$ is denoted by $P$, to avoid confusing it with a prefix factor.
The summed prime sequence below is restricted to $(n,N)=1$ for its
analytic estimates. The pointwise weight inequality itself needs no
coprimality condition.

Define functions $Y_i(n)$ as follows:
\begin{align*}
 Y_1(n)&=Q_1(N^a;n),&Y_2(n)&=Q_1(N^h;n),\\
 Y_3(n)&=\sum_{N^a\le p<N^{1/3}}Q_p(N^a;n),&
 Y_4(n)&=\sum_{N^a\le p<N^d}Q_p(N^a;n),\\
 Y_5(n)&=\sum_{N^a\le p_1<p_2<N^h}Q_{p_1p_2}(N^a;n),&
 Y_6(n)&=\sum_{N^a\le p_1<N^h\le p_2<N^d}Q_{p_1p_2}(N^a;n).
\end{align*}
The five switching terms have the explicit tuple definitions
\begin{align*}
 Y_7(n)&=\sum_{N^d\le p_1<p_2<(N/p_1)^{1/2}}
           Q_{p_1p_2}(p_2;n),\\
 Y_8(n)&=\sum_{N^a\le p_1<N^{1/3}\le p_2<(N/p_1)^{1/2}}
           Q_{p_1p_2}(p_2;n),\\
 Y_9(n)&=\sum_{N^h\le p_1<N^d\le p_2<(N/p_1)^{1/2}}
           Q_{p_1p_2}((N/(p_1p_2))^{1/2};n),\\
 Y_{10}(n)&=\sum_{N^a\le p_1<p_2<p_3<p_4<N^h}
           Q_{p_1p_2p_3p_4}(p_2;n),\\
 Y_{11}^{\mathrm{full}}(n)&=
   \sum_{N^a\le p_1<p_2<p_3<N^h\le p_4<N^d}
           Q_{p_1p_2p_3p_4}(p_2;n).
\end{align*}
In particular, the upper endpoint in the last line is $N^d$.
It is not replaced by $N^{1/2-2a}/p_3$.
Put $\Upsilon_i=\sum_{P\le N,\ P\text{ prime},\ (P,N)=1}Y_i(N-P)$,
with the superscript ``full'' retained for the last term.

\paragraph{The elementary two-threshold weight.}
For completeness, the basic Chen weight also has a direct quotient
interpretation. Let $0<\kappa<\sigma\le1/3$, and let $n\le N$ be
squarefree and $N^\kappa$-rough. Define $S_1,S_2,S_3,S_4$ pointwise
by the four tuple sums of Wu's Lemma~2.1, replacing every sieve condition
by the corresponding $Q_b$ condition. If $k=\Omega(n)\ge3$ and
$r$ factors of $n$ are below $N^\sigma$, then
\[
 S_1=r,\qquad S_2=\mathbf1_{r=0},\qquad
 S_3=\mathbf1_{r=1},\qquad S_4=\max(r-2,0).
\]
Indeed, roughness at the second specified prime forces the first two
specified primes to be the smallest two factors of $n$. If these factors
are $q_1<q_2$, then $q_1q_2^2<n\le N$ for $k\ge3$.
The third specified factor in $S_4$ can be any subsequent factor below
$N^\sigma$. Consequently
\[
 2-S_1-2S_2-S_3+S_4=0\qquad(k\ge3).
\]
For $k\le2$ the same expression is at most $2$.
If $n$ is not $N^\kappa$-rough every term is zero. This proves the
basic two-threshold weight after the square-factor exceptions are removed.

\paragraph{A direct proof of the repaired eleven-term weight.}
We prove the pointwise inequality
\begin{align*}
 W(n):={}&3Y_1(n)+Y_2(n)-Y_3(n)-Y_4(n)+Y_5(n)+Y_6(n)\\
         &-2Y_7(n)-Y_8(n)-Y_9(n)-Y_{10}(n)
           -Y_{11}^{\mathrm{full}}(n)
       \le4\mathbf1_{\Omega(n)\le2}
\end{align*}
for every squarefree integer $1\le n\le N$.
Every nonzero term requires $n$ to be $N^a$-rough. For $Y_9$, in
particular, its quotient cutoff exceeds $N^{(1-d)/4}>N^a$.
For $n=1$, the weight is $W(1)=4$. Otherwise we assume this roughness
and list the factors as
$q_1<\cdots<q_k$. Write
\[
 r=\#\{q_i<N^h\},\qquad s=\#\{q_i<N^d\},\qquad
 t=\#\{q_i<N^{1/3}\};\qquad 0\le r\le s\le t.
\]
The first six terms satisfy
\[
 Y_1=1,\quad Y_2=\mathbf1_{r=0},\quad Y_3=t,\quad Y_4=s,
 \quad Y_5=\binom r2,\quad Y_6=r(s-r).
\]
For $r\ge2$, quotient roughness makes the first two specified factors
in $Y_{10}$ and $Y_{11}^{\mathrm{full}}$ equal to $q_1,q_2$. Hence
\[
 Y_{10}=\binom{r-2}{2},\qquad
 Y_{11}^{\mathrm{full}}=(r-2)(s-r),
\]
where $\binom j2=0$ for $j=0,1$. Discarding the remaining nonpositive
terms therefore gives $W(n)\le s-t\le0$.

Suppose now that $k\ge3$ and $r\le1$.
If $r=1$, then $W(n)\le2-t-Y_8$, and $Y_8=1$ when $t=1$;
thus $W(n)\le0$.
If $r=0$ and $s\ge2$, then $W(n)\le4-s-t\le0$.
If $r=s=0$, then $Y_7=1$ and $t\ge1$, since three factors at least
$N^{1/3}$ would have product greater than $N$. Also $Y_8=1$ when
$t=1$, so $W(n)\le2-t-\mathbf1_{t=1}\le0$.
It remains to treat $r=0,s=1$. If $t\ge3$, then
$W(n)\le3-t\le0$. If $t\le2$, the inequalities
\[
 q_1\ge N^h,\qquad q_2\ge N^d,\qquad
 q_3,q_4\ge N^{1/3},\qquad h+d+2/3>1
\]
exclude $k\ge4$. Therefore $k=3$, and the same factor-size bounds give
\[
 q_1q_2q_3^2\ge N^{h+d+2/3}>N.
\]
The strict switching endpoint holds because
$q_1q_2^2<q_1q_2q_3=n\le N$, and the displayed inequality gives
the quotient cutoff in $Y_9$. Thus $Y_9\ge1$; no uniqueness of its
tuple is asserted. When $t=1$ we also have $Y_8=1$, proving
$W(n)\le3-t-1-\mathbf1_{t=1}\le0$.
Finally, if $k\le2$, either $Y_5+Y_6=0$, in which case $W(n)\le4$,
or $Y_5+Y_6=1$, in which case $Y_2=0$ and $Y_3,Y_4\ge2$, giving
$W(n)\le0$. This completes the pointwise proof.

\paragraph{Exceptional integers and the summed inequality.}
The pointwise proof applies also to $n<N^{1-a}$, so this range need
not be discarded. On the entire nonzero rough support every integer
has at most $\lfloor1/a\rfloor=11$ prime factors counted with
multiplicity. Bounding each tuple by the corresponding subset of
these factors gives
\[
 |W(n)|\le 4+2\binom{11}{1}+6\binom{11}{2}
                     +2\binom{11}{4}=1016.
\]
The number of integers having a square factor $q^2$ with prime
$q\ge N^a$ is at most
\[
 N\sum_{q\ge N^a}q^{-2}\le2N^{1-a}\qquad(N^a\ge2),
\]
where summing over all integers instead of primes proves the last
bound. Summing over the primes defining $\Upsilon_i$, and observing
that their good representations are a subset of those counted by
$D_{1,2}(N)$, therefore gives
\begin{align*}
4D_{1,2}(N)\ge{}&3\Upsilon_1+\Upsilon_2-\Upsilon_3-\Upsilon_4
 +\Upsilon_5+\Upsilon_6\\
 &-2\Upsilon_7-\Upsilon_8-\Upsilon_9-\Upsilon_{10}
   -\Upsilon_{11}^{\mathrm{full}}-2032N^{1-a}.
\end{align*}
This proof does not use the final contraction of the fourth-prime range
in Wu's Lemma~2.2.

\paragraph{The corresponding full-range switching mass.}
The last tuple domain is exactly
$a<t_1<t_2<t_3<h<t_4<d$, with quotient roughness threshold $N^{t_2}$.
Its Buchstab argument is bounded below by
\[
 \frac{1-3h-d}{h}=\frac{1-d}{h}-3>\frac53.
\]
Using $\omega(u)=1/u$ on $[1,2]$ and $\omega(u)\le4/7$ for $u\ge2$
therefore permits the uniform bound $\omega(u)\le3/5$ throughout this
domain. Integrating the three other variables gives the following upper
bound for the associated coefficient with switching factor $8$:
\[
 U_{11}^{\mathrm{full}}\le
 \frac{24}{5}\log\frac dh
 \int_a^h \frac{\log(t/a)\log(h/t)}{t^2}\,dt.
\]
The switching estimate itself, and the prime-prefix distribution inputs
for the other terms, are supplied below.

\subsection{The analytic main terms}
For brevity write $U_{11}$ for the full-range upper constant.
Define the following lower and upper main-term constants:
\[
 L_1=\frac{2e^{-\gamma}}a f\left(\frac1{2a}\right),\qquad
 L_2=\frac{2e^{-\gamma}}h f\left(\frac1{2h}\right),
\]
\[
 U(l,u;w)=\frac{2e^{-\gamma}}a\int_l^u
       F\left(\frac{w(t)-t}{a}\right)\frac{dt}{t}.
\]
Write $w_0(t)=1/2$ and set
\[
 U_3=U(a,e;v)+U(e,1/3;w_0),\qquad U_4=U(a,d;v).
\]
The two additional lower constants are
\begin{align*}
 L_5&=\frac{2e^{-\gamma}}a\int_a^h\int_a^t
       f\left(\frac{v(t)-t-u}{a}\right)\frac{du\,dt}{tu},\\
 L_6&=\frac{2e^{-\gamma}}a\int_h^d\int_a^h
       f\left(\frac{v(t)-t-u}{a}\right)\frac{du\,dt}{tu}.
\end{align*}
The corresponding classical-switching upper constants are
\begin{align*}
 U_7&=8\int_d^{1/3}\int_t^{(1-t)/2}
       \frac{\omega((1-t-u)/u)}{tu^2}\,du\,dt,\\
 U_8&=8\int_a^{1/3}\frac{\log(2-3t)}{t(1-t)}\,dt,\\
 U_9&=8\int_h^d\frac{\log((1-t-d)/d)}{t(1-t)}\,dt,\\
 U_{10}&=8\int_{a<t_4<t_3<t_2<t_1<h}
       \frac{\omega((1-t_1-t_2-t_3-t_4)/t_3)}
            {t_1t_2t_3^2t_4}\,dt_1\,dt_2\,dt_3\,dt_4,\\
 U_{11}&=8\int_{a<t_1<t_2<t_3<h<t_4<d}
       \frac{\omega((1-t_1-t_2-t_3-t_4)/t_2)}
            {t_1t_2^2t_3t_4}\,dt_1\,dt_2\,dt_3\,dt_4.
\end{align*}
The sieve estimates proved below give the combined coefficient
\[
 c_* = \frac14\left(3L_1+L_2-U_3-U_4+L_5+L_6
                  -2U_7-U_8-U_9-U_{10}-U_{11}\right).
\]
The rigorous enclosure below is applied after further elementary
overestimates for the upper terms and underestimates for the lower terms.

\paragraph{Boundary and distribution conditions.}
For a one-prime prefix the required square condition is $2t<v(t)$,
valid for $t<e$. At $e$, trim a fixed small terminal strip, use the
classical upper sieve on that strip, and only afterwards send its width
to zero. For a two-prime prefix, the further condition is
$t+2u<v(t)$. It holds throughout $L_5$ and in the part of $L_6$ with
$t\le m$. In the remaining part, if the lower-sieve function is positive,
then $v(t)-t-u>2a>h\ge u$, which proves the required condition.
Where the function vanishes only nonnegativity is needed.

% Integration fragment: no documentclass and no bibliography environment.
% Requires amsmath, amssymb, amsthm, hyperref, and a lemma environment.
% Source numbering below refers to arXiv:2109.02851v2 and
% arXiv:2309.08522v1. This is a derived interface, not a verbatim theorem.

% Integration fragment: no documentclass.
% Insert before Lemma lem:semiprime55 (or before the first HB interface).
% Uses existing bibliography key Lichtman2023 and ordinary lemma/proof.
% No distribution estimate for primes to large moduli is used in this lemma.

\subsection{Siegel--Walfisz stability of the individual rough factors}
\label{sec:rough-sw-stability}

The fact that an individual Heath--Brown variable has coefficient
$1$ or $\mu$ does not by itself check the additional roughness
restriction. The following elementary stability argument supplies
that check, including the auxiliary coprimality and the Mellin twists.

\begin{lemma}[Rough Siegel--Walfisz stability]
\label{lem:rough-sw}
Fix $\xi>0$ and $B\ge0$. Let
\[
 x^\xi\le U\le x,\qquad z_0=x^{1/(\log\log x)^3},
\]
and let $I\subset[U,2U]$ be any interval.
Let $e(n)$ be either $1$ or $\mu(n)$.
Suppose $h$ is a smooth function on $[U,2U]$ such that
\[
 |h(v)|\le1,\qquad
 \int_U^{2U}|h'(v)|\,dv\ll(\log x)^B,
 \qquad |t|\le(\log x)^B.
\]
Set
\[
 \alpha(n)=e(n)\mathbf1_{P^-(n)>z_0}\mathbf1_I(n)
                  h(n)n^{it}.
\]
For every $A>0$, all integers $D,q\ge1$, and every residue $r$
with $(r,q)=1$, one has
\begin{equation}\label{eq:rough-sw-full}
 \left|
 \sum_{\substack{n\sim U\\n\equiv r\pmod q\\(n,D)=1}}\alpha(n)
 -\frac1{\varphi(q)}
       \sum_{\substack{n\sim U\\(n,Dq)=1}}\alpha(n)
 \right|
 \ll_{A,B,\xi}\frac{U\,\tau(D)}{(\log x)^A}.
\end{equation}
The implied constant is uniform in $I,D,q,r,t$ and in the stated
bounds for $h$. Consequently these sequences meet
\cite[Definition 2.1 and Proposition 5.2]{Lichtman2023}'s
Siegel--Walfisz hypothesis. Fixed powers of $\log n$ are also
allowed after normalization by the corresponding power of
$\log(2U)$.
\end{lemma}

\begin{proof}
We first prove the assertion for $D=1$, $h=1$, and $t=0$,
uniformly for every interval $I\subset[U,2U]$.
The classical inputs are the ordinary Siegel--Walfisz estimate
for the unrough sequence $\mu$ and the fundamental lemma of
the one-dimensional sieve. The needed auxiliary-coprimality
version of the former can itself be reduced to the ordinary
version as follows.
For any integer $E\ge1$ the exact convolution identity is
\[
 \mu(n)\mathbf1_{(n,E)=1}
       =\sum_{\substack{v\mid n\\v\mid E^\infty}}\mu(n/v).
\]
On a scale $V$, split this at $v=V^{1/4}$.
For the small $v$, apply ordinary Siegel--Walfisz to $n/v$;
the sum of the factors $1/v$ is at most
$\prod_{p\mid E}(1-1/p)^{-1}\le\tau(E)$.
For the tail use
\[
 \sum_{\substack{v>V^{1/4}\\v\mid E^\infty}}\frac1v
 \le V^{-1/8}\prod_{p\mid E}(1-p^{-1/2})^{-1}
 \le V^{-1/8}\tau(E)^2.
\]
Both its progression contribution and its normalized coprime
mean are therefore $O(V^{7/8}\tau(E)^2)$, by bounding them
with the total absolute mass. Terms $(v,q)>1$ contribute zero
in the primitive residue and its coprime mean.
Thus the auxiliary-coprimality estimate has the fixed divisor
factor $\tau(E)^2$, uniformly on subintervals of $[V,2V]$
and in each fixed logarithmic range of $q$.
This uses no rough Siegel--Walfisz assertion.

\paragraph{Small moduli and the fundamental-lemma error.}
Fix a constant $B_0$ and first take $q\le(\log x)^{B_0}$.
Put $L=U^{1/4}$. For large $x$ one has $z_0<L$, and
\[
 s_0=\frac{\log L}{\log z_0}
       \ge\frac{\xi}{4}(\log\log x)^3.
\]
Take upper and lower fundamental-lemma weights $\lambda^\pm_\ell$
of level $L$, supported on squarefree divisors of $P(z_0)$,
with $|\lambda^\pm_\ell|\le1$, and write
\[
 \rho^\pm(n)=\sum_{\ell\mid n}\lambda^\pm_\ell.
\]
They satisfy
$\rho^-(n)\le\mathbf1_{P^-(n)>z_0}\le\rho^+(n)$.
Endpoint conventions at $z_0$ can be made identical throughout.

In a primitive residue class $r\bmod q$ the local density at a
sifting prime is $1/p$ if $p\nmid q$ and zero otherwise.
Deleting the primes dividing $q$ preserves the uniform
one-sided dimension-one hypothesis of the fundamental lemma.
The main-term difference of the two weights is therefore
$O(V_q(z_0)e^{-s_0})$, with $V_q(z_0)\le1$.
Counting multiples of $\ell$ in $I\cap\{n\equiv r\pmod q\}$
has an error $O(1)$ for each $(\ell,q)=1$, and none occur
when $(\ell,q)>1$. Hence
\begin{equation}\label{eq:rough-beta-gap}
 \sum_{\substack{n\in I\\n\equiv r\pmod q}}
              (\rho^+(n)-\rho^-(n))
 \ll \frac Uq e^{-s_0}+L
 \ll_C U(\log x)^{-C}
\end{equation}
for every fixed $C>0$.
The same bound holds for the average over the primitive residue
classes. This controls the normalized coprime main sum as well.
It is uniform in the endpoints of $I$ and in $q$ in the indicated
range.

Since $|e(n)|\le1$,
\[
 \big|e(n)\mathbf1_{P^-(n)>z_0}-e(n)\rho^+(n)\big|
       \le\rho^+(n)-\rho^-(n).
\]
Thus \eqref{eq:rough-beta-gap} justifies replacing the rough
indicator by $\rho^+$ in both terms of the desired discrepancy.
This absolute-error argument is necessary when $e=\mu$;
no positivity of $\mu$ is asserted.

\paragraph{The ordinary M枚bius Siegel--Walfisz main sum.}
For $e=\mu$, the replacement progression sum becomes
\[
 \sum_{\substack{\ell\le L\\(\ell,q)=1}}
   \lambda^+_\ell\mu(\ell)
 \sum_{\substack{m\in I/\ell\\
             m\equiv r\bar\ell\pmod q\\(m,\ell)=1}}\mu(m),
\]
using the exact identity
$\mu(\ell m)=\mu(\ell)\mu(m)\mathbf1_{(m,\ell)=1}$
for squarefree $\ell$.
The normalized primitive-class average is the identical outer
sum with the inner sum replaced by
\[
 \frac1{\varphi(q)}
       \sum_{\substack{m\in I/\ell\\(m,\ell q)=1}}\mu(m).
\]
Here $U/\ell\ge U^{3/4}\ge x^{3\xi/4}$, so
$\log(U/\ell)$ is comparable with $\log x$ and
$q\le(\log x)^{B_0}$ lies in a fixed logarithmic
Siegel--Walfisz range for every inner sum.

Use the unrough auxiliary-coprimality estimate just derived:
with the fixed divisor exponent $C_0=2$ and any requested
logarithmic saving $C$ its inner discrepancy is
\[
 \ll_{C,B_0,\xi}
       \frac{(U/\ell)\tau(\ell)^{C_0}}{(\log x)^C}.
\]
The same estimate holds on arbitrary subintervals of
$[U/\ell,2U/\ell]$ by subtracting endpoint estimates.
Summing over $\ell$ and using the fixed divisor moment
\[
 \sum_{\ell\le L}\frac{\tau(\ell)^{C_0}}{\ell}
                   \ll(\log L)^{C_1}
\]
gives $O_A(U(\log x)^{-A})$ after choosing
$C>A+C_1+1$. The divisor exponent here is the fixed exponent
in the classical auxiliary-coprimality input; it is not
allowed to grow with $\ell$.
This calculation explains explicitly the auxiliary condition
$(m,\ell)=1$ introduced by the identity for $\mu(\ell m)$.

For $e=1$, expand $\rho^+$ in the same way without $\mu$.
For each $(\ell,q)=1$, counts of $m\in I/\ell$ in any two
residue classes modulo $q$ differ by $O(1)$.
Thus their difference from the primitive-class average
is also $O(1)$, and summing the weights costs $O(L)$.
Together with \eqref{eq:rough-beta-gap}, this proves the
unweighted assertion for all small $q$ in both cases.

\paragraph{Large moduli.}
For every sequence under consideration with absolute value
at most one, the discrepancy without auxiliary $D$ is trivially
\[
 \ll \frac Uq+1+\frac U{\varphi(q)}
       \ll Uq^{-1/2}+1.
\]
Choose $B_0>2A+4$. For $q>(\log x)^{B_0}$ this is
$O_A(U(\log x)^{-A})$, since $U\ge x^\xi$.
Together with the small-modulus proof, this gives the estimate
for all $q\ge1$, uniformly on every interval $I$.

\paragraph{An arbitrary auxiliary integer.}
Return to any $D\ge1$.
Primes at most $z_0$ dividing $D$ cannot divide a rough $n$.
The number of remaining integers of $[U,2U]$ removed by
$(n,D)=1$ is at most
\[
 \sum_{\substack{\ell\mid D\\\ell>z_0}}
           \left(\frac U\ell+1\right)
       \le\left(\frac U{z_0}+1\right)\omega(D).
\]
The change in the progression term plus its normalized
coprime mean is at most twice this bound, because the
coefficients are bounded by one and $1/\varphi(q)\le1$.
For every fixed $A$,
\[
 \left(\frac U{z_0}+1\right)\omega(D)
        \ll_{A,\xi}\frac{U\,\tau(D)}{(\log x)^A}.
\]
This is uniform for arbitrarily large $D$, and proves the
complete auxiliary-coprimality assertion. It does not
impose a hidden size restriction on that auxiliary integer.

\paragraph{Mellin twists and smooth coefficients.}
The estimates just obtained hold for every subinterval
of $[U,2U]$.
On this interval the function $h(v)v^{it}$ has total variation
\[
 \ll \int_U^{2U}|h'(v)|\,dv
                 +|t|\int_U^{2U}\frac{dv}{v}
 \ll(\log x)^B.
\]
Partial summation therefore costs at most a fixed power of
$\log x$. Apply the untwisted estimate with a correspondingly
larger arbitrary saving to obtain \eqref{eq:rough-sw-full}.
The interval indicator has only two endpoint jumps.
A normalized fixed power of $\log v$ has bounded variation
here, so the final assertion follows in the same manner.
\end{proof}

\paragraph{Use in the Heath--Brown arguments.}
In the preceding or following factorization arguments, take
$\xi=2\zeta$ for a selected individual variable with
$x^{2\zeta}\le U\le x^{1/3+\zeta/2}$.
Its coefficient after Mellin separation is precisely one of
the sequences in Lemma~\ref{lem:rough-sw}, after harmless
normalization of logarithms.
One must select an individual $1$ or $\mu$ slot; this lemma
does not assert Siegel--Walfisz for an arbitrary grouped
convolution.

Before Mellin separation, discard dyadic Heath--Brown boxes
which do not meet the original prime-variable interval.
If an original factor lies in $[R,2R]$, the product of its
at most six dyadic lower endpoints is then at least
$R/2^6$ and at most $2R$.
These group-scale inequalities remain true for the retained
boxes after Mellin separation.
They justify the assertion that, in the absence of medium
slots, each original group of size at least $N^a/2$
must contain a long slot when $12\zeta<a/2$.
At most two long slots can occur globally because three
would have product exceeding $x^{1+3\zeta/2}$.

\subsection{The interface for a fixed prime prefix}
\label{sec:prefix-interface}

There are two different weights in this argument. The weight used to
sieve $\mathcal A_b$ is a weight on the remaining divisor $q$, of level
approximately $D/b$. It is not identified with the restriction to
$d=bq$ of a single global weight $\widetilde\lambda_D^{\pm}(d)$.
In particular, the signs and the prefix conditions in the ordinary
Rosser weights change when a prime is removed. For example, for
$q\mid P(z)$ and $z<p$, the relation
$-\lambda_D^-(pq)=\lambda_{D/p}^+(q)$ requires $p^2<D$;
the opposite-sign relation involves $p^3<D$.
For two prefix primes a corresponding lower-weight restriction
can require $p_1p_2^3<D$, which is stronger than
$p_1p_2^2<D$. Thus the assertion below is derived at the level of
the individual box atoms in the sieve construction.

For precise references write $L_{21}$ for
\href{https://arxiv.org/pdf/2109.02851v2}{Lichtman,
\emph{A modification of the linear sieve, and the count of twin primes}},
and $L_{23}$ for
\href{https://arxiv.org/pdf/2309.08522v1}{Lichtman,
\emph{Primes in arithmetic progressions to large moduli, and Goldbach
beyond the square-root barrier}}.
The assertion below is a consequence of their proofs, with the
additional convolution argument supplied here; it is not the literal
statement of $L_{23}$, Proposition~6.6.

Put $\sigma=7/32$ and
\begin{equation}\label{eq:prefix-v}
 v(t)=\min\left\{\frac{1+t}{2},
       \frac{1-(3/2-2\sigma)t}{1+\sigma},
       1-\frac{\sigma+3t}{2}\right\}.
\end{equation}
For $N\in\mathbb N$ let
\[
 E_N(d)=\pi(N;d,N)-\frac{\pi(N)}{\varphi(d)}
 \qquad ((d,N)=1).
\]

\begin{lemma}[Derived prime-prefix sieve interface]
\label{lem:prefix-interface}
Fix $r\in\{1,2\}$ and positive constants $\delta,h,u,t_1,\ldots,t_r$,
and let
\[
 I_i=(N^{t_i},N^{t_i+h}],\qquad
 T_+=\sum_{i=1}^r(t_i+h),\qquad z=N^u.
\]
Assume that the intervals are ordered and disjoint, with every member
of $I_i$ larger than every member of $I_{i+1}$, and $u<t_r$.
Let $\nu$ satisfy
\begin{align}
 \nu+20\delta&\le\inf_{t\in[t_1,t_1+h]}v(t),
       \label{eq:prefix-level-margin}\\
 \sum_{i<j}(t_i+h)+2(t_j+h)&\le\nu-20\delta
       \quad(1\le j\le r),
       \label{eq:prefix-vector-margin}\\
 \nu-T_+-20\delta&>0.
       \label{eq:prefix-residual-positive}
\end{align}
Set $R=N^{\nu-T_+-20\delta}$. Choose the standard modified
linear-sieve weights $\rho_R^{\pm}$ on the remaining divisor $q$,
with sieve precision sufficiently small in terms of the fixed margins.
Use weights independent of the particular primes inside each box.
Then, for every $A>0$,
\begin{equation}\label{eq:prefix-remainder}
 \left|\sum_{\substack{p_i\in I_i\ (1\le i\le r)\\
                         b=p_1\cdots p_r, (b,N)=1}}
       \sum_{\substack{q\mid P(z)\\(q,N)=1}}
       \rho_R^{\pm}(q)E_N(bq)\right|
 \ll_{A,\delta,h,u,\boldsymbol t}N(\log N)^{-A}.
\end{equation}
The constant is independent of the even integer $N$.
The weights can be chosen with the usual linear-sieve main terms
$F(\log R/\log z)$ and $f(\log R/\log z)$, up to the error tending
to zero with the sieve precision. Here it suffices to use the lemma
when $\log R/\log z\ge2$; a zero lower bound requires no distribution
estimate.
\end{lemma}

\begin{proof}
We keep $\delta,h,u,\boldsymbol t$ fixed while $N$ tends to infinity.
Choose the construction precision $\eta>0$ and the associated exponent
mesh sufficiently small compared with $\delta$ and $u$.
All losses below are fixed powers $N^{O(\eta)}$, not factors that are
silently absorbed into an asymptotic error with $\eta$ fixed.
They are paid for by the explicit margins in
\eqref{eq:prefix-level-margin}--\eqref{eq:prefix-residual-positive}.

We use the construction in $L_{21}$, equations~(5.14)--(5.17),
with residual level $R$, rather than imposing a restriction on a
global weight of level $N^\nu$. Its remainder is a finite linear
combination of atoms of the form
\begin{equation}\label{eq:prefix-atom}
 \chi_{J_1}*\cdots*\chi_{J_m}*\beta,
\end{equation}
where $\chi_J$ denotes the prime indicator of a construction interval,
all the displayed primes exceed a cutoff $w$, the small-prime weight
$\beta$ is supported on $k\mid P(w)$ and $k\le R^{O(\eta)}$, and
$|\beta(k)|\le1$. The cutoff is a fixed positive power of $N$,
depending on $\eta$. The number of atoms and the sum of the absolute
values of their scalar coefficients are bounded in terms of $\eta$.
The construction may shrink the rough-part level before adding the
small-prime weight; this changes the residual exponent by $O(\eta)$.
The vectors of rough-prime intervals belong to
$D^{\mathrm{well}}(R^{1+O(\eta)})$. These assertions follow from
the ordinary modified linear sieve, not from its further modification
using $F^*$ in $L_{21}$.

Convolve each atom with $\chi_{I_1}*\cdots*\chi_{I_r}$.
Because $z<\inf I_r$, the prefix remains the largest $r$ primes.
The first $r$ conditions defining $D^{\mathrm{well}}$ hold by
\eqref{eq:prefix-vector-margin}. For a later rough prime $q_j$,
the corresponding condition follows from
\[
 (p_1\cdots p_r)(q_1\cdots q_{j-1}q_j^2)
 \le N^{T_+}R^{1+O(\eta)}
 \le N^{\nu-19\delta}
\]
after making $\eta$ small enough. Refining or rounding all intervals
to the common exponent mesh costs another $N^{O(\eta)}$.
Consequently every concatenated rough vector is in
$D^{\mathrm{well}}(N^{\nu-c\delta})$ for an absolute $c>0$.
Its largest prime has exponent in $[t_1,t_1+h+O(\eta)]$.

Apply $L_{23}$, Proposition~6.1, to a representative of each such
rough box. The partition of the prime indices furnished by that
proposition is kept fixed throughout that box. The exponent mesh
then enlarges each of the three factor supports by only
$N^{O(\eta)}$; it does not require choosing a partition separately
for every integer in the box. This is the step implemented in
$L_{21}$, Lemma~5.3. Its proof uses the factorization of a rough
box and the size of the small-prime weight, not membership of the
concatenated vector in a global $D^+$ or $D^-$.

For completeness, the treatment of signs, repeated intervals, and
coefficients can be made explicit. Write an ordered-tuple rough atom
as $T_J=\chi_{J_1}*\cdots*\chi_{J_m}$, with the squarefree restriction
on the rough product understood. Retain exactly its scalar
coefficient from the sieve construction, including the original
factorials for repeated interval labels in $L_{21}$, (5.10).
For an explicit convention, let $m_I$ be the multiplicity of each
interval label and put $M_J=\prod_I m_I!$. If $A_J$ is the indicator
of the corresponding sorted squarefree prime-box product, then
\begin{equation}\label{eq:box-symmetry-normalization}
 \mathbf1_{\rm squarefree}T_J=M_J A_J.
\end{equation}
Indeed each squarefree product has exactly $m_I!$ assignments to
the slots carrying label $I$. With the tuple-convolution convention
of $L_{21}$, (5.10), the scalar already includes
$(-1)^m/M_J$, in addition to any other original scalar factors.
Keep this reciprocal factorial exactly once; it is not inserted a
second time during the three-group normalization below.
If $\mathcal I_1,\mathcal I_2,\mathcal I_3$ is the fixed partition
of its prime slots (including the prefix), put
\[
 K_i=|\mathcal I_i|!,\qquad
 \gamma_i=K_i^{-1}\mathop{*}_{j\in\mathcal I_i}\chi_{J_j}.
\]
The empty convolution is the point mass at $1$. Unique prime
factorization bounds the number of ordered representations by
$K_i$, even on products containing repeated primes, so
$|\gamma_i(n)|\le1$ for every integer $n$. The unrestricted
convolution identity is exactly
\[
 T_J=(K_1K_2K_3)\,\gamma_1*\gamma_2*\gamma_3.
\]
Absorb the displayed, fixed factorial product into the original
scalar coefficient. The extra nonsquarefree rough products will
be removed by the error estimate below; no claim of a unique
ordered prime-slot representation is needed.
The small-prime factor is combined with the third rough factor.
The factorization of an integer into its primes
at most $w$ and its primes greater than $w$ is unique. Hence, if
$\gamma_3$ is the normalized third rough factor,
\[
 |(\gamma_3*\beta)(n)|
 =|\gamma_3(n_{>w})\beta(n_{\le w})|\le1.
\]
This argument allows either sign of $\beta$ and does not assume
positivity. Its support is enlarged by at most $R^{O(\eta)}$;
the degree of every factor-size inequality in Proposition~6.1 is
bounded, so that the resulting loss is again $N^{O(\eta)}$.

One can also implement the interval normalization using unrestricted
prime convolutions and remove collisions afterwards. This gives
exact convolution identities away from integers divisible by
$p^2$ for a rough prime $p>w$. The omitted discrepancy is at most
\[
 \sum_{\substack{d\le N^{\nu+O(\eta)}\\
                  p^2\mid d\ \mathrm{for\ some}\ p>w}}
   \tau_C(d)\left(\frac{N}{d}+1+\frac{N}{\varphi(d)}\right)
 \ll_C (N+N^{\nu+O(\eta)})w^{-1}(\log N)^{C'},
\]
where $C,C'$ are fixed with the precision. Since $w$ is a fixed
positive power of $N$ and $\nu<1$, this is
$O_A(N(\log N)^{-A})$ for every fixed $A$. Thus the squarefree
convention in the sieve causes no unaccounted convolution error.

We spell out the source of uniformity in the residue, since the
displayed constant in $L_{23}$, Proposition~6.6, is written with
a subscript $a$. The input used for a box atom is $L_{23}$,
Proposition~5.2, not an inference from that subscript. Proposition~5.2
has the explicit supremum $0<|a|<x^{1+\epsilon}$ and permits arbitrary
complex, $1$-bounded coefficients on its three modulus factors.
With its Type~II length denoted here by $U$, its size conditions are
\begin{align*}
 Q_1&<Ux^{-\epsilon},&
 U^2Q_2Q_3^2&<x^{1-8\epsilon},\\
 |a|^\sigma UQ_1Q_2^5Q_3^2&<x^{2-15\epsilon},&
 (UQ_3/Q_2)^{2\sigma}UQ_1Q_2^5Q_3^2
     &<x^{2-15\epsilon}.
\end{align*}
Here is a quantitative check of the parameters, including the
endpoints of the Type~II range. Fix a factorization margin
$0<\zeta<10^{-5}$, sufficiently small in terms of $\delta$, and take
$\eta\ll\epsilon\ll\zeta$. For
\[
 x^{2\zeta}\le U\le x^{1/3+\zeta/2},\qquad A_0=Ux^{-\zeta},
\]
we have $x^\zeta\le A_0\le x^{1/3-\zeta/2}$, exactly the allowed
range of the parameter $A$ in Proposition~6.1. Denote its three
factors by $a_0,b_0,c_0$. Apart from $x^{O(\eta)}$ box enlargements,
put $Q_1=a_0$, $Q_2=b_0$, and $Q_3=c_0$ and then split into dyadic
intervals. In particular $Q_1\le A_0<Ux^{-\epsilon}$.
Replacing $A_0$ by $U$ in the three monomials in (6.3) costs,
respectively,
\[
 x^{2\zeta},\qquad x^\zeta,\qquad
 x^{(1+2\sigma)\zeta}.
\]
Their original savings are $3\zeta,3\zeta,2\zeta$.
Thus the smallest remaining saving is
$(1-2\sigma)\zeta=9\zeta/16>0$.
Choosing, for example, $\epsilon<\zeta/1000$ and then the construction
mesh sufficiently small makes all four displayed conditions strict.
If a dyadic value of $Q_2$ is smaller than its support endpoint,
the last monomial is still bounded by its endpoint value: its power
of $Q_2$ is $5-2\sigma>0$. This justifies the same conclusion on
every dyadic subbox.

There are two additional bookkeeping points in passing to primes.
First, the primes in $\pi(N)$ cannot all be treated on the artificial
scale $x=N$. Remove the initial range $p\le N^{1-\kappa}$, with a
fixed $0<\kappa\ll\epsilon$. Its progression term is bounded by
\[
 \sum_{p\le N^{1-\kappa}}\tau_C(N-p)
       \ll N^{1-\kappa+o(1)},
\]
and its mean term is $O(N^{1-\kappa}(\log N)^{C'})$.
On the remaining dyadic intervals $[x,2x]$, use
$\alpha=\log N/\log x=1+O(\kappa)$ in (6.3).
The exponent coordinates, total level, and the function
$v_\alpha(t)$ from (6.2) differ from their $N$-scale counterparts
by $O(\kappa)$. These changes are paid for by the fixed margins.
Also $N<x^{1+\epsilon}$ after making $\kappa$ small enough.
Thus Proposition~5.2 applies with residue $a=N$ on every interval;
it is not enough merely to set $\alpha=1$ on the first interval.
If the full support is at most $N^{1/2-c\delta}$, ordinary
Bombieri--Vinogradov suffices for each bounded atom. Only the
remaining case needs Proposition~6.1, whose section assumes its
level parameter is at least $1/2$.

Second, apply the $k=3$ Heath--Brown identity of $L_{23}$,
Lemma~3.3, as in the proof of Theorem~1.7, but with middle interval
$[x^{2\zeta},x^{1/3+\zeta/2}]$. The six dyadic variables are
$N_1,N_2,N_3,M_1,M_2,M_3$, with $M_i\le(2x)^{1/3}$.
If a variable lies in that interval, take it as the Type~II variable
$U$ and combine the other five. This uses exactly the coefficient
sequences of the cited Heath--Brown argument, rather than a new
Siegel--Walfisz assertion for an arbitrary convolution. The chosen
sequence is $1$ or $\mu$, with the roughness cutoff
$P^-(n)>x^{1/(\log\log x)^3}$ used there; the logarithm is handled
by partial summation. Lemma~\ref{lem:rough-sw} verifies the full
Siegel--Walfisz condition~(2.2), including its auxiliary coprimality
and the rough cutoff. Divisor-bounded
coefficients on the other factor are reduced to $1$-bounded ones
by the truncation argument in the proof of Proposition~3.1, using
Lemmas~4.1 and~4.3; the resulting powers of $\log x$ are absorbed
by increasing the arbitrary exponent in the error estimate.

If no variable lies in the middle interval, every $M_i$ is less
than $x^{2\zeta}$, and at most two of the $N_i$ exceed
$x^{1/3+\zeta/2}$. If just one does, the product of the other five
is at most $x^{10\zeta}$, so the long variable is at least
$x^{1-10\zeta+o(1)}$ and Lemma~3.4 applies. If two do, the
product of the remaining four is at most $x^{8\zeta}$; apply
Proposition~3.2, which has an explicit supremum over all $a\ne0$.
The total modulus support is at most $x^{3/5}$, after making the
fixed margins sufficiently small: indeed $\max_{t\ge0}v(t)
=153/256<3/5$. For small $\zeta,\epsilon$ this is less than
$(x/x^{8\zeta})^{2/3-3\epsilon}$, as required by that proposition,
and less than the long variable times $x^{-\epsilon}$ as required
by Lemma~3.4. Having no large variable is impossible because the
product of all six variables is comparable with $x$.
These arguments also allow the sharp product cutoff on a dyadic
interval, by the same separation and partial summation used in
the proof of Proposition~3.1. Restrictions $(d,N)=1$ split into
$(q_i,N)=1$ on the modulus factors and preserve their bounds.
No fixed-residue estimate from $L_{21}$, Theorem~2.5, is used here.

We obtain a remainder $O_A(N(\log N)^{-A})$ for each atom, uniformly
in $N$. Summing the fixed number of normalized atoms proves
\eqref{eq:prefix-remainder}. Notice that the signs of the atoms were
never used to deduce a distribution estimate. They only enter when
reassembling the residual upper or lower sieve weights.

Finally, the main term is evaluated before convolving in the prefix.
For each $b=p_1\cdots p_r$ the residual density is
$g_N(\ell)=1/(\ell-1)$ for $\ell\nmid N$ and $g_N(\ell)=0$
for $\ell\mid N$, because $(b,q)=1$. Uniformity in $N$ here
requires a minor precaution: the two-sided condition (2.3) of
$L_{21}$ need not hold with an absolute constant all the way down
to $2$, since $N$ can contain many small primes.
Use the small-prime presieve already present in (5.15), with
$w=R^{\eta^2}$ and small-factor level $R^{O(\eta)}$.
Below $w$ the one-sided dimension-one estimate is uniform, since
deleting prime factors from the usual product only decreases its
ratio. The fundamental lemma therefore gives the full factor
$V_N(w)$, with a relative error tending to zero with $\eta$.
Above $w$ the two-sided estimate is uniform: for $w\le y<z$,
\[
 \sum_{\substack{\ell\mid N\\\ell\ge w}}
       \frac1{\ell-1}
 \ll\frac{\log N}{w\log w},\qquad
 \frac{V_N(y)}{V_N(z)}
 =\frac{\log z}{\log y}
   \left(1+O\!\left(\frac1{\log y}
                    +\frac{\log N}{w\log w}\right)\right).
\]
Consequently the usual rough-box evaluation in Theorem~2.10
applies with constants independent of $N$. Reattaching the
presieve gives the usual $F(s),f(s)$, with
$s=\log R/\log z$, and an error tending to zero with the precision.
The entire product $V_N(z)$ is retained; it is not replaced by a
product with the prime divisors of $N$ restored. On summing over
the prefix box the main factor is
\[
 \pi(N)V_N(z)
 \sum_{\substack{p_i\in I_i\ (1\le i\le r)\\(p_1\cdots p_r,N)=1}}
       \frac1{\varphi(p_1\cdots p_r)},
 \qquad
 V_N(z)=\prod_{\substack{\ell<z\\\ell\nmid N}}
       \left(1-\frac1{\ell-1}\right).
\]
Its upper and lower multipliers are respectively
$F(s)+O(\eta)$ and $f(s)+O(\eta)$, with the stated precision
error interpreted on a fixed compact range of $s\ge2$.
This evaluation uses weights constructed for the residual sequence;
it does not replace their main term by a global-prefix weight sum.
\end{proof}

\paragraph{Order of limits and use in the integrals.}
For each fixed positive level margin $\delta$, first choose a finite
exponent partition and sieve precision so that all box enlargements
are smaller than that margin, and then let $N\to\infty$.
Next let the partition width and sieve precision tend to zero,
still within the fixed margin, and finally let $\delta\downarrow0$.
Continuity of $F,f$ on the relevant compact intervals gives the
arguments $(v(t_1)-t_1)/u$ and
$(v(t_1)-t_1-t_2)/u$. Prefix boxes meeting ordering boundaries
can be discarded for lower bounds, and bounded with a classical
sieve for upper bounds; their main contribution tends to zero with
the partition width. A boundary on which the prefix
$D^{\mathrm{well}}$ condition is an equality is treated by a thin
strip before taking the limit, not by applying the lemma at equality.

\paragraph{The middle branch in the level.}
The two-branch form of $v$ used in the numerical calculation is exact.
Put $A(t)=(1+t)/2$, $B(t)=(32-34t)/39$, and
$C(t)=57/64-3t/2$. Then
\[
 B-A=\frac{25-107t}{78},\qquad
 C-A=\frac{25}{64}-2t,\qquad
 B-C=\frac{1568t-175}{2496}.
\]
For $t\le25/128$ the minimum is $A$, whereas for $t\ge25/128$
the minimum is $C$ (the zero of $B-C$ is $25/224<25/128$).
Thus for every $t\ge0$,
\[
 v(t)=\begin{cases}
 (1+t)/2,&t\le25/128,\\
 57/64-3t/2,&t\ge25/128.
 \end{cases}
\]

\paragraph{Status of this interface.}
The lemma supplies an adaptation of the box-atom proof, rather than
quoting a published prime-prefix theorem verbatim. The distribution
and linear-sieve theorems cited above are external inputs; the extra
prefix, coefficient, and uniformity arguments have been supplied here.
\subsection{Completing the conservative calculation}

Classical Bombieri--Vinogradov and the lower linear sieve, first at
each fixed level $N^{1/2-\delta}$ and then with $\delta\downarrow0$,
give $\Upsilon_i\ge(L_i-o(1))\HN$ for $i=1,2$.
The uniform presieve argument in Lemma~\ref{lem:prefix-interface}
also applies to these classical estimates.
For the prime-prefix terms, that lemma, the boundary conditions already
checked, and partial summation give
\[
 \Upsilon_3\le(U_3+o(1))\HN,\quad
 \Upsilon_4\le(U_4+o(1))\HN,\quad
 \Upsilon_5\ge(L_5-o(1))\HN,\quad
 \Upsilon_6\ge(L_6-o(1))\HN.
\]
Here the one-prime density $1/(p-1)$ can be replaced by $1/p$ at
negligible cost. Omitting primes dividing $N$ changes these compact
positive-exponent integrals by $o(1)$ uniformly, since there are only
boundedly many such primes above $N^a$. The band $e\le t\le1/3$
in $U_3$ uses only the classical upper sieve. If its residual sieve
parameter $s=(1/2-t)/a$ is below two, lower the sifting point to the
square root of the residual level. Since $F(s)=2e^\gamma/s$ for
$1\le s\le3$, this gives the same limiting upper coefficient
$4/(1/2-t)$ used in $U_3$. All level margins are sent to zero after
the large-$N$ limit.

\paragraph{Why the switching coefficient is eight.}
In each of $\Upsilon_7,\ldots,\Upsilon_{11}^{\mathrm{full}}$, retain the multiplicity
of the factorization after switching. Terms with a repeated specified
prime, or with an original prime dividing $N$, cost
$O(N^{1-a}(\log N)^{O(1)})$. Every remaining factor of a switched
integer is at least $N^a$, and the number of factors is bounded in terms
of $a$. Discard integers below $N/(\log N)^K$, with $K$ arbitrarily
large, at negligible cost. Split into the finitely many possible numbers
of prime factors, and use the multiplicative grid argument from
Lemma~\ref{lem:bv}, preserving multiplicities. In each rectangular
piece, one prime-interval indicator supplies the Siegel--Walfisz factor;
the convolution of all other indicators is divisor-bounded.
This gives generalized Bombieri--Vinogradov for any fixed level
$N^\theta$, $\theta<1/2$.

Here the exact coprime mass can differ from the total switched mass $X$.
Replacing it incurs, after summing over the squarefree moduli, at most
\[
 \ll (\log N)\sum_j w(j)\sum_{\ell\mid j}\frac1{\ell-1}
 \ll N^{1-a}(\log N)^{O(1)},
\]
where $w(j)$ is the bounded multiplicity. For the first inequality use
$\varphi(\ell q)=\varphi(\ell)\varphi(q)$ on the squarefree moduli.
Thus the upper linear sieve at $z=N^{1/5}$ gives
$(4/\theta+o(1))\CN X/\log N$. First take the large-$N$ limit,
then $\theta\uparrow1/2$. This gives the coefficient eight in each
displayed switched mass, without the stronger distribution lemma or
the computed constant $G_1$ in \cite{Li2026}.

To evaluate those masses, sum over the specified prime prefixes and
apply the standard Buchstab rough-number estimate to the remaining
quotient. The case in which the quotient equals one is negligible:
the five possible prefix products are respectively at most
$N^{2/3},N^{2/3},N^{(1+d)/2},N^{4h},N^{3h+d}$, all fixed powers
smaller than $N$. For $U_7,U_8$, first trim a fixed-width exponent
strip at Buchstab argument one. The upper estimate
$\Phi(X,y)\ll X/\log y$ bounds its mass by $O(\eps)N/\log N$;
apply the asymptotic away from that strip and finally let its width
tend to zero. The other domains stay away from that endpoint.
Partial summation now gives the displayed masses, with all tuple
multiplicities retained. In $U_8$ the Buchstab argument is at most
two, so $\omega(u)=1/u$ gives its elementary one-dimensional form.
For $U_9$ the quotient sifted to its square root is prime, apart from
the already removed unit case; its prime-number-theorem mass gives
the displayed logarithmic integral. Consequently
$\Upsilon_i\le(U_i+o(1))\HN$ for all five switching terms.

\paragraph{Elementary numerical enclosures.}
The four-dimensional integrals admit convenient overestimates. For
$s\ge2$ one has $\omega(s)\le4/7$. On $[2,3]$, this follows from
$\omega(s)=(1+\log(s-1))/s$: the maximum of
$1+\log(s-1)-4s/7$ is $\log(7/4)-4/7<0$. For instance the positive
logarithm series with argument $3/11$ proves the last strict inequality.
The relation $(s\omega(s))'=\omega(s-1)$ then propagates the bound
over each successive unit interval. All arguments of $\omega$ in
$U_{10}$ exceed two, since $6h<1$. For $U_{11}$ the argument is
at least $(1-d)/h-3>5/3$, so the bound $\omega\le3/5$ applies,
using $1/u$ below two and $4/7$ above two.
Integrating the elementary factors gives
\begin{align*}
 U_{10}&\le\frac{16}{7}\int_a^h
          \frac{\log^2(h/t)\log(t/a)}{t^2}\,dt,\\
 U_{11}&\le\frac{24}{5}\log\frac dh\int_a^h
      \frac{\log(t/a)\log(h/t)}{t^2}\,dt.
\end{align*}
Similarly,
\[
 U_7\le8\int_d^{1/3}\frac{\log((1-2t)/t)}{t(1-t)}\,dt
       +\frac{32}{3d^2}(1/4-d)^2\log(1/d-3).
\]
Only the triangular region $d<t<1/4$, $t<u<(1-t)/3$ has a
Buchstab argument above two. Its area is $2(1/4-d)^2/3$;
there $tu(1-t-u)\ge d^2/2$ and the extra logarithm is at most
$\log(1/d-3)$, proving the last bound.

No numerical differential-equation solution is needed. Define
\[
 J(s)=\begin{cases}\displaystyle\int_2^{s-1}
                   \frac{\log(t-1)}t\,dt,&s>3,\\0,&s\le3.
       \end{cases}
\]
For $1\le s\le5$, $e^{-\gamma}F(s)=2(1+J(s))/s$;
for $s>5$ the value at five is an upper bound. Also
$e^{-\gamma}f(s)\ge2\log(s-1)/s$ for $s>2$, with zero below two.
For $s_1=1/(2a)<6$, the exact formula is
\[
 L_1=8\left(\log(s_1-1)+\int_3^{s_1-1}\frac{J(v)}v\,dv\right),
 \qquad L_2=8\log(1/(2h)-1).
\]
Consequently an outward-rounded Darboux calculation with explicit
logarithm bounds encloses every term. The certificate uses fixed-point
integer intervals with denominator $2^{80}$; logarithms are reduced to
$\log 2$ and $2\operatorname{artanh}(u)$ with $0\le u\le1/3$.
After 32 terms the omitted tail is at most
$2u^{65}/(65(1-u^2))$. The elementary operations are rounded outwards
using integer floor and ceiling division. A 20,000-cell monotone
prefix table encloses $J$, 2,000 cells enclose each one-dimensional
integral, and $256^2$ rectangular cells enclose each two-dimensional
integral after mapping its domain to a rectangle. No floating-point
quadrature contributes to these bounds.

The resulting exact lower endpoint for the elementary replacement
functional is
\[
 \frac{1475894317448524132649014}{1208925819614629174706176}
 >\frac{61}{50}.
\]
For a shorter check, outward rounding of the individual certificates
gives the following rational bounds (every terminating decimal is exact):
\begin{center}
\begin{tabular}{ll@{\qquad}ll}
\toprule
Lower term&Certified lower bound&Upper term&Certified upper bound\\
\midrule
$L_1$&12.8990&$U_3$&19.3178\\
$L_2$&5.8973&$U_4$&13.1277\\
$L_5$&2.3441&$U_7$&1.5373\\
$L_6$&2.2784&$U_8$&4.5967\\
&&$U_9$&3.3285\\
&&$U_{10}$&0.2553\\
&&$U_{11}$&0.6338\\
\bottomrule
\end{tabular}
\end{center}
These weaker rounded entries already give
$c_*\ge48824/40000>61/50$. The preceding analytic estimates and
the directly proved quotient weight therefore imply
$D_{1,2}(N)\ge(c_*-o(1))\HN\ge(61/50)\HN$ eventually.
The numerical certificate establishes the integral inequality; the
number-theoretic premises are supplied by the proofs above.

For the imbalanced conclusion, set $b=43/50$, $\theta=4999/10000$,
and $\vartheta=9999/20000$. For $u=7/93$,
\[
 \log(50/43)
 \le2\left(u+\frac{u^3}{3}+\frac{u^5}{5(1-u^2)}\right)
 =\frac{7824846421}{51881026500}.
\]
Hence
\[
 \frac{61}{50}-\frac{40000}{4999}\log(50/43)
 \ge\frac{341710995767}{25935325147350}>\frac1{100}.
\]
The strict gap absorbs the asymptotic errors in the transfer theorem,
yielding $D_{0.86}(N)\ge\HN/100$ for sufficiently large even $N$.
This completes Theorem~\ref{thm:conservative}.
The result is stronger than the explicitly stated $b=0.894$
comparison, but weaker than the consequence obtained by accepting the
published $1.9728$ input. It is therefore not presented as fulfilling
the requested advance beyond all existing results.

% Integration fragment: no documentclass. Uses the main document's
% lemma/proof environments, \CN, \HN and bibliography keys.
% Derived argument, not a quotation of a published semiprime theorem.

\section{A derived uniform upper sieve for balanced semiprimes}
\label{sec:semiprime-upper}

This section supplies a separate derivation for the product sequence.
It does not apply a prime distribution theorem directly to two prime
variables of lengths approximately $N^{1/2}$. Write $L_{21}$ for the
author version arXiv:2109.02851v2 of \cite{Lichtman2025}, and $L_{23}$
for arXiv:2309.08522v1 of \cite{Lichtman2023}. Numbered propositions
below refer to those versions.

\begin{lemma}[Derived smooth semiprime distribution and upper bound]
\label{lem:semiprime55}
Fix $1/2<b<1$ and $0<\varepsilon_0<1/100$, and put
\[
 \nu=\frac{11}{20}-\varepsilon_0,\qquad z=N^{1/5}.
\]
Let $\rho_D^+$ be the ordinary Iwaniec modified upper linear-sieve
weight of level $D=N^\nu$, restricted to divisors of $P(z)$.
Its construction precision is sufficiently small in terms of
$\varepsilon_0$.

The following distribution assertion holds for smooth nonnegative
weights $W(r,q)$, localized in dyadic rectangles with
$r,q\ge N^c$, $rq\asymp x$, $N^{1-\kappa}\le x\ll N$, where $c>1/3$
is fixed and $\kappa>0$ is sufficiently small. Assume the normalized
smooth derivatives of $W$ are bounded by fixed powers of $\log N$.
Writing
\[
 X_W=\sum_{r,q\ {\rm prime}}W(r,q),
\]
one has, uniformly in the even integer $N$, for every fixed $A>0$,
\begin{equation}\label{eq:smooth-semiprime-error}
 \sum_{\substack{d\mid P(z)\\(d,N)=1}}\rho_D^+(d)
 \left\{\sum_{\substack{r,q\ {\rm prime}\\rq\equiv N\pmod d}}
       W(r,q)-\frac{X_W}{\varphi(d)}\right\}
 \ll_{A,W,\varepsilon_0,c}N(\log N)^{-A}.
\end{equation}
Here the dependence on $W$ means the fixed derivative exponents and
bounds, not its dyadic location. A polylogarithmic number of such
rectangles is permitted by increasing $A$.

Consequently, if
\[
 B_b(N)=\#\{(p,r,q):\ p,r,q\ {\rm prime},\
                     N=p+rq,\ q^b<r\le q\},
\]
then
\begin{equation}\label{eq:semiprime-new-upper}
 B_b(N)\le
 \left(\frac{80}{11}\log\frac1b+o(1)\right)\HN .
\end{equation}
The distribution assertion is for the displayed smooth weights.
The sharp-region upper bound follows by positive smooth outer
approximation; no unproved sharp-boundary distribution assertion
is inserted.
\end{lemma}

\begin{proof}
We give the coefficient, length, and residue checks separately.
All small positive parameters are fixed before $N$ tends to infinity.
One may take
\[
 \eta\ll\epsilon\ll\zeta\ll\varepsilon_0,\qquad
 \kappa\ll\epsilon,
 \qquad 0<\zeta<10^{-5},
\]
where $\eta$ is the sieve/box precision, $\epsilon$ is the
distribution precision, and $\zeta$ is the factorization margin.
Their choices can also depend on the fixed lower bound $c$.

\paragraph{The actual upper-sieve atoms.}
The construction of $L_{21}$, Section~5, equations~(5.6), (5.10),
and (5.14)--(5.17), expands the modified upper weight into finitely
many rough-prime box atoms followed by a small-prime $\beta$-sieve
factor. The number and sum of absolute scalar coefficients are
bounded in terms of $\eta$. The rough boxes have lower endpoints
$D_1\ge\cdots\ge D_j$. The definition of $\mathbf D_j^+$ is
\[
 D_1\cdots D_{m-1}D_m^3<D_0
       \quad(m\le j,\ m\ {\rm odd})
\]
when $j$ is even, and with right side $D_0^{1/(1+\tau)}$
when $j$ is odd, where $\tau=\eta^9$ and
$D_0=D^{1+O(\eta)}$ denotes the rough construction level.
It follows directly from this modified-weight construction that
\[
 D_1^3<D_0,\qquad
 t_1:=\frac{\log p_1}{\log N}\le\frac{\nu}{3}+O(\eta).
\]
This fact is not an identification of the modified weight with
the ordinary Rosser weight. Its rough vector also belongs to
$\mathcal D^{\rm well}(N^{\nu+O(\eta)})$.
The appended small-prime factor has support at most $N^{O(\eta)}$
and primes at most a cutoff $w=N^{c_\eta}$, $c_\eta>0$.

For an ordered box atom $T_J=\chi_{J_1}*\cdots*\chi_{J_j}$,
retain its original tuple-convolution scalar coefficient, including
the reciprocal symmetry factor $1/M_J$ in
\eqref{eq:box-symmetry-normalization}, exactly once. Given a partition of the prime slots
$\mathcal I_1,\mathcal I_2,\mathcal I_3$, put
\[
 K_i=|\mathcal I_i|!,\qquad
 \gamma_i=K_i^{-1}\mathop{*}_{k\in\mathcal I_i}\chi_{J_k}.
\]
Then $|\gamma_i|\le1$ and, for unrestricted convolutions,
$T_J=K_1K_2K_3\,\gamma_1*\gamma_2*\gamma_3$.
The fixed factorial is absorbed into the scalar coefficient.
Attach the small-prime $\beta$ factor to $\gamma_3$.
The unique decomposition into primes at most $w$ and greater
than $w$ preserves its bound by one.
The squarefree restriction can be restored with total discrepancy
\[
 \ll (N+N^{\nu+O(\eta)})w^{-1}(\log N)^{C_\eta}
      \ll_A N(\log N)^{-A}.
\]
Indeed, on the added moduli with a rough prime square $p^2\mid d$,
$p>w$, use the elementary progression bound $N/d+1$, the mean
bound $N/\varphi(d)$, and divisor moments. This is the same
coefficient argument as in the prime-prefix interface.

\paragraph{A uniform factorization margin for every upper atom.}
Set $\sigma=7/32$ and $\mu=25/128$. The cube condition implies
$t_1<\mu$ when $\eta$ is sufficiently small.
Use the fixed division point $T=21/200$.
If $t_1\ge T$, the single-variable branch in $L_{23}$,
Proposition~6.1, gives
\[
 v_1(t_1)=\frac{1+t_1}{2}\ge\frac{221}{400}>\frac{11}{20}.
\]
If $t_1<T$ and there are at least three rough primes, its
three-variable branch $w(t_1,t_2,t_3)$ has inner minimum at least
\[
 \min\left\{
   \frac{(5-4\sigma)-(3-4\sigma)T}{8-6\sigma},
   \frac{5-\sigma-3T}{8},1-2T\right\}
 =
 \min\left\{\frac{6243}{10700},\frac{3573}{6400},
                  \frac{79}{100}\right\}
 =\frac{3573}{6400}.
\]
The threshold defining $\psi$ in that branch is at most
$(1+T)/2=221/400$, strictly smaller than this minimum.
Thus the branch is present and exceeds $11/20$.
The condition $t_1\le\mu$ for using that branch has already
been checked.

If there are zero, one, or two rough primes, the entire atom has
support at most $N^{2\nu/3+O(\eta)}<N^{1/2-c_0}$ for a fixed
$c_0>0$. The ordinary convolution Bombieri--Vinogradov argument
of Lemma~\ref{lem:bv}, including smooth separated weights,
handles these atoms. Their coefficients are bounded by a constant
depending only on $\eta$.

For all other atoms, Proposition~6.1 therefore supplies the
factorizations with fixed positive margins. The same partition
of prime slots is used throughout each sufficiently fine box.
Changing representatives and attaching the small-prime factor
costs only $N^{O(\eta)}$ in support.
This constructs the necessary three coefficient factors;
the complete upper weight is not assumed triply well-factorable.

\paragraph{The Type II inequalities.}
Work temporarily on a local scale $x\asymp rq$.
For a Type II length $U$ in
\[
 x^{2\zeta}\le U\le x^{1/3+\zeta/2},
\]
use $A_0=Ux^{-\zeta}$ in Proposition~6.1.
Then $x^\zeta\le A_0\le x^{1/3-\zeta/2}$.
Write its factorization as $d=a_0b_0c_0$.
After the box construction and dyadic subdivision, the modulus
factor lengths $Q_1,Q_2,Q_3$ satisfy the four hypotheses of
$L_{23}$, Proposition~5.2:
\[
 Q_1<Ux^{-\epsilon},\quad
 U^2Q_2Q_3^2<x^{1-8\epsilon},\quad
 N^\sigma UQ_1Q_2^5Q_3^2<x^{2-15\epsilon},
\]
\[
 (UQ_3/Q_2)^{2\sigma}UQ_1Q_2^5Q_3^2<x^{2-15\epsilon}.
\]
For clarity, replacing $A_0$ by $U$ in the three monomials
of Proposition~6.1 costs $x^{2\zeta}$, $x^\zeta$,
and $x^{(1+2\sigma)\zeta}$, respectively. Their available
savings are $3\zeta,3\zeta,2\zeta$.
The smallest remaining saving is $(1-2\sigma)\zeta=9\zeta/16$.
The conditions remain strict after all $x^{O(\eta)}$ losses.
Monotonicity on dyadic subboxes follows by rewriting the last
monomial with powers $Q_2^{5-2\sigma}Q_3^{2+2\sigma}$.

On scales $x\ge N^{1-\kappa}/O(1)$, take the fixed parameter
$\alpha=1+3\kappa$. For all sufficiently large $N$ one has
$N<x^\alpha$.
The coordinates and the envelope $v_\alpha$ differ from their
$N$-scale values by $O(\kappa)$, absorbed by the margins above.
Moreover $N<x^{1+\epsilon}$, precisely the uniform residue
range in Proposition~5.2. Alternatively its $\alpha=1$
factorization can be used and the extra factor
$N^\sigma/x^\sigma=x^{O(\kappa)}$ paid from the same margin.
This is an explicit moving-residue argument.

\paragraph{Double Heath--Brown decomposition.}
Start with $\Lambda(r)\Lambda(q)$ and apply the $k=3$
Heath--Brown identity to each factor, at its own scale.
There are at most twelve variables, six from each factor,
of which at most six carry $\mu$ coefficients.
The $\mu$ variables in the $r$ group have length at most
$O(R^{1/3})$, and those in the $q$ group at most $O(Q^{1/3})$.
In the balanced region $R=x^{t+o(1)}$, $Q=x^{1-t+o(1)}$;
these bounds are $x^{t/3+o(1)}$ and $x^{(1-t)/3+o(1)}$,
respectively, both strictly below $x^{1/3}$.
They need not all be at most $x^{1/6}$.

Retain $P^-(rq)>z_0=x^{1/(\log\log x)^3}$ before applying
the identities; it is automatic on the original primes.
Every variable in the identities then has this same roughness
condition. Put the twelve variables into dyadic intervals.
If one individual variable lies in
$[x^{2\zeta},x^{1/3+\zeta/2}]$, choose it as $U$.
Its coefficient is $1$ or $\mu$, with a possible logarithm.
Lemma~\ref{lem:rough-sw}, with $\xi=2\zeta$, verifies its
rough Siegel--Walfisz condition; no Siegel--Walfisz property
of an arbitrary grouped convolution is asserted.
Group the other eleven variables, whose coefficient is
divisor-bounded. The truncation and interval-separation argument
in the proof of Proposition~3.1 reduces that coefficient to the
$1$-bounded setting of Proposition~5.2.

If there is no such variable, every $\mu$ variable is shorter
than $x^{2\zeta}$. There are at most two variables longer than
$x^{1/3+\zeta/2}$, since the product of all twelve is comparable
with $x$. Each original prime group must contain one such
variable: otherwise its product would be at most
$x^{12\zeta+o(1)}$, contradicting its fixed positive exponent.
Thus there are exactly two long variables, both with coefficient
$1$ or a logarithm, and the product $M$ of the other ten is at
most $x^{20\zeta+o(1)}$.
Proposition~3.2 applies, since
\[
 Q_{\rm modulus}\le x^{3/5}
   <(x/M)^{2/3-3\epsilon}
\]
for the fixed small choices of $\zeta,\epsilon$.
The two long lengths exceed $x^{3\epsilon}$.
The zero-long and one-long alternatives are excluded by the
two separate original product lengths, rather than left as
unhandled balanced Type II blocks.

\paragraph{Product cutoffs, smooth twists, and prime powers.}
The separate original $r$ and $q$ interval conditions cannot be
discarded when grouping the twelve variables.
For a smooth weight on a dyadic rectangle, use its two-variable
Mellin transform. Truncating at
$|\tau_1|,|\tau_2|\le(\log x)^L$ gives separated powers on every
variable in its original prime group.
For the chosen Type II variable this gives $n^{i\tau}\mu(n)$
or $n^{i\tau}$, times a smooth cutoff.
Partial summation in the Siegel--Walfisz bound costs only a
fixed power of $\log x$, absorbed by increasing its arbitrary
logarithmic saving. The same applies to logarithmic factors.

For the two-long-variable case, the required partial-interval
uniformity is explicit in the proof of Proposition~5.2 in
Maynard,
\href{https://arxiv.org/pdf/2006.07088v1}
{arXiv:2006.07088v1}, the input cited for $L_{23}$,
Proposition~3.2. That proof permits arbitrary subintervals
$\mathcal I_1\subset[N_1,2N_1]$ and
$\mathcal I_2\subset[N_2,2N_2]$.
Its Lemma~7.1 also gives the underlying smooth estimate with
polylogarithmic derivative bounds and uniformity in the residue.
Two partial summations therefore permit the two Mellin twists,
at cost $O((1+|\tau_1|)(1+|\tau_2|))$.
This does not assume an arbitrary-coefficient extension of
Proposition~3.2.

The tails of the Mellin integrals are harmless with arbitrary
logarithmic saving. For $m\ne N$, sum the progression contribution
over moduli by $d\mid(N-m)$. Divisor bounds and Cauchy--Schwarz
give
\[
 \sum_{\substack{m\ll N\\m\ne N}}
    \tau_C(m)\tau_C(|N-m|)\ll_C N(\log N)^{C'},
\]
using the usual second moments of fixed divisor functions.
The mean contribution has the same type of bound.
The case $m=N$, if introduced at an intermediate step, is
$O(N^{\nu+o(1)})$. Rapid Mellin decay now allows $L$ and
the number of integrations by parts to give any requested
logarithmic saving.

Prime powers introduced by $\Lambda(r)\Lambda(q)$ have total
mass $N^{1-c/2+o(1)}$: for one proper prime power $v\ge N^c$,
sum the elementary bound $O(N/v)$ for the other factor, and use
$\sum_{v\ge N^c,\ v\ {\rm proper\ prime\ power}}v^{-1}
 \le N^{-c/2+o(1)}$.
Their summed progression error is bounded by the same
$d\mid(N-m)$ divisor argument and hence remains
$O(N^{1-c/2+o(1)})$; their mean error uses
$\sum_d|\rho_D^+(d)|/\varphi(d)\ll_\eta(\log N)^{C_\eta}$.
In particular, their mass is not multiplied by $D$.
Remove the logarithmic prime weights by smooth partial summation.
Finally, since the original primes are larger than $z$,
$(rq,d)=1$ for every $d\mid P(z)$. Thus the coprime main mass
is exactly $X_W$, as in \eqref{eq:smooth-semiprime-error}.
This proves that assertion.

\paragraph{Smooth outer approximation and the sieve main term.}
Discard $rq\le N^{1-\kappa}$ in the original counting problem;
its contribution is $O(N^{1-\kappa+o(1)})$.
For the remaining curved region, use nonnegative smooth outer
approximations on dyadic rectangles. For each fixed approximation
precision $\delta_1>0$ these may be chosen with polylogarithmic
normalized derivative bounds, and with total prime-product mass
\[
 X_{\rm outer}
 \le \left(\log(1/b)+O_b(\delta_1)+o(1)\right)\frac{N}{\log N}.
\]
This follows from the prime number theorem and the convergent
integral
\[
 \int_{b/(1+b)}^{1/2}\frac{dt}{t(1-t)}=\log(1/b);
\]
the boundary strips have mass tending to zero.
More explicitly, choose one global smooth majorant in the variables
$\log r/\log N$, $\log q/\log N$, and $rq/N$, expanding the two
exponent boundaries by $O(\delta_1)$ and the upper product boundary
by a factor $1+O(\delta_1)$. Take $\delta_1,\kappa$ sufficiently
small in terms of $b-1/2$ so its two prime variables still exceed
$N^c$ for some fixed $c>1/3$.
Multiply this single majorant by nonnegative smooth dyadic
partitions of unity in $r$ and $q$. Their products have bounded
overlap and sum back to the same global majorant exactly.
The $O(\delta_1)$ excess is therefore a bound for the total mass,
not an error paid separately on each of the $O((\log N)^2)$ boxes.
All normalized derivatives have uniform bounds over these boxes.
The analytic errors are summed by first requesting two additional
logarithmic powers in \eqref{eq:smooth-semiprime-error}.
For each fixed $\delta_1$ the prime number theorem evaluates the
global mass; dominated convergence in the displayed compact
exponent integral then permits $\delta_1\downarrow0$.
Positivity is used only to compare the original prime count with
the smooth weighted sifted count. It is not used to sandwich a
signed sieve-weight remainder for the sharp curved region.
If the outer approximation extends slightly beyond $rq=N$,
apply the upper sieve to $|N-rq|$. Divisibility is unchanged
and all required positive-prime representations remain included.

The density is $g_N(\ell)=1/(\ell-1)$ for $\ell\nmid N$
and zero for $\ell\mid N$.
As in the prime-prefix interface, presieve to $w=N^{c_\eta}$,
use the uniform one-sided dimension bound below $w$, and the
uniform two-sided Mertens estimate above $w$. This preserves
the complete factor
\[
 V_N(z)=\prod_{\substack{\ell<z\\\ell\nmid N}}
             \left(1-\frac1{\ell-1}\right)
       =(2e^{-\gamma}+o(1))\frac{\CN}{\log z}.
\]
The ordinary modified upper sieve has multiplier $F(s)+o(1)$,
where $s=\log D/\log z=5\nu\in(1,3)$.
Since $F(s)=2e^\gamma/s$ on this interval, it yields
\[
 S\le\left(\frac4\nu+o(1)\right)
                   \frac{\CN X_{\rm outer}}{\log N}.
\]
Representations with $p\le z$ cost only $O(N^{1/5+o(1)})$.
First let $N\to\infty$, then let the smoothing and construction
precisions tend to zero within the fixed margins, and finally
let $\varepsilon_0\downarrow0$.
This gives \eqref{eq:semiprime-new-upper}.
\end{proof}

\paragraph{Scope of this derivation.}
The residue-uniform analytic inputs are the cited Type II and
two-long-variable estimates. The double Heath--Brown expansion,
smooth product-cutoff argument, and uniform factorization of
the actual upper-sieve atoms are the additional deductions here.
The result is not quoted as a previously published theorem for
semiprimes, nor does this section establish its priority.

% Integration fragment: no documentclass.
% Uses the main manuscript's notation and Lemma lem:semiprime55.
% Do not interpret bounded multiplicity as arbitrary arithmetic weights.

\subsection{The same upper sieve for finitely many prime factors}
\label{sec:primeproducts55}

The upper-weight factorization established in
Lemma~\ref{lem:semiprime55} is a property of the modulus coefficients.
It can also be used for the switched product sequences below.
We record the additional argument, particularly the coprime main
mass and the multiplicity of the factorizations.

Put $a=100/1149$ and $K_0=\lfloor1/a\rfloor=11$.
An admissible product multiset here means a finite positive sum
of labelled prime-tuple families
\[
 m=p_1\cdots p_k,\qquad 1\le k\le K_0,
 \qquad p_j\ge N^a,\qquad m\le N,
\]
whose domains are specified by finitely many ordering and monomial
inequalities in the $p_j$ and $N$. Labels distinguish different
specified-prefix choices for the same product.
Their coefficients are fixed bounded nonnegative numbers.
In particular, the resulting multiplicity $w(m)$ is bounded
independently of $N$. This does not mean that an arbitrary
bounded arithmetic function $w(m)$ is permitted: the tuple-domain
description is needed for separation of variables.

\begin{lemma}[Derived upper sieve for finite prime-product multisets]
\label{lem:primeproducts55}
Let $\mathcal M_N$ be an admissible product multiset as above,
and write
\[
 X_N=\sum_m w(m).
\]
Assume its mass has a limit
$X_N=(J+o(1))N/\log N$, evaluated with its stated labels and
multiplicities. Then
\begin{equation}\label{eq:primeproduct-prime-upper}
 \sum_{\substack{m<N\\N-m\ {\rm prime}}}w(m)
 \le\left(\frac{80}{11}J+o(1)\right)\HN .
\end{equation}
More generally the same conclusion applies to a finite positive
sum of such domains and to positive smooth outer approximations,
with their own total mass $X_N$.

This lemma applies to all five switched multisets defined by
$Y_7,Y_8,Y_9,Y_{10},Y_{11}^{\rm full}$ in
Section~\ref{sec:repaired-chen-weight}. If $U_i=8J_i$ denotes
the corresponding upper constant computed there with coefficient
eight, then
\begin{equation}\label{eq:switch-55}
 \Upsilon_i\le\left(\frac{10}{11}U_i+o(1)\right)\HN
 \quad(7\le i\le10),\qquad
 \Upsilon_{11}^{\rm full}
       \le\left(\frac{10}{11}U_{11}^{\rm full}+o(1)\right)\HN.
\end{equation}
Here the last mass uses the complete range
$N^h\le p_4<N^d$.
\end{lemma}

\begin{proof}
Fix $0<\varepsilon_0<1/100$ and put
$D=N^{11/20-\varepsilon_0}$ and $z=N^{1/5}$.
Use precisely the Iwaniec modified upper weight $\rho_D^+$
and its normalized box-atom expansion constructed in
Lemma~\ref{lem:semiprime55}. In particular
\[
 |\rho_D^+(d)|\le C_\eta,\qquad d\mid P(z),
\]
where $\eta$ is the fixed construction precision.
Its factorization works for all medium lengths
$U\in[x^{2\zeta},x^{1/3+\zeta/2}]$.
The positive margins from the two cases
$t_1\ge21/200$ and $t_1<21/200$ depend only on modulus
factorization, not on the sequence being counted.
Use
\[
 \eta\ll\epsilon\ll\zeta\ll\varepsilon_0,\qquad
 12\zeta<a/2,\qquad \kappa\ll\epsilon .
\]
All parameters and $K_0$ are fixed before $N\to\infty$.

\paragraph{Small products, repeated primes, and smooth domains.}
Products below $N^{1-\kappa}$ cost
$O(N^{1-\kappa})$ in the prime count because $w$ is bounded.
On $m\le CN$, the products having a repeated prime at least
$N^a/2$ cost at most
\[
 C' N\sum_{\ell\ge N^a/2,\ \ell\ {\rm prime}}\ell^{-2}
       \ll N^{1-a}
\]
for bounded multiplicity. Thus actual repeated primes can be
removed before using strictly ordered prime tuples, with an
additive error of this size in the desired upper bound.
This removal does not identify distinct labelled prefix choices.

\paragraph{Prime-product mass lemma and radial boundaries.}
The mass calculation requires both exponent coordinates and the
radial coordinate $y=m/N$. For a fixed number $k$ of prime slots put
\[
 t_j=\frac{\log p_j}{\log N}\quad(j<k),\qquad
 t_k^*=1-\sum_{j<k}t_j,\qquad y=\frac{p_1\cdots p_k}{N}.
\]
Fix $0<a_0<a$, and work on an ambient domain
$t_j,t_k^*\ge a_0$, $0<y\le C$, with $C$ fixed.
The labelled prime-product counting measure, multiplied by $\log N/N$,
converges, on continuous
test weights and subsequently on domains with null boundary in these
coordinates, to
\[
 d\mu_k=\frac{dt_1\cdots dt_{k-1}\,dy}
 {t_1\cdots t_{k-1}t_k^*}.
\]
For $k=1$ this means $d\mu_1=dy$. Fixed bounded constants attached
to the finite labels are retained. This is a bounded density on the
displayed simplex times a radial interval; it is not a density in
the full space of the $k$ prime exponents.

Here is the requisite proof and uniform error estimate. First take
$\eta\le y\le C$ with $\eta>0$ fixed. With the first $k-1$ primes
held fixed and $P=p_1\cdots p_{k-1}$, write $p_k=Ny/P$. The prime
number theorem and partial summation give the last-prime measure
\[
 \frac{dp_k}{\log p_k}
 =\frac{N}{P\log N}\frac{dy}{t_k^*+\log y/\log N},
\]
with error $O_A((N/P)(\log N)^{-A})$ for fixed smooth seminorms.
All nonempty relevant last-prime intervals have their endpoints at
least a fixed positive power of $N$. Since
\[
 \sum_{N^{a_0}/2\le p\le CN}\frac1p=O_{a_0,C}(1),
\]
the sum of these errors over the other primes is
$O_{A,k,a_0}(N(\log N)^{-A})$.
The reciprocal-prime measures in the remaining coordinates converge
to $dt_j/t_j$, again by prime-number-theorem partial summation.
On $y\ge\eta$ the displayed denominator tends uniformly to $t_k^*$,
proving the assertion first for smooth and then continuous weights.
The elementary upper bound for primes also gives
\[
 \#\{(p_1,\ldots,p_k):p_j\ge N^{a_0}/2,\ 
                  p_1\cdots p_k\le\eta N\}
 \ll_{k,a_0}\frac{\eta N}{\log N}.
\]
Indeed fix $P$, bound its last-prime count by
$O_{a_0}(\eta N/(P\log N))$, and sum reciprocal products.
This removes the radial restriction by first taking $N\to\infty$
and then $\eta\downarrow0$.

For completeness, each fixed monomial constraint
$\prod_{j=1}^kp_j^{\beta_j}\le cN^\gamma$ becomes exactly
\[
 \sum_{j<k}(\beta_j-\beta_k)t_j+\beta_k-\gamma
 \le\frac{\log c-\beta_k\log y}{\log N}.
\]
If the $\beta_j$ are not all equal, its limiting boundary is a
nonconstant affine hyperplane in the simplex coordinates and has
$\mu_k$-measure zero. If all $\beta_j=\lambda$ but
$\gamma\ne\lambda$, its truth value is eventually constant on
$\eta\le y\le C$. If all $\beta_j=\gamma=\lambda\ne0$, it is
the exact radial condition $y^\lambda\le c$ and must be smoothed
in $y$. When $\lambda=0$ its constant truth value $1\le c$ is
kept exactly. This classification covers ordering constraints,
equalities, and finite intersections and unions. In particular it
does not discard a pure total-product boundary by using the full
exponent-space measure. For example the empty one-prime domain
$p=N$ at even $N$ has exponent outer strips
$N^{1-\delta}\le p\le N$ of normalized limiting mass one, whereas
radial strips around $p/N=1$ have mass tending to zero with their width.

\paragraph{A positive outer approximation before sieving.}
For each fixed error budget $\tau>0$, choose $\eta$ sufficiently
small. On $\eta\le y\le C$ smooth the nonconstant affine walls
outward by small fixed exponent widths, and the critical pure-product
walls by small fixed radial widths. Retain a nonnegative smooth
majorant of the whole strip $0<y\le\eta$, supported in $y<2\eta$;
its mass is $O(\eta)N/\log N$. The preceding classification and
mass lemma then give one nonnegative bounded global outer weight
$W_{N,\tau}$ with
\[
 W_{N,\tau}\ge\mathbf1_{\mathcal D_N},\qquad
 \limsup_{N\to\infty}\frac{\log N}{N}
       \sum_{p_1,\ldots,p_k}
        (W_{N,\tau}-\mathbf1_{\mathcal D_N})\le\tau.
\]
Keep a separate lower-cutoff factor equal to one for $p_j\ge N^a$
and zero for $p_j\le N^a/2$. Here this factor-two transition has
vanishing excess mass: for $k\ge2$ its limiting face is a
nonconstant simplex wall, and for $k=1$ the cutoff $a<1$ is
noncritical. For a general critical pure-product cutoff the radial
rule is required instead of a fixed factor-two transition.
Thus every prime in the actual outer family remains at least $N^a/2$.

The operators $p_j\partial_{p_j}$ send $t_i$ to
$\delta_{ij}/\log N$ and send $y$ to $y$. Hence all normalized
derivatives on dyadic charts are bounded for each fixed smoothing
precision, uniformly in $N$ and the chart location. A smooth dyadic
partition of unity gives $O((\log N)^k)$ rectangles and sums back
to this one global majorant. The excess mass is therefore counted
once; only distribution errors acquire the harmless logarithmic
number of rectangles, absorbed by a larger requested saving.

The order of use is essential. First dominate the desired positive
prime count by its count with $W_{N,\tau}$. Then apply the smooth
exact-coprime distribution statement and the upper sieve below.
Small excess total mass alone is not asserted to control a difference
of signed progression remainders. After these estimates, take
$N\to\infty$ for the fixed family, widths, and factorization margins,
and only then take $\tau\downarrow0$. No averaging over $N$ or
uniformity in a cutoff tending to zero is used.

\paragraph{Distribution with the exact coprime mass.}
For a smooth $k$-tuple weight $W$ on one rectangle, put
\[
 X_W^{(d)}
   =\sum_{\substack{p_1,\ldots,p_k\ {\rm prime}\\
                    (p_1\cdots p_k,d)=1}}
                         W(p_1,\ldots,p_k).
\]
We first prove, on scales $N^{1-\kappa}/O(1)\le x\ll N$,
\begin{equation}\label{eq:primeproduct-coprime-error}
 \sum_{\substack{d\mid P(z)\\(d,N)=1}}\rho_D^+(d)
 \left\{
  \sum_{\substack{p_1,\ldots,p_k\ {\rm prime}\\
            p_1\cdots p_k\equiv N\pmod d}}W(p_1,\ldots,p_k)
       -\frac{X_W^{(d)}}{\varphi(d)}
 \right\}
 \ll_A N(\log N)^{-A}.
\end{equation}

Apply the $k_{\rm HB}=3$ Heath--Brown identity to each of the
$k$ prime slots, at that slot's own scale, initially with
$\Lambda$ weights. There are at most $6k\le66$ variables,
of which at most $3k$ carry a $\mu$ coefficient.
Retain the condition
$P^-(p_1\cdots p_k)>x^{1/(\log\log x)^3}$ before expansion.
Consequently every nonunit variable is rough at that cutoff.
Each $\mu$ variable has length at most the cube root of
its own original factor scale, and in particular at most
$x^{1/3+o(1)}$.
Before Mellin separation, discard every dyadic block that cannot
intersect its original prime-slot interval. If the original slot has
scale $R_i$, the product of that group's at most six dyadic lower
endpoints is at least $R_i/2^6\ge N^a/O(1)$. This group-scale bound
persists after separating the smooth weights, including at points
outside the original curved domain.

If an individual variable has length
\[
 x^{2\zeta}\le U\le x^{1/3+\zeta/2},
\]
take that individual variable as the Type II factor and group
all the others. Its coefficient is $1$ or $\mu$, with a possible
logarithm. Lemma~\ref{lem:rough-sw}, with $\xi=2\zeta$,
verifies its complete rough Siegel--Walfisz input.
The other coefficient is divisor-bounded, with fixed order
depending on $K_0$. The normalized three-factor modulus
coefficients satisfy $L_{23}$, Proposition~5.2, by the
factorization already proved in Lemma~\ref{lem:semiprime55}.
The truncation of divisor-bounded coefficients and their
logarithmic losses are handled exactly as in that lemma.

Suppose no individual variable is in this middle interval.
All $\mu$ variables are then shorter than $x^{2\zeta}$.
Every original prime group has at most six variables and
original product at least $N^a/2$. It must therefore contain
a variable longer than $x^{1/3+\zeta/2}$: otherwise its
product would be at most $x^{12\zeta+o(1)}<N^a/2$.
On the other hand there are at most two such long variables
globally, because the entire product is comparable with $x$.
It follows that this case is impossible for $k\ge3$.

For $k=2$, there is exactly one long variable in each original
group. Both have coefficient $1$ or a logarithm, and the
remaining ten variables have product at most $x^{20\zeta+o(1)}$.
Apply $L_{23}$, Proposition~3.2:
\[
 N_1,N_2>x^{1/3+\zeta/2}>x^{3\epsilon},
 \qquad
 Q_{\rm modulus}\le x^{3/5}
       <(x/x^{20\zeta+o(1)})^{2/3-3\epsilon}.
\]
For $k=1$, either there is one long variable and the other
five have product at most $x^{10\zeta+o(1)}$, so Lemma~3.4
applies, or there are two long variables and the remaining
four have product at most $x^{8\zeta+o(1)}$, so
Proposition~3.2 applies. These include all alternatives.

The separate original product cutoffs are retained using
the $k$-variable Mellin transform of $W$, truncated at
polylogarithmic frequencies. A chosen individual Type II
variable acquires only a twist $n^{i\tau}$; partial summation
preserves its Siegel--Walfisz estimates with a polylogarithmic
loss. The two long variables, when present, are handled by
two partial summations, using the arbitrary-subinterval form
explicit in the proof of Maynard's Proposition~5.2,
\href{https://arxiv.org/pdf/2006.07088v1}
{arXiv:2006.07088v1}, which is the input for $L_{23}$,
Proposition~3.2.
The one-long-variable case allows the same partial summation
in Lemma~3.4. The reconciliation of its printed unrestricted mean
with our exact coprime mean is given explicitly in
Section~\ref{sec:one-long-coprime-mean}.
No Siegel--Walfisz property of an arbitrary convolution or
arbitrary-coefficient extension of the two-long-variable
estimate is assumed.

For every local scale use the fixed residue parameter
$\alpha=1+3\kappa$. Then $N<x^\alpha<x^{1+\epsilon}$
for sufficiently small $\kappa$ and sufficiently large $N$.
Changes of exponent coordinates are $O(\kappa)$ and are
absorbed by the fixed factorization margins.
Rough modulus atoms with at most two large-prime slots have
support below $N^{2\nu/3+O(\eta)}$, $\nu=11/20-\varepsilon_0$.
They are covered by ordinary Bombieri--Vinogradov for $k=1$.
For $k\ge2$, select one original prime-interval factor and
group the other $k-1$; both have lengths between fixed
positive powers of the product scale. The generalized
Bombieri--Vinogradov input \eqref{eq:geh}, with its verified
Siegel--Walfisz factor, applies. Smooth Mellin twists cost
only logarithmic powers.

As in Lemma~\ref{lem:semiprime55}, Mellin tails are controlled
after summing over divisors of $N-m$, using
\[
 \sum_{\substack{m\ll N\\m\ne N}}
      \tau_C(m)\tau_C(|N-m|)
        \ll N(\log N)^{C'},
\]
which follows by Cauchy--Schwarz and divisor second moments.
Any isolated $m=N$ term is $O(DN^{o(1)})$.
Proper prime powers introduced by the $\Lambda$ identities
have total mass at most $N^{1-a/2+o(1)}$; after the divisor
argument their cumulative progression error still has
this power saving. Their mean error uses
$\sum_d|\rho_D^+(d)|/\varphi(d)\ll_\eta(\log N)^{C_\eta}$.
Remove logarithmic prime weights by partial summation.
This proves \eqref{eq:primeproduct-coprime-error}.
Increasing $A$ absorbs the number of dyadic rectangles,
Mellin losses, and the fixed number of prime-factor patterns.

\paragraph{Replacing the coprime mass: a power-saving bound.}
Here $X_W^{(d)}$ cannot be identified with $X_W$:
some prime factors may be below $z$.
The exact correction is nevertheless negligible even after
the modulus weights have been summed in absolute value.
For a prime $\ell$,
\[
 \sum_{\substack{d\mid P(z)\\\ell\mid d}}
     \frac{|\rho_D^+(d)|}{\varphi(d)}
 \le \frac{C_\eta}{\ell-1}
          \sum_{e\le D/\ell}\frac1{\varphi(e)}
 \ll_\eta \frac{\log N}{\ell-1}.
\]
Indeed $d$ is squarefree, so $d=\ell e$ with $(e,\ell)=1$
and $\varphi(d)=(\ell-1)\varphi(e)$; dropping the residual
support restrictions only increases this positive bound.
After summing the smooth tuples, let $w_{\rm out}(m)$ denote
their nonnegative bounded total multiplicity and $X_{\rm out}$
their total mass. The union bound therefore gives
\begin{align*}
 &\sum_{\substack{d\mid P(z)\\(d,N)=1}}
    \frac{|\rho_D^+(d)|}{\varphi(d)}
                      (X_{\rm out}-X_{\rm out}^{(d)})\\
 &\qquad\ll_\eta
    (\log N)\sum_m w_{\rm out}(m)
          \sum_{\ell\mid m}\frac1{\ell-1}\\
 &\qquad\ll_{\eta,K_0}
    X_{\rm out}N^{-a}\log N
       \ll N^{1-a}(\log N)^{C}.
\end{align*}
Every such $\ell$ is at least $N^a/2$ and each product has at
most $K_0$ distinct factors. This proves the claimed uniform
power saving. It is valid for positive bounded labelled
multiplicities and does not lose a factor $D$.
If the original switched set excludes products sharing a prime
factor with $N$, the raw mass of those products is separately bounded by
\[
 O\left(N\sum_{\substack{\ell\mid N\\\ell\ge N^a/2}}\frac1\ell\right)
       =N^{1-a+o(1)}.
\]
This deletion is therefore also harmless; it is distinct from the
coprime-main-mass correction just estimated.

We may now replace $X_W^{(d)}$ by the total mass throughout
\eqref{eq:primeproduct-coprime-error}.
Apply the upper linear sieve with density
$g_N(\ell)=1/(\ell-1)$ for $\ell\nmid N$ and zero otherwise.
The uniform small-prime presieve and the full product $V_N(z)$
are those established in Lemma~\ref{lem:semiprime55}.
For $s=5\nu\in(1,3)$ one has $F(s)=2e^\gamma/s$, and hence
\[
 S\le\left(\frac4\nu+o(1)\right)
                    \frac{\CN X_{\rm out}}{\log N}.
\]
If the outer approximation reaches products just above $N$,
sieve $|N-m|$; every desired representation is still counted.
The cases $0<N-m\le z$ in the desired prime count contribute
$O(N^{1/5})$ for bounded
multiplicity. First let $N\to\infty$, then remove the smooth
outer approximation and the sieve precision within the fixed
margins, and finally let $\varepsilon_0\downarrow0$.
This proves \eqref{eq:primeproduct-prime-upper}.

\paragraph{The five switched multisets and the rough quotient.}
In each of $Y_7,Y_8,Y_9,Y_{10},Y_{11}^{\rm full}$, keep the
specified primes with exactly their original labels and ordering.
Factor the rough quotient uniquely as
\[
 u=q_1\cdots q_j,\qquad q_1<\cdots<q_j,
\]
after the repeated-prime exception above is removed.
Split into the finitely many possible values of $j$.
This is a canonical ordered factorization of the quotient;
there is no additional division by $j!$ and no deletion of
different specified-prefix choices for the same full product.
All quotient primes are at least the prescribed roughness
threshold, and in these five domains that threshold is at least
$N^a$. Thus the entire product has at most $K_0$ factors.
The resulting domains satisfy the finite monomial-inequality
description used above, including the moving quotient threshold
$(N/(p_1p_2))^{1/2}$ in $Y_9$.

The case $u=1$ has negligible mass. The five possible specified
prefix products are bounded respectively by
\[
 N^{2/3},\quad N^{2/3},\quad N^{(1+d)/2},
                  \quad N^{4h},\quad N^{3h+d},
\]
all smaller than $N$ by a fixed power.
Likewise, specified factors dividing $N$ or any repeated
prime affect the original prime count only by a power-saving
error. No restriction of the fourth-prime range is introduced:
for $Y_{11}^{\rm full}$ it remains
$N^a\le p_1<p_2<p_3<N^h\le p_4<N^d$.

Reassembling the canonical quotient factorizations gives exactly
the original rough-number mass, with its prefix multiplicities.
Its Buchstab/prime-number-theorem evaluation is therefore
unchanged. In the notation of the main calculation,
\[
 X_i=(J_i+o(1))\frac{N}{\log N},
 \qquad J_i=\frac{U_i}{8}.
\]
The treatment of the Buchstab argument-one boundary by a thin
strip is preserved; smoothing does not remove its mass silently.
Applying the just-proved upper sieve to each of these masses
gives \eqref{eq:switch-55}.
\end{proof}

\paragraph{Scope.}
This is a derived extension of the modulus-atom and smooth
product argument in Lemma~\ref{lem:semiprime55}.
It retains the full domains and tuple multiplicities of the
five switching terms. It changes only their analytic upper
coefficient from $8$ to $80/11$.
It makes no assertion of distribution for arbitrary bounded
arithmetic reweightings or of bibliographic priority.

% Integration fragment only: no documentclass.
% Uses the main manuscript's lemma/proof environments, \CN, \HN,
% sec:prefix-interface, sec:semiprime-upper, and L_1,L_2 notation.
% Source numbering: arXiv:2109.02851v2 and arXiv:2309.08522v1.

\subsection{A higher level for the two restricted lower sieves}
\label{sec:restricted-lower-level}

The upper-sieve cube condition is not needed for this additional
step. Here the restriction on the sifting primes supplies the
necessary bound for the largest prime in every lower-sieve atom.
Throughout this subsection, $L_{21}$ and $L_{23}$ have the meanings
specified in Section~\ref{sec:semiprime-upper}.

\begin{lemma}[Derived restricted lower-sieve level]
\label{lem:restricted-lower-level}
Put $h=50/309$. Fix $0<u\le h$ and $0<\varepsilon_0<1/100$, and set
\[
 \nu=\frac{11}{20}-\varepsilon_0,\qquad
 D=N^\nu,\qquad z=N^u.
\]
The ordinary Iwaniec modified lower linear sieve of level $D$
can be constructed, with sufficiently small fixed precision, so
that its signed remainder on
$\mathcal A_N=\{N-p:p<N,\ p\text{ prime}\}$ satisfies
\begin{equation}\label{eq:restricted-lower-remainder}
 \left|\sum_{\substack{d\mid P(z)\\(d,N)=1}}
  \rho_D^-(d)
   \left\{\pi(N;d,N)-\frac{\pi(N)}{\varphi(d)}\right\}\right|
 \ll_{A,u,\varepsilon_0,\eta}N(\log N)^{-A}
\end{equation}
for every fixed $A>0$, uniformly over even $N$. Here $\eta$ is the
fixed sieve precision, and $\rho_D^-$ denotes the signed remainder
weight with the convention of the lower-sieve inequality.
Consequently, with $P(N)=\{\ell\text{ prime}:\ell\nmid N\}$,
\begin{equation}\label{eq:restricted-lower-main}
 S(\mathcal A_N;P(N),N^u)
 \ge \left\{\frac{2e^{-\gamma}}{u}
        f\left(\frac{\nu}{u}\right)+o(1)\right\}\HN.
\end{equation}
The $o(1)$ conclusion is obtained by first taking $N\to\infty$
at each fixed precision and then letting that precision tend to zero.
This is a deduction from the sieve construction and uniform
dispersion estimates, rather than a quotation of a published
restricted lower-sieve theorem.
\end{lemma}

\begin{proof}
Choose fixed parameters
\[
 \eta\ll\epsilon\ll\zeta\ll\varepsilon_0,
 \qquad \kappa\ll\epsilon,
 \qquad 0<\zeta<10^{-5},
\]
also sufficiently small in terms of $u$. The ordinary construction
is the one with main term $f$, not the additional upper-sieve
modification with main term $F^*$.

\paragraph{All lower atoms and their signs.}
The expansion in $L_{21}$, (5.17), explicitly applies to both
modified signs; (5.18) states
$\mathbf D_j^\pm\subset\mathbf D_j^{\rm well}$.
Thus the restricted lower remainder is a finite sum of atoms
\[
 c_J^-\,\chi_{J_1}*\cdots*\chi_{J_j}*\beta_J^-,
\]
with the original signed scalar $c_J^-$, the original parity-dependent
small-prime sieve factor $\beta_J^-$, and all rough prime slots
restricted to primes below $z$. The small-prime factor has absolute
value at most one, uses primes at most $w=N^{c_\eta}$, and has support
at most $N^{O(\eta)}$. The number and total absolute scalar mass
of these atoms are bounded in terms of $\eta$.
Intersecting a prime box with $[w,z)$ only restricts that slot's
coefficient; it does not alter any scalar sign or substitute an
upper weight for a lower weight. The rough vectors lie in
$\mathcal D^{\rm well}(N^{\nu+O(\eta)})$ and satisfy
\[
 t_1\le u+O(\eta)\le h+O(\eta)<\mu=\frac{25}{128}.
\]
In particular, no assertion $p_1^3<D$ is made for the lower sieve.

Repeated prime boxes, their factorial scalars, and the small-prime
factor are handled exactly by the explicit convolution identities
in Section~\ref{sec:prefix-interface}. Those identities allow either
sign of $\beta_J^-$. After normalizing each group of prime slots by
its factorial, each modulus-factor coefficient is bounded by one.
The small factor is attached to the third group by the unique
small-prime/rough-prime decomposition. Added rough-square collisions
have total discrepancy
\[
 \ll (N+N^{\nu+O(\eta)})w^{-1}(\log N)^{C_\eta}
 =O_A(N(\log N)^{-A}).
\]
Hence applying the distribution estimate to each atom and summing
absolute errors is legitimate irrespective of its sign.

\paragraph{The complete factorization range.}
Set $T=21/200$ and $\sigma=7/32$. For an atom with at least three
rough primes, the cases below exhaust all its possible largest
prime exponents. If $t_1\ge T$, then $t_1<\mu$ gives
\[
 v_1(t_1)=\frac{1+t_1}{2}\ge\frac{221}{400}
 >\frac{11}{20}.
\]
If $t_1<T$, all of $t_1,t_2,t_3$ are less than $T$.
In $L_{23}$, Proposition~6.1, (6.4), the inner minimum defining
$w(t_1,t_2,t_3)$ is at least
\[
 \min\left\{
 \frac{(5-4\sigma)-(3-4\sigma)T}{8-6\sigma},
 \frac{5-\sigma-3T}{8},1-2T\right\}
 =\frac{3573}{6400}>\frac{11}{20}.
\]
Its $\psi$ threshold is at most
$(1+T)/2=221/400<3573/6400$. This branch is therefore present.
Its additional hypothesis $t_1\le\mu$ has been checked above.
This use of the three-variable branch is entirely inside the
domain of Proposition~6.1, and needs no conjectural extension.

For $j=0,1,2$, the whole atom, including its small factor, has
support at most
\[
 N^{2u+O(\eta)}\le N^{2h+O(\eta)}<N^{1/2-c_0}
\]
for a fixed $c_0>0$. Ordinary Bombieri--Vinogradov, with its
maximum over primitive residue classes, handles these bounded
atoms. Thus no undefined $t_3$ is used in the small-rank cases.

For the remaining atoms, choose the level parameter in
Proposition~6.1 as $\vartheta=\nu+\varepsilon_0/2$; make $\eta$
and $\zeta$ small enough that the rough support belongs to
$\mathcal D^{\rm well}(N^{\vartheta-2\zeta})$.
The strictly positive factorization margins just proved persist
under box inflation, the small factor, and passage to local scales.

\paragraph{Uniform prime distribution.}
Here there is only one prime variable to decompose, so the six-variable
Heath--Brown argument already supplied in
Section~\ref{sec:prefix-interface} applies directly to these atoms.
For clarity its parameter choices are
\[
 x^{2\zeta}\le U\le x^{1/3+\zeta/2},
 \qquad A_0=Ux^{-\zeta}.
\]
The three savings of Proposition~6.1 after replacing $A_0$ by $U$
are respectively $\zeta$, $2\zeta$, and
$(1-2\sigma)\zeta=9\zeta/16$. After all sufficiently small
box losses, they give all four inequalities of $L_{23}$,
Proposition~5.2, including $Q_1<Ux^{-\epsilon}$.
Its supremum over $0<|a|<x^{1+\epsilon}$ supplies the required
moving residue $a=N$, rather than a fixed-residue implied constant.
Indeed, discard primes at most $N^{1-\kappa}$ at cost
$N^{1-\kappa+o(1)}$ using $d\mid(N-p)$ and divisor bounds.
For the remaining dyadic prime scales use the fixed
$\alpha=1+3\kappa$: for large $N$, $N<x^\alpha<x^{1+\epsilon}$.
All exponent changes are $O(\kappa)$ and fit the same margins.

The selected middle variable has the rough $1$ or $\mu$ coefficient.
Lemma~\ref{lem:rough-sw} verifies its Siegel--Walfisz condition,
including auxiliary coprimality; partial summation handles a logarithm.
If there is no middle variable, the one-long-variable and
two-long-variable cases use Lemma~3.4 and Proposition~3.2 of
$L_{23}$ respectively, as checked in the prime-prefix proof.
The total modulus support is below $x^{3/5}$, and the remaining
short-variable products are at most $x^{10\zeta}$ and
$x^{8\zeta}$ respectively. The restrictions $(d,N)=1$ split
between the three modulus coefficients and preserve their bounds.
These arguments establish \eqref{eq:restricted-lower-remainder}.

\paragraph{The lower main term uniformly in $N$.}
The local density is $g_N(\ell)=1/(\ell-1)$ if $\ell\nmid N$,
and zero otherwise. Presieve to $w=N^{c_\eta}<z$.
The uniform one-sided dimension estimate below $w$ permits both
signs in the fundamental lemma and retains the full factor $V_N(w)$.
More explicitly, in the source construction take
$w=D_0^{\eta^2}$ and the small-prime sieve level $D_0^\eta$,
where $D_0$ is the nominal level before the harmless box inflation.
For both signs the small-factor main sum is
$V_N(w)\{1+O(e^{-c/\eta})\}$, uniformly in even $N$.
Above $w$, the missing primes dividing $N$ have total density at most
\[
 \sum_{\substack{\ell\mid N\\\ell\ge w}}\frac1{\ell-1}
 \ll\frac{\log N}{w\log w},
\]
so the two-sided Mertens condition used in the rough-box main-term
calculation is uniform. Reattaching the presieve therefore gives
the ordinary lower multiplier $f(\nu/u)+O_u(\eta)$ and
\[
 V_N(N^u)=(2e^{-\gamma}+o(1))\frac{\CN}{u\log N}.
\]
Indeed, let $\delta_N$ denote the displayed sum over missing primes.
At fixed $\eta$, the finite total absolute mass of the rough atoms
is $C_\eta$. Replacing their rough densities by
$g_0(\ell)=1/(\ell-1)$ changes the normalized main sum by at most
$C_\eta\delta_N=o_{N;\eta}(1)$, while
$V_N(z)/V_N(w)=\{V_0(z)/V_0(w)\}(1+O(\delta_N))$.
Thus the full main term is
$\pi(N)V_N(z)\{f(\nu/u)+O_u(\eta)+o_{N;\eta}(1)\}$.
No uniform two-sided density identity below $w$ is being assumed.
The signed convention in the displayed lower inequality of
$L_{21}$, Theorem~2.10, is immaterial after bounding the absolute
remainder in \eqref{eq:restricted-lower-remainder}; explicitly the
result is at least
$\pi(N)V_N(z)\{f(\nu/u)-O_u(\eta)}-O_A(N(\log N)^{-A})$.
Take $A>2$, first let $N\to\infty$, and then let $\eta$ tend to zero
within the fixed factorization margins. This proves
\eqref{eq:restricted-lower-main}.
\end{proof}

\paragraph{Replacement of the two lower constants.}
For $a=100/1149$ and $h=50/309$, the two permissible replacements are
\[
 L_1(\nu)=\frac{2e^{-\gamma}}a f\left(\frac\nu a\right),
 \qquad
 L_2(\nu)=\frac{2e^{-\gamma}}h f\left(\frac\nu h\right).
\]
After taking the limits in the lemma, one may let
$\varepsilon_0\downarrow0$ and use continuity of $f$.
No other term in the pointwise Chen weight is changed by this step.

There is a useful numerical distinction from the earlier level
$1/2$: at $\nu=11/20$ one has $\nu/a=6.3195>6$.
The earlier truncated delay-differential formula for $L_1$ must
therefore no longer be called exact. A sufficient rigorous lower
bound follows by putting
\[
 J_0(v)=\begin{cases}
 \displaystyle\int_2^{v-1}\frac{\log(t-1)}t\,dt,&v>3,\\
 0,&v\le3,
 \end{cases}
\]
and using the nonnegative successive terms in the sieve differential
equations to obtain, for $s>2$,
\[
 e^{-\gamma}f(s)\ge\frac2s
 \left\{\log(s-1)+\int_3^{s-1}\frac{J_0(v)}v\,dv\right\}.
\]
Here the last integral is zero for $s\le4$.
For completeness, this follows successively from
$e^{-\gamma}F(v)\ge2/v$, then
$e^{-\gamma}f(v)\ge2\log(v-1)/v$, then
$e^{-\gamma}F(v)\ge2(1+J_0(v))/v$, and a further integration
of $(sf(s))'=F(s-1)$. Thus
\[
 L_1(\nu)\ge\frac4\nu
 \left\{\log(\nu/a-1)
       +\int_3^{\nu/a-1}\frac{J_0(v)}v\,dv\right\}.
\]
The second parameter obeys $2<\nu/h<4$, so its elementary formula
remains exact:
\[
 L_2(\nu)=\frac4\nu\log(\nu/h-1).
\]
This subsection uses neither a claimed external Chen constant,
nor the double-sieve correction functions $H,h$.

\subsection{The enhanced transfer and its numerical consequences}
\label{sec:enhanced-corollaries}

\begin{theorem}[Enhanced transfer]\label{thm:enhanced-transfer}
With the normalization of Theorem~\ref{thm:transfer}, suppose
$D_{1,2}(N)\ge(c-o(1))\HN$ for every sufficiently large even $N$.
Then for each fixed $b\in(1/2,1)$,
\[
 \liminf_{\substack{N\to\infty\\2\mid N}}
 \frac{D_b(N)}{\HN}
 \ge c-\frac{80}{11}\log(1/b).
\]
\end{theorem}
\begin{proof}
The representations excluded from $D_b$ have
$N-p=rq$ with $q^b<r\le q$. The part with
$rq<N/(\log N)^{10}$ contributes $o(\HN)$, as do the possible
unit and small-prime endpoint terms. Apply
Lemma~\ref{lem:semiprime55} to the remaining excluded representations
and subtract their upper bound from $D_{1,2}(N)$.
\end{proof}

For the published input $c=1.9728=1233/625$, the exact threshold in
this theorem is
\[
 b>\exp(-13563/50000)=0.762418241890\ldots.
\]
In particular, the result gives
\begin{equation}\label{eq:enhanced-763}
 D_{0.763}(N)\ge\frac{1}{200}\HN
\end{equation}
eventually, \emph{with the external-input dependency on
\cite[Theorem 1.1]{Li2026} retained}. This numerical implication
does not resolve the independent input audit in
Section~\ref{sec:input-audit}.

For completeness, the strict margin in \eqref{eq:enhanced-763}
has a short rational certificate. For $0<u<1$,
\[
 \log\frac{1+u}{1-u}
 \le 2\left(u+\frac{u^3}{3}+
                 \frac{u^5}{5(1-u^2)}\right).
\]
Set $u=237/1763$, corresponding to $b=763/1000$. Then
\[
 \log(1000/763)
 \le\frac{11309527869744957}{41810125855610000},
\]
and hence
\[
 \frac{1233}{625}-\frac{80}{11}\log(1000/763)
 \ge\frac{159434349239058}{28744461525731875}
 >\frac1{200}.
\]
The strict inequality absorbs the asymptotic errors.

There is also a consequence which does not use the published
$1.9728$ input. Combining Theorems~\ref{thm:conservative}
and~\ref{thm:enhanced-transfer} gives
\begin{equation}\label{eq:enhanced-846}
 D_{0.846}(N)\ge\frac1{300}\HN
\end{equation}
for sufficiently large even $N$. Indeed, $u=77/923$ gives
\[
 \log(500/423)\le
 \frac{278128212904157}{1663088937705000},
\]
and therefore
\[
 \frac{61}{50}-\frac{80}{11}\log(500/423)
 \ge\frac{310893234857}{83154446885250}>
 \frac1{300}.
\]
This consequence still depends on the distribution and sieve theorems
cited in the new lemma and the conservative proof. Exact arithmetic
certifies these displayed numerical implications, not the analytic
inputs or the priority of the resulting statements.

% Integration fragment; all analytic dependencies are named explicitly.
\subsection{Combining the improved lower and switching estimates}
\label{sec:enhanced-chen}

\begin{theorem}[Enhanced conservative Chen calculation]
\label{thm:enhanced-chen}
The preceding derived sieve lemmas and the full-range quotient weight give,
for every sufficiently large even $N$,
\begin{equation}\label{eq:enhanced-chen-constant}
 D_{1,2}(N)\ge\frac85\HN.
\end{equation}
Consequently,
\begin{equation}\label{eq:enhanced-805}
 D_{0.805}(N)\ge\frac1{50}\HN.
\end{equation}
This deduction does not use the numerical Chen constant in
\cite{Li2026}. It still depends on the published distribution and sieve
inputs and on the adaptations proved in this draft.
\end{theorem}

\begin{proof}
Keep the pointwise quotient weight and the constants
$U_3,U_4,L_5,L_6,U_7,\ldots,U_{11}^{\rm full}$ from
Section~\ref{sec:conservative}. Write
$\widetilde U_i=(10/11)U_i$ for the last five constants.
Lemma~\ref{lem:primeproducts55} supplies these reduced upper constants,
and Lemma~\ref{lem:restricted-lower-level} supplies
$L_1(\nu),L_2(\nu)$ with $\nu\uparrow11/20$.
The same pointwise inequality therefore gives
\[
 \liminf_{\substack{N\to\infty\\2\mid N}}
       \frac{D_{1,2}(N)}{\HN}\ge c_{55},
\]
where
\[
 c_{55}=\frac14\left\{
 3L_1(11/20)+L_2(11/20)-U_3-U_4+L_5+L_6
 -2\widetilde U_7-\widetilde U_8-\widetilde U_9
 -\widetilde U_{10}-\widetilde U_{11}^{\rm full}\right\}.
\]
Every level is first used with a fixed positive margin; continuity
is used only after taking the large-$N$ limit.

Replace $L_1(11/20)$ by the explicit lower functional in
Section~\ref{sec:restricted-lower-level}. Reuse the outward enclosures
for the unchanged terms and multiply each of the five switching
enclosures by the exact rational number $10/11$.
The new one-dimensional enclosure uses 4,000 cells, a 30,000-cell
monotone prefix table for $J_0$, the same denominator $2^{80}$,
and the same 32-term logarithm series with its explicit tail.
The resulting exact lower endpoint is
\[
 c_{55}\ge
 \frac{1945721329817687310960304}{1208925819614629174706176}
 >\frac85.
\]
In particular, the excess over $8/5$ is at least
\[
 \frac{3575005760712697322007}{377789318629571617095680}>0.
\]
The certificate is \texttt{certify\_enhanced\_chen.json}; its source
calculation is \texttt{certify\_enhanced\_chen.py}, using the valid
full-range certificate \texttt{certify\_repaired\_chen.json}.
The legacy narrow-range certificate is not a mathematical premise.
All upper endpoints for the new $L_1$ calculation enclose only the
chosen lower functional, not the exact value of $L_1$.
For a short arithmetic check, the following outward-rounded bounds
suffice; the last five upper entries already include the factor $10/11$.
Every terminating decimal in the table is an exact rational number.
\begin{center}
\begin{tabular}{ll@{\qquad}ll}
\toprule
Lower term&Lower bound&Upper term&Upper bound\\
\midrule
$L_1(11/20)$&12.9014&$U_3$&19.3178\\
$L_2(11/20)$&6.3640&$U_4$&13.1277\\
$L_5$&2.3441&$\widetilde U_7$&1.3976\\
$L_6$&2.2784&$\widetilde U_8$&4.1788\\
&&$\widetilde U_9$&3.0259\\
&&$\widetilde U_{10}$&0.2321\\
&&$\widetilde U_{11}^{\rm full}$&0.5762\\
\bottomrule
\end{tabular}
\end{center}
Substitution in the displayed formula gives
$c_{55}\ge6437/4000=8/5+37/4000$ already from this shorter table.
The strict gap proves \eqref{eq:enhanced-chen-constant} after absorbing
the asymptotic error.

Now apply Theorem~\ref{thm:enhanced-transfer} with $c=8/5$ and
$b=161/200$. For $u=39/361$ the elementary positive logarithm series gives
\[
 \log(200/161)
 \le 2\left(u+\frac{u^3}{3}+\frac{u^5}{5(1-u^2)}\right)
 =\frac{3285966272199}{15148773682000}.
\]
Thus
\[
 \frac85-\frac{80}{11}\log(200/161)
 \ge\frac{46763937841}{2082956381275}
 =\frac1{50}+\frac{10209620431}{4165912762550}
 >\frac1{50}.
\]
This second strict gap absorbs the transfer error and proves
\eqref{eq:enhanced-805}.
\end{proof}

The numerical calculation certifies the displayed functional and
logarithm inequalities. It does not certify the analytic lemmas,
the external theorems, or the priority of the resulting statement.
The result remains a research proof draft for independent review.

% Integration fragment. Insert after the enhanced conservative calculation.
% This strengthens the upper level only; it does not alter the lower sieve.
\subsection{Using the third-prime support condition}
\label{sec:sharper-upper-support}

The uniform level $11/20$ above used only the first cube condition
of the actual modified upper weight. Retaining its third-prime
condition gives a further, entirely rational, improvement. Put
\[
 \theta_0=\frac{89}{160},\qquad \mu=\frac{25}{128}.
\]
For this subsection let $\Phi(t_1,t_2,t_3)$ denote the complete
eight-branch envelope in $L_{23}$, Proposition~6.1, equation~(6.4),
with residue parameter $\alpha=1$ and spectral parameter $\sigma=7/32$.
In particular its branches include
\[
 v(t_1),\qquad w(t_1,t_2,t_3),\qquad
 \psi\bigl(v(t_2+t_3),2t_1+t_2+t_3\bigr),
\]
where
\[
 v(t)=\min\left\{\frac{1+t}{2},\frac{32-34t}{39},
                         \frac{57}{64}-\frac{3t}{2}\right\},
 \qquad \psi(x,y)=x\mathbf1_{x\ge y},
\]
and
\[
 w(t_1,t_2,t_3)=\psi\left(
  \min\left\{\frac{66-34t_3}{107},
                 \frac{153}{256}-\frac{3t_3}{8},1-2t_2\right\},
  \frac{1+t_1}{2}\right).
\]

\begin{lemma}[An upper-support envelope bound]
\label{lem:upper-support-89}
If
\[
 0\le t_3\le t_2\le t_1,\qquad 3t_1\le\theta_0,
 \qquad t_1+t_2+3t_3\le\theta_0,
\]
then $\Phi(t_1,t_2,t_3)\ge\theta_0$.
\end{lemma}
\begin{proof}
The first support condition implies $t_1\le89/480<\mu$, so
$v(t_1)=(1+t_1)/2$. Write $B(t)=153/256-3t/8$.
The inner minimum in $w$ is exactly $B(t_3)$, since
\[
 \frac{66-34t}{107}-B(t)=\frac{525}{27392}+\frac{49t}{856}>0
 \quad(t\ge0),
\]
and $1-2t_2\ge151/240>153/256\ge B(t_3)$.
Consequently $\max\{v(t_1),w(t_1,t_2,t_3)\}
=\max\{v(t_1),B(t_3)\}$.

Suppose for a contradiction that $\Phi<\theta_0$.
The two preceding branches force
\[
 t_1<\frac9{80}=\frac{54}{480},\qquad
 t_3>\frac{53}{480}.
\]
Set $s=t_2+t_3$. Then
\[
 \frac{106}{480}<s<\frac9{40},\qquad
 2t_1+t_2+t_3\le4t_1<\frac9{20}<\theta_0.
\]
If $v(s)\ge\theta_0$, its displayed $\psi$ branch is therefore
present and contradicts $\Phi<\theta_0$. Hence $v(s)<\theta_0$.
On the indicated interval $s>\mu$, the minimum defining $v(s)$
is $57/64-3s/2$. Indeed the other two branches minus this
selected branch are nonnegative there. Thus
\[
 s>\frac{107}{480}.
\]
Finally, using $t_1\ge t_2$,
\[
 \theta_0\ge t_1+t_2+3t_3\ge2s+t_3
       >\frac{214+53}{480}=\theta_0,
\]
which is impossible.
\end{proof}

\begin{proposition}[A sharper derived upper level]
\label{prop:upper-89}
Lemma~\ref{lem:semiprime55} and Lemma~\ref{lem:primeproducts55}
remain valid with upper-sieve level
$D=N^{89/160-\varepsilon_0}$ and upper coefficient $640/89$
in place of $80/11$. The same admissible product domains,
bounded multiplicities, and exact coprime-mass convention are retained.
No stronger lower-sieve level is asserted by this proposition.
\end{proposition}
\begin{proof}
The construction of the actual modified upper atoms used in
Lemma~\ref{lem:semiprime55} imposes, for at least three rough slots,
both
\[
 D_1^3<D_0,\qquad D_1D_2D_3^3<D_0
\]
(with a smaller right side for the odd-rank case).
After box inflation their logarithmic exponents thus obey
\[
 3t_1\le\nu+O(\eta),\qquad
 t_1+t_2+3t_3\le\nu+O(\eta),
 \qquad \nu=89/160-\varepsilon_0.
\]
These are properties of the actual Iwaniec upper-weight expansion,
not assumptions that the ordinary Rosser weight is factorable.

First work at residue parameter $\alpha=1$.
For small enough $\eta$, the exponent vector lies in the closed
domain of Lemma~\ref{lem:upper-support-89}, which gives the
level $\theta_0$ for every such vector. In applying
$L_{23}$, Proposition~6.1, choose instead the strictly smaller
parameter $\vartheta=\nu+\varepsilon_0/2$.
The rough vector belongs to
$\mathcal D^{\rm well}(N^{\vartheta-2\zeta})$ after choosing
$\eta\ll\zeta\ll\varepsilon_0$.
On a local scale write $\rho=\log N/\log x=1+O(\kappa)$.
The precise support requirement becomes
\[
 \rho(\nu+C\eta)+2\zeta<\vartheta.
\]
It holds with the full parameter hierarchy
$\eta,\kappa\ll\epsilon_{\rm disp}\ll\zeta\ll\varepsilon_0$;
the same hierarchy also ensures $N<x^{1+\epsilon_{\rm disp}}$.

The passage to local scales and to the fixed residue parameter
$\alpha=1+3\kappa$ needs a quantitative margin, since $\psi$
alone is discontinuous. The proof of the lemma supplies one.
At any vector, either $v(t_1)\ge\theta_0$, or
$B(t_3)\ge\theta_0$ with its gate no larger than $v(t_1)$,
or the final pair branch supplies the bound.
In the last case its gate is at most $4(2\theta_0-1)=9/20$,
strictly below $\theta_0$. In the middle case the maximum
with $v(t_1)$ equals the continuous function
$\max\{v(t_1),B(t_3)\}$.
The single-variable branches change by $O(\kappa)$, as do
the inner minimum and the gates after changing scale and
$\alpha$. Thus the applicable maximum stays at least
$\theta_0-O(\kappa+\eta)$.
Choose $\kappa,\eta\ll\varepsilon_0$; this still exceeds
$\vartheta$ with a fixed margin, and $t_1<\mu_\alpha$ remains
strict because $89/480<25/128$.

All subsequent steps of the two earlier lemmas are unchanged:
normalized modulus-factor coefficients, rough Siegel--Walfisz
stability, the single- and multi-group Heath--Brown alternatives,
uniformity for residue $N$, Mellin separation, and the coprime
main-mass correction. There is ample space in the two-long-variable
estimate since $\theta_0<3/5<2/3$.
Atoms with at most two rough slots have full support
$N^{2\nu/3+O(\eta)}<N^{1/2-c_0}$ and are still handled by
ordinary Bombieri--Vinogradov or its convolution form.

Finally at $z=N^{1/5}$ the linear-sieve parameter satisfies
$1<5\nu<3$, so $F(5\nu)=2e^\gamma/(5\nu)$.
The upper coefficient is $4/\nu$. Taking the limits in the
same order as before gives $4/\theta_0=640/89$.
\end{proof}

\begin{corollary}[Sharper transfer]
\label{cor:transfer-89}
If $D_{1,2}(N)\ge(c-o(1))\HN$ for every sufficiently large even $N$,
then for each fixed $b\in(1/2,1)$,
\[
 \liminf_{\substack{N\to\infty\\2\mid N}}\frac{D_b(N)}{\HN}
 \ge c-\frac{640}{89}\log(1/b).
\]
\end{corollary}
\begin{proof}
Apply Proposition~\ref{prop:upper-89} to the excluded semiprimes
and make the same subtraction as in Theorem~\ref{thm:enhanced-transfer}.
\end{proof}

\paragraph{Exact numerical consequences.}
Replace all five switching upper constants in $c_{55}$ by
$\widehat U_i=(80/89)U_i$, retaining $L_1(11/20)$,
$L_2(11/20)$ and every other term. Proposition~\ref{prop:upper-89}
covers exactly these same labelled, full-range product multisets.
Call the resulting lower functional $c_{89}$. Exact propagation
of the preceding numerical enclosures gives
\[
 c_{89}\ge
 \frac{705742503700722615189406077}
      {430377591782807986195398656}
 >\frac{1639}{1000}.
\]
Thus $D_{1,2}(N)\ge(1639/1000)\HN$ eventually. The strict
inequality absorbs the Chen asymptotic error. For $b=797/1000$,
the positive atanh-series bound already used above gives
\[
 \frac{1639}{1000}-\frac{640}{89}\log\frac{1000}{797}
 >\frac{9082360163231797}{1234848885445827000}
 >\frac1{200}.
\]
Consequently the separate route gives
\begin{equation}\label{eq:sharper-797}
 D_{0.797}(N)\ge\frac1{200}\HN
\end{equation}
for every sufficiently large even $N$.

Retaining the published numerical input $c=1233/625$ instead,
the corresponding arithmetic is
\[
 \frac{1233}{625}-\frac{640}{89}\log\frac{1000}{761}
 >\frac{6086593811525371}{693512037060091875}
 >\frac1{200}.
\]
It gives the external-input consequence
\begin{equation}\label{eq:sharper-761}
 D_{0.761}(N)\ge\frac1{200}\HN,
\end{equation}
retaining exactly the limitation discussed in
Section~\ref{sec:input-audit}.

These rational calculations and the source hashes are recorded in
\texttt{certify\_sharper\_transfer.json}. The certificate also gives
a strictly sub-endpoint sieve level and explicit allocations of the
positive slack to the asymptotic errors. It is a certificate of the
numerical implications, not a verification of the analytic inputs.

% Independent integration/research fragment: no documentclass.
% Does not edit or supersede the main manuscript automatically.
% Uses the existing v, a, h, d, H_N and prime-prefix interface notation.

\subsection{A stronger uniform level for the two-prime lower sieve}
\label{sec:double-prefix-level}

This refinement concerns the two-prime lower terms $L_5,L_6$.
It retains the residual sieve construction of
Lemma~\ref{lem:prefix-interface}; in particular it does not identify
that construction with a prefix restriction of a global lower weight.
Set
\[
 a=\frac{100}{1149},\qquad h=\frac{50}{309},\qquad
 d=\frac12-3a,\qquad \sigma=\frac7{32},\qquad
 \mu=\frac{25}{128},
\]
and define
\[
 K_a=\frac{5-\sigma-3a}{8}
       =\frac{153}{256}-\frac{3a}{8}
       =0.56501917428\ldots .
\]
For $a\le u\le\min(t,h)$ and $t\le d$, put
\begin{equation}\label{eq:double-prefix-level}
 \mathcal V_2(t,u)=
 \begin{cases}
 \max\{v(t),v(t+u),K_a\},&t\le\mu,\\
 v(t),&t>\mu.
 \end{cases}
\end{equation}
The second line is essential: the apparent constant-branch gain
beyond the largest-prime domain of the factorization theorem
must not be used.

\begin{lemma}[Derived two-prime level]
\label{lem:double-prefix-level}
On the $L_5$ and $L_6$ domains, the two-prime version of
Lemma~\ref{lem:prefix-interface} remains valid with $v(t)$ replaced
by $\mathcal V_2(t,u)$, using fixed positive margins and sufficiently
fine exponent boxes. Its residual level has exponent
\[
 \mathcal V_2(t,u)-t-u,
\]
and its lower-sieve multiplier at cutoff $N^a$ is therefore
\[
 f\left(\frac{\mathcal V_2(t,u)-t-u}{a}\right).
\]
The assertion uses the ordinary modified lower sieve and the
uniform prime-dispersion argument already proved in that interface.
It introduces no double-sieve correction.
\end{lemma}

\begin{proof}
Keep all margins and exponent boxes fixed while $N\to\infty$.
Write the two prefix primes in decreasing order as
$p_1=N^{t+O(\Delta)}$, $p_2=N^{u+O(\Delta)}$, where
$u>a$ after trimming a boundary strip. Construct the lower weight
on the residual divisor $q\mid P(N^a)$ at level
\[
 R=N^{\nu-T_+-20\delta},
 \qquad T_+=(t+\Delta)+(u+\Delta),
\]
where $\nu$ is below the infimum of $\mathcal V_2$ on the
prefix box by a fixed positive margin. Its rough atom is a product
of primes smaller than $N^a$, followed by the usual small-prime
factor of support $N^{O(\eta)}$.
Thus the specified prefix supplies the two largest prime slots
of every concatenated rough atom.

The two initial $D^{\rm well}$ inequalities are
\[
 2(t+\Delta)<\nu-20\delta,
 \qquad (t+\Delta)+2(u+\Delta)<\nu-20\delta.
\]
Every later inequality follows by multiplying the residual
$D^{\rm well}(R^{1+O(\eta)})$ condition by $p_1p_2$.
Consequently the full rough vector belongs to
$D^{\rm well}(N^{\nu-c\delta})$ after all box losses.
This explains why the total analytic level is $\nu$, while the
sieve level available after the prefix is $\nu-t-u$.

If the residual rough atom is empty, the full atom consists of
$p_1p_2$ and the small-prime factor. Its support is at most
\[
 N^{t+u+O(\eta)}\le N^{d+h+O(\eta)}<N^{1/2-c_0}
\]
for a fixed $c_0>0$. Ordinary Bombieri--Vinogradov handles this
case, with bounded normalized atom coefficients. There is no
third-prime exponent to which a three-variable formula must be
artificially assigned.

Otherwise the concatenated atom has at least three rough slots.
Its third exponent, denoted $r$, satisfies
$0<r\le a+O(\eta)$. First take $t<\mu$ with a fixed margin.
Proposition~6.1, (6.4), of
\href{https://arxiv.org/pdf/2309.08522v1}
{Lichtman, arXiv:2309.08522v1}, explicitly permits the
subsum branch $v(t+u)$ as well as $v(t)$.
The condition on the largest prime is checked on the actual
concatenated vector, rather than on a reordered or artificial one.

Here is a uniform lower envelope for the permitted $w$ branch.
At residue exponent $\alpha=1$, its inner minimum is
\[
 \min\left\{A(r),B(r),1-2u\right\},
 \qquad A(r)=\frac{66-34r}{107},\quad
 B(r)=\frac{153-96r}{256}.
\]
One has
\[
 A(r)-B(r)=\frac{525+1568r}{27392}>0,
 \qquad 1-2u\ge1-2h>\frac{153}{256}\ge B(r).
\]
Thus this inner minimum is exactly $B(r)$, which is at least
$B(a)=K_a$ up to the fixed construction losses.
Its $\psi$ threshold is $(1+t)/2=v(t)$ because $t<\mu$.
If $K_a>v(t)$, the branch is present with a strict margin and
supplies level at least $K_a$. If $K_a\le v(t)$, the existing
$v(t)$ branch suffices. It follows that every nonempty residual
atom admits the uniform level in \eqref{eq:double-prefix-level}.
The alternative displayed branch $w(t,r,u)$ cannot improve
this worst-case bound: its inner minimum is $B(u)\le B(a)$.
This statement only concerns the branches used here, not the
other $\psi$ branches in the complete formula (6.4).

For $t>\mu$ we retain the already established single-variable
level. In particular, although $K_a$ can exceed $v(t)$ near the
right end of the $L_6$ interval, the argument identifying the
$\psi$ threshold with $v(t)$ is then false. Such a substitution
would give an unjustified gain.

For each residual box the prime-slot partition supplied by
Proposition~6.1 is fixed throughout that box. The original signs,
factorial normalizations, and small-prime factor are preserved
exactly as in Lemma~\ref{lem:prefix-interface}.
No atom-dependent level is substituted into a fixed sieve main
term: the same chosen $R$ is used on every residual atom in a
given prefix box, and all those atoms have just been covered.
Also
\[
 \mathcal V_2(t,u)\le\frac{153}{256}<\frac35.
\]
Therefore the Type~II and both Heath--Brown endpoint checks of
the prime-prefix proof continue to apply without increasing
their total-modulus bound. Its moving-residue proof, rough
Siegel--Walfisz lemma, and uniform density calculation are
unchanged. In passing to local scales use the fixed positive
margins to absorb $\alpha=1+O(\kappa)$ and exponent perturbations.

For $L_5$, both prefix inequalities hold throughout the domain
with room to spare already at the old level: $2t\le2h<v(t)$
and $t+2u\le3h<(1+a)/2\le v(t)$.
For $L_6$, where the lower function is positive,
$\nu-t-u>2a>h\ge u$ implies $t+2u<\nu$.
The other inequality $2t<v(t)$ holds for every $t\le d$.
Where the lower function is zero, nonnegativity suffices.
Trim ordering boundaries, $t=\mu$, and any equality boundary;
take $N\to\infty$, then the mesh and construction precision
to zero, and finally the margins to zero. The stated lower
multiplier follows by continuity of $f$.
\end{proof}

\paragraph{What this refinement actually improves.}
The legitimate replacement is
\[
 L_5^{(2)}=\frac{2e^{-\gamma}}a
   \int_a^h\int_a^t
       f\left(\frac{\mathcal V_2(t,u)-t-u}{a}\right)
                   \frac{du\,dt}{tu}.
\]
It is strictly larger than $L_5$. For instance, on the nonempty
open triangle $a<u<t$ and $t+u<\mu$, one has
\[
 \mathcal V_2(t,u)=v(t+u)=\frac{1+t+u}{2}
                 >\frac{1+t}{2}=v(t).
\]
The sieve lower function is strictly increasing at these arguments.
At the lower corner, the limiting arguments change from
$4.245$ to $4.745$; no gain estimate is inferred from that
single point alone. More generally the subsum branch improves
the old level when either $t+u\le\mu$, or
\[
 t+u>\mu,\qquad 4t+3u<\frac{25}{32}.
\]
The constant branch gives additional level improvement when
$t<2K_a-1=0.13003834856\ldots$ and is larger than the other branches.

For this particular uniform envelope, $L_6^{(2)}=L_6$.
Indeed, for $h\le t\le\mu$,
\[
 K_a<v(h)\le v(t),\qquad
 v(t+u)\le v(h+a)<v(h)\le v(t),
\]
because $h+a>\mu$ and $v$ decreases above $\mu$.
For $t>\mu$ the definition deliberately retains $v(t)$.
Thus this route supplies a genuine gain to the combined Chen
coefficient of exactly $(L_5^{(2)}-L_5)/4$, pending a certified
evaluation, and does not claim an improvement to $L_6$.

The full maximum in Proposition~6.1 contains other subsums and
conditional $\psi$ branches. Optimizing those would require a
separate uniform envelope over the actual third exponent and
checking every associated threshold. They are not silently
included in the present formula.

\subsection{Certified consequence with the two-prime improvement}
\label{sec:double-prefix-corollary}

\begin{theorem}[Separate calculation with the two-prime improvement]
\label{thm:double-prefix-chen}
Combining the full-range quotient weight with the derived estimates
above gives, for every sufficiently large even integer $N$,
\[
 D_{1,2}(N)\ge\frac{1643}{1000}\HN,
 \qquad D_{0.796}(N)\ge\frac1{500}\HN.
\]
This route does not use the numerical Chen bound in \cite{Li2026}.
It retains the published analytic inputs and the derived interfaces
proved in the present manuscript.
\end{theorem}
\begin{proof}
Replace $L_5$ by $L_5^{(2)}$ in the functional $c_{89}$ from
Section~\ref{sec:sharper-upper-support}, keeping all other terms.
Denote the resulting functional by $c_{89,2}$. The preceding lemma
justifies this replacement in the same pointwise quotient weight.

The sieve argument in $L_5^{(2)}$ obeys
\[
 2<s=\frac{\mathcal V_2(t,u)-t-u}{a}
 \le\frac{153}{256a}-2<6.
\]
For the lower inequality, already at the old level
$v(t)-t-u\ge(1-3h)/2>2a$.
Consequently the exact formula
\[
 \frac{s f(s)}{2e^\gamma}
   =\log(s-1)+K(s),\qquad
 K(s)=\begin{cases}
 0,&s\le4,\\
 \displaystyle\int_3^{s-1}\frac{J_0(v)}v\,dv,&4<s\le6
 \end{cases}
\]
applies throughout this integral. Here $J_0$ is defined in
Section~\ref{sec:restricted-lower-level}.

Map the triangle to $t\in[a,h]$, $y\in[0,1]$ by
$u=a+(t-a)y$, with Jacobian $t-a$.
Use $512^2$ rectangle cells, the existing 80-bit integer interval
arithmetic and explicit 32-term logarithm tail, and 30,000-cell
prefix enclosures for each of $J_0$ and $K$.
Interval minima and maxima enclose the affine branches in
$\mathcal V_2$ even across their intersections.
The calculation gives
\[
 L_5^{(2)}\ge
 \frac{2853747390748858418997292}{1208925819614629174706176}
 >2.3605.
\]
The numerical source is \texttt{certify\_double\_prefix.py} and
its exact endpoints are recorded in
\texttt{certify\_double\_prefix.json}.
The following weaker rational table suffices for a short check.
The upper switching entries include the factor $80/89$.
\begin{center}
\begin{tabular}{ll@{\qquad}ll}
\toprule
Lower term&Lower bound&Upper term&Upper bound\\
\midrule
$L_1(11/20)$&12.9014&$U_3$&19.3178\\
$L_2(11/20)$&6.3640&$U_4$&13.1277\\
$L_5^{(2)}$&2.3605&$\widehat U_7$&1.3819\\
$L_6$&2.2784&$\widehat U_8$&4.1318\\
&&$\widehat U_9$&2.9919\\
&&$\widehat U_{10}$&0.2294\\
&&$\widehat U_{11}^{\rm full}$&0.5697\\
\bottomrule
\end{tabular}
\end{center}
Substitution gives, by exact rational arithmetic,
\[
 c_{89,2}\ge\frac{263}{160}
       =\frac{1643}{1000}+\frac3{4000}.
\]
The strict gap absorbs the Chen asymptotic error and proves the
first assertion.

For the second, apply Corollary~\ref{cor:transfer-89} with
$c=1643/1000$ and $b=199/250$.
Set $u=51/449$ in the positive atanh-series bound. It gives
\[
 \frac{1643}{1000}-\frac{640}{89}\log\frac{250}{199}
 >\frac{3730223387621}{1603179334639000}>\frac1{500}.
\]
The certificate also exhibits a strictly smaller sieve level than
$89/160$ and a positive allocation for every asymptotic error.
This proves the asserted eventual counting bound without using an
unproved endpoint distribution statement.
\end{proof}

All decimals in the short table are terminating rational bounds.
The numerical certificate verifies these integral and logarithm
inequalities; it does not replace the analytic proofs or establish
priority over the literature.

\section{Audit of the numerical input}\label{sec:input-audit}
The following is a limitation of this draft's independent verification,
not a disproof of the theorem in \cite{Li2026}.
Its equations (16) and (23)--(25) combine linear-sieve terms evaluated at
an improved distribution level with double-sieve corrections evaluated
at the older level $1/2$. Schematically the upper combination has the form
\[
 F(s_{\rm new})-\frac{2e^\gamma}{s_{\rm old}}
                      H_{1/2}(s_{\rm old}),
 \qquad s_{\rm new}>s_{\rm old}.
\]
In contrast, the explicitly checked extension in
\cite[Lemma 6.4]{Lichtman2025} uses
$F(s_{\rm new})-2e^\gamma H_{\vartheta}(s_{\rm new})/s_{\rm new}$,
with additional distribution hypotheses on all extended prime vectors
and the switched sets. Monotonicity
$H_{\vartheta}(s)\ge H_{1/2}(s)$ at fixed $s$ does not by itself justify
the displayed mixed-parameter correction. The analogous issue affects
the lower-sieve correction $h$. We have not supplied the further
argument that would close this interface.

A second interface concerns the distribution function with three
arguments. The improvement in \cite[Lemmas 2.1--2.2]{Li2026} is stated
under the restriction that the largest prime exponent is at most
$25/128$. Equations (25) and (34) contain arguments beyond this range.
For example, $(t_1,t_2,u)=(0.22,0.15,1/11.49)$ occurs in (25);
the raw three-variable maximum is at least $0.575$, whereas the stated
single-variable level at $0.22$ is $0.560625$. If a piecewise extension
back to the single-variable level is intended, that convention and its
use in the numerical calculation need to be confirmed. An independent
recalculation must respect the stated range explicitly.

Consequently Theorem~\ref{thm:transfer} remains a complete transfer
statement, while its $1.9728$ applications retain the external-input
dependency stated above. One possible independent repair is to set all
double-sieve savings to zero, restrict all distribution parameters to
their proved domains, and certify the remaining numerical integrals.
Section~\ref{sec:conservative} supplies a separate conservative
replacement with its own derived sieve interface, full-range quotient
weight, and exact numerical enclosure. It yields $c=1.22$.
The additional upper- and lower-sieve derivations then strengthen
this separate route to $c=1.6$ in Theorem~\ref{thm:enhanced-chen},
without importing the mixed-parameter double-sieve corrections.
The support and two-prime refinements give $c=1.643$
in Theorem~\ref{thm:double-prefix-chen}; the retained-base construction
raises this separate calculation to $c=1.731$ in
Theorem~\ref{thm:retained-base-chen}, and the variable-prefix
construction raises it to $c=1.783$ in
Theorem~\ref{thm:variable-prefix-chen}. The extended first stop
then gives $c=1.785$ in Theorem~\ref{thm:extended-prefix-chen}.
The restricted-lower and Buchstab improvement gives $c=1.817$
in Theorem~\ref{thm:buchstab-lower-chen}, still without any
mixed-level double-sieve input.
The older constant $c=0.899$ from \cite{Wu2008} can be
used without this particular mixed-level step, but merely substituting
that constant does not establish the requested advance beyond the
existing benchmark.

% Source-audit fragment, 2026-10-07. No documentclass; not inserted in main.
% This identifies missing interfaces and gives valid replacement formulas.
% It neither refutes Li's theorem nor certifies its numerical constant.

\subsection{What the cited double-sieve bounds actually transfer}
\label{sec:external-chen-input-resolution}

Let $d=N^{T+o(1)}$ be a specified ordered prime prefix and let
$z=N^u$. Write
\[
 s_0=\frac{1/2-T}{u},\qquad s_1=\frac{\nu-T}{u},
 \qquad \nu>1/2.
\]
All statements below concern an averaged positive prefix box,
with the common main factor denoted by $M>0$. They do not assert
an estimate for each individual prefix prime.
The improved linear estimate at level $N^\nu$ must itself be
justified on that box.

In Wu's \href{https://arxiv.org/pdf/0709.3764}{2008 author version},
Section 3, the total level is $Q=N^{1/2-\delta}$, the residual
level is $Q/d$, and $A(s)=sF(s)/(2e^\gamma)$,
$a(s)=sf(s)/(2e^\gamma)$.
The functions $H,h$ are defined through the defect from these
same-level normalized functions in (3.4)--(3.7).
Their prefix class satisfies
$p_1\cdots p_{j-1}p_j^2\le Q$.
This is also the normalization underlying Section 3 and Lemmas
5.1--5.2 of
\href{https://www.impan.pl/shop/en/publication/transaction/download/product/82961}
{Wu's 2004 paper}; the latter's functional inequalities iterate
the same total level, rather than adding an independent saving
to a different linear-sieve baseline.

Thus, after removing fixed margins, the two separately established
upper estimates have the form
\[
 S/M\le F(s_1)+o(1),\qquad
 S/M\le F(s_0)-D_H+o(1),\qquad
 D_H=\frac{2e^\gamma}{s_0}H(s_0).
\]
The analogous two lower estimates use $f(s_1)$ and
$f(s_0)+D_h$, where $D_h=2e^\gamma h(s_0)/s_0$.
Whenever both inputs are valid on the same positive box, their
rigorous combination is
\begin{align}
 U_{\rm safe}
 &=\min\{F(s_1),F(s_0)-D_H\}\notag\\
 &=F(s_1)-[D_H-\{F(s_0)-F(s_1)\}]_+,
 \label{eq:external-safe-upper}\\
 L_{\rm safe}
 &=\max\{f(s_1),f(s_0)+D_h\}\notag\\
 &=f(s_1)+[D_h-\{f(s_1)-f(s_0)\}]_+.
 \label{eq:external-safe-lower}
\end{align}
Here $[a]_+=\max(a,0)$. One can use certified lower bounds for
$H,h$ in these formulas. A finite common partition permits the
choice of the better bound on each box before taking the
large-$N$ limit. On boxes outside Wu's prefix class, omit the
classical double-sieve option.

These formulas follow by taking a minimum or maximum of bounds
for the same quantity. They do not require a new double-sieve
theorem. Relative to the mixed expressions
$U_{\rm mix}=F(s_1)-D_H$ and $L_{\rm mix}=f(s_1)+D_h$, the exact
local corrections needed are
\begin{align*}
 U_{\rm safe}-U_{\rm mix}
   &=\min\{D_H,F(s_0)-F(s_1)\},\\
 L_{\rm mix}-L_{\rm safe}
   &=\min\{D_h,f(s_1)-f(s_0)\}.
\end{align*}
They quantify the loss that a corrected integral calculation
must include. Merely imposing the two separate bounds cannot
remove that loss: for example, when
$F(s_1)\le F(s_0)-D_H$ and $D_H>0$, the hypothetical value
$S/M=F(s_1)$ satisfies both bounds but violates $U_{\rm mix}$.
This is a logical non-implication, not a counterexample to a
prime-counting theorem.

\paragraph{The additional statement needed for the full mixed saving.}
Lichtman,
\href{https://arxiv.org/html/2109.02851v2}{arXiv:2109.02851v2},
Lemma 6.4 and (6.23), give the upper correction
\[
 F(s_1)-\frac{2e^\gamma}{s_1}H_\nu(s_1),
 \qquad s_1\in[2,3],
\]
subject to distribution at the same level for every extended
prime vector and the corresponding switched sets.
The largest specified prime is preserved during the iteration
in the application following that lemma. The stated monotonicity
is $H_\nu(s)\ge H_{1/2}(s)$ at a fixed argument $s$.
To deduce the full old-argument saving used in Li's formulas
from this statement one would additionally need
\begin{equation}\label{eq:external-missing-cross-level}
 H_\nu(s_1)\ge\frac{s_1}{s_0}H_{1/2}(s_0),
 \qquad
 h_\nu(s_1)\ge\frac{s_1}{s_0}h_{1/2}(s_0),
\end{equation}
or a direct joint sieve construction proving the desired mixed
bounds. Fixed-argument monotonicity does not imply these
inequalities; $H_{1/2}$ decreases as its argument increases.
For the lower correction, the final paragraph of Lichtman's
Section 6.3 explicitly leaves the analogous formula derivations
to the reader. Extension of its stated $H$ lemma beyond $[2,3]$
also requires an argument at the new level.

The checked passages in
\href{https://imar.ro/journals/Mathematical_Reports/Pdfs/2026/1_2/3.pdf}
{Li II}, (16), (19), (21), and (23)--(25), use the old arguments
$(1/2-T)/u$ in $H,h$ alongside improved linear terms. The
paragraph following (16) cites fixed-argument monotonicity.
The earlier \href{https://arxiv.org/html/2405.05727v3}{2024 v3}
already has the same combination and does not supply a further
cross-level lemma. Consequently the checked citations do not
close \eqref{eq:external-missing-cross-level}. This audit has
not proved those inequalities false.

\paragraph{The classical prefix domains are consistent.}
Put $a=100/1149$, $h=50/309$, $\mu=25/128$, and
$d_*=1/2-3a$. The old double-sieve domains in Li (23)--(25)
are compatible with the old support class. For (23),
$t_1+2t_2\le3h<1/2$. In (24), the split at $1/2-2h$
and the later bound $t_2\le39/256=(1/2-\mu)/2$ ensure
$t_1+2t_2\le1/2$. In the correction of (25),
$t_1\le d_*$ and $t_2\le3a/2$ give the same inequality.
Also $2t_1<1/2$ throughout these regions.
Thus the old-level terms are usable as separate bounds;
the issue is their addition to the new baseline.
Where only tabulated values for $s\ge2$ are supplied, using
$h(s)\ge0$ for smaller $s$ is a valid conservative choice.

\subsection{The largest-prime restriction is active}
\label{sec:external-largest-prime-resolution}

Let $\Phi$ denote the raw eight-branch maximum of L23 (6.4),
and let $v(t)$ be its single-variable level. Li's Section 2
refers to these functions and states the three-variable
improvement only when the largest prime exponent is at most
$\mu$. No piecewise extension of the three-variable function
was located in the checked Li II text or its cited 2024 v3.
In contrast, L23 itself uses the three-variable expression in
(7.5) only for $c\le\mu$ and explicitly returns to $v(t)$
in (7.6) for $c>\mu$.

A conservative level supported by the stated distribution
theorems, with their usual strict margins, is therefore
\[
 \Phi_{\rm safe}(x,y,z)=
 \begin{cases}
  \Phi(x,y,z),&x\le\mu,\\
  v(x),&x>\mu,
 \end{cases}
 \qquad x\ge y\ge z.
\]
This remains a valid conservative choice. The additional construction
in Section~\ref{sec:selected-subproduct-factorization} proves some
strict-gate branches also for $x>\mu$; it does not extend the entire
raw eight-branch maximum. The full prefix interface is still required
wherever the first case is invoked; this notation does not bypass
its hypotheses.
For (34), membership in $D^{\rm well}$ must also be recomputed
at the selected safe level; when it fails, retain only the
switched estimate.

This restriction is not rendered moot by vanishing integrands.
In (25), take
\[
 (t_1,t_2,u)=(11/50,3/20,100/1149).
\]
This lies strictly inside its two-dimensional integration
domain. Here
\[
 \Phi=\frac{23}{40},\quad v(t_1)=\frac{897}{1600},\quad
 s_{\rm raw}=\frac{47109}{20000}>2,\quad
 s_{\rm safe}=\frac{70089}{32000}>2.
\]
Since both arguments lie between 2 and 4, the lower-sieve
function is strictly increasing there. Hence the raw and safe
lower terms differ on an open neighborhood.

For (34), the earlier comparison used the ascending tuple
\[
 (t_1,t_2,t_3,t_4)=(11/125,9/100,23/250,101/500).
\]
It is strictly inside the displayed integration domain, with
$t_4>\mu$. The full four-prime square-prefix exponents are
\[
 101/250,\quad193/500,\quad237/500,\quad14/25.
\]
They are all less than both
\[
 \Phi(t_4,t_3,t_2)=\frac{591}{1000},\qquad
 v(t_4)=\frac{4701}{8000}.
\]
The raw maximizing pair branch has gate surplus $1/200$.
Writing $S=59/125$, the two upper-sieve arguments are
$119/88$ and $925/704$, respectively. Both exceed 1, so neither
the support indicator nor a zero-density convention removes
this example. With the published switched coefficient cap
$G_1=6.06932$, the linear option is also smaller than the
switched option here: their ratio is
\[
 \frac{F((\nu-S)/t_1)t_2}
 {G_1\omega((1-S)/t_2)t_1}
 =\frac{2e^\gamma t_2}{G_1\omega((1-S)/t_2)(\nu-S)}<1
\]
for either displayed $\nu$, using
$\omega(u)\ge1/2$ for $u\ge2$, $e^\gamma<1.782$, and
$\nu-S\ge37/320$. Thus the minimum does not erase the difference
between these two numerical choices. However, the selected-subproduct
construction now proves the pair branch at this point and on a
neighborhood. This example is no longer an unresolved factorization
obstruction or evidence that a fallback is necessary there.
The Li (25) example above remains outside the branches proved so far.

The conservative fallback and the new branch-specific construction
provide rigorously specified options. The surviving task is to recompute
the affected functionals, with
\eqref{eq:external-safe-upper}--\eqref{eq:external-safe-lower}
or a proved stronger joint construction. No new numerical
Chen coefficient is certified in this fragment. In particular,
these findings neither independently establish nor disprove
the published value $1.9728$.

\section{Comparison and limitations}
In the same ``$1+a$'' convention, $a=1+b$ and smaller $a$ is stronger.
The explicitly stated comparison in \cite[Theorem 1.1 and its following
remark]{LiLiu2026} is $1+1.9$, with an indicated improvement to
$1+1.894$. That printed result is not the only relevant benchmark.
Even the classical transfer in this manuscript combined with the
published $c=1.9728$ gives $b=0.782$.

A stronger published-method comparison comes from
\cite[Lemma 2.2 and (27)--(29)]{Li2026}. The outer number of prime
factors $k$ in that lemma may equal two; its sieve-prefix rank $r$
is a different parameter. The displayed Buchstab functional has
$G_1<6.06932$, with the same normalization giving
\[
 K_{G_1}=2e^{-\gamma}G_1<\frac{852}{125}=6.816.
\]
If its uniform shifted local-product version is supplied, applying
that functional to the removed semiprime pairs gives
\[
 \liminf\frac{D_b(N)}{\HN}
 \ge\frac{1233}{625}-\frac{852}{125}\log(1/b).
\]
In particular the right side is greater than $1/100$ at $b=3/4$
and greater than $1/500$ at $b=749/1000$. These rational margin
checks take the published $G_1$ integral bound as an input; they do
not independently certify that integral or the Chen theorem.

The localization qualification matters. The literal Lemma 2.2
uses $N$ as both residue and product scale and normalizes the outer
exponents to sum to one. Products $rq<N$ require local scale $x$
with the residue still $N$, including variable-level prefixed sieve
weights. The source itself uses localization for its switched
three- and four-prime sequences. Proposition 5.2 of
\cite{Lichtman2023} supplies the needed moving-residue range at
the convolution-estimate level, but that alone is not a verbatim
statement of the full localized prime-product lemma.
Section~\ref{sec:g1-localized-interface} supplies a derived interface
for the small-initial-cutoff limit, using ordinary BV for its first
three terms. This avoids reliance on unverified virtual-prefix
levels in reconstructing the exact printed finite-cutoff formula.
Neither $b=0.749$ nor the optimized $b=0.74865$ is attributed as
a theorem printed in Li's paper.

Theorem~\ref{thm:variable-prefix-chen} gives $b=0.770$ using its
separate $c=1.783$ calculation, and $b=0.74865$ using the external
$c=1.9728$ input. The exact comparison certificate gives
\[
 6.815354545604\le2e^{-\gamma}(6.06932)
                  \le6.815354545605,
\]
whereas the new coefficient is strictly below $6.814$.
At $b=0.74865$, subtraction using that printed decimal has a
negative margin, while the new coefficient has margin greater
than $1/4000$. This compares two available numerical upper bounds;
it is not a lower bound for the actual $G_1$ integral or a claim
that the printed method cannot be optimized.
Section~\ref{sec:g1-limit-comparison} proves the precise relation
to its optimized limiting functional. The separate $b=0.770$
conclusion remains weaker than the external-input conclusion.
The extended first stop subsequently certifies $K<6.810$, giving
$b=0.7695$ on the separate route and $b=0.7485$ with the same external
input, as stated in Theorem~\ref{thm:extended-prefix-chen}. The new
first-stop range is covered by the single-variable source assertion;
its noncompact third-prefix tail is bounded explicitly. These changes
do not establish superiority to every known optimization of this method.
The separate count is then improved to $c=1.817$ and $b=0.766$ by
Theorem~\ref{thm:buchstab-lower-chen}. Its starting lower level is
$249/448$ and its next step is an exact least-prime decomposition;
the stronger external-input implication remains unchanged.
Uncertified floating explorations are excluded from all theorems.

A separate check of Pascadi's newer distribution results in
\cite[Theorems 1.3 and 5.4, Proposition 4.4]{Pascadi2025}
does not supply a direct stronger bound for this problem. Those
statements retain an error constant depending on the residue $a$;
the introduction explicitly discusses additional obstacles when
$a$ grows with the product scale. Their improved fixed-residue
levels therefore cannot simply replace the moving-residue input
with $a=N$. This does not rule out a future independent extension.
That paper also identifies the capacity issue in the L23
factorization; the repairs in
Section~\ref{sec:two-factor-greedy-repair} are included explicitly.

A search not finding the same corollary does not prove priority.
No world record, exhaustive literature coverage, or independent
validation of the cited Chen theorem is asserted. The full Goldbach
conjecture would correspond to $b=0$; the present transfer theorem
only concerns $b>1/2$ and proves nothing at that endpoint.

% Integrated proof fragment; numerical arithmetic is certified separately.
\section{An upper sieve with a third-prefix stopping rule}
\label{sec:pruned-upper-sieve}

The construction below changes the upper weight itself. It retains
only third-prefix boxes for which the available distribution estimates
have a fixed positive margin, and pays explicitly for the stopped boxes.
It does not use the multiplier attached to the different fourth/sixth-prefix
construction in \cite{Lichtman2025}.

\begin{proposition}[Derived upper sieve with third-prefix stopping]
\label{prop:pruned-upper}
Put $K_{\rm cap}=3549/500$. For every admissible product multiset
in Lemma~\ref{lem:primeproducts55}, of mass
$X_N=(J+o(1))N/\log N$, one has
\[
 \sum_{\substack{m<N\\N-m\in\PP}}w(m)
 \le (K_{\rm cap}J+o(1))\HN.
\]
In particular the five original switching constants satisfy
\begin{equation}\label{eq:pruned-switching}
 \Upsilon_i\le(K_{\rm cap}U_i/8+o(1))\HN
 \quad(i=7,8,9,10,11),
\end{equation}
where $i=11$ denotes the full fourth-prime range. Also, for every
fixed $1/2<b<1$,
\begin{equation}\label{eq:pruned-bad-semiprimes}
 B_b(N)\le(K_{\rm cap}\log(1/b)+o(1))\HN.
\end{equation}
Consequently a lower bound $D_{1,2}(N)\ge(c-o(1))\HN$ implies
\begin{equation}\label{eq:pruned-transfer}
 \liminf_{\substack{N\to\infty\\2\mid N}}D_b(N)/\HN
 \ge c-K_{\rm cap}\log(1/b).
\end{equation}
All these assertions retain the domains and labelled multiplicities
of the preceding product lemmas.
\end{proposition}
\begin{proof}
\paragraph{Parameters and a finite grid.}
Fix $\theta_0=89/160<\theta\le57/100$ and a sufficiently small
$e>0$, with $\theta-2e>\theta_0$. Write
\[
 \theta_R=\theta-2e,\qquad \theta_A=\theta-e,
 \qquad \vartheta=\theta-3e/2,
 \qquad D_0=N^{\theta_R},\quad z=N^{1/5}.
\]
The raw weight uses $D_0$; $\theta_A$ is the envelope threshold;
$\vartheta$ is the intermediate factorization parameter.
Choose $0<\zeta\ll e$, the distribution precision
$\epsilon_{\rm disp}\ll\zeta$, and a small presieve precision
$\eta\ll\epsilon_{\rm disp}$. Put
\[
 w=D_0^{\eta^2}=N^\omega,\quad \omega=\theta_R\eta^2,
 \qquad L=D_0^\eta=N^\ell,\quad \ell=\theta_R\eta.
\]
Choose an exponent mesh $\xi>0$ so small that
$\xi\omega^{-5}\to0$ with $\eta$; for example $\xi\le\omega^7$,
with a further decrease if needed for the fixed analytic margins.
Partition $[w,z)$ into disjoint intervals $J$ with exponent endpoints
$l(J),u(J)$ and $0<u(J)-l(J)\le\xi$.
Use $P_N(y)=\prod_{p<y,\ p\nmid N}p$ throughout; the rough window is $[w,z)$ and the small-prime window is $p<w$. Take the intervals left-closed and right-open. The last interval is
truncated at $z$. The number of boxes is fixed before $N\to\infty$.
Let $M=\lceil\theta_R/\omega\rceil+2$.

All primes in the following construction belong to the allowed sifting
set $\mathcal P_N=\{p:p\nmid N\}$. This restriction does not affect
the ordering or any upper estimate for their density sums.

\paragraph{The third-box rule, independent of rank.}
Define
\[
 v_A=2\theta_A-1,\quad
 k_A=(153/32-8\theta_A)/3,\quad u_A=v_A+k_A.
\]
An ordered triple of interval labels $(J_1,J_2,J_3)$ is called
\emph{bad} if its closed ordered exponent box meets
\[
 t_1\le v_A,\qquad t_3\ge k_A,\qquad t_2+t_3\ge u_A.
\]
Otherwise it is good. Crucially, this label decision is identical
in every later rank; it is not made separately for even and odd ranks.
For a full prime tuple whose third prime is less than $w$, declare no
extra stop. For small enough $\eta,\xi$, no bad box can contain
such a triple anyway, since $k_A-\xi>\omega$.

For every point of a good box satisfying the raw or slightly enlarged
cube conditions, the full envelope $\Phi$ in Lichtman 2023,
Proposition 6.1, is at least $\theta_A$. Here is the explicit check.
The first cube bounds $t_1<\mu=25/128$, with a fixed gap.
If $t_1\ge v_A$, the branch $v(t_1)=(1+t_1)/2$ suffices.
Otherwise, if $t_3\le k_A$, then $B(t_3)=153/256-3t_3/8\ge\theta_A$.
The inner minimum of $w(t_1,t_2,t_3)$ is $B(t_3)$, since
$(66-34t_3)/107>B(t_3)$ and $1-2t_2>B(t_3)$.
Its gate is $v(t_1)<\theta_A$, so this branch is available.
In the remaining case $t_1<v_A$, $t_3>k_A$, and
$s=t_2+t_3\le u_A$. Then
\[
 v_A<2k_A<s\le u_A<2v_A,
 \qquad 2t_1+s<4v_A<\theta_A.
\]
Each of the three affine branches defining $v(s)$ is at least
$\theta_A$: the first uses $s>v_A$, the third uses $s\le u_A$,
and the middle follows from
$(32-34(2v_A))/39>\theta_A$ for $\theta_A\le.57<4/7$.
Thus the pair branch is available. These checks also show stability
under an $O(\kappa)$ change of scale and residue exponent:
the last gate has a fixed gap, and the middle case is controlled by
the continuous maximum $\max\{v(t_1),B(t_3)\}$.
The surplus $\theta_A-\vartheta=e/2$ absorbs these perturbations.

\paragraph{Raw support and exact first-exit identities.}
Let $\mathcal R$ consist of the squarefree divisors of $P_N(z)$
$d=p_1\cdots p_r\in\mathcal D^+(D_0)$ for which $r<3$ or the
third prefix has a good box label (using the preceding convention
when $p_3<w$). Set $\lambda_{\mathcal R}(d)=\mu(d)\mathbf1_{d\in\mathcal R}$.
The support is prefix-closed. Every first excluded prefix has odd
length: either an ordinary odd cube fails or a bad third prefix is
encountered. Consequently, for every squarefree $Q$ dividing $P_N(z)$
and multiplicative $0\le h(p)\le1$,
\begin{equation}\label{eq:odd-first-exit-identity}
 \sum_{d\mid Q}\lambda_{\mathcal R}(d)h(d)
 =\prod_{p\mid Q}(1-h(p))+
 \sum_{\substack{P=p_1\cdots p_j\mid Q\text{ first exit}\\j\text{ odd}}}
 h(P)\prod_{\substack{p\mid Q\\p<p_j}}(1-h(p)).
\end{equation}
In particular the sum on the left is nonnegative, and with $h=1$
it majorizes the indicator of avoiding all primes of $Q$.

There is a useful exact density comparison. With
$Q=P_N(z)/P_N(w)$, multiply \eqref{eq:odd-first-exit-identity}
by $V_N(w)$ and put $h=g_N$. Then
\begin{align}
 V_N(w)\sum_{d\mid Q}\lambda_{\mathcal R}(d)g_N(d)
 &=V_N(z)+\sum_{\substack{\mathcal R\text{ first exits }P\\p_j\ge w}}
                  g_N(P)V_N(p_j)\notag\\
 &\le\sum_{d\mid P_N(z)}\lambda_{\mathcal R}(d)g_N(d)
        =:V_{\mathcal R}(D_0,z).\label{eq:rough-density-one-sided}
\end{align}
Rough first exits are exactly the full first exits with their last
prime above $w$. All omitted full exits are nonnegative. Thus no
unproved equality between the rough and full main terms is needed.
Comparing first exits with ordinary $\mathcal D^+(D_0)$ also gives
\begin{equation}\label{eq:raw-third-budget}
 V_{\mathcal R}(D_0,z)\le V^+(D_0,z)+
 \sum_{\substack{p_1>p_2>p_3>w\\p_1p_2p_3\in\mathcal D^+(D_0)\\
                         (J(p_1),J(p_2),J(p_3))\text{ bad}}}
                    g_N(p_1p_2p_3)V_N(p_3).
\end{equation}
New third exits add the displayed terms; removed later ordinary exits
can only decrease this difference.

\paragraph{Parity-dependent cube boxing.}
For a nonincreasing rank-$r$ vector of interval labels
$\boldsymbol J=(J_1,\ldots,J_r)$, keep the good-third-label condition
when $r\ge3$. For even $r$ retain the vector if
\[
 l(J_1)+\cdots+l(J_{m-1})+3l(J_m)\le\theta_R
                      \quad(m\le r,\ m\text{ odd});
\]
for odd $r$ retain it if the same inequalities hold with every $l$
replaced by $u$ and with a strict sign $<\theta_R$.
The raw cube convention is strict $<D_0$; this strict odd-box test
therefore includes no negative coefficient at an excluded raw endpoint.
The even outer-box test remains non-strict. The empty vector is retained.
Let $c(d)=(-1)^r$ when the distinct prime factors of a rough squarefree
$d$ have a retained vector, and $c(d)=0$ otherwise. Set $c=0$ at ranks
larger than $M$. No raw supported tuple has such a rank, since
$d\le D_0$ and every rough prime is at least $w$.
For every squarefree divisor $d$ of $Q$ one has
\[
 c(d)\ge\lambda_{\mathcal R}(d).
\]
For even rank this follows because every raw supported tuple passes
the lower-endpoint test; for odd rank the upper-endpoint test ensures
every tuple in the box is raw supported. The third-label condition
does not change under either test. Therefore, for every integer $n$,
\begin{equation}\label{eq:rough-box-pointwise}
 \sum_{\substack{d\mid n\\d\mid Q}}c(d)
 \ge\sum_{\substack{d\mid n\\d\mid Q}}\lambda_{\mathcal R}(d)
 \ge\mathbf1_{(n,Q)=1}.
\end{equation}
This comparison is coefficientwise and does not rely on the boxed
support being prefix-closed.

The density difference is supported only near ordinary cube boundaries.
If $c(d)\ne\lambda_{\mathcal R}(d)$, some odd $m\le M$ obeys
\[
 \left|t_1+\cdots+t_{m-1}+3t_m-\theta_R\right|
                         \le(M+2)\xi.
\]
Uniformly for an exponent interval of width $H$ above $\omega$,
$\sum g_N(p)\ll H/\omega+o_N(1)$ by the prime number theorem;
also $\prod_{w\le p<z}(1+g_N(p))\ll\omega^{-1}$.
Delete the tagged prime $p_m$, sum it over this interval with all
other primes fixed, and then sum over their squarefree subsets.
There are at most $M$ choices of the tagged index and an interval
width $O(M\xi)$. This proves
\[
 \sum_{d\mid Q}(c(d)-\lambda_{\mathcal R}(d))g_N(d)
 \ll M^2\xi\omega^{-2}+o_N(1).
\]
As $V_N(w)/V_N(z)\ll\omega^{-1}$ and $M\ll\omega^{-1}$,
the normalized cost is
\begin{equation}\label{eq:cube-box-density-error}
 V_N(w)\sum_{d\mid Q}(c(d)-\lambda_{\mathcal R}(d))g_N(d)
 \ll V_N(z)\xi\omega^{-5}+o_N(V_N(z)).
\end{equation}
All parameters are fixed in the $N\to\infty$ error. The bound is
deliberately coarse. No third-prefix boundary term is hidden here,
because the raw support was defined using the same box labels.

\paragraph{Small-prime weights and a pointwise complete upper bound.}
Choose fundamental-lemma weights $\beta^+,\beta^-$ of level $L$,
on divisors of $P_N(w)$, with $|\beta^\pm|\le1$, such that
\[
 B^-(n):=\sum_{b\mid n}\beta^-(b)
 \le\mathbf1_{(n,P_N(w))=1}\le
 B^+(n):=\sum_{b\mid n}\beta^+(b).
\]
Their density sums equal $V_N(w)(1+O(e^{-c/\eta}))$ for an absolute
$c>0$. The uniform one-sided dimension hypothesis suffices here;
deleting prime divisors of $N$ preserves it.
Explicitly set $\beta_r=\beta^+$ for even $r$ and
$\beta_r=\beta^-$ for odd $r$, always retaining the outer sign
$(-1)^r$. This convention is fixed here rather than inferred from
the parity wording in a statement of the source.
Write $B_r(n)=\sum_{b\mid n}\beta_r(b)$. Define the complete weight through
\[
 \rho(bd)=c(d)\beta_{\Omega(d)}(b),
 \qquad b\mid P_N(w),\quad d\mid Q,
\]
where the plus choice is made for even rank and the minus choice
for odd rank; zero coefficients are omitted. The smooth/rough
decomposition $bd$ is unique. In particular $|\rho(q)|\le1$, and
$\rho$ is supported on squarefree
$q\le N^{\theta_R+M\xi+\ell}$. Then
\begin{align*}
 \sum_{q\mid n}\rho(q)
 &=\sum_{\substack{d\mid n\\d\mid Q}}c(d)B_{\Omega(d)}(n)\\
 &\ge\mathbf1_{(n,P_N(w))=1}\sum_{\substack{d\mid n\\d\mid Q}}c(d)
 \ge\mathbf1_{(n,P_N(z))=1}.
\end{align*}
Thus the construction supplies a genuine pointwise upper bound,
even if a lower fundamental-lemma polynomial is negative.
It follows for every nonnegative sequence with local density $g_N$
that its sifted mass is at most
$X\sum_q\rho(q)g_N(q)+\sum_q\rho(q)r(q)$.

Since $|c(d)|\le1$, the fundamental-lemma density error is bounded by
\[
 e^{-c/\eta}V_N(w)\prod_{w\le p<z}(1+g_N(p))
 \ll V_N(z)e^{-c/\eta}\omega^{-2}.
\]
Combining \eqref{eq:rough-density-one-sided},
\eqref{eq:raw-third-budget}, and \eqref{eq:cube-box-density-error}
proves the explicit complete-weight main-term estimate
\begin{equation}\label{eq:complete-third-main}
 \sum_q\rho(q)g_N(q)
 \le V^+(D_0,z)+\sum_{\text{raw bad third prefixes }P}g_N(P)V_N(p_3)
 +O\!\left(V_N(z)\{\xi\omega^{-5}+e^{-c/\eta}\omega^{-2}+o_N(1)\}\right).
\end{equation}
This is the requested $V^+$ plus third-exit budget, now for the
actual complete weight rather than only an ordinary Rosser sum.

\paragraph{Actual atoms and factorization margins.}
For each retained vector let $\gamma(\boldsymbol J)$ be the product
of factorials of its repeated interval-label multiplicities.
The rough coefficient has the exact squarefree ordered-tuple expansion
\[
 c(d)=\sum_{r\le M}\sum_{\boldsymbol J\text{ retained}}
 \frac{(-1)^r}{\gamma(\boldsymbol J)}
 \sum_{\substack{d=p_1\cdots p_r\text{ squarefree}\\p_i\in J_i}}1.
\]
Each complete atom appends the corresponding $\beta^+$ or $\beta^-$.
Every actual rough vector belongs to
$\mathcal D^{\rm well}(N^{\theta_R+O(M\xi)})$:
the lower-endpoint odd cube conditions imply all lower-endpoint
square conditions, and rounding costs at most $O(M\xi)$.
Its total rough support is at most $N^{\theta_R+M\xi}$.
If $r\ge3$, its first three labels are good by definition;
the envelope check above therefore applies to every actual vector.
If $r\le2$, the first cube implies total support, including the
small factor, at most $N^{2\theta_R/3+2\xi+\ell}<N^{1/2-c_0}$.
These latter atoms use the classical distribution estimate.

For $r\ge3$, use a representative of each box and apply
Lichtman 2023, Proposition 6.1, at parameter $\vartheta$ with
factorization margin $\zeta$ and $A_0=Ux^{-\zeta}$,
$x^{2\zeta}\le U\le x^{1/3+\zeta/2}$.
Choose the hierarchy so that
\[
 \theta_R+O(M\xi)+2\zeta<\vartheta<\theta_A,
 \qquad \ell+O(M\xi)\ll\zeta,
 \qquad \kappa\ll\epsilon_{\rm disp}\ll\zeta\ll e.
\]
The source supplies a partition of the prime slots into three factors.
Keep that partition for the whole box. Rounding changes each support
by $N^{O(M\xi)}$; appending the small factor changes it by $N^\ell$.
The strict monomial inequalities retain their savings. In particular
the smallest original Type II saving after replacing $A_0$ by $U$
is $(1-2\sigma)\zeta=9\zeta/16$, and the displayed losses are
chosen smaller than it.
Use a fixed $\alpha=1+3\kappa$ on scales
$N^{1-\kappa}/O(1)\le x\ll N$. The good-box argument and its gate
checks lose only $O(\kappa)$, absorbed by $e/2$.

To obtain bounded factor coefficients, divide each grouped prime
convolution by the factorial of its number of slots and retain that
factorial in the scalar coefficient. Combining the third group with
$\beta^\pm$ preserves its bound by one, by the unique smooth/rough
decomposition. Restoring squarefreeness costs at most
$(N+N^{\theta_R+\ell+O(M\xi)})w^{-1}(\log N)^{C_\eta}$,
by the rough-square collision estimate already used in the manuscript.
There are finitely many atoms, and their total absolute scalar mass
depends only on the fixed precisions. Thus the existing individual
atom distribution proof applies, without an assumption that the
whole new weight is triply well-factorable.

\paragraph{Distribution and the exact product mass.}
Here the previously proved distribution argument is applied atom by
atom, with the new factorization just constructed. In particular,
for every separated smooth tuple weight $W$, it gives
\[
 \sum_{\substack{q\mid P_N(z)}}\rho(q)
 \left\{\sum_{\substack{\boldsymbol p\ \mathrm{prime}\\
               p_1\cdots p_k\equiv N\pmod q}}W(\boldsymbol p)
       -\frac{X_W^{(q)}}{\varphi(q)}\right\}
 \ll_A N(\log N)^{-A}.
\]
For a medium individual Heath--Brown variable this is the
three-factor estimate with its verified rough Siegel--Walfisz input.
Without such a variable, preservation of each original group scale
leaves at most two long variables, as in
Lemma~\ref{lem:primeproducts55}. All new modulus supports are below
$N^{3/5-c}$, so the one-long and two-long estimates, their arbitrary
subinterval versions, and the divisor treatment of Mellin tails and
prime powers are unchanged. Low-rank atoms were handled above.
The number of atoms is fixed, and each loss is absorbed by increasing $A$.

The exact coprime mass is still necessary. Since $|\rho(q)|\le1$,
\[
 \sum_{\substack{q\mid P_N(z)\\p\mid q}}
       \frac{|\rho(q)|}{\varphi(q)}
 \ll\frac{\log N}{p-1}.
\]
Every prime factor in a smooth outer product family is at least
$N^a/2$, and there are at most eleven. Summing this inequality
against that family's nonnegative weights gives
\[
 \sum_q\frac{|\rho(q)|}{\varphi(q)}
              (X_{\rm out}-X_{\rm out}^{(q)})
 \ll X_{\rm out}N^{-a}\log N
 \ll N^{1-a}(\log N)^C.
\]
Thus the local mean can be replaced by the total mass. The global
nonnegative outer approximation and its dyadic partition are those
of Lemma~\ref{lem:primeproducts55}; its excess mass is counted once.
The original labels, quotient factorizations, full fourth-prime range,
unit-quotient errors and Buchstab argument-one boundary strips are
retained. Repeated-prime and small-product exceptions still have a
fixed power saving. Finally $0<N-m\le z$ contributes only $O(z)$
for bounded multiplicity. These facts connect the complete pointwise
weight and its density estimate to the desired prime-complement count.

\paragraph{Limiting coarse coefficient.}
Every raw bad third prefix lies within $O(\xi)$ of the three-plane
bad set at $\theta_A$, and its third exponent is at least
$k_A-\xi>k_A/2$. Its cube conditions are at the smaller level
$\theta_R$. Prime-number-theorem integration is therefore on a compact
domain with a bounded density. Uniformly for $t$ in a fixed compact subinterval
of $(0,1)$, the local product obeys
\[
 V_N(N^t)=(2e^{-\gamma}+o(1))\frac{\CN}{t\log N}.
\]
Indeed the primes dividing $N$ above $N^t$ have total reciprocal
mass $O(N^{-t}\log N)$; below $N^t$ their full local factors are
retained in $\CN$. Also $g_N(p)=p^{-1}(1+O(p^{-1}))$ on the
allowed large primes. The prime number theorem therefore evaluates
the third-prefix sum uniformly; omitting primes dividing $N$ can
only reduce it, and their large-prime mass is negligible.
Letting the mesh tend to zero and then
$e\downarrow0$, the third-prefix contribution in \eqref{eq:complete-third-main},
divided by $V_N(z)$, has upper limit at most
\[
 \frac15\int_{\substack{0<t_3\le t_2\le t_1\\
 3t_1\le\theta,\ t_1+t_2+3t_3\le\theta\\
 t_1\le v,\ t_3\ge k,\ t_2+t_3\ge v+k}}
 \frac{dt_1dt_2dt_3}{t_1t_2t_3^2},
\]
where $v=2\theta-1$, $k=(153/32-8\theta)/3$.
Here the elementary geometric upper bound is as follows. Set
$\delta=\theta-89/160$ and, on the displayed integration domain,
\[
 x=v-t_1,\quad y=v-t_2,\quad z'=t_3-k.
\]
The order and three stopping inequalities give
$0\le x\le y\le z'$. The third cube gives
$3z'-x-y\le5\delta$, because $\theta-(2v+3k)=5\delta$.
Writing $A=x$, $B=y-x$, $C=z'-y$ maps the domain into
\[
 A,B,C\ge0,\qquad A+2B+3C\le5\delta,
\]
with unit absolute Jacobian. This simplex has volume
$(5\delta)^3/36$. Moreover $t_2t_3\ge kv$: for
$k\le t_3\le(v+k)/2$, use
$t_2t_3\ge t_3(v+k-t_3)\ge kv$; for larger $t_3$, use
$t_2\ge t_3$ and $(v+k)^2/4\ge kv$.
It follows that $t_1t_2t_3^2\ge k^2v^2$, and hence
\[
 I\le\frac{(5\delta)^3}{36k^2v^2}.
\]
The bound $2e^{-\gamma}<6/5$ follows from
$\gamma>H_{10}-\log11>\log(5/3)$: the exact logarithm
certificate gives $\log(55/3)<291/100<7381/2520=H_{10}$.

For the ordinary Rosser density sum use the standard one-sided
linear upper-sieve bound
\[
 V^+(D_0,z)\le V_N(z)\{F(5\theta_R)+o(1)\}.
\]
Its dimension constant is uniform in $N$: removing primes dividing
$N$ only decreases every inverse density-product ratio. This is
the one-sided hypothesis for the ordinary upper sieve noted in
\cite[Remark 2.9]{Lichtman2025}; no full-range two-sided condition
uniform in $N$ is being imposed. Here $2<5\theta_R<3$, so
$F(5\theta_R)=2e^\gamma/(5\theta_R)$.
Together with this ordinary upper-sieve density bound,
this is the coefficient budget
\[
 \frac4\theta+2e^{-\gamma}I
 \le\frac4\theta+\frac65\frac{(5(\theta-89/160))^3}{36k^2v^2}.
\]
First let $N\to\infty$ at fixed parameters. Then take
$\eta\downarrow0$ and $\xi\omega^{-5}\downarrow0$ within the
analytic margins, and finally $e\downarrow0$.


At $\theta=17/30$ the exact rational budget is
\[
 \delta=1/96,\quad k=119/1440,\quad v=2/15,\qquad
 \overline K=\frac{25730385}{3625216},\qquad
 K_{\rm cap}-\overline K=\frac{174771}{453152000}>0.
\]
Choose the fixed construction and distribution parameters sufficiently
small within this strict gap. The complete density bound and the
preceding distribution estimate then give
\[
 S\le(K_{\rm cap}+o(1))\frac{\CN X_{\rm out}}{\log N}.
\]
Removing the smooth outer approximation proves the first assertion.
For the five switched masses, $J_i=U_i/8$ with exactly their original
labels, giving \eqref{eq:pruned-switching}. For bad semiprimes,
Lemma~\ref{lem:mass} gives $J=\log(1/b)$, and their prime factors
exceed a fixed power greater than $N^{1/3}$ after discarding small
products. Thus the same construction proves
\eqref{eq:pruned-bad-semiprimes}. Subtracting these representations
from $D_{1,2}$, with at most $O(1)$ from the unit convention, proves
\eqref{eq:pruned-transfer}. No endpoint distribution theorem is used.
\end{proof}

% Final-form numerical consequence; see the preceding analytic proposition.
\subsection{Certified consequence of the stopped upper sieve}
\label{sec:pruned-upper-corollary}
Put $K_{\rm cap}=3549/500=7.098$. The original upper switching
constants are rescaled by $K_{\rm cap}/8=3549/4000$; all lower
terms, $U_3,U_4$ and the full-range eleven-term weight remain as
in the preceding two-prime-prefix calculation.

\begin{theorem}[Separate calculation with the stopped upper sieve]
\label{thm:pruned-chen}
Every sufficiently large even $N$ satisfies
\[
 D_{1,2}(N)\ge\frac{839}{500}\HN,
 \qquad D_{79/100}(N)\ge\frac1{250}\HN.
\]
This first route does not use the numerical Chen bound in \cite{Li2026}.
If, separately, the external lower bound
$D_{1,2}(N)\ge(1233/625-o(1))\HN$ is assumed, then
\[
 D_{379/500}(N)\ge\frac1{200}\HN.
\]
The last assertion retains that separate external-input dependency.
\end{theorem}

\begin{proof}
Proposition~\ref{prop:pruned-upper} supplies both required upper estimates. Write $U_i^{\rm cap}=(3549/4000)U_i$ for the five switching terms.
Reusing the certified $L_5^{(2)}$ integral and all the other existing
intervals gives the following outward bounds. Each terminating
decimal in this table denotes an exact rational number.
\begin{center}
\begin{tabular}{ll@{\qquad}ll}
\toprule
Lower term&Lower bound&Upper term&Upper bound\\
\midrule
$L_1(11/20)$&12.9014&$U_3$&19.3178\\
$L_2(11/20)$&6.3640&$U_4$&13.1277\\
$L_5^{(2)}$&2.3605&$U_7^{\rm cap}$&1.3640\\
$L_6$&2.2784&$U_8^{\rm cap}$&4.0784\\
&&$U_9^{\rm cap}$&2.9532\\
&&$U_{10}^{\rm cap}$&0.2265\\
&&$U_{11}^{{\rm cap},{\rm full}}$&0.5623\\
\bottomrule
\end{tabular}
\end{center}
The eleven-term functional consequently has lower bound
\[
 \frac14\left(3L_1+L_2-U_3-U_4+L_5^{(2)}+L_6
 -2U_7^{\rm cap}-U_8^{\rm cap}-U_9^{\rm cap}
 -U_{10}^{\rm cap}-U_{11}^{{\rm cap},{\rm full}}\right)
 \ge\frac{16783}{10000}
 =\frac{839}{500}+\frac3{10000}.
\]
The strict gap absorbs the Chen asymptotic error.
Subtracting \eqref{eq:pruned-bad-semiprimes} then gives
$\liminf D_b(N)/\HN\ge c-K_{\rm cap}\log(1/b)$ for either
lower constant used in the statement.

For a short rational check of the two transfer margins, use
\[
 \log\frac{1+u}{1-u}
 <2\left(u+\frac{u^3}{3}+\frac{u^5}{5(1-u^2)}\right)
 \qquad(0<u<1).
\]
With $u=21/179$ and $u=121/879$, respectively, this yields
\begin{align*}
 \log(100/79)&<\frac{106803868101}{453091781000},\\
 \log(500/379)&<\frac{1069768337033803}{3860975930715000}.
\end{align*}
Hence
\begin{align*}
 \frac{839}{500}-K_{\rm cap}\log(100/79)
 &>\frac1{250}
      +\frac{190892806551}{226545890500000},\\
 \frac{1233}{625}-K_{\rm cap}\log(500/379)
 &>\frac1{200}
      +\frac{735463365840551}{643495988452500000}.
\end{align*}
These positive rational gaps absorb the remaining asymptotic errors,
using the uniform estimates proved above.
\end{proof}

All numerical comparisons above are reproduced by
\texttt{certify\_pruned\_upper.py}; its JSON output
retains both the sharper $\overline K$ calculation and this
$K_{\rm cap}$ variant, with source hashes and exact interval endpoints.
This computation certifies numerical inequalities only. The analytic
construction is proved in the preceding proposition; the published
external Chen input retains the audit limitations already stated.
No bibliographic priority is established by either calculation.

% Research proof fragment, 7 October 2026. Not an assertion of priority.
% This file is intentionally separate from the main manuscript until reviewed.
\section{Retaining the old support in the bad third-prefix boxes}
\label{sec:retained-base-upper}

Set $\theta_0=89/160$ and $\theta=29/50$. Let $\Phi$ be the full
eight-branch envelope of Lichtman 2023, Proposition 6.1, (6.4), at
$\alpha=1$, $\sigma=7/32$, as in Section~\ref{sec:sharper-upper-support}.
Write $S=t_1+t_2+t_3$ and
\[
 E_\theta=\{0<t_3\le t_2\le t_1<1/5:\quad
 3t_1\le\theta,\quad t_1+t_2+3t_3\le\theta,\quad
 \Phi(t_1,t_2,t_3)<\theta\}.
\]
Extend the ordinary linear lower-sieve function by $f(s)=0$ for
$s\le2$. Define
\begin{equation}\label{eq:retained-base-functional}
 K(\theta)=\frac4\theta+2e^{-\gamma}
 \int_{E_\theta}
 \frac{f((\theta-S)/t_3)-f((\theta_0-S)/t_3)}
 {t_1t_2t_3^2}\,dt_1dt_2dt_3.
\end{equation}
No numerical assertion is part of this definition.

\begin{proposition}[Retained-base upper-sieve interface]
\label{prop:retained-base-upper}
For every $K_*>K(29/50)$, the conclusions of
Proposition~\ref{prop:pruned-upper} hold with $K_*$ in place of
$3549/500$. The product families, multiplicities, moving residue $N$,
and full fourth-prime range are exactly those in that proposition.
\end{proposition}

We first justify the asymptotic equality needed for the residual lower
sieves. A one-sided lower-sieve bound alone would not justify subtracting
two residual density sums.

\begin{lemma}[Uniform residual lower density]
\label{lem:uniform-residual-lower-density}
Let $R=N^\beta$, $y=N^t$, where $(\beta,t)$ ranges over a compact
subset of $\{\beta>0,t>0,\beta\ge2t\}$. Put
$g_N(p)=\mathbf1_{p>2,\,p\nmid N}/(p-1)$ and
$V_N(y)=\prod_{p<y}(1-g_N(p))$. With the ordinary Rosser support
\[
 \mathcal D^-(R)=\{d=q_1\cdots q_r:\ q_1>\cdots>q_r,
 q_1^2<R,\quad q_1\cdots q_{j-1}q_j^3<R
 \text{ for even }j\le r\},
\]
including the empty product, one has uniformly
\[
 \frac{\sum_{d\mid P_N(y),\ d\in\mathcal D^-(R)}\mu(d)g_N(d)}
 {V_N(y)}=f(\beta/t)+o(1).
\]
The strict or non-strict cube convention gives the same asymptotic.
\end{lemma}
\begin{proof}
Here is the precise classical input used, including convergence. For
dimension one and the cubic Rosser support, the normalized positive
first-exit integrals $a_j(s)$ obey
\[
 a_1(s)=(3-s)_+/s,\qquad
 a_j(s)=\frac1s\int_{\max(s,\,2+\epsilon_j)}^{2+j}
                  a_{j-1}(u-1)\,du\quad(j\ge2),
\]
where $\epsilon_j=0$ for even $j$ and $1$ for odd $j$, and an
empty integral is zero. The classical Rosser calculation gives
\begin{equation}\label{eq:classical-positive-exit-series}
 \sum_{j\ge2,\ j\ {
m even}}a_j(s)=1-f(s)\quad(s\ge2),
\end{equation}
with uniform convergence on compact $s$ intervals. This is a
convergence input, not a conclusion drawn merely from finite-depth
asymptotics. Explicit sources are
\href{https://www.suzuyu1729.com/notes/NagoyaSeminar2022/RosserIwaniec.pdf}
{Y. Suzuki, \emph{The Rosser--Iwaniec sieve}}, Sections 7--9
(the exit recurrence and (9.1)--(9.2)), Propositions 9.4 and 11.8
(compact convergence and its existence), and Section 12
(the optimal linear parameter $\beta=2$).
The normalization to $f$ is also given by (3.2.10) in
\href{https://mathweb.tifr.res.in/Documents/Publications/Lectures/tifr72.pdf}
{Y. Motohashi, \emph{Lectures on Sieve Methods and Prime Number Theory}},
TIFR 72 (1983); its Section 3.3, Lemma 14, supplies the convergence
estimate in the classical linear case. These inputs concern the
ordinary support displayed here, not a factorable modified weight.

Since $y\le\sqrt R$, the initial square condition never excludes
a prime below $y$. All first exclusions therefore have even length.
The finite first-exit identity gives exactly
\begin{equation}\label{eq:residual-lower-positive-exits}
 \frac{V_N^-(R,y)}{V_N(y)}
 =1-\sum_{j\ge2,\ j\ {
m even}} T_{j,N}(R,y),\qquad
 T_{j,N}=\sum_{\substack{P=q_1\cdots q_j\text{ first exit}\\q_1<y}}
 g_N(P)\frac{V_N(q_j)}{V_N(y)}\ge0.
\end{equation}
Fix a rank $j$. Every preceding prefix belongs to
$\mathcal D^{\rm well}(R)$: the even cube condition implies its
own square condition and, by ordering, the next square condition;
the first square is part of the lower support definition.
If $A_i=q_1\cdots q_i$, then $A_i^2<R A_{i-1}$, and induction gives
$A_i<R^{1-2^{-i}}$. Failure of the $j$th cube consequently forces
\[
 q_j^3\ge R/A_{j-1}>R^{2^{1-j}}.
\]
Thus all primes in this fixed-depth sum exceed $N^{c_j}$ for some
$c_j>0$ uniform in the compact parameter set.

For $g_0(p)=\mathbf1_{p>2}/(p-1)$, prime-number-theorem integration
on this compact exponent domain gives
$T_{j,0}(R,y)=a_j(\beta/t)+o(1)$ uniformly. The recurrence above
can alternatively be obtained by first summing the largest prime.
All boundary hyperplanes have zero limiting measure.
Deleting $p\mid N$ changes this finite-depth asymptotic by $o(1)$:
\[
 \sum_{p\mid N,\ p\ge N^{c_j}}\frac1{p-1}
 \ll N^{-c_j}\log N,
\]
and the other prime reciprocal sums and the product ratio are bounded
in terms of $j$ and the compact parameter set. More explicitly, each
surviving normalized summand satisfies
\[
 g_N(P)\frac{V_N(q_j)}{V_N(y)}
 =g_0(P)\frac{V_0(q_j)}{V_0(y)}
   \prod_{\substack{q_j\le p<y\\p\mid N}}(1-g_0(p))
 =\bigl(1+o(1)\bigr)
   g_0(P)\frac{V_0(q_j)}{V_0(y)}.
\]
Consequently $T_{j,N}=a_j(\beta/t)+o(1)$ for each fixed $j$.

There is no interchange of an uncontrolled infinite prime sum here.
The standard one-sided linear lower-density bound applies uniformly
to $g_N$, since deleting primes decreases every inverse local-product
ratio. It gives $V_N^-/V_N\ge f(\beta/t)-o(1)$.
In the opposite direction, discard the positive exits of rank $>J$
from \eqref{eq:residual-lower-positive-exits}. The fixed-depth
asymptotics then give
\[
 \frac{V_N^-}{V_N}\le
 1-\sum_{2\le j\le J,\ j\ {
m even}}a_j(\beta/t)+o_J(1).
\]
First let $N\to\infty$, then $J\to\infty$, using the compact-uniform
convergence in \eqref{eq:classical-positive-exit-series}. This proves
the lemma, including $\beta/t=2$ by continuity. In particular no
two-sided dimension hypothesis with a constant independent of $N$
has been silently assumed.
\end{proof}

\begin{proof}[Proof of Proposition~\ref{prop:retained-base-upper}]
\paragraph{Levels and the full envelope.}
Fix $e>0$ sufficiently small and use the two raw exponents
\[
 \nu_0=\theta_0-2e,\qquad \nu_1=\theta-2e,
 \qquad R_i=N^{\nu_i},\qquad \theta_A=\theta-e.
\]
Take the small-prime and rough windows, beta weights, exponent grid,
rank cutoff and hierarchy of Section~\ref{sec:pruned-upper-sieve},
with $D_0=R_1$. Thus $w=N^\omega$, the beta level is $N^\ell$,
the grid mesh is $\xi\ll\omega^5$, and every parameter is fixed
before $N\to\infty$.

Here is a buffered box rule which uses all eight branches and handles
their gates. Set $v(t)$ as in Section~\ref{sec:sharper-upper-support},
and write $W_1,W_2$ for the inner minima in
$w(t_1,t_2,t_3)$ and $w(t_1,t_3,t_2)$, respectively.
A closed ordered third-prefix box is good if throughout it at least
one fixed member of the following finite list supplies its indicated
test:
\[
 \begin{gathered}
 v(t_1),\ v(t_2),\ v(t_1+t_2),\ v(S)\ \ge\theta_A;\\
 \max\{v(t_1),W_i\}\ge\theta_A\quad(i=1,2);\\
 v(t_1+t_3)\ge\theta_A,\quad
 v(t_1+t_3)-(t_1+2t_2+t_3)\ge e;\\
 v(t_2+t_3)\ge\theta_A,\quad
 v(t_2+t_3)-(2t_1+t_2+t_3)\ge e.
 \end{gathered}
\]
Otherwise the box is bad. The label decision is made once and is
independent of the eventual rank and its parity. Third primes below
$w$ have the good label. For small enough $\omega,\xi$, this agrees
with the continuous $W_1$ test on all relevant cube-supported vectors.

The first raw cube gives $t_1\le\theta/3+O(\xi)<25/128$ with a
fixed gap. Hence $v(t_1)=(1+t_1)/2$, which is exactly the gate of
both $W_i$. The maximum with that gate removes the discontinuity
of each $w$ branch. The pair branches instead have the explicit
positive gate margin $e$; their discontinuities are not treated as
continuous. Scale changes and changing the residue exponent to
$\alpha=1+3\kappa$ perturb their affine pieces by $O(\kappa)$.
They therefore retain an envelope value above
$\theta-3e/2$ when $\kappa\ll e$. If desired, for each pair atom
one can choose its factorization parameter a little below its
$v$-value and above its gate, rather than use a common parameter
below the gate. Both choices have surplus over its actual raw
support. At a $W_i$ gate crossing the $v(t_1)$ branch supplies the
same surplus. This explicitly preserves the hypotheses in the proof
of Lichtman 2023, Proposition 6.1.

Bad boxes satisfy, up to $O(\xi)$,
\[
 t_1\le2\theta_A-1,\qquad
 t_3\ge(153/32-8\theta_A)/3.
\]
For the latter assertion, on the entire cube-supported domain
$1-2t_2>153/256$ and
$(66-34t_3)/107>153/256-3t_3/8$, so $W_1$ is its middle
affine piece. In particular, in a bad box
\[
 3t_1\le3(2\theta_A-1)+O(\xi)<\nu_0.
\]
This strict old first-cube margin is important below.

\paragraph{The actual mixed support and its upper property.}
Let $\mathcal R$ be the union of $\mathcal D^+(R_0)$ and
the subset of $\mathcal D^+(R_1)$ whose third label is good,
with ranks below three always kept at the new level. Equivalently,
good descendants use $R_1$ and bad descendants use $R_0$.
The support is prefix-closed. A first exclusion is odd: an ordinary
new odd cube fails before the third choice, or a bad third prefix
fails the old third cube, or its descendant first fails a later old
odd cube. An old first-cube failure cannot first appear at a bad
third prefix because of the strict margin just proved.
Thus the first-exit identity \eqref{eq:odd-first-exit-identity}
holds for $\lambda_{\mathcal R}(d)=\mu(d)\mathbf1_{d\in\mathcal R}$.
It proves the pointwise upper majorant and the rough/full one-sided
density comparison \eqref{eq:rough-density-one-sided} for this support.

The exact full density identity is
\begin{equation}\label{eq:retained-base-exact-density}
 V_{\mathcal R}-V_N^+(R_1,z)
 =\sum_{\substack{P=p_1p_2p_3\in\mathcal D^+(R_1)\\
        p_1>p_2>p_3,\ p_1<z\\\text{third label bad}}}
 g_N(P)\left\{V_N^-(R_1/P,p_3)
 -\mathbf1_{P\in\mathcal D^+(R_0)}V_N^-(R_0/P,p_3)\right\}.
\end{equation}
Indeed, the two supports differ only in the disjoint trees with a
bad third prefix. For either raw level $R_i$ at which that prefix is
present, its total descendant contribution is
$-g_N(P)V_N^-(R_i/P,p_3)$: three prefix primes reverse the
parity, so the later odd cubes become the residual even cubes.
The residual initial-square condition is automatic, since
$Pp_3^2<R_i$. This proves the identity with its stated sign.
The old prefix is absent precisely when its third cube fails,
equivalently when $\log(R_0/P)/\log p_3\le2$ (apart from harmless
endpoint conventions). The initial old cube has already been checked.

\paragraph{Boxing, beta weights, and distribution.}
Use exactly the parity boxing in Section~\ref{sec:pruned-upper-sieve},
but with the effective exponent $\nu_1$ on good third labels and
$\nu_0$ on bad labels. Ranks below three use $\nu_1$.
At even ranks impose all relevant lower-endpoint odd cubes with
non-strict inequality; at odd ranks impose the upper-endpoint
odd cubes with strict inequality. Hence $c(d)\ge\lambda_{\mathcal R}(d)$
coefficientwise. Any discrepancy lies within $O(M\xi)$ of an odd
cube boundary at one of the two levels. The same tagged-prime
argument bounds its normalized density cost by $O(\xi\omega^{-5})$.
There is no additional third-label discrepancy.

Append $\beta^+$ to even ranks and $\beta^-$ to odd ranks, with
the outer sign retained. The unique small/rough decomposition gives
$|\rho(q)|\le1$ and the complete pointwise upper bound, exactly as
in \eqref{eq:rough-box-pointwise} and the following beta construction.
Its additional normalized density error is
$O(e^{-c/\eta}\omega^{-2})$.

An actual good atom belongs to
$\mathcal D^{\rm well}(N^{\nu_1+O(M\xi)})$ and has the full
buffered envelope surplus proved above. An actual bad atom belongs
to $\mathcal D^{\rm well}(N^{\nu_0+O(M\xi)})$ and retains its
old first and third cube conditions. Apply
Lemma~\ref{lem:upper-support-89} and the margin argument of
Proposition~\ref{prop:upper-89} to the latter. Choose separate
intermediate factorization parameters for the old and new atoms,
then choose
\[
 \ell+O(M\xi),\ \kappa\ll\epsilon_{\rm disp}\ll\zeta\ll e
\]
small enough for both. Keep one representative slot partition on
each box; the rounded support and appended beta factor cost
$N^{O(M\xi)+\ell}$. This is smaller than the minimum Type II
margin $9\zeta/16$. The reciprocal repeated-label factorial occurs
once, and normalizing each grouped convolution by its slot factorial
gives the bounded coefficients used by the earlier atom argument.
Ranks at most two have complete support below $N^{1/2-c}$.
All other complete moduli have support below $N^{3/5-c}$, since
$\nu_1<.58$. Thus the prime-product distribution proof, its
medium-variable and at-most-two-long-variable alternatives,
squarefree collision estimate, divisor Mellin tails, and uniformity
for residue $N$ apply without alteration. The number of atoms is
fixed before the large-$N$ limit.

In particular the exact coprime mass is replaced by the total mass
with the same bound
$O(X_{\rm out}N^{-a}\log N)$, using $|\rho(q)|\le1$.
The nonnegative outer approximation is global, its excess is counted
once, and the labels and full fourth-prime range are retained.
No distribution estimate for arbitrary bounded outer arithmetic
weights, and no assertion that the whole new weight is triply
well-factorable, is used.

\paragraph{Evaluation and order of limits.}
On all bad boxes, $t_3$ is bounded below by a fixed positive constant:
at the limiting level it is at least $113/2400$. Thus the residual
parameters in \eqref{eq:retained-base-exact-density} range over
compact positive sets wherever a residual term is present. Apply
Lemma~\ref{lem:uniform-residual-lower-density} to both residual sums.
Across the old third-cube boundary its omitted contribution agrees
continuously with $f(s)=0$ at $s=2$.
Uniform Mertens estimates for $V_N(N^{t_3})$ and prime-number-theorem
integration then evaluate the correction in
\eqref{eq:retained-base-exact-density}. Missing prime divisors of
$N$ in these three fixed positive-power prime slots have total
reciprocal mass $o(1)$ uniformly.

The bad sets of the buffered box rules converge in measure to
$E_\theta$. All branches are piecewise affine, and every additional
level or gate equality is contained in a finite union of proper
hyperplanes. Away from these planes, a strict good branch is
eventually certified on its whole box. The $W_i$ rescue and the
explicit pair gate margins above justify this statement despite
the jumps in $\psi$. The correction integrand is bounded on the
compact bad domain, so the extra boundary strips have vanishing cost.

Finally, the ordinary upper density is at most
$V_N(z)(F(5\nu_1)+o(1))$, and $5\nu_1<3$.
Combining all the estimates gives the coefficient
\eqref{eq:retained-base-functional}. First take $N\to\infty$ at
fixed construction parameters; then take the beta and mesh errors
to zero within the analytic margins; finally let $e\downarrow0$.
Given the strict slack $K_*>K(\theta)$, one can choose fixed
parameters whose complete coefficient is below $K_*$.
Applying the resulting upper bound to the five switched masses
$J_i=U_i/8$ and to the bad-semiprime mass $\log(1/b)$ proves
the proposition. This establishes an analytic interface for a
separately certified value of the integral; it makes no claim of
priority or of comparison with every previously published method.
\end{proof}

% Numerical consequence, conditional on the retained-base analytic proposition.
\subsection{Numerical consequence of the retained-base upper sieve}
\label{sec:retained-base-corollary}
The integer interval calculation in
\texttt{certify\_retained\_base\_candidate.py} gives
\[
 K_{\rm integral}\le
 \frac{121917146575524836424592375}{17529424384412123033239552}
 <\frac{1391}{200}=6.955.
\]
It encloses the integral in Proposition~\ref{prop:retained-base-upper};
it does not certify the analytic proposition. The source code received
a separate numerical audit, including an algebraic verification of
the six affine good regions against all eight envelope branches.
The final exact run used 2,856,205 adaptive subdivisions.

\begin{theorem}[Consequence of the retained-base construction]
\label{thm:retained-base-chen}
Under the analytic conclusion of Proposition~\ref{prop:retained-base-upper},
every sufficiently large even $N$ satisfies
\[
 D_{1,2}(N)\ge\frac{1731}{1000}\HN,
 \qquad D_{39/50}(N)\ge\frac1{500}\HN.
\]
This calculation does not use the numerical Chen bound in \cite{Li2026}.
If the external input
$D_{1,2}(N)\ge(1233/625-o(1))\HN$ is used separately, then
\[
 D_{377/500}(N)\ge\frac1{200}\HN.
\]
\end{theorem}

\begin{proof}
Put $K_{\rm cap}=1391/200$ and multiply the five original switching
upper constants by $K_{\rm cap}/8=1391/1600$.
The remaining terms are the same as in
Theorem~\ref{thm:pruned-chen}. Outward rational rounding gives
\begin{center}
\begin{tabular}{ll@{\qquad}ll}
\toprule
Lower term&Lower bound&Upper term&Upper bound\\
\midrule
$L_1$&12.9014&$U_3$&19.3178\\
$L_2$&6.3640&$U_4$&13.1277\\
$L_5^{(2)}$&2.3605&$U_7^{\rm cap}$&1.3365\\
$L_6$&2.2784&$U_8^{\rm cap}$&3.9962\\
&&$U_9^{\rm cap}$&2.8937\\
&&$U_{10}^{\rm cap}$&0.2219\\
&&$U_{11}^{{\rm cap},{\rm full}}$&0.5510\\
\bottomrule
\end{tabular}
\end{center}
The eleven-term functional is at least
\[
 \frac{34629}{20000}
 =\frac{1731}{1000}+\frac9{20000}.
\]
Thus its positive slack absorbs the asymptotic error in the Chen count.
The same universal upper bound for the removed semiprimes gives
\[
 \liminf_{\substack{N\to\infty\\2\mid N}}
 \frac{D_b(N)}{\HN}\ge c-K_{\rm cap}\log(1/b).
\]
For a rational margin check, apply
\[
 \log\frac{1+u}{1-u}
 <2\left(u+\frac{u^3}{3}+\frac{u^5}{5(1-u^2)}\right)
\]
at $u=11/89$ and $u=123/877$. This gives
\begin{align*}
 \log(50/39)&<\frac{3415573051}{13746895500},\\
 \log(500/377)&<\frac{359019451176843}{1271481760705000},
\end{align*}
and therefore
\begin{align*}
 \frac{1731}{1000}-K_{\rm cap}\log(50/39)
 &>\frac1{500}+\frac{201103843}{211490700000},\\
 \frac{1233}{625}-K_{\rm cap}\log(500/377)
 &>\frac1{200}+\frac{77561935082399}{19561257857000000}.
\end{align*}
Both gaps are strictly positive and absorb the remaining asymptotic errors.
\end{proof}

The exact interval recombination and dependency hashes are recorded by
\texttt{certify\_retained\_base\_transfer.py}. These are numerical
certificates, with all analytic dependencies stated separately.
Neither consequence establishes priority: the published $G_1$ method
in \cite[(27)--(29)]{Li2026} gives a potentially stronger comparison
after its shifted local-product interface is supplied.

% Fragment for inclusion in the main manuscript; no standalone preamble.
% This repairs the C<1 capacity case in the use of L23, Section 6.
\subsection{A two-factor repair for a capacity below one}
\label{sec:two-factor-greedy-repair}

In this subsection $\vartheta$ denotes the distribution exponent and
$\sigma$ the spectral parameter (called $\theta$ in L23).
The argument of L23, Lemma~6.3 and Proposition~6.5, uses a capacity
$C=D/(A b_0)$ after selecting a factor $b_0\le x^u$.
Lemma~6.3 itself assumes an existing factorization with $c\le C$;
its algebra is valid.  The subsequent greedy construction requires
care when $C<1$, because an integer factor $c\ge1$ cannot satisfy
$c\le C$.  Replacing this capacity by $\max(1,C)$ does not by itself
justify the three inequalities in L23~(6.3).  The following argument
handles that case by reallocating every prime factor of $d$.
It is the analogue, with the additional moving-residue constraints,
of the two-factor repair in Pascadi, arXiv:2505.00653v2,
Lemma~5.6 and the proof of Lemma~5.7.

\begin{lemma}[Two-factor repair]\label{lem:two-factor-greedy-repair}
Let $x>1$, $0<\delta<10^{-5}$, $1/2\le\vartheta\le5/8$,
$0\le\sigma\le7/32$, and $\alpha\ge0$.  Put
\[
 D=x^{\vartheta-2\delta},\qquad
 x^\delta\le A\le x^{1/3-\delta/2},\qquad B=D/A,
\]
and
\[
 u=\min\left\{
 \frac{2-(2+2\sigma)\vartheta}{3-4\sigma},
 \frac{2-2\vartheta-\alpha\sigma}{3}\right\}.
\]
Suppose $B\le x^u$ and
\[
 d=p_1\cdots p_r\in\mathcal D^{\rm well}(D),\qquad
 p_1\ge\cdots\ge p_r,
\]
where the defining condition is
$p_1\cdots p_{j-1}p_j^2\le D$ for every $j\le r$.
Then there is a factorization $d=a'b'c'$ with
$a'\le A$, $b'\le B$, $c'=1$, and
\begin{align}
 A^2b'c'^2&\le x^{1-3\delta},\label{eq:two-factor-repair-first}\\
 x^{\alpha\sigma}Aa'b'^5c'^2&\le x^{2-4\delta}
                       \le x^{2-3\delta},\label{eq:two-factor-repair-second}\\
 (Ac'/b')^{2\sigma}Aa'b'^5c'^2
   &\le x^{2-(4+4\sigma)\delta}
                       \le x^{2-2\delta}.
                       \label{eq:two-factor-repair-third}
\end{align}
These are the three constraints in L23~(6.3), with slack.
\end{lemma}

\begin{proof}
We have $A\ge1$ and
\[
 B\ge x^{\vartheta-1/3-3\delta/2}
       \ge x^{1/6-3\delta/2}>1.
\]
Start with $a'=b'=1$.  Immediately before assigning $p_j$, the
assigned product equals $a'b'=p_1\cdots p_{j-1}$.  The support
condition gives
\[
 p_j^2\le\frac{D}{a'b'}=\frac{A}{a'}\frac{B}{b'}.
\]
Consequently at least one of $p_j\le A/a'$ and $p_j\le B/b'$
holds.  Assign $p_j$ to that factor.  Iterating over all primes
preserves $a'\le A$, $b'\le B$ and yields $a'b'=d$.
This construction also covers $d=1$, by taking both factors equal
to one.  Set $c'=1$.

First,
\[
 A^2b'\le A^2B=AD
 \le x^{\vartheta+1/3-5\delta/2}
 \le x^{23/24-5\delta/2}
 \le x^{1-3\delta},
\]
where the last inequality follows from $\delta<10^{-5}<1/12$.
Next, using the second upper bound in the definition of $u$,
\[
 x^{\alpha\sigma}Aa'b'^5
 \le x^{\alpha\sigma}A^2B^5
 =x^{\alpha\sigma}D^2B^3
 \le x^{\alpha\sigma+2\vartheta-4\delta+3u}
 \le x^{2-4\delta}.
\]
Finally, $5-2\sigma>0$ and $3-4\sigma>0$, so the first upper
bound in the definition of $u$ gives
\begin{align*}
 (A/b')^{2\sigma}Aa'b'^5
 &\le A^{2+2\sigma}b'^{5-2\sigma}
 \le A^{2+2\sigma}B^{5-2\sigma}\\
 &=D^{2+2\sigma}B^{3-4\sigma}\\
 &\le x^{(2+2\sigma)\vartheta-(4+4\sigma)\delta
                         +(3-4\sigma)u}
 \le x^{2-(4+4\sigma)\delta}.
\end{align*}
\end{proof}

To apply the repair in the capacity case above, retain only the
information $b_0\le x^u$ and $C=D/(A b_0)<1$.  These imply
$B=D/A<b_0\le x^u$, which is the hypothesis of
Lemma~\ref{lem:two-factor-greedy-repair}.  The resulting $b'$ need
not equal $b_0$ or remain in $[x^v,x^u]$; the required conclusion is
the existence of some factorization satisfying L23~(6.3), already
proved directly.  If $C\ge1$, the original greedy construction
has two legitimate capacities and the usual Lemma~6.3 applies.
Thus this repair covers the missing case without reducing the
admissible distribution exponent, including the earlier level
$\vartheta=11/20$.

\paragraph{Correct initialization and capacities in the three-factor case.}
There is a separate issue in the printed proof of L23, Lemma~6.4:
the initialization must use $\{1,p_1,p_2\}$, and must append $p_3$
only in the subsequent greedy process.  Moreover, the printed
choice $D_2=D^2/x=x^{v-4\delta}$ does not equal $x^v$ and does not
give the claimed lower bound $b>x^v$ at the first failed step.
The capacities below avoid that discrepancy and give exact
versions of all three estimates.  We use the strict margins
already imposed on compact prefix boxes; no assertion at a gate
boundary without such margins is needed.

\begin{lemma}[Buffered three-factor construction]
\label{lem:buffered-three-factor-greedy-repair}
Use the notation and ranges of
Lemma~\ref{lem:two-factor-greedy-repair}, without assuming
$B\le x^u$, and let $v=2\vartheta-1$.
Let $d=p_1\cdots p_r\in\mathcal D^{\rm well}(D)$ with
$p_1\ge\cdots\ge p_r$.
Suppose $r\ge3$, $p_3\le x^{u-v}$, and, for some
$\{i,j\}=\{1,2\}$,
\[
 p_i\le x^v,\qquad p_j^2\le x^{1-\vartheta-2\delta}.
\]
Then $d$ admits a factorization satisfying L23~(6.3).
\end{lemma}

\begin{proof}
Set
\[
 D_1=A,\qquad D_2=x^v,\qquad D_3=\frac{D}{Ax^v}.
\]
Their product is exactly $D$.  All three capacities are at least
one: for the last one,
\[
 D_3\ge x^{2/3-\vartheta-3\delta/2}
       \ge x^{1/24-3\delta/2}>1.
\]
Place $p_i$ in the second factor.  Since
$p_j^2\le D_1D_3$, place $p_j$ in either the first or third
factor that has sufficient capacity, leaving the other equal to
one.  The assigned factors are thus a permutation of
$\{1,p_1,p_2\}$ with the stated placement of $p_i$.
Now append $p_3,p_4,\ldots$ whenever any factor has capacity.

If every prime is appended, write the resulting factors as
$a,b,c$ with $a\le D_1$, $b\le D_2$, $c\le D_3$.
Since $p_3>1$ and $p_3\le x^{u-v}$, we have $u\ge v$.
The two inequalities defining $u\ge v$ imply, respectively,
\[
 (8-6\sigma)\vartheta+4\sigma\le5,
 \qquad \alpha\sigma+8\vartheta\le5.
\]
Consequently
\begin{align*}
 A^2bc^2
 &\le A^2D_2D_3^2
   =x^{1-4\delta},\\
 x^{\alpha\sigma}Aab^5c^2
 &\le x^{\alpha\sigma}A^2D_2^5D_3^2
   =x^{\alpha\sigma+8\vartheta-3-4\delta}
   \le x^{2-4\delta},\\
 (Ac/b)^{2\sigma}Aab^5c^2
 &\le A^{2+2\sigma}D_2^{5-2\sigma}D_3^{2+2\sigma}\\
 &=x^{(8-6\sigma)\vartheta+4\sigma-3-(4+4\sigma)\delta}
   \le x^{2-(4+4\sigma)\delta}.
\end{align*}
These imply the three required inequalities.

Otherwise let $p_k$, with $k\ge3$, be the first prime that cannot
be appended, and denote the current factors by $a_0,b_0,c_0$.
Thus $a_0b_0c_0=p_1\cdots p_{k-1}$ and
$a_0p_k>A$, $b_0p_k>x^v$, $c_0p_k>D_3$.
Put $b=b_0p_k$.  Then
\[
 x^v<b\le x^v p_3\le x^u.
\]
The support condition and $a_0p_k>A$ give
\[
 A b c_0<(a_0p_k)(b_0p_k)c_0\le D.
\]
It follows that $C=D/(Ab)>c_0\ge1$, so both remaining
capacities $A,C$ are legitimate.  For every remaining prime
$p_l$ the support condition gives
\[
 p_l^2\le
 \frac{D}{a_0bc_0p_{k+1}\cdots p_{l-1}}
 =\frac{AC}{a_0c_0p_{k+1}\cdots p_{l-1}}.
\]
The two-factor greedy argument, starting from $a_0,c_0$, therefore
assigns all remaining primes and produces $d=abc$ with
$a\le A$, $c\le C=D/(Ab)$, and $x^v<b\le x^u$.
Directly,
\begin{align*}
 A^2bc^2&\le D^2/b\le x^{1-4\delta},\\
 x^{\alpha\sigma}Aab^5c^2
 &\le x^{\alpha\sigma}D^2b^3\le x^{2-4\delta},\\
 (Ac/b)^{2\sigma}Aab^5c^2
 &\le D^{2+2\sigma}b^{3-4\sigma}
   \le x^{2-(4+4\sigma)\delta}.
\end{align*}
This also covers a failed step at $k=r$, when no primes remain
to be reassigned.
\end{proof}

For a strict $W$ branch, its inequalities state that the
designated prefix exponent is below $v$, that the other prefix
prime has square below $x^{1-\vartheta}$, and that
$t_3<u-v$.  On a fixed compact box in this strict region, choose
$\delta$ small enough that the second inequality still holds
with exponent $1-\vartheta-2\delta$; shrink the other margins to
cover the box widths and the fixed perturbation of $\alpha$.
Lemma~\ref{lem:buffered-three-factor-greedy-repair} then applies.
As these positive buffers tend to zero, the limiting eight-branch
envelope $\Phi$ is unchanged.  This is a buffered use of the
source envelope, not a claim that its printed boundary argument
or its printed initialization was exact.

% Derived factorization fragment, 2026-10-08; separately saved for inclusion.
% Requires the existing two-factor-greedy-repair.tex.
% This is a branch-specific extension, not an extension of the full raw Phi.

\subsection{Selecting a subproduct without a largest-prime restriction}
\label{sec:selected-subproduct-factorization}

Write $\sigma$ for the spectral parameter and $\vartheta$ for the
distribution exponent.  Set
\[
 \frac12\le\vartheta<\frac58,\quad 0\le\sigma\le\frac7{32},\quad
 \alpha\ge0,\quad 0<\delta<10^{-5},\quad D=x^{\vartheta-2\delta},
 \qquad x^\delta\le A\le x^{1/3-\delta/2},
\]
and define
\[
 v_{\rm crit}=2\vartheta-1,\qquad
 u=\min\left\{
 \frac{2-(2+2\sigma)\vartheta}{3-4\sigma},
 \frac{2-2\vartheta-\alpha\sigma}{3}\right\}.
\]
The conclusion needed below is the existence of $d=abc$, with $a\le A$,
satisfying the three inequalities
\begin{equation}\label{eq:selected-subproduct-target}
 A^2bc^2\le x^{1-3\delta},\quad
 x^{\alpha\sigma}Aab^5c^2\le x^{2-3\delta},\quad
 (Ac/b)^{2\sigma}Aab^5c^2\le x^{2-2\delta}.
\end{equation}
These are L23 (6.3).  The
\href{https://arxiv.org/html/2309.08522v1}
{algebraic criterion in L23, Lemma 6.3}, does not assume
$p_1\le x^u$.  That extra assumption occurs in the subsequent
Proposition 6.5.  We prove a separate construction and do not quote
that proposition with its hypothesis omitted.

\begin{lemma}[A selected-subproduct criterion]
\label{lem:selected-subproduct-factorization}
Let $d=p_1\cdots p_r\in\mathcal D^{\rm well}(D)$ with the primes
in decreasing order.  Select a product $b_0$ of some of these prime
slots, and list the remaining slots in their inherited order as
$q_1\ge\cdots\ge q_m$.  Suppose
\begin{equation}\label{eq:selected-subproduct-assumptions}
 x^{v_{\rm crit}}\le b_0\le x^u,
 \qquad b_0q_1\cdots q_{j-1}q_j^2\le D\quad(1\le j\le m).
\end{equation}
Then \eqref{eq:selected-subproduct-target} holds for some factorization
$d=abc$.  No condition $p_1\le x^u$ or $t_1\le\mu$ is required.
\end{lemma}

\begin{proof}
First suppose $C=D/(A b_0)\ge1$.  Start with $a=c=1$ and keep
$b=b_0$ fixed.  Just before appending $q_j$, the assigned product is
$ac=q_1\cdots q_{j-1}$, and the second condition in
\eqref{eq:selected-subproduct-assumptions} gives
\[
 q_j^2\le\frac{D}{b_0ac}=\frac A a\frac C c.
\]
At least one of $a q_j\le A$ and $c q_j\le C$ holds.  Choose such
a factor and continue.  This assigns every remaining prime and gives
$d=ab_0c$, $a\le A$, $c\le D/(A b_0)$.
Directly, since $5-2\sigma>0$ and $3-4\sigma>0$,
\begin{align*}
 A^2b_0c^2&\le D^2/b_0\le x^{1-4\delta},\\
 x^{\alpha\sigma}Aa b_0^5c^2
   &\le x^{\alpha\sigma}D^2b_0^3\le x^{2-4\delta},\\
 (Ac/b_0)^{2\sigma}Aa b_0^5c^2
   &\le D^{2+2\sigma}b_0^{3-4\sigma}
    \le x^{2-(4+4\sigma)\delta}.
\end{align*}
These imply \eqref{eq:selected-subproduct-target}.

If $C<1$, put $B=D/A$.  Then $B<b_0\le x^u$.
Lemma~\ref{lem:two-factor-greedy-repair}, applied to the full original
vector, gives $d=a'b'$, $a'\le A$, $b'\le B$, $c'=1$, and all three
target inequalities.  This case deliberately reallocates the primes
of $b_0$ as well.  The final $b'$ need not lie in the critical interval;
the repaired two-factor argument proves the target directly.
Its hypotheses are the original $\mathcal D^{\rm well}(D)$ condition
and $B\le x^u$, both already checked.  In particular it does not
introduce a bound $p_1\le x^u$.
\end{proof}

\paragraph{The pair branch in question.}
For $b_0=p_2p_3$, the remaining list is $p_1,p_4,p_5,\ldots$.
The initial condition in \eqref{eq:selected-subproduct-assumptions}
is precisely
\[
 p_1^2p_2p_3\le D.
\]
Before every later remaining prime $p_j$, $j\ge4$, the product of
$b_0$ and the already assigned remaining primes is
$p_1p_2\cdots p_{j-1}$.  The required square condition is therefore
exactly the original support condition, with no additional factor.
Thus the proposed two-capacity start and every subsequent step are valid.

More generally, the following selected products have exactly the
listed additional requirements; a product of the first $k$ slots may
be used whenever it is defined.
\begin{center}
\begin{tabular}{c|c}
Selected product $b_0$ & Additional square condition \\
\hline
$p_1\cdots p_k$ & none beyond $d\in\mathcal D^{\rm well}(D)$ \\
$p_2$ & $p_1^2p_2\le D$ \\
$p_1p_3$ & $p_1p_2^2p_3\le D$ \\
$p_2p_3$ & $p_1^2p_2p_3\le D$
\end{tabular}
\end{center}
For $p_2$ the inherited conditions start at $p_3$; for the two
nonconsecutive pairs they start at $p_4$.  This table does not assert
that an arbitrary nonconsecutive subproduct needs only one extra test.
Such a product must satisfy every remaining-square condition in
\eqref{eq:selected-subproduct-assumptions}.

\paragraph{Inverting the critical interval.}
For $p_i=x^{t_i}$, write $s=\log b_0/\log x$ and set
\[
 V_{\alpha,\sigma}(s)=\min\left\{
 \frac{1+s}{2},\quad
 \frac{2-(3-4\sigma)s}{2+2\sigma},\quad
 1-\frac{\alpha\sigma+3s}{2}\right\}.
\]
Direct rearrangement, using $3-4\sigma>0$, gives the exact equivalence
\[
 v_{\rm crit}\le s\le u
 \quad\Longleftrightarrow\quad
 \vartheta\le V_{\alpha,\sigma}(s).
\]
At $\alpha=1$, $\sigma=7/32$, this becomes the same single-variable
function already used in the manuscript:
\[
 V_{1,7/32}(s)=v(s)
 =\min\left\{\frac{1+s}{2},\frac{32-34s}{39},
                       \frac{57}{64}-\frac{3s}{2}\right\}
 =\min\left\{\frac{1+s}{2},\frac{57}{64}-\frac{3s}{2}\right\}.
\]
In particular the paired branch is valid, with fixed margins, if
\begin{equation}\label{eq:selected-pair-gate}
 2t_1+t_2+t_3+2\delta<\vartheta<v(t_2+t_3).
\end{equation}
This proves its strict-gate version without $t_1\le25/128$.
The other nonconsecutive pair has gate
$t_1+2t_2+t_3+2\delta<\vartheta$; the $p_2$ branch has gate
$2t_1+t_2+2\delta<\vartheta$.
The prefix-product branches require only the support hypothesis and
the corresponding critical-interval test.

The distinction between an endpoint envelope and a fixed level matters.
For a branch with endpoint $V$ and gate $G$, its construction requires
$G+2\delta<\vartheta<V$, not merely $\vartheta< V$.
When $G<V$ one may approach the endpoint from within that interval.
An atom of a smaller level can also be embedded in a slightly larger
support level in this interval before applying the construction.
No positive margin is asserted when $G=V$; such gate boundaries must
be trimmed before an integral limiting argument.

\paragraph{Uniformity on actual sieve boxes.}
In this paragraph take $\sigma=7/32$.
These assertions are stable under the boxes and small factors used
by the existing prefix interface.  To see this concretely, fix
$\varepsilon>0$ such that $\vartheta\le v(s)-\varepsilon$ and every
required extra square exponent is at most $\vartheta-\varepsilon$.
At $\alpha=1$ the selected logarithmic product lies above
$v_{\rm crit}$ by at least $2\varepsilon$; its distance below the two
upper endpoints defining $u$ is at least $39\varepsilon/34$ and
$2\varepsilon/3$, respectively.  If $0\le\alpha-1\le\varepsilon$,
the latter gap decreases by at most $7\varepsilon/96$, and so remains
at least $19\varepsilon/32$.

Choose box widths, the support inflation, and $\delta$ sufficiently
small compared with these fixed margins.  Carry out the greedy
construction on the ordered upper endpoints of the prime intervals.
The support and extra square conditions hold for that entire endpoint
vector.  The selected product has an upper endpoint $b_+$ and lower
endpoint $b_-$ satisfying $x^{v_{\rm crit}}<b_-\le b_+<x^u$.
If $D/(A b_+)\ge1$, the slot assignment obtained at the upper
endpoints remains valid for every prime tuple in the box: in
particular the resulting actual $c$ is at most $D/(A b_+)$ and hence
at most $D/(A b_0)$.  If $D/(A b_+)<1$, apply the full-vector
two-factor repair to those upper endpoints.  Its proof uses only
their ordered positive sizes, so the resulting slot assignment again
works throughout the box.  This furnishes a fixed partition of prime
slots for every tuple in the box.

The original signs, multiplicity factorials and bounded small-prime
factor are consequently retained by the existing box convolution
interface.  Reattach a factor $e\le x^{O(\eta)}$ to $c$, preserving $a\le A$.
The three target left sides are multiplied by $e^2$, $e^2$, and
$e^{2+2\sigma}$, respectively; choose $\eta\ll\delta$ to absorb them
in the slack.  In the two-factor case the first inequality has the
additional slack $1/24-\delta/2$, as follows from
$A^2b'\le AD\le x^{\vartheta+1/3-5\delta/2}$;
the other two retain their displayed positive multiples of $\delta$.  For moving residues, the local-scale perturbation
$\alpha=1+O(\kappa)$ is absorbed by choosing $\kappa$ below the same
fixed margins.  No new Type II or long-variable distribution theorem
is inferred from this algebra: their use remains the existing
interface, once these factorization hypotheses have been supplied.

\paragraph{The previous Li (34) test point is repaired.}
Take the first three decreasing exponents
\[
 (t_1,t_2,t_3)=\left(\frac{101}{500},\frac{23}{250},\frac9{100}\right)
                  =(0.202,0.092,0.090).
\]
Then $s=t_2+t_3=91/500$ and
\[
 v(s)=\frac{591}{1000},\qquad
 2t_1+s=\frac{293}{500}=\frac{586}{1000}.
\]
For example choose $\vartheta=59/100$, and any sufficiently small
fixed $\delta<10^{-5}$.  The selected product lies in the critical
interval, whose endpoints at $\alpha=1$ are
\[
 v_{\rm crit}=\frac9{50},\qquad
 u=\min\left\{\frac{899}{3400},\frac{481}{2400}\right\}
   =\frac{481}{2400}.
\]
Here $s-v_{\rm crit}=1/500$ and $u-s=221/12000$.
Indeed $t_1>u$, so the unmodified hypothesis of Proposition 6.5
would not have covered even this buffered example.
The new criterion does cover it.  Moreover
\[
 v(t_1)=\frac{4701}{8000}<\frac{59}{100}<\frac{591}{1000}.
\]
The fourth specified exponent $11/125$ in the earlier Li (34)
test has square-prefix exponents $101/250,193/500,237/500,14/25$,
all strictly below $59/100$.  In that ordinary upper-sieve option
the residual sifting cutoff is the smallest specified prime, so the
specified first three slots persist in the actual concatenated atom.
The first four square conditions and the residual well support
therefore imply the full support condition at a buffered total level.
This repairs the largest-prime factorization obstruction at that
point and on a neighborhood.  The former audit should no longer
use that point as an example requiring fallback to $v(t_1)$.
It does not by itself verify the integral or all of Li (34).

\paragraph{The Li (25) test is not repaired by these branches.}
For $t_1=11/50$, $t_2=3/20$ and a third exponent
$0<t_3\le a=100/1149$, the raw improvement $v(t_2)=23/40$
has the new gate
\[
 2t_1+t_2=\frac{59}{100}>\frac{23}{40}.
\]
It therefore fails.  The leading-product branches and the $p_1p_3$
branch are at most $v(t_1)=897/1600$, because $v$ decreases above
$\mu=25/128$.  For $s=t_2+t_3$, the $p_2p_3$ gate also fails:
if $s\le\mu$, its condition would give
$2(11/50)+s<(1+s)/2$, hence $s<3/25$, contrary to $s>3/20$;
if $s\ge\mu$, then $2(11/50)+s>v(s)$ already at $s=\mu$,
and the left side increases while the right side decreases.
The two raw $W$ branches are absent here as well: their inner values
are below $153/256$, but their gate is $(1+t_1)/2=61/100$.
Thus this construction leaves that specific discrepancy unresolved.
It does not assert nonexistence of a different factorization argument.

For $t_1>\mu$ at $\alpha=1$, the same observations show why one must
not extend the whole raw eight-branch maximum automatically.
The leading larger subsums and the $p_1p_3$ branch are dominated by
$v(t_1)$, and both $W$ branches fail their gates.  A proved additional
choice is $v(t_2)$ only with its extra gate, or $v(t_2+t_3)$ with its
displayed pair gate.  Other points require their own verified branch.
The separate cross-level $H/h$ issue is unaffected by this factorization
extension.

% Research proof fragment, 7 October 2026. No claim of priority.
% The analytic coefficient is defined by an integral; numerical certification
% is separate. Root alone decides whether to integrate this fragment.
\section{A first-prefix stop and variable third-prefix upper level}
\label{sec:variable-prefix-upper}

Put
\[
 \Theta=\frac{597}{1000},\qquad \mu=\frac{25}{128},\qquad
 k=\frac7{4000},\qquad B(t)=\frac{153}{256}-\frac{3t}{8}.
\]
Let $v$ and $\Phi$ be the single-variable and full eight-branch
envelopes in Section~\ref{sec:sharper-upper-support}, at
$\alpha=1$ and $\sigma=7/32$. Set
\[
 L(\boldsymbol t)=\min\{\Theta,\Phi(t_1,t_2,t_3)\},\qquad
 S=t_1+t_2+t_3,
\]
and
\[
 \mathcal E=\{0<t_3\le t_2\le t_1\le\mu:\quad
          3t_1\le\Theta,\quad t_1+t_2+3t_3\le\Theta\}.
\]
Write $\bar f=f/(2e^\gamma)$, extending $\bar f(s)=0$ for $s\le2$.
The coefficient for this construction is
\begin{align}
 K_{\rm var}={}&\frac4\Theta+
 4\int_\mu^{\Theta/3}
       \frac{\bar f((\Theta-t)/t)}{t^2}\,dt\notag\\
 &+4\int_{\mathcal E}
 \frac{\bar f((\Theta-S)/t_3)
       -\bar f((L(\boldsymbol t)-S)/t_3)}
 {t_1t_2t_3^2}\,d\boldsymbol t.
 \label{eq:variable-prefix-upper-coefficient}
\end{align}
The third integrand vanishes for $t_3\le k$, as proved below;
its apparent singularity at zero therefore causes no improper-limit
issue. Both integrands are nonnegative. The first correction can
also be written exactly as
\[
 \frac4\Theta\int_2^{6427/3125}\frac{\log(s-1)}s\,ds,
\]
using $\bar f(s)=\log(s-1)/s$ for $2\le s\le4$.
Equivalently, the sum of the base and first correction is
\[
 \frac4\mu\,\bar F(\Theta/\mu),\qquad
 \bar F=F/(2e^\gamma).
\]
Indeed $(s\bar F(s))'=\bar f(s-1)$ and
$3\bar F(3)=1$, while $\Theta/\mu=9552/3125\in(3,4)$.
This is the ordinary upper density obtained by using the shorter
sieve cutoff $N^\mu$ from the outset. We retain the explicit first
correction to match the separate numerical integral.

\begin{proposition}[Variable-prefix upper-sieve interface]
\label{prop:variable-prefix-upper}
Let $K_*>K_{\rm var}$. For every admissible product multiset in
Lemma~\ref{lem:primeproducts55}, with its original labelled mass
$X_N=(J+o(1))N/\log N$, one has
\[
 \sum_{\substack{m<N\\N-m\in\PP}}w(m)
       \le(K_*J+o(1))\HN.
\]
Consequently the five original switched multisets satisfy
\[
 \Upsilon_i\le(K_*U_i/8+o(1))\HN\quad(i=7,8,9,10),\qquad
 \Upsilon_{11}^{\rm full}
       \le(K_*U_{11}^{\rm full}/8+o(1))\HN.
\]
For every fixed $1/2<b<1$,
\[
 B_b(N)\le(K_*\log(1/b)+o(1))\HN.
\]
In particular a lower bound $D_{1,2}(N)\ge(c-o(1))\HN$ implies
\[
 \liminf_{\substack{N\to\infty\\2\mid N}}D_b(N)/\HN
       \ge c-K_*\log(1/b).
\]
No numerical bound on $K_{\rm var}$ is asserted in this proposition.
\end{proposition}

\begin{proof}
\paragraph{Uniform margins and the compact third-prefix correction.}
On $0\le t_1\le\mu$ one has $v(t_1)=(1+t_1)/2$, and hence
\[
 v(t_1)-3t_1\ge\frac3{256},\qquad
 \Theta-3t_1\ge\Theta-3\mu=\frac{177}{16000}.
\]
As $\Phi\ge v(t_1)$, these inequalities give the strict uniform
first-cube surplus
\begin{equation}\label{eq:variable-cap-first-cube-margin}
 L(\boldsymbol t)-3t_1\ge\frac{177}{16000}.
\end{equation}
For the inner minimum $W_1$ of $w(t_1,t_2,t_3)$,
\[
 \frac{66-34t_3}{107}>B(t_3),\qquad
 1-2t_2\ge1-2\mu=\frac{39}{64}>\frac{153}{256},
\]
so $W_1=B(t_3)$. Its gate is exactly $v(t_1)$.
Therefore $\max\{v(t_1),w(t_1,t_2,t_3)\}
=\max\{v(t_1),B(t_3)\}\ge B(t_3)$.
Since $B(k)=\Theta$, one has $L=\Theta$ whenever $t_3\le k$.
This proves the compactness assertion following
\eqref{eq:variable-prefix-upper-coefficient}.

\paragraph{Coarse labels, robust caps, and a separate fine grid.}
Fix a small rational $e>0$, and put
\[
 \nu=\Theta-2e,\qquad D=N^\nu,\qquad z=N^{1/5},\qquad
 r_e=\mu-4e.
\]
Partition the exponent interval $[0,1/5]$ into finitely many
intervals of width at most $\chi$, where
$\chi\ll\min(e,k)$; use left-closed, right-open intervals below
the final endpoint. These are the \emph{coarse labels}.
A first label is kept if its closed upper endpoint is at most
$r_e$. This decision is identical in all subsequent ranks.
The cutoff represented by these labels lies between
$r_e-\chi$ and $r_e$.

For a closed ordered coarse third box $B$ with a kept first label,
define $\Gamma_B$ as follows. Take the maximum of the infima over
$B$ of
\[
 v(t_1),\ v(t_2),\ v(t_1+t_2),\ v(S),\qquad
 \max\{v(t_1),W_1\},\quad\max\{v(t_1),W_2\},
\]
where $W_2$ is the inner minimum of $w(t_1,t_3,t_2)$.
In addition include $\inf_B v(t_1+t_3)$ if
\[
 \inf_B\{v(t_1+t_3)-(t_1+2t_2+t_3)\}\ge e,
\]
and include $\inf_B v(t_2+t_3)$ if
\[
 \inf_B\{v(t_2+t_3)-(2t_1+t_2+t_3)\}\ge e.
\]
All these infima are taken on the closed ordered box. Put
\begin{equation}\label{eq:variable-prefix-coarse-cap}
 \lambda_B=\min\{\nu,\Gamma_B-2e\}.
\end{equation}
Empty ordered boxes are omitted. The number of caps is finite.
Since $W_1=B(t_3)$, one also has
$\Gamma_B\ge B(\mu)=537/1024>1/2$; thus the intermediate
factorization parameters below lie in the source's permitted range.
For sufficiently small $e,\chi$, every actual vector in a kept
box satisfies $3t_1<\lambda_B$, by
\eqref{eq:variable-cap-first-cube-margin} and
$\Gamma_B\ge\inf_Bv(t_1)$.

Choose the presieve precision $\eta$ later, and put
\[
 w=D^{\eta^2}=N^\omega,\quad \omega=\nu\eta^2,
 \qquad L_{\rm small}=D^\eta=N^\ell,\quad \ell=\nu\eta.
\]
Take $\omega<k/4$. Partition $[w,z)$ into \emph{fine} exponent
intervals of width at most $\tau$, refining the coarse labels and
clipped at $w,z$. Every fine label inherits its coarse parent.
In particular the third cap is unchanged when the fine grid is
refined. Let $M=\lceil\nu/\omega\rceil+2$.
The fine width $\tau$ is chosen after the coarse grid and $\eta$;
it will satisfy all analytic margins and the density estimate below.
There is no assertion uniform in the number of coarse labels.

\paragraph{Raw support and a complete pointwise majorant.}
All primes are restricted to $p\nmid N$, and
$P_N(y)=\prod_{p<y,\ p\nmid N}p$.
Inside $\mathcal D^+(D)$ define $\mathcal R$ by keeping only
tuples whose first prime has a kept coarse label. If a tuple has
at least three primes and its third prime is at least $w$, also
require that it belongs to $\mathcal D^+(N^{\lambda_B})$,
where $B$ is its first three coarse labels. If its third prime is
below $w$, impose no extra cap. The empty product is kept.

The support is prefix-closed. Every first exclusion has odd rank.
The additional first stop has rank one. A capped third tree first
fails its third cube or a later odd cube; its first cube cannot fail
by the strict cap margin. Thus
$\lambda_{\mathcal R}(d)=\mu(d)\mathbf1_{d\in\mathcal R}$ obeys
the positive odd-first-exit identity
\eqref{eq:odd-first-exit-identity}. In particular it gives a pointwise
upper bound on avoiding the sifting primes.
For $Q=P_N(z)/P_N(w)$, the same identity gives
\begin{equation}\label{eq:variable-prefix-rough-density}
 V_N(w)\sum_{d\mid Q}\lambda_{\mathcal R}(d)g_N(d)
 =V_N(z)+\sum_{\substack{\mathcal R\text{ first exits }P\\
                         \text{last prime}\ge w}}
                   g_N(P)V_N(\text{last prime})
 \le V_{\mathcal R}.
\end{equation}
Here $V_{\mathcal R}$ is the full raw density on $P_N(z)$.
All rough first exits are full first exits with last prime at least
$w$, since their earlier primes are larger. Omitted full exits are
positive. This is a one-sided comparison, not an assumed equality.

For a nonincreasing fine label vector of rank $j\le M$, retain
the first-label rule. Use effective level $\nu$ if $j<3$ and
the fixed coarse cap $\lambda_B$ if $j\ge3$.
At even rank impose all odd cube conditions on lower endpoints,
with non-strict inequality. At odd rank impose them on upper
endpoints, with strict inequality. Let $c(d)=(-1)^j$ on the
resulting squarefree rough support and $c(d)=0$ elsewhere.
Then
\[
 c(d)\ge\lambda_{\mathcal R}(d)
\]
for every rough squarefree $d$: a positive raw coefficient passes
the outer test, and a negative boxed coefficient is certainly raw
supported. The rank cutoff omits no raw supported tuple.
Consequently $\sum_{d\mid n,\,d\mid Q}c(d)\ge\mathbf1_{(n,Q)=1}$.

Let $C_\chi$ be one plus the number of coarse caps. A discrepancy
between $c$ and $\lambda_{\mathcal R}$ lies within $O(M\tau)$
of an odd cube hyperplane at one of these $C_\chi$ levels.
For each level and tagged prime, drop the coarse-label restrictions
and use the prime number theorem above $N^\omega$; the tagged
interval has reciprocal prime mass $O(M\tau/\omega)+o_N(1)$.
Summing the remaining subsets using
$\prod_{w\le p<z}(1+g_N(p))\ll\omega^{-1}$ gives
\begin{equation}\label{eq:variable-prefix-box-error}
 V_N(w)\sum_{d\mid Q}(c(d)-\lambda_{\mathcal R}(d))g_N(d)
 \ll V_N(z)C_\chi\tau\omega^{-5}+o_N(V_N(z)).
\end{equation}
This argument is legitimate even when the tagged prime is in the
first three positions: the union over the \emph{frozen} finite cap
list is an upper bound. Choose $\tau$ so that
$C_\chi\tau\omega^{-5}$ is as small as required.

Take fundamental-lemma weights $\beta^\pm$ on primes below $w$,
of level $L_{\rm small}$, with $|\beta^\pm|\le1$ and density
$V_N(w)(1+O(e^{-c/\eta}))$. Append $\beta^+$ at even rough rank
and $\beta^-$ at odd rough rank, retaining the sign of $c$:
\[
 \rho(bd)=c(d)\beta_{\Omega(d)}(b),\qquad
 \beta_j=\begin{cases}\beta^+,&j\text{ even},\\
                         \beta^-,&j\text{ odd}.
          \end{cases}
\]
The small/rough decomposition is unique, so $|\rho(q)|\le1$.
For every $n$ its complete divisor sum satisfies
\[
 \sum_{q\mid n}\rho(q)
 \ge\mathbf1_{(n,P_N(w))=1}
        \sum_{\substack{d\mid n\\d\mid Q}}c(d)
 \ge\mathbf1_{(n,P_N(z))=1}.
\]
The first inequality holds term by term with the appropriate sign,
even if a lower beta polynomial is negative. Its normalized density
error is $O(e^{-c/\eta}\omega^{-2})$. Combining this fact with
\eqref{eq:variable-prefix-rough-density} and
\eqref{eq:variable-prefix-box-error} yields
\begin{equation}\label{eq:variable-prefix-complete-density}
 \sum_q\rho(q)g_N(q)
 \le V_{\mathcal R}+
 O\left(V_N(z)\{C_\chi\tau\omega^{-5}
                     +e^{-c/\eta}\omega^{-2}+o_N(1)\}\right).
\end{equation}

\paragraph{Actual atoms, gate margins, and their redistribution.}
The fine support has its exact ordered-slot expansion. A repeated
fine label of multiplicity $m$ contributes one factor $1/m!$;
the product of these factorials is used exactly once. Thus each
rough coefficient is a finite scalar combination of squarefree
prime-slot convolutions, followed by its assigned $\beta^\pm$.
At rank at least three its actual prime vector satisfies
\[
 d\in\mathcal D^{\rm well}
       (N^{\lambda_B+O(M\tau)}),\qquad
 d\le N^{\lambda_B+O(M\tau)}.
\]
Indeed an odd cube implies its own square condition and the next
square condition by ordering; the lower-endpoint tests lose only
$O(M\tau)$ when evaluated on an actual tuple.

Every coarse box supplies $\Gamma_B\ge\lambda_B+2e$.
For an ordinary $v$ branch this is a continuous affine-piece margin.
For a rescued $W_i$ branch, the gate is $v(t_1)$ on all these boxes.
At a representative vector, if $W_i-v(t_1)\le e/4$, the $v(t_1)$
branch itself retains a value at least $\Gamma_B-e/4$; otherwise
the $W_i$ branch has a positive gate gap. In the latter case choose
its factorization parameter below its value and above its gate.
For the pair branches the definition of $\Gamma_B$ gives a gate
gap at least $e$ throughout the coarse box, so one can likewise
choose the factorization parameter below the branch value and
above the gate. In every case the chosen parameter remains more
than a fixed multiple of $e$ above the actual cap. This does not
treat a jump in $\psi$ as a continuous function.

On a local product scale
$N^{1-\kappa}/O(1)\le x\ll N$, use $\alpha=1+3\kappa$.
All exponent coordinates and affine branches change by $O(\kappa)$.
The first-prime bound $t_1\le\mu-4e$ ensures the required
$t_1\le\mu_\alpha$ after this change, and the preceding gate or
rescue margins survive. Choose the hierarchy
\[
 \ell+O(M\tau),\ \kappa\ll\epsilon_{\rm disp}\ll\zeta\ll e,
 \qquad \ell+O(M\tau)<(3/5-\Theta)/10.
\]
For each required $U$, apply Lichtman 2023, Proposition 6.1, to
a representative fine box, at its selected parameter and with
$A_0=Ux^{-\zeta}$ for
$x^{2\zeta}\le U\le x^{1/3+\zeta/2}$.
The source's three-factor partition has strict monomial margins.
Keep this partition of the prime slots for the whole fine box;
rounding multiplies its factor bounds by $N^{O(M\tau)}$ and
appending the beta factor multiplies one by $N^\ell$.
These losses are chosen smaller than the least Type II saving
$9\zeta/16$ in the previously checked redistribution argument.

Normalize each grouped prime convolution by its slot factorial,
retaining that factorial in its scalar coefficient. The three
group coefficients are then bounded by one. Combining the small
beta factor with one rough group preserves this bound by the unique
small/rough decomposition. Removing and restoring the squarefree
restriction costs
\[
 (N+N^{\nu+\ell+O(M\tau)})w^{-1}(\log N)^C,
\]
by the same rough-square collision and divisor argument used in
Section~\ref{sec:pruned-upper-sieve}. The number of atoms and their
total scalar mass are finite at fixed parameters; the exponent $C$
may depend on those parameters. All complete moduli satisfy
$q\le N^{\nu+\ell+O(M\tau)}<N^{3/5-c_0}$ for some $c_0>0$.
Atoms of rank at most two instead have complete support below
$N^{2\mu+\ell+O(M\tau)}<N^{1/2-c_1}$.
They use the classical/generalized Bombieri--Vinogradov cases
already proved in Lemma~\ref{lem:primeproducts55}.

\paragraph{Distribution for the unchanged outer product families.}
The preceding construction checks all modulus-atom inputs to the
proof of Lemma~\ref{lem:primeproducts55}. To spell out the remaining
dependency, apply the $k_{\rm HB}=3$ Heath--Brown decomposition to
each original prime slot of a separated smooth outer tuple weight.
Every original group still has its original positive-power scale.
An individual variable of length between $x^{2\zeta}$ and
$x^{1/3+\zeta/2}$ uses the verified rough Siegel--Walfisz lemma
and the three-factor distribution estimate with the atom partition
just constructed. If no such variable occurs, every original
group contains a long variable, so there are at most two groups;
the one- and two-long-variable estimates apply because the complete
modulus support is below $N^{3/5-c_0}$. Their partial-interval forms,
Mellin twists, and divisor treatment of the tails are the same as
in that lemma. Prime powers, repeated outer primes, and products
below $N^{1-\kappa}$ retain a power saving. No arbitrary bounded
arithmetic outer weight, or Siegel--Walfisz claim for an arbitrary
convolution, has been introduced.

Thus, for each separated smooth tuple weight $W$, the complete
coefficient $\rho$ satisfies
\[
 \sum_q\rho(q)\left\{
   \sum_{\substack{\boldsymbol p\ \mathrm{prime}\\
                p_1\cdots p_j\equiv N\pmod q}}W(\boldsymbol p)
       -\frac{X_W^{(q)}}{\varphi(q)}\right\}
       \ll_A N(\log N)^{-A}.
\]
The implied constant may depend on the fixed grids and precisions.
The exact coprime mass can be replaced by total mass: $|\rho|\le1$
gives, for each outer prime $p$,
\[
 \sum_{q:\,p\mid q}\frac{|\rho(q)|}{\varphi(q)}
       \ll\frac{\log N}{p-1},
\]
and every outer prime is at least $N^a/2$, with at most eleven
factors. The total correction is
$O(X_{\rm out}N^{-a}\log N)$, without a factor equal to the level.
Use the single global nonnegative smooth outer approximation from
that lemma. Its excess mass is counted once; its dyadic partition,
labels and bounded multiplicities are preserved.
These facts attach \eqref{eq:variable-prefix-complete-density} to
an actual prime-complement upper bound, rather than just an ordinary
Rosser density calculation.

\paragraph{Exact density of the first and third modifications.}
The stopped rank-one trees and the capped rank-three trees are
disjoint. Their exact full-density identity is
\begin{align}
 V_{\mathcal R}-V_N^+(D,z)
 ={}&\sum_{\substack{p<z,\ p^3<D\\
                       \text{first coarse label stopped}}}
          g_N(p)V_N^-(D/p,p)\notag\\
 &+\sum_{\substack{P=p_1p_2p_3\in\mathcal D^+(D)\\
                   p_1>p_2>p_3\ge w\\
                   \text{first coarse label kept}}}
 g_N(P)\left\{V_N^-(D/P,p_3)
 -\mathbf1_{P\in\mathcal D^+(N^{\lambda_B})}
                     V_N^-(N^{\lambda_B}/P,p_3)\right\}.
 \label{eq:variable-prefix-upper-exact-density}
\end{align}
For either modification the prefix is odd, so its descendant
contribution in an ordinary upper support is minus its density
times a residual lower density. This proves the signs in the
identity. Where the capped third prefix is present, its cube gives
$Pp_3^2<N^{\lambda_B}$, making the residual lower initial-square
condition automatic. Where it is absent, its first cube still
holds, so precisely the third cube fails. The omitted residual
term is consequently represented by $f(s)=0$ for $s\le2$.

In the first sum the prime exponent lies in a fixed compact
interval near $[\mu,\Theta/3]$, and its residual parameter is at
least two. For the second sum, caps equal the global exponent
$\nu$ throughout $t_3\le k/2$, provided $\chi<k/8$ and the
other fixed parameters are sufficiently small. Indeed the entire
coarse box then has third exponent below $k$, so
$\Gamma_B\ge\inf_BB(t_3)>\Theta$ and
\eqref{eq:variable-prefix-coarse-cap} gives $\lambda_B=\nu$.
Such terms cancel exactly. All nonzero third corrections therefore
have $t_3>k/2$, independently of the small presieve cutoff $\omega$.

Apply Lemma~\ref{lem:uniform-residual-lower-density} to every
present residual term. Its parameter pairs lie in compact positive
sets, uniformly over the finite cap list; $f(2)=0$ handles a capped
third prefix as it disappears. The uniform formula
\[
 V_N(N^t)=(2e^{-\gamma}+o(1))\frac{\CN}{t\log N}
\]
and prime-number-theorem integration evaluate both corrections.
All primes in these corrections exceed fixed powers of $N$,
so omitted primes dividing $N$ have uniformly negligible reciprocal
mass. First let $N\to\infty$ at fixed construction parameters.

Now remove the beta error and fine-grid error within the displayed
hierarchy, keeping the coarse labels fixed when refining cube tests.
Then refine the coarse grid and let $e\downarrow0$.
Off a finite union of level and gate hyperplanes, the robust cap
$\lambda_B$ tends to $L(\boldsymbol t)$: a strict pair gate is
eventually admitted, a failed gate remains omitted, and the two
rescued $W_i$ maxima are continuous. The excluded first strip tends
to the plane $t_1=\mu$. All these boundaries have measure zero.
The third correction is confined to the fixed compact set
$t_3\ge k/2$ and has a bounded integrand, so dominated convergence
applies. There is no appeal to a distribution theorem at a gate,
at $t_1=\mu$, or at a limiting level.

For the ordinary global density use the uniform one-sided upper
bound $V_N^+(D,z)\le V_N(z)(F(5\nu)+o(1))$.
Here $5\nu<3$, giving the limiting coefficient $4/\Theta$.
The preceding two correction limits give exactly
\eqref{eq:variable-prefix-upper-coefficient}, since
$2e^{-\gamma}f=4\bar f$.
The strict gap $K_*>K_{\rm var}$ allows all grids and precision
parameters to be fixed so that their total coefficient is below
$K_*$. Removing the smooth outer approximation gives the first
assertion of the proposition. The desired cases $0<N-m\le z$
cost only $O(z)$ for bounded multiplicity.

\paragraph{The five switched masses and bad semiprimes.}
For each of $Y_7,Y_8,Y_9,Y_{10},Y_{11}^{\rm full}$, the original
specified primes and their labels remain in place. Factor the rough
quotient canonically into increasing primes and split into its
finitely many possible factor counts, as in
Lemma~\ref{lem:primeproducts55}. There is no additional quotient
factorial and no identification of different specified-prefix
choices. The moving quotient cutoff in $Y_9$, the negligible unit
quotient cases, and the full fourth-prime interval
$N^h\le p_4<N^d$ in $Y_{11}^{\rm full}$ are unchanged.
Reassembling gives exactly the five masses $J_i=U_i/8$.

For $B_b(N)$ use the ordered two-prime family $m=rq$, $r\le q$,
$r>q^b$. After removing $m<N^{1-\kappa}$, both prime factors
exceed $N^{1/3+\delta_b}$ for a suitable fixed $\delta_b>0$ and
sufficiently small $\kappa$. The repeated-prime case has a power
saving. Lemma~\ref{lem:mass} gives its mass
$(\log(1/b)+o(1))N/\log N$, with multiplicity one.
The same complete weight therefore bounds this family with
coefficient $K_*\log(1/b)$. Subtracting it from $D_{1,2}$, with
the previously recorded $O(1)$ unit-convention adjustment, proves
the final assertion. This completes the analytic interface; a
strict numerical coefficient requires a separate certificate.
\end{proof}

% Algebraic integration fragment; no documentclass.
% Valid only on the ordered three-prime domain stated below.

\subsection{A continuous six-branch expression on the ordered domain}
\label{sec:variable-envelope-simplification}

Put $\mu=25/128$, $\Lambda=153/256$, and
\[
 A(t)=\frac{1+t}{2},\qquad
 C(t)=\frac{57}{64}-\frac32t,\qquad
 B(t)=\Lambda-\frac38t.
\]
At residue exponent $\alpha=1$ and spectral exponent $7/32$,
the definitions in Lichtman,
\href{https://arxiv.org/html/2309.08522v1}{arXiv:2309.08522v1},
(6.2) and (6.4), become
\begin{align*}
 v(t)&=\min\left\{A(t),\frac{32-34t}{39},C(t)\right\},\\
 \psi(a,b)&=a\mathbf1_{a\ge b},\\
 w(a,b,c)&=\psi\left(
   \min\left\{\frac{66-34c}{107},B(c),1-2b\right\},A(a)\right),\\
 \Phi(x,y,z)&=\max\big\{
 v(x),v(y),v(x+y),v(x+y+z),w(x,y,z),w(x,z,y),\\
 &\hspace{38mm}\psi(v(x+z),x+2y+z),
             \psi(v(y+z),2x+y+z)\big\}.
\end{align*}
Thus all eight original branches, including both gates of each
type, are retained in this starting definition.

\begin{lemma}[Ordered-domain envelope identity]
\label{lem:variable-envelope-simplification}
For $0<z\le y\le x\le\mu$,
\begin{equation}\label{eq:variable-envelope-six-branches}
 \Phi(x,y,z)=\max\{v(x),v(x+y),v(x+y+z),B(z),v(x+z),v(y+z)\},
 \qquad v(t)=\min\{A(t),C(t)\}.
\end{equation}
If either of the pair branches $v(s)$, with
$(s,r)=(x+z,y)$ or $(s,r)=(y+z,x)$, satisfies $v(s)\ge v(x)$,
then its original gate has the uniform surplus
\begin{equation}\label{eq:variable-envelope-pair-gate-margin}
 v(s)-(s+2r)\ge\frac3{256}.
\end{equation}
\end{lemma}

\begin{proof}
The functions $A$ and $C$ meet at $t=\mu$.
For $t\le\mu$,
\[
 \frac{32-34t}{39}-A(t)=\frac{25-107t}{78}>0,
\]
whereas for $t\ge\mu$,
\[
 \frac{32-34t}{39}-C(t)=\frac{-175+1568t}{2496}>0.
\]
The displayed inequalities hold at the shared endpoint as well.
Consequently the middle affine branch is redundant, in fact for
every $t\ge0$. In particular $v(x)=A(x)$ and $v(y)\le v(x)$
on the stated ordered domain.

Write $W_1,W_2$ for the inner minima in $w(x,y,z)$ and
$w(x,z,y)$, respectively. For $u\ge0$,
\[
 \frac{66-34u}{107}-B(u)
 =\frac{525}{27392}+\frac{49u}{856}>0.
\]
Also $1-2y\ge39/64>\Lambda\ge B(z)$ and
$1-2z\ge39/64>\Lambda\ge B(y)$.
Thus $W_1=B(z)$ and $W_2=B(y)\le B(z)$.
Both original $w$ gates equal $A(x)=v(x)$, so
\[
 \max\{v(x),w(x,y,z),w(x,z,y)\}
     =\max\{v(x),B(z)\}.
\]
This identity uses the baseline branch when a gate fails;
it does not assert $w(x,y,z)=B(z)$ at those points.

It remains to inspect the two pair gates. In either case
$r\le x\le\mu$. Suppose $v(s)\ge v(x)$.
If $s\le\mu$, then
\[
 v(s)-(s+2r)=\frac12-\frac{s}{2}-2r
      \ge\frac12-\frac52\mu=\frac3{256}.
\]
If $s\ge\mu$, comparison with $v(x)=A(x)$ gives
\[
 \frac{57}{64}-\frac32s\ge\frac{1+x}{2},
 \qquad s\le\frac{25/32-x}{3}.
\]
Hence
\[
 v(s)-(s+2r)=\frac{57}{64}-\frac52s-2r
 \ge\frac{23}{96}-\frac76x
 \ge\frac{23}{96}-\frac76\mu=\frac3{256}.
\]
These are precisely the gates $x+2y+z$ and $2x+y+z$.
A pair branch below $v(x)$ cannot change the maximum whether
or not its gate succeeds; a branch at least $v(x)$ has just
been shown to pass its gate. This completes the comparison
of all eight branches with the six in
\eqref{eq:variable-envelope-six-branches}.
\end{proof}

The resulting envelope is continuous on this ordered domain.
In particular
\[
 B(z)\le\Phi\le\Lambda,
 \qquad \Phi-3x\ge\frac12-\frac52x\ge\frac3{256}.
\]
For $\Theta<\Lambda$, it follows that $\Phi\ge\Theta$ whenever
$z\le(8/3)(\Lambda-\Theta)$.
At $\Theta=597/1000$ this cutoff is exactly $7/4000$.
These observations justify the compact third-correction domain
and supply strict margins for pair branches which can affect
the envelope. At a $w$ gate equality the baseline branch supplies
the same value; strict distribution margins must still be chosen
before perturbing the residue exponent or exponent boxes.

The six-branch formula is asserted only for
$0<z\le y\le x\le25/128$. It must not be applied to an
unordered triple without first sorting its coordinates, nor
to the perturbed residue exponent without its separate margin
argument.

% Research integration fragment; no documentclass. Not yet in the main file.
% This supplies an algebraic comparison and a derived local interface.
% It is not an independent verification of Li's Chen constant or a priority claim.

\subsection{The small-initial-cutoff limit of the published four-term sieve}
\label{sec:g1-limit-comparison}

Write $\mu=25/128$, $\Lambda=153/256$, and let
$\Phi(t_1,t_2,t_3)$ be the eight-branch expression in
\href{https://arxiv.org/html/2309.08522v1}{Lichtman,
arXiv:2309.08522v1}, (6.4), with residue exponent $1$ and
spectral exponent $7/32$. In this subsection
\[
 \mathcal T_\mu=\{0<t_3<t_2<t_1<\mu\},\qquad
 S=t_1+t_2+t_3,\qquad d\omega=\frac{dt_1dt_2dt_3}{t_1t_2t_3^2}.
\]
Extend the lower linear-sieve function by $f(s)=0$ for $s\le2$,
and put $\bar f=f/(2e^\gamma)$. Its usual normalization is
$F(s)\to1$, $f(s)\to1$, with $F\ge1\ge f\ge0$.

The published formula (29) in
\href{https://imar.ro/journals/Mathematical_Reports/Pdfs/2026/1_2/3.pdf}
{Runbo Li, \emph{On Chen's theorem, Goldbach's conjecture and almost
prime twins II}} uses the fixed initial exponent $\sigma=1/500$
and gives the numerical upper bound $G_1<6.06932$.
The limiting functional below is an algebraic consequence of that
four-term architecture. It is not a numerical theorem explicitly
printed in that source.

\paragraph{A limit independent of the first three levels.}
For $0<\sigma<\mu$, consider
\begin{align*}
 G_\sigma={}&\frac1\sigma F(\nu_0/\sigma)
 -\frac1\sigma\int_\sigma^\mu
          f((\nu_1(t)-t)/\sigma)\frac{dt}{t}\\
 &+\frac1\sigma\int_{\sigma<t_2<t_1<\mu}
          F((\nu_2(t_1,t_2)-t_1-t_2)/\sigma)
                         \frac{dt_1dt_2}{t_1t_2}\\
 &-\int_{\sigma<t_3<t_2<t_1<\mu}
          f((\Phi(\boldsymbol t)-S)/t_3)\,d\omega.
\end{align*}
Suppose the first three residual exponents have a positive uniform
lower bound:
\[
 \nu_0\ge c,\qquad \nu_1(t)-t\ge c,\qquad
 \nu_2(t_1,t_2)-t_1-t_2\ge c>0.
\]
They may depend on $\sigma$. Then
\begin{equation}\label{eq:g1-limit-identity}
 \lim_{\sigma\downarrow0}G_\sigma
 =\frac1\mu+\int_{\mathcal T_\mu}
       \{1-f((\Phi(\boldsymbol t)-S)/t_3)\}\,d\omega
 =:\mathcal G_\Phi.
\end{equation}
This covers Li's first three levels, since
$\Phi(t_1,t_2,t_3)\ge(1+t_1)/2$ on this domain. It also covers
the more conservative choice of total level $1/2-\delta$ for
\emph{all} of the first three terms: then one may take any
$c<7/64-\delta$, since $1/2-2\mu=7/64$.

Here is a direct convergence proof, avoiding an unquantified
cancellation of divergent integrals. Put $h=F-f$. The
delay equations give $(sh(s))'=-h(s-1)$ for $s>3$.
Thus $h$ is decreasing and
\[
 h'(s)=-\frac{h(s-1)+h(s)}s\le-\frac{2h(s)}s,
 \qquad h(s)\le\frac{9h(3)}{s^2}\quad(s\ge3).
\]
In particular $|F-1|+|1-f|\ll s^{-2}$ suffices here.
The first three terms can therefore be replaced by
\[
 T_\sigma=\frac1\sigma
    \left\{1-\log(\mu/\sigma)+\frac12\log^2(\mu/\sigma)\right\}
\]
with error $O_c(\sigma(1+\log^2(\mu/\sigma)))$.
Direct integration gives
\[
 \int_{\sigma<t_3<t_2<t_1<\mu}d\omega=T_\sigma-\frac1\mu.
\]
Finally
\[
 \Phi-t_1-t_2\ge\frac12-\frac32t_1
       \ge\frac{53}{256}>0.
\]
For sufficiently small $t_3$, the remaining sieve argument is
at least $c_0/t_3$. Consequently the integrand's numerator is
$O(t_3^2)$, and the integrated small-$t_3$ tail is bounded by a
constant times $\int_0^\sigma\log^2(\mu/t_3)\,dt_3$.
This proves both absolute convergence and
\eqref{eq:g1-limit-identity}. The proof only uses the standard
linear-sieve functions and their displayed delay equations;
it does not assume that Li's fixed-$\sigma$ numerical value equals
its limit.

\paragraph{Exact comparison with the compact candidate.}
Take $3\mu<\Theta<3/5$, and set
$L_\Theta(\boldsymbol t)=\min(\Theta,\Phi(\boldsymbol t))$.
The compact first-stop candidate has functional
\begin{equation}\label{eq:compact-g1-functional}
 K_\Theta=
 \frac{2e^{-\gamma}}\mu F(\Theta/\mu)
 +2e^{-\gamma}\int_{\mathcal T_\mu}
       \left\{f((\Theta-S)/t_3)
                    -f((L_\Theta-S)/t_3)\right\}\,d\omega.
\end{equation}
The integrand is zero outside the ordinary third-cube domain
$S+2t_3<\Theta$. Also the ordinary term can be written as
\begin{equation}\label{eq:first-stop-base-identity}
 \frac{2e^{-\gamma}}\mu F(\Theta/\mu)
 =\frac4\Theta+
   4\int_\mu^{\Theta/3}
       \frac{\bar f((\Theta-t)/t)}{t^2}\,dt.
\end{equation}
Indeed, for $3<s<4$,
$sF(s)/(2e^\gamma)=1+\int_2^{s-1}\bar f(u)\,du$.
Thus a first-prefix stop at $\mu$ is simply a change of the
sifting limit to $N^\mu$ at the density-functional level.

Apply the same four-term identity at a constant level $\Theta$.
It follows either by separating an ordinary Rosser upper weight
according to its first zero, one, two, or three primes above
$N^\sigma$, or by the corresponding delay identities. Since
$\Theta>3\mu$, all first-prefix cubes are admitted. Taking
$\sigma\downarrow0$ as above gives
\[
 \frac1\mu F(\Theta/\mu)
 =\frac1\mu+\int_{\mathcal T_\mu}
          \{1-f((\Theta-S)/t_3)\}\,d\omega.
\]
Combining this with \eqref{eq:compact-g1-functional} proves
\begin{equation}\label{eq:g1-cap-exact-comparison}
 K_\Theta=2e^{-\gamma}\mathcal G_\Phi+
 4\int_{\mathcal T_\mu}
       \left\{\bar f((\Phi-S)/t_3)
             -\bar f((\min(\Theta,\Phi)-S)/t_3)\right\}\,d\omega.
\end{equation}
The correction is nonnegative. Every branch of $\Phi$ is at most
$\Lambda=153/256$: the $v$ branches have maximum $\Lambda$, and
each $w$ branch is at most $\Lambda-3t_3/8$ or
$\Lambda-3t_2/8$. Hence equality holds in
\eqref{eq:g1-cap-exact-comparison} whenever $\Theta\ge\Lambda$.
At $\Theta=.597<\Lambda$ the inequality is strict. For example,
near $(t_1,t_2,t_3)=(.11,\mu-.11,.08)$ the ungated
$v(t_1+t_2)$ branch is larger than $.597$, and both lower-sieve
arguments are larger than $2$ on an open subregion.

For the proposed certification values $.596$ and $.597$, all
nonzero third-correction points have
\[
 t_3\ge k_\Theta:=\frac{153/32-8\Theta}{3}>0.
\]
To see this, either $v(t_1)\ge\Theta$, or
$(1+t_1)/2<\Theta$. In the latter case $t_3\le k_\Theta$
makes every inner entry of $w(t_1,t_2,t_3)$ at least $\Theta$;
its gate is also satisfied. Thus $\Phi\ge\Theta$.
The two concrete lower bounds are $53/12000$ and $7/4000$,
respectively. This explains why the capped functional is convenient
for a compact exact enclosure even though the uncapped limit is
numerically a little smaller.

\paragraph{What an exact improvement would mean.}
Li prints a fixed-$\sigma$ bound, not an optimality assertion.
A new certificate $K_\Theta<K_*$ would imply
$2e^{-\gamma}\mathcal G_\Phi<K_*$ by
\eqref{eq:g1-cap-exact-comparison}. Together with the local
distribution interface below, it would give an improved explicit
consequence of the published sieve method. This does not make
that newly certified consequence pre-existing prior art, and it
does not establish a new analytic method or world priority.
An uncomputed possible optimization must not be substituted for
an established numerical comparison bound.

The shifted prime-product proof and the order of limits are written
in the separate fragment \texttt{g1-optimized-transfer.tex}.

% Derived research integration fragment; no documentclass.
% Requires the algebraic definitions in g1-limit-comparison.tex.
% Numeric K certification is pending; the Li Chen count is an external input.

\subsection{A moving-residue interface sufficient for the limiting functional}
\label{sec:g1-localized-interface}

\begin{proposition}[Derived optimized prime-product upper sieve]
\label{prop:g1-optimized-transfer}
Let $\mathcal M_N$ be an admissible finite positive labelled
prime-product multiset in the sense of the existing prime-product
lemma, with a fixed positive lower exponent for each prime factor,
bounded tuple multiplicity, and total mass
$X_N=(J+o(1))N/\log N$, for sufficiently large even $N$. For any
$K_*>2e^{-\gamma}\mathcal G_\Phi$, the shifted prime count satisfies
\[
 \sum_{\substack{m<N\\N-m\ {\rm prime}}}w(m)
       \le (K_*J+o(1))\frac{C(N)N}{\log^2N}.
\]
Consequently, if the external published input
$D_{1,2}(N)\ge1.9728\,C(N)N/\log^2N$ is accepted, then
\[
 D_b(N)\ge\{1.9728-K_*\log(1/b)+o(1)\}
                         \frac{C(N)N}{\log^2N}
 \qquad(1/2<b<1).
\]
The proposition is a derived consequence of the analytic and
box-construction inputs detailed below. No numeric assertion
$K_*<6.815354546$ is made here: that requires a separate exact
certificate. In particular, the floating candidate values are
not premises of this proposition.
\end{proposition}

\begin{proof}
The literal normalization $\sum\epsilon_i=1$ and residue $N$ in
Li's Lemma~2.2 does not itself permit replacing $N$ by a product
scale $x<N$ while keeping the residue $N$. We instead use the
uniform analytic inputs and the already written box-atom argument.
The following is a \emph{derived} interface; it is not attributed
verbatim to any printed theorem.

Fix $b\in(1/2,1)$ and, initially, a fixed $\sigma>0$.
Let the outer sequence consist of ordered prime pairs $r\le q$,
$r>q^b$, $rq<N$, with bounded positive smooth outer weights.
Delete $rq<N^{1-\kappa}$ at cost $O(N^{1-\kappa})$ and choose
the outer approximation to preserve a lower factor exponent
$a>\mu$. This is possible because
$(1-\kappa)b/(1+b)>\mu$ for small $\kappa$.
For a local smooth rectangle write $x\asymp rq$ and let $X_W$
be its prime-pair mass. Every sieve modulus below is composed of
primes $<N^\mu$, so its coprime mass is exactly $X_W$.

For a closed ordered third-prefix exponent box
$\mathcal B\subset\{\sigma\le t_3\le t_2\le t_1\le\mu-\delta\}$,
choose a constant total exponent $\nu_\mathcal B$ which is below
one valid branch of $\Phi$ throughout the box by a fixed margin.
Every gate used by that branch has a fixed positive surplus.
Use a residual lower sieve of level
\[
 R_\mathcal B=N^{\nu_\mathcal B-T_+-C\delta},
 \qquad T_+=\sup_{\mathcal B}(t_1+t_2+t_3).
\]
It is only needed on boxes where
$\nu_\mathcal B-T_+>2\sup_{\mathcal B}t_3+C\delta$;
the other lower terms may be set to zero before the boundary
strip is removed. Then, for every $A>0$, its averaged remainder
is $O_{A,\sigma,\delta,\eta}(N(\log N)^{-A})$, uniformly in $N$,
including the outer positive smooth prime-pair domain.

Here are the details which make this a shifted, variable-prefix
statement rather than a substitution in a fixed-residue theorem.

\begin{enumerate}
\item
Use the residual lower Iwaniec weight at the constant level
$R_\mathcal B$, not a claimed restriction of one global weight.
Retain the inequality ``residual primes $<p_3$'' as an ordering
condition. Its ordinary modified-sieve expansion has finitely
many rough-prime box atoms at every fixed construction precision
$\eta$, followed by a small-prime factor. Concatenate the three
specified prefix slots with the sorted residual slots.
The first two $D^{\rm well}$ conditions have a uniform margin,
because $\Phi-3t_1\ge3/256$. The active third cube provides
the remaining prefix condition. Later conditions follow from
the residual $D^{\rm well}(R_\mathcal B^{1+O(\eta)})$ support.
Thus each complete rough atom belongs to
$D^{\rm well}(N^{\nu_\mathcal B-c\delta})$.

\item
The varying cutoff $p_3$ is not frozen in the wrong direction.
It is handled by sorting all the rough slots before boxing.
When prefix and residual primes occupy the same interval label,
the prefix slots are the largest required ranks in that label.
On squarefree products, the sorted indicator is the full tuple
convolution divided by the factorial of each full label
multiplicity. Keep this scalar exactly once. This is the same
repeated-box normalization used in the existing prefix and
upper-weight fragments; it does not require disjoint labels
across the prefix/residual boundary. Rough square collisions
cost $O((N+D)w^{-1}\log^C N)$, where $w$ is the fixed positive
power presieving cutoff.

\item
Apply L23 Proposition~6.1 to a representative of each full
rough box, retain that slot partition throughout the box, and
append the small-prime factor. Dividing grouped prime
convolutions by their slot factorials makes each of the three
modulus coefficients $1$-bounded. The small/rough decomposition
is unique, so appending the small-prime weight preserves this
bound. The number and total absolute scalar mass of atoms are
fixed before $N\to\infty$.

\item
On every local scale use $\alpha=1+3\kappa$, not $\alpha=1$.
For $N^{1-\kappa}/O(1)\le x\ll N$, one has
$N<x^\alpha<x^{1+\epsilon_{\rm disp}}$ after choosing
$\kappa\ll\epsilon_{\rm disp}$. All exponent coordinates
change by $O(\kappa)$. The largest-prime restriction becomes
$t_{1,x}\le\mu_\alpha$ because the removed $\delta$ strip
dominates this change. The buffered branch and gate inequalities
survive as well. Boundary strips are removed only after taking
the large-$N$ limit; no continuity across a discontinuous gate
is assumed.

\item
Apply the order-three Heath--Brown identity to each of the two
original prime factors at its own scale, retaining its original
product group and the roughness cutoff. If one individual
variable has length $U\in[x^{2\zeta},x^{1/3+\zeta/2}]$, take
it as the Type~II factor. Its coefficient is $1$ or $\mu$,
up to logarithms and Mellin twists, and the existing rough
Siegel--Walfisz lemma applies. The other coefficient is
divisor-bounded; the usual truncation absorbs logarithmic
losses. In Proposition~6.1 take $A_0=Ux^{-\zeta}$.
Replacing $A_0$ by $U$ leaves the minimum saving
$(1-2\cdot7/32)\zeta=9\zeta/16$.
Choose the mesh and beta support costs much smaller than this
margin. L23 Proposition~5.2 now applies with its explicit
supremum over $0<|a|<x^{1+\epsilon_{\rm disp}}$, taking $a=N$.

\item
If there is no such medium variable, each original prime group
must contain a long variable. There are exactly two long
variables, one in each group, and the remaining ten variables
have product at most $x^{20\zeta+o(1)}$. The two-long-variable
estimate L23 Proposition~3.2 has a supremum over all nonzero
residues. It applies because every complete modulus is below
$x^{3/5-c}$, whereas
$(x/x^{20\zeta+o(1)})^{2/3-3\epsilon_{\rm disp}}$ is larger.
This checks the case which cannot be handled by pretending
that the original balanced primes themselves have Type~II
length at most $x^{1/3}$.

\item
Use one global positive outer approximation, then a smooth
dyadic partition and Mellin separation. The divisor moment
bound for $\tau_C(rq)\tau_C(|N-rq|)$ controls all Mellin tails
after summation; no factor $D$ is charged per boundary error.
Proper prime powers in the Heath--Brown expansion have a
power-saving total mass. The exact coprime mean is $X_W$,
since $r,q>N^a>N^\mu$. Increasing the arbitrary logarithmic
error exponent absorbs the dyadic rectangles and Mellin losses.
\end{enumerate}

These are precisely the coefficient, local-scale, medium-variable,
two-long-variable, and smoothing steps proved in
\texttt{prefix-sieve-interface.tex},
\texttt{uniform-prime-products-upper.tex}, and
\texttt{rough-sw-stability.tex}; only the full three-variable
branch replaces the single-variable branch in the modulus
factorization step. Their validity depends on the supplied
strict full-atom margins, not on a fixed global value $.55$.
The uniform main density of each residual lower sieve is the
$f$ density established in the retained-base density argument.
At fixed $\sigma$, all relevant residual parameters are compact;
the zero boundary agrees with $f(2)=0$.

\paragraph{The stated finite-prime-product scope.}
For a general admissible labelled product $m=p_1\cdots p_k$,
retain the lower bound $p_i\ge N^a/2$ in the common positive
outer approximation. Apply the order-three identity at each
of its at most eleven original prime groups. A medium variable
is handled exactly as above. In the absence of a medium
variable, each original group must contain a variable exceeding
$x^{1/3+\zeta/2}$: its at most six shorter variables could
otherwise have product at most $x^{12\zeta+o(1)}<N^a/2$.
There are at most two long variables altogether. Thus this
case is impossible for $k\ge3$; for $k=2$ it is the case
already checked. For $k=1$ the existing one-long-variable
Lemma~3.4 or two-long-variable Proposition~3.2 applies.
All dyadic blocks retain their original prime-group scale
restrictions before Mellin separation.

Here outer primes can be smaller than $N^\mu$, so the coprime
mass cannot be silently replaced by $X_W$. For the full
coefficient sequence $\lambda_j(d)$ of each of the four
Buchstab terms, keep its bounded coefficient norm
$|\lambda_j(d)|\le C_{\sigma,\eta}$ and its squarefree support.
The first $j$ largest rough primes identify the specified
prefix, even for the third term with the variable cutoff.
The same correction as in the prime-product fragment gives
\begin{align*}
 &\sum_d\frac{|\lambda_j(d)|}{\varphi(d)}
                       (X_{\rm out}-X_{\rm out}^{(d)})\\
 &\qquad\ll_{\sigma,\eta}(\log N)
       \sum_mw_{\rm out}(m)\sum_{\ell\mid m}\frac1{\ell-1}
 \ll_{\sigma,\eta}X_{\rm out}N^{-a}\log N.
\end{align*}
Indeed the contribution from moduli divisible by any fixed
$\ell$ is $O(\log N/(\ell-1))$. No factor equal to the
modulus level is lost. Repeated outer primes and outer primes
dividing $N$ are removed with the same power savings as in
the earlier finite-product proof. This supplies the stated
scope; it does not permit arbitrary bounded arithmetic
reweightings of the outer products.

\paragraph{A complete route to the limit without virtual-prefix levels.}
Use one common nonnegative outer prime-product weight $W$ for
\emph{all} terms of the Buchstab identity. It is chosen before
that identity is applied, and its total excess mass is counted
once. All four terms therefore refer to the same weighted
sequence of positive integers $|N-m|$; the isolated $m=N$
term is omitted with the usual power-saving error. Changing
the positive approximation separately after a negative
Buchstab term would not justify the subtraction.

Use the Buchstab identity (27) on this common sequence,
with initial cutoff $N^\sigma$ and final cutoff $N^\mu$.
For its rank-zero, rank-one and rank-two terms use total modulus
level $N^{1/2-\delta}$. Their full coefficients are bounded,
and their largest zero, one or two primes identify the specified
prefix; generalized Bombieri--Vinogradov therefore controls
their remainders uniformly in residue $N$.
Their residual exponents are bounded below by $7/64-\delta$.
Use the just-derived lower interface for the negative rank-three
term. All branch, ordering, cube, and $t_1=\mu$ boundary strips
have vanishing main mass at fixed $\sigma$.

First let $N\to\infty$, then remove smoothing, box and sieve
errors within fixed margins, then let the margins tend to zero,
and finally let $\sigma\downarrow0$.
Formula~\eqref{eq:g1-limit-identity} yields
for the common admissible product sequence
\[
 \sum_{\substack{m<N\\N-m\ {\rm prime}}}w(m)
 \le\{2e^{-\gamma}\mathcal G_\Phi J+o(1)\}
             \frac{C(N)N}{\log^2N}.
\]
For the bad-semiprime specialization this is
\[
 B_b(N)\le
  \{2e^{-\gamma}\mathcal G_\Phi\log(1/b)+o(1)\}
             \frac{C(N)N}{\log^2N}.
\]
Here the prime-pair mass is
$\{\log(1/b)+o(1)\}N/\log N$.
Prime squares, products sharing a prime factor with $N$, and
the desired prime $N-rq<N^\mu$ contribute only a power-saving
or $O(\omega(N))$ exception, as in the existing transfer proof.
For a strict numerical cap one chooses all fixed parameters
before the final sufficiently-large-$N$ assertion; no effective
threshold is claimed.

This route closes the required moving-residue/local-scale issue
\emph{as a derived consequence of the already stated analytic
and box-construction inputs}. It does not claim that Li's literal
Lemma~2.2 alone has this formulation. In particular it avoids
having to independently justify the virtual remaining-prime
levels $\Phi(t,\sigma,\sigma)$ and $\Phi(t_1,t_2,\sigma)$ for
the first two nonempty prefixes. Reproducing the exact finite
$\sigma=1/500$ printed functional solely from literal source
statements would still require that separate interface check.
Subtracting the bad-semiprime estimate from the stated external
Chen count proves the last assertion. A separate exact
certificate for $K_\Theta<K_*$, together with
\eqref{eq:g1-cap-exact-comparison}, would be a sufficient
numerical premise for using this proposition with that $K_*$.
\end{proof}

% Numerical consequence of the separately proved variable-prefix interface.
\subsection{Certified consequence of the variable-prefix construction}
\label{sec:variable-prefix-corollary}
The exact interval calculation in
\texttt{certify\_variable\_prefix.py} gives
\[
 K_{\rm var}\le
 \frac{4917859459084088392523169865}{721728714309933617299587072}
 <\frac{3407}{500}=6.814.
\]
The gap below $3407/500$ is exactly
\[
 \frac{27974909469527055451}{90216089288741702162448384000}>0.
\]
This is a certificate for the integral in
\eqref{eq:variable-prefix-upper-coefficient}, not a computer proof
of Proposition~\ref{prop:variable-prefix-upper}.
It uses 80-bit outward integer arithmetic, curvature bounds for
the delay equations of the sieve functions, weighted Jensen and
chord inequalities on interior cells, and exact rectangular
overestimates on boundary cells. All subdivision decisions are exact.
The final run used 995,180 subdivisions. A separate driver replayed
the first 50,000 subdivisions with identical rational endpoints.
The source, arithmetic dependency, and review hashes accompany the
certificate. Floating values are used only for display.

\begin{theorem}[Consequence of the variable-prefix construction]
\label{thm:variable-prefix-chen}
Under the analytic conclusion of
Proposition~\ref{prop:variable-prefix-upper}, every sufficiently
large even $N$ satisfies
\[
 D_{1,2}(N)\ge\frac{1783}{1000}\HN,
 \qquad D_{77/100}(N)\ge\frac1{1000}\HN.
\]
This calculation does not use Li's numerical Chen constant.
Using the external published input
$D_{1,2}(N)\ge(1233/625-o(1))\HN$ separately gives
\[
 D_{14973/20000}(N)\ge\frac1{10000}\HN.
\]
Thus its exponent is $14973/20000=0.74865$.
The external-input review restrictions in
Section~\ref{sec:input-audit} continue to apply.
\end{theorem}

\begin{proof}
Put $K_*=3407/500$ and rescale precisely the five switching
upper constants by $K_*/8=3407/4000$. Retain the complete
eleven-term quotient weight, including the full fourth-prime
range in $Y_{11}^{\rm full}$. The lower terms and the other
upper terms are unchanged. Outward rational rounding gives
\begin{center}
\begin{tabular}{ll@{\qquad}ll}
\toprule
Lower term&Lower bound&Upper term&Upper bound\\
\midrule
$L_1$&12.9014&$U_3$&19.3178\\
$L_2$&6.3640&$U_4$&13.1277\\
$L_5^{(2)}$&2.3605&$U_7^{\rm var}$&1.3094\\
$L_6$&2.2784&$U_8^{\rm var}$&3.9152\\
&&$U_9^{\rm var}$&2.8350\\
&&$U_{10}^{\rm var}$&0.2174\\
&&$U_{11}^{{\rm var},{\rm full}}$&0.5398\\
\bottomrule
\end{tabular}
\end{center}
The resulting eleven-term functional is at least
\[
 \frac{35677}{20000}
 =\frac{1783}{1000}+\frac{17}{20000}.
\]
Its positive slack absorbs the asymptotic error in $D_{1,2}$.
The same coefficient $K_*$ bounds the bad semiprimes, so
\[
 \liminf_{\substack{N\to\infty\\2\mid N}}
 \frac{D_b(N)}{\HN}\ge c-K_*\log(1/b).
\]
The elementary bound
\[
 \log\frac{1+u}{1-u}
 <2\left(u+\frac{u^3}{3}+\frac{u^5}{5(1-u^2)}\right)
 \qquad(0<u<1)
\]
at $u=23/177$ and $u=5027/34973$ gives respectively
\begin{align*}
 \frac{1783}{1000}-K_*\log(100/77)
 &>\frac1{1000}+\frac{678900927197}{640474411500000},\\
 \frac{1233}{625}-K_*\log(20000/14973)
 &>\frac1{10000}
    +\frac{3353708303566756189617}{21349427496285714700000000}.
\end{align*}
Both rational gaps are positive and absorb the remaining
asymptotic errors. The exact interval recombination and the
short rational checks are recorded by
\texttt{certify\_variable\_prefix\_transfer.py}.
\end{proof}

The comparison with the small-initial-cutoff limit of the published
$G_1$ construction explains this result as an optimized numerical
consequence of that sieve architecture. The certified coefficient
improves the bound obtained from the printed value $G_1<6.06932$;
it does not establish a new analytic method or priority over all
published or unpublished work. The separate $b=0.770$ conclusion
also remains weaker than the external-input $b=0.74865$ conclusion.

\subsection{The exact coprime mean in the one-long-variable case}
\label{sec:one-long-coprime-mean}
Here is an explicit reconciliation needed when applying L23,
Lemma~3.4, in the one-long-variable alternatives above.
Write $z_0=x^{1/(\log\log x)^3}$ and let $R(n)$ be the indicator
that all prime factors of $n$ exceed $z_0$. The printed main
sum in that lemma is unrestricted, whereas our discrepancy uses
the exact coprime sum. For $q\le x^{3/5}$ and a long-variable
endpoint $t\gg x^{1-10\zeta}$, the difference is bounded by
\[
 \frac{1}{\varphi(q)}
 \sum_{\substack{p\mid q\\p>z_0}}(t/p+1)
 \ll \frac tq\left\{\frac{\log^C x}{z_0}
                         +\frac{\log^C x}{t}\right\}
 \ll_A\frac{t}{q\log^A x}
\]
for every fixed $A$. Indeed, a rough integer not coprime to
$q$ is divisible by one of the displayed primes; use the union
bound, $q/\varphi(q)\ll\log^C x$, and
$\omega(q)\ll\log x$. The bound applies at both endpoints
of any long-variable interval. Partial summation allows the
logarithm and Mellin twist already present in the proof.

For a fixed short product $m$, if $(m,q)>1$, the progression
and exact coprime terms both vanish, since $(N,q)=1$.
Otherwise the long-variable residue is $N\overline m\pmod q$
and is primitive. Summation over the short products and
moduli costs only powers of the logarithm, by the fixed
divisor moment estimates for $\sum_{m\asymp M}\tau(m)^C/m$
and $\sum_{q\le Q}\tau(q)^C/q$. Increasing $A$ therefore
absorbs this correction in the stated remainder. This
reconciliation changes neither a distribution level nor a
sieve constant.
% Research extension fragment. Numerical certification and independent review
% are separate; do not infer a new theorem constant from floating probes.
\section{Extending the first-prime stop with a single-variable fallback}
\label{sec:extended-prefix-upper}

Set
\[
 \rho=\frac{19}{96},\qquad \mu=\frac{25}{128},\qquad
 \Theta=\frac{299}{500}.
\]
Retain the normalization $\bar f=f/(2e^\gamma)$ and extend
$\bar f(s)=0$ for $s\le2$. On ordered triples define
\[
 A(x,y,z)=
 \begin{cases}
 \Phi(x,y,z),&0<z\le y\le x\le\mu,\\
 v(x),&0<z\le y\le x,\quad\mu<x\le\rho,
 \end{cases}
 \qquad L=\min\{\Theta,A\},\quad S=x+y+z.
\]
Here the full envelope is used only within its largest-prime
range. The fallback uses the first assertion of L23,
Proposition~6.1, which requires the level to be below $v(x)$
and does not impose $x\le\mu$. All factorization calls retain
the repairs in Section~\ref{sec:two-factor-greedy-repair}.

Put
\[
 \mathcal E_\rho=\{0<z\le y\le x\le\rho:\ x+y+3z\le\Theta\}.
\]
The first raw cube $3x\le\Theta$ is automatic because
$3\rho=19/32<\Theta$. Define
\begin{align}
 K_{\rm ext}={}&\frac4\Theta+
 4\int_\rho^{\Theta/3}
          \frac{\bar f((\Theta-t)/t)}{t^2}\,dt\notag\\
 &+4\int_{\mathcal E_\rho}
 \frac{\bar f((\Theta-S)/z)-\bar f((L-S)/z)}{xyz^2}
                  \,dx\,dy\,dz.
 \label{eq:extended-prefix-functional}
\end{align}
The first correction is equivalently
\[
 \frac4\Theta\int_2^{4801/2375}\frac{\log(s-1)}s\,ds.
\]
Unlike the preceding compact construction, this third correction
can extend to arbitrarily small $z$. The tail estimate below is
therefore part of both the analytic proof and any numerical certificate.

\begin{proposition}[Extended first-stop upper sieve]
\label{prop:extended-prefix-upper}
For any $K_*>K_{\rm ext}$ and any admissible labelled prime-product
multiset of Lemma~\ref{lem:primeproducts55}, of mass
$X_N=(J+o(1))N/\log N$, one has
\[
 \sum_{\substack{m<N\\N-m\in\PP}}w(m)
              \le(K_*J+o(1))\HN.
\]
The five switched multisets therefore have their original upper
constants multiplied by $K_*/8$, and
\[
 B_b(N)\le\{K_*\log(1/b)+o(1)\}\HN\qquad(1/2<b<1).
\]
The proposition asserts no particular numerical value for $K_*$.
\end{proposition}

\begin{proof}
We use the actual upper support and the two-grid construction of
Proposition~\ref{prop:variable-prefix-upper}. The required changes
are given explicitly below; the outer-product families, their
labels and their mass normalization are unchanged.

\paragraph{Margins at the larger first stop.}
For $0\le t\le\rho$, one has $v(t)\ge3t$. Indeed,
on $[0,\mu]$,
\[
 v(t)-3t=\frac12-\frac52t\ge\frac3{256},
\]
and on $[\mu,\rho]$,
\[
 v(t)-3t=\frac{57}{64}-\frac92t
                   =\frac92(\rho-t).
\]
Fix a small rational $e>0$, put $\nu=\Theta-2e$, and keep
first coarse labels only when their closed upper endpoint is
at most $\rho-4e$. On these labels the latter margin is at least
$18e$. Also
\[
 \nu-3(\rho-4e)=\frac{17}{4000}+10e>0.
\]
Freeze a coarse exponent grid of width $\chi\ll e$. For a
third coarse box wholly below $x=\mu-4e$, use the robust full
envelope $\Gamma_B$ from
\eqref{eq:variable-prefix-coarse-cap}. For every other third box
use $\Gamma_B=\inf_B v(x)$. Define
\[
 \lambda_B=\min\{\nu,\Gamma_B-2e\}.
\]
By the preceding margins and the bounded slopes, after decreasing
$\chi/e$ every retained first box has
$\lambda_B>3x$ throughout its closed ordered portion.
The finite cap list is fixed before choosing the small-prime
precision and the refining fine grid.

\paragraph{Integer support, weights, and distribution.}
Inside $\mathcal D^+(N^\nu)$ retain the common first-label rule;
at ranks at least three whose third prime exceeds the rough
cutoff $w=N^\omega$, impose the ordinary upper support at
level $N^{\lambda_B}$. Impose no additional third cap below $w$.
The cap is fixed by the first three coarse labels and its first
cube cannot fail. Thus the support is prefix closed and all
first exclusions are odd, including the first-prime stop.

Apply the same fine-label parity rounding and parity-dependent
small beta weights as in the preceding proposition. In particular
the exact rough/full one-sided density comparison, the complete
pointwise majorant, and the bound on the coefficient norm remain
valid. If $C_\chi$ is one plus the number of frozen caps, the
normalized error remains
\[
 O(C_\chi\tau\omega^{-5}
                 +e^{-c/\eta}\omega^{-2}+o_N(1)).
\]
The fine width $\tau$ is selected after the coarse grid and
$\eta$; no uniformity in the number of coarse caps is asserted.

For actual atoms from the lower first-prime region, the buffered
full-envelope construction is unchanged. In the other region the
single-variable $v$ branch has a fixed surplus above the actual
cap. On a local scale $N^{1-\kappa}/O(1)\le x_0\ll N$, take
$\alpha=1+3\kappa$. The coordinate and branch perturbations are
$O(\kappa)$ and are absorbed by choosing $\kappa\ll e$.
Only boxes wholly below $\mu-4e$ use the full three-prime
formula, so the perturbed largest-prime restriction is respected.
The transition strip uses the single-variable assertion throughout.

Keep the same hierarchy of Type II savings, small beta support
and fine-box losses. The complete moduli have support below
$N^{3/5-c}$ for a fixed $c>0$, since $\Theta=0.598<3/5$.
Ranks at most two instead have support below
$N^{2\rho+o(1)}$, and $2\rho=19/48<1/2$.
The existing one- and two-long-variable alternatives, rough
Siegel--Walfisz estimates and residue-uniform Type II estimate
therefore apply. The label factorials, unique small/rough
decomposition, exact coprime-mass correction and single positive
outer approximation are the same as before. These operations
preserve the original five switched masses and the bad-semiprime
mass, including the full range of $Y_{11}^{\rm full}$.

\paragraph{Density at fixed cutoff.}
The exact density identity
\eqref{eq:variable-prefix-upper-exact-density} holds with the
new first labels and caps: first stopped trees and third capped
trees are disjoint, and their odd parity produces residual lower
densities with the same signs. At each fixed $\omega>0$ their
arguments range over compact sets. The uniform residual-density
lemma therefore evaluates them after taking $N\to\infty$.
The absence of a capped third prefix is encoded by $f(s)=0$
when its third cube fails. No uniform asymptotic statement for
unbounded sieve arguments is used at this stage.

\paragraph{Uniform integrability as the rough cutoff decreases.}
On the full ordered domain $x\le\rho$ one has
\begin{equation}\label{eq:extended-prefix-residual-margin}
 L-x-y\ge\rho.
\end{equation}
Indeed $\Theta-x-y\ge\Theta-2\rho\ge\rho$. For the other
branch, $A\ge v(x)$, and each piece of $v(x)-2x$ is at least
$\rho$: on $[0,\mu]$ its minimum is $53/256>\rho$, while on
$[\mu,\rho]$ it is decreasing to
$v(\rho)-2\rho=\rho$.
The buffered cap obeys the same inequality with an
$O(e+\chi)$ loss. Once these parameters are small enough,
$\lambda_B-x-y\ge\rho/2$ uniformly in all boxes and in
the rough cutoff.

Let $D_0(s)=\bar F(s)-\bar f(s)$. Standard monotonicity and
the delay equations give
\[
 (sD_0(s))'=-D_0(s-1),\qquad
 D_0(s)\le D_0(s_0)(s_0/s)^2\quad(s\ge s_0\ge3).
\]
For ordered levels the difference of the two lower sieve values
is nonnegative and at most $D_0((L-S)/z)$, because their common
limit lies between $\bar f$ and $\bar F$.
If $d_0\ge D_0(s_0)$ and $0<h\le\rho/(s_0+1)$, the whole
limiting contribution with $z<h$ is at most
\begin{equation}\label{eq:extended-prefix-tail-bound}
 \frac{2d_0s_0^2}{(\rho-h)^2}\,h
 \left\{\log^2(\rho/h)+2\log(\rho/h)+2\right\}.
\end{equation}
To verify it, use $(L-S)/z\ge(\rho-z)/z$ and integrate
\[
 \int_{z\le y\le x\le\rho}\frac{dx\,dy}{xy}
                       =\tfrac12\log^2(\rho/z).
\]
The remaining integral is
$\int_0^h\log^2(\rho/z)\,dz
=h\{\log^2(\rho/h)+2\log(\rho/h)+2\}$.
For buffered caps the same argument with $\rho/2$ in the
residual bound supplies a common integrable majorant. The upper
limit $\rho$ in the logarithm is unchanged. Thus the contribution
from $z<\omega$ tends to zero uniformly in the sufficiently small
buffers, even though it is not identically zero.

First take $N\to\infty$ at fixed parameters. Then let the beta
and fine-grid errors tend to zero with the coarse caps fixed;
the tail bound just proved controls the accompanying decrease
of $\omega$. Finally refine the coarse grid and send $e$ to zero.
The limiting fallback is $v(x)$ above $\mu$, and the full envelope
is recovered below it. The strip at $\mu$ has vanishing integral
by the same domination. In fact the two definitions agree
continuously there, since $\Phi(\mu,y,z)=v(\mu)=153/256$.
All gate and level boundary sets have measure zero.

The ordinary upper density contributes $4/\Theta$, since
$5\nu<3$. The two density corrections converge to those in
\eqref{eq:extended-prefix-functional}. Strict slack
$K_*>K_{\rm ext}$ permits all construction parameters to be
fixed before the final sufficiently-large-$N$ statement. Applying
this bound to the unchanged masses proves the proposition.
\end{proof}
\subsection{Certified consequence of the extended first stop}
\label{sec:extended-prefix-corollary}
The exact enclosure in \texttt{certify\_extended\_prefix.py} gives
\[
 K_{\rm ext}\le
 \frac{1230679674488580033396846301}{180734410032387061618573312}
 <\frac{681}{100}=6.810.
\]
The exact gap below $681/100$ is
\[
 \frac{3041445799396405640948843}{4518360250809676540464332800}>0.
\]
This run used two million subdivisions and includes the entire
small-third-prime tail. It attempted the stronger numerical target
$6.809$ without certifying it; the saved certificate preserves
that unsuccessful target and its false success flag. The assertion
above follows from a separate exact rational comparison with
$6.810$. No floating value is used to decide either comparison.
The root's independent adaptive driver reproduces the first
$50{,}000$ subdivisions' rational endpoints exactly, using the
separately reviewed enclosure primitives.

In particular, the tail below $z=1/100$ contributes at most
\[
 \frac{1028753216356588863}{1208925819614629174706176}.
\]
It follows from \eqref{eq:extended-prefix-tail-bound} with
$s_0=10$, using the outward table value of $D_0(10)$.
For the buffered analytic domination, one may explicitly take
$h\le\rho/[2(s_0+1)]$ and use the residual lower bound $\rho/2$.
The numerical certificate concerns the unbuffered limiting
integral; the analytic proof supplies the passage to that limit.

\begin{theorem}[Consequence of the extended first stop]
\label{thm:extended-prefix-chen}
Under Proposition~\ref{prop:extended-prefix-upper}, for every
sufficiently large even $N$,
\[
 D_{1,2}(N)\ge\frac{357}{200}\HN,
 \qquad D_{1539/2000}(N)\ge\frac1{2000}\HN.
\]
Thus $b=1539/2000=0.7695$ is an admissible exponent in this
research derivation, without using Li's numerical Chen constant.
If the external published input
$D_{1,2}(N)\ge(1233/625-o(1))\HN$ is accepted, then separately
\[
 D_{1497/2000}(N)\ge\frac1{50000}\HN,
 \qquad 1497/2000=0.7485.
\]
The latter implication retains the external-input restrictions in
Section~\ref{sec:input-audit}.
\end{theorem}

\begin{proof}
Use $K_*=681/100$ in Proposition~\ref{prop:extended-prefix-upper}.
Exactly the same five upper switching terms are multiplied by
$K_*/8=681/800$. Their outward upper bounds become
\[
 U_7\le1.3086,\quad U_8\le3.9129,\quad U_9\le2.8334,
 \quad U_{10}\le0.2173,\quad U_{11}^{\rm full}\le0.5395.
\]
The remaining entries of the table in
Section~\ref{sec:variable-prefix-corollary} are unchanged.
The complete eleven-term functional is therefore at least
\[
 \frac{71413}{40000}
 =\frac{357}{200}+\frac{13}{40000}.
\]
This strict slack absorbs the asymptotic error in the Chen count.
The same upper coefficient bounds the bad semiprimes, giving
$\liminf D_b(N)/\HN\ge c-K_*\log(1/b)$.

For $0<u<1$, the short upper bound
\[
 \log\frac{1+u}{1-u}
 <2\left(u+\frac{u^3}{3}+\frac{u^5}{5(1-u^2)}\right)
\]
at $u=461/3539$ and $u=503/3497$ proves
\begin{align*}
 \frac{357}{200}-\frac{681}{100}\log\frac{2000}{1539}
 &>\frac1{2000}
   +\frac{8273727872308673}{45476711094294000000},\\
 \frac{1233}{625}-\frac{681}{100}\log\frac{2000}{1497}
 &>\frac1{50000}
   +\frac{1298487768167459}{42679314784054000000}.
\end{align*}
Both positive rational gaps absorb the remaining asymptotic errors.
The source-linked exact recombination and these short checks are
saved by \texttt{certify\_extended\_prefix\_transfer.py}; a separate
review also recomputed the five-term rescaling without its helper
functions.
\end{proof}

This improves the preceding numerical consequences within the
same general sieve construction. It remains a research proof
draft relative to the analytic and published inputs already named;
neither formal verification nor superiority to every existing
result follows from the numerical certificate.

% Analytic extension. Numerical consequences are certified separately.
\section{A restricted lower sieve followed by Buchstab decomposition}
\label{sec:buchstab-restricted-lower}

Put
\[
 \nu_* =\frac{249}{448},\qquad w_* =\frac{25}{224},
 \qquad h=\frac{50}{309},\qquad a=\frac{100}{1149}.
\]
Thus $\nu_*=(1+w_*)/2$, $\nu_*/w_*=249/50$, and
$a<w_*<h<\mu=25/128$. This section improves only the two
initial lower terms of the complete Chen functional. All
upper switching terms and quotient domains are unchanged.

\begin{lemma}[Uniform level for a restricted lower sieve]
\label{lem:restricted-lower-sharp-level}
For every fixed $0<u\le h$, the lower-sieve conclusion in
Lemma~\ref{lem:restricted-lower-level} holds at every total
level $\nu_* -\varepsilon$ with sufficiently small fixed $\varepsilon>0$.
Consequently its limiting lower coefficient is
\[
 \mathcal L_*(u)=\frac{2e^{-\gamma}}u f(\nu_*/u)
               =\frac4u\bar f(\nu_*/u),
 \qquad \bar f=f/(2e^\gamma).
\]
The statement retains the source factorization repairs in
Section~\ref{sec:two-factor-greedy-repair}.
\end{lemma}

\begin{proof}
Only the factorization bound in the proof of
Lemma~\ref{lem:restricted-lower-level} needs strengthening.
For a rough atom of rank at least three write its first
three exponents as $0<t_3\le t_2\le t_1\le h+O(\eta)$.
They lie strictly below $\mu$ at sufficiently small precision.
Lemma~\ref{lem:variable-envelope-simplification} gives
\[
 \Phi(t_1,t_2,t_3)
 \ge\max\left\{\frac{1+t_1}{2},
                      \frac{153}{256}-\frac{3t_3}{8}\right\}
 \ge\max\left\{\frac{1+t_1}{2},
                      \frac{153}{256}-\frac{3t_1}{8}\right\}
 \ge\frac{249}{448}.
\]
The two affine expressions intersect at $t_1=w_*$.
The branch-rescue argument in that lemma is essential: the
second affine expression need not be an active raw branch
by itself, but the maximum with the single-variable branch
is a valid lower bound for the full envelope.

Choose the source factorization level strictly between
$\nu_* -\varepsilon$ and $\nu_*$, and then choose all atom,
Type II, local-scale and small-prime precisions smaller than
its surplus. The $O(\kappa)$ perturbation to the moving-residue
parameter $\alpha=1+3\kappa$ is absorbed in the same surplus.
The complete support remains below $N^{3/5-c}$.
Ranks zero through two have support at most
$N^{2h+O(\eta)}<N^{1/2-c_0}$ and use ordinary BV as before.
All signed atoms, their scalar factors, residual coefficients,
and the uniform main-density argument are unchanged.
Thus the cited proof supplies the lower multiplier at the new
level. Take $N\to\infty$ first, then the fixed construction
errors to zero, and finally $\varepsilon\downarrow0$.
Continuity of $f$ gives the displayed limiting coefficient.
\end{proof}

The uniform pointwise envelope bound cannot be increased over
this entire ordered cube: at
$t_1=t_2=t_3=w_*$ every branch in the simplified envelope is
at most $\nu_*$, and the first, pair and affine branches attain
it. The square-prefix conditions hold strictly there. This
observation concerns that envelope bound, not an impossibility
theorem for other sieve constructions.

\begin{proposition}[Buchstab improvement of the second lower term]
\label{prop:buchstab-restricted-lower}
For fixed $w_*\le u\le h$, the shifted prime sequence has the
limiting lower coefficient
\begin{equation}\label{eq:buchstab-restricted-coefficient}
 \mathcal B_*(u)
 =2e^{-\gamma}\left\{
    (1/u+1)f\big((1/u+1)/2\big)-f(249/50)\right\}.
\end{equation}
In particular the permissible lower constants for the first
two terms of the Chen weight are
\begin{align}
 L_1&\ge\frac4a\bar f\left(\frac{286101}{44800}\right),
 \label{eq:buchstab-new-L1}\\
 L_2&\ge8\log(259/100)-4\bar f(249/50).
 \label{eq:buchstab-new-L2}
\end{align}
Each displayed inequality denotes a lower main coefficient;
the corresponding counts allow an $o(1)$ loss. No double-sieve
function $H$ or $h$ is used here.
\end{proposition}

\begin{proof}
Let $\mathcal P_N$ be the primes not dividing $N$, and define
\[
 T_N(z)=\#\{P<N:P\in\PP,\ (N-P,P_N(z))=1\},
 \qquad P_N(z)=\prod_{\substack{\ell<z\\\ell\in\mathcal P_N}}\ell.
\]
Let $T_N(\ell;z)$ be the same count with $\ell\mid N-P$
and the quotient $(N-P)/\ell$ free of primes in $P_N(z)$.
For $z\le\ell$ these are the usual prefixed sifted counts.
The exact least-prime decomposition gives
\begin{equation}\label{eq:buchstab-restricted-identity}
 T_N(N^u)=T_N(N^{w_*})
  -\sum_{\substack{N^{w_*}\le\ell<N^u\\\ell\in\mathcal P_N}}
                       T_N(\ell;\ell).
\end{equation}
No squarefree assumption is needed for this identity.

The first term has lower coefficient
$2e^{-\gamma}f(\nu_*/w_*)/w_*$ by the preceding lemma.
For the sum, use the averaged one-prime prefix interface
of Lemma~\ref{lem:prefix-interface}, approaching the total level
$v(t)=(1+t)/2$ for $w_*\le t\le h$.
Here the residual cutoff is the prefixed prime itself;
all residual primes are smaller than the named largest prime.
The necessary first square-prefix condition has fixed surplus:
\[
 v(t)-2t=\frac{1-3t}{2}\ge\frac{1-3h}{2}>0.
\]
Later square-prefix conditions follow from the residual
modified sieve as in the existing prefix proof. Empty residual
rough atoms have support $N^{t+O(\eta)}<N^{1/2-c}$ and use BV.
The single-variable source assertion imposes no additional
largest-prime domain beyond the checked range of its formula.

For clarity, this is an estimate after averaging a positive
prime interval, not a distribution theorem for each fixed
$\ell$. Partition $[w_*,u]$ into finitely many closed exponent
boxes. Fix also $0<\xi\ll w_*$. In a box $[t,t+\Delta]$,
replace the cutoff $\ell$ by $N^{t-\xi}$ to obtain an upper
bound. This enforces the strict cutoff inequality $u<t_1$
in the statement of the existing interface. Choose the total level below
the box infimum of $v$ by a fixed margin, and subtract the
largest possible prefix exponent to define its residual
level. The actual prefixed atoms then meet precisely the
strict hypotheses above. The uniform moving-residue estimate,
positive prime-box normalization and the Mertens factor in
the existing interface give the following limit: first take
$N\to\infty$ with all parameters fixed, next let the prefix
mesh and sieve precision tend to zero within the fixed margins,
then let those margins tend to zero, and finally let $\xi\downarrow0$.
Summing the boxes yields
\[
 \limsup_{\substack{N\to\infty\\2\mid N}}
 \frac1{\HN}
 \sum_{\substack{N^{w_*}\le\ell<N^u\\\ell\in\mathcal P_N}}
 T_N(\ell;\ell)
 \le 2e^{-\gamma}
       \int_{w_*}^{u}\frac{F((1-t)/(2t))}{t^2}\,dt.
\]
Before $\xi$ is removed, the main-density denominator is
$t(t-\xi)$ and the sieve argument has denominator $t-\xi$.
They converge uniformly to the displayed integrand. Throughout
these limits the sieve arguments remain greater than $2$,
since their limiting minimum is $(1-h)/(2h)=259/100$.
All arguments range over a compact interval bounded away
from zero. The sum of reciprocal masses of omitted primes
dividing $N$ is $o(1)$ at these fixed positive cutoffs.
Thus there is no singular-cutoff limit or interchange of an
unbounded number of boxes with $N$ in this step.

Put $s(t)=(1/t+1)/2$. The delay equation
$(s f(s))'=F(s-1)$ implies the exact antiderivative
\[
 \frac{d}{dt}\{(1/t+1)f(s(t))\}
       =-\frac{F((1-t)/(2t))}{t^2}.
\]
Since $s(w_*)=\nu_*/w_*=249/50$, subtraction in
\eqref{eq:buchstab-restricted-identity} yields
\eqref{eq:buchstab-restricted-coefficient}.
For $u=h$, $s(h)=359/100\in(2,4)$ and
$\bar f(s)=\log(s-1)/s$, giving
\eqref{eq:buchstab-new-L2}. The first lower term uses the
preceding uniform lemma directly at $u=a$, giving
\eqref{eq:buchstab-new-L1}.
\end{proof}

The choice of $w_*$ also has an exact variational explanation
within this particular combination of inputs. Starting at
$w\le u$ with the uniform-level coefficient
$2e^{-\gamma}f(\nu_*/w)/w$ and making the same subtraction,
the derivative of the resulting lower coefficient with
respect to $w$, divided by $2e^{-\gamma}$, is
\[
 \frac{F((1-w)/(2w))-F(\nu_*/w-1)}{w^2}.
\]
It is nonnegative for $w\le2\nu_*-1$ and nonpositive for
$w\ge2\nu_*-1$, by the monotonicity of $F$.
Thus $w_*$ maximizes this derived bound in its admissible
interval. This optimization uses the same two already
justified estimates, rather than a mixed double-sieve saving.
\subsection{Certified consequence of the lower-sieve improvement}
\label{sec:buchstab-lower-corollary}

The integer interval calculation in
\texttt{certify\_lower\_prefix\_buchstab.py} encloses the two
expressions in \eqref{eq:buchstab-new-L1}--\eqref{eq:buchstab-new-L2}.
In particular,
\[
 L_1\ge
 \frac{15597450959322408038477303}{1208925819614629174706176}
 >12.9019,
\]
and
\[
 L_2\ge
 \frac{1962215427253960434935021}{302231454903657293676544}
 >6.4924.
\]
The certificate uses the previously reviewed outward sieve-function
table at the exact arguments $286101/44800$ and $249/50$.
The logarithm is enclosed by integer interval arithmetic.
A separate check recomputes it with a rational $50$-term
atanh series and explicit tail, and reproduces the sieve-table
endpoints. These computations certify real-valued expressions;
their interpretation as lower bounds uses
Proposition~\ref{prop:buchstab-restricted-lower}.

\begin{theorem}[Separate Chen count and imbalanced representation]
\label{thm:buchstab-lower-chen}
Under the derived lower and upper interfaces in
Propositions~\ref{prop:buchstab-restricted-lower} and
\ref{prop:extended-prefix-upper}, every sufficiently large even
$N$ satisfies
\[
 D_{1,2}(N)\ge\frac{1817}{1000}\HN,
 \qquad D_{383/500}(N)\ge\frac1{1000}\HN.
\]
Thus the separate research derivation gives
$N=p+rq$ with $r=1$ or prime and
$r\le q^{383/500}=q^{0.766}$.
It does not use the external Chen coefficient $1.9728$ or
any double-sieve correction.
\end{theorem}

\begin{proof}
Replace only the first two lower rows in the table of
Theorem~\ref{thm:extended-prefix-chen}. All other rows and
the full eleven-term pointwise quotient weight are retained.
The increase in its lower coefficient is exactly
\[
 \frac{3(12.9019-12.9014)+(6.4924-6.3640)}4
 =\frac{1299}{40000}.
\]
The resulting table therefore supplies
\[
 \frac{71413}{40000}+\frac{1299}{40000}
 =\frac{9089}{5000}
 =\frac{1817}{1000}+\frac1{1250}.
\]
The positive slack absorbs the asymptotic errors in the Chen
count. The existing bad-semiprime bound has coefficient
$681/100$. At $u=117/883$, the elementary upper bound
\[
 \log\frac{500}{383}
 <2\left(u+\frac{u^3}{3}+\frac{u^5}{5(1-u^2)}\right)
 =\frac{351453088200357}{1318411216105000}
\]
gives
\[
 \frac{1817}{1000}-\frac{681}{100}\log\frac{500}{383}
 >\frac1{1000}
  +\frac{83923780224883}{131841121610500000}.
\]
This positive rational gap absorbs the remaining asymptotic
errors in the subtraction. The complete source-linked numerical
recombination is saved in \texttt{certify\_buchstab\_transfer.json}.
\end{proof}

The stronger $b=0.7485$ implication using the external
coefficient $1.9728$ is unchanged. The present improvement
strengthens the separate route, but does not establish that
it exceeds that external-input implication or every known
result of the same quantifier. All conclusions remain in a
research proof draft subject to the explicitly named inputs.

% Optional classical double-sieve increment, separated from the quoted theorem.
The additional increment below applies because neither of the preceding
lower-term improvements changes $U_{10}$. It is not used in the numerical
constant or exponent of Theorem~\ref{thm:buchstab-lower-chen}; that theorem
retains its derivation without double-sieve corrections. The comparisons
with the older lower bounds below also exclude improvement over their
stronger replacements.

% Integration fragment, 2026-10-08.  Main manuscript and state not edited.
% The arithmetic certificate does not certify the cited analytic inputs.

\subsection{A quantitative use of the classical double sieve}
\label{sec:safe-double-sieve-gain}

The safe minimum/maximum construction has a concrete, but small,
consequence for the present eleven-term functional.  The published
step tables do not improve its already enhanced one- and two-prefix
linear-sieve terms.  They do improve a positive-volume part of the
four-prefix upper term $U_{10}$, which presently uses only switching.
No old-level correction is added to a new-level linear function here.

Put
\[
 a=\frac{100}{1149},\quad h=\frac{50}{309},\quad
 d_*=\frac12-3a,\quad K=\frac{681}{100},\quad
 b=\frac{12}{125},\quad H_0=\frac{223939}{10^7}.
\]
Here $h$ without an argument denotes the cutoff.  The double-sieve
functions with arguments are those of
\href{https://www.impan.pl/shop/en/publication/transaction/download/product/81935}
{Wu, Acta Arithmetica 131 (2008), Section 3 and p.~376, Tables 1--2}.
In particular,
\[
 H(2)\ge H(2.2)\ge H_0.
\]
The table is used as a published input, rather than claimed to have
been independently rederived.  The published entries
$h(3)\ge0.0077162$ and $h(4.8)\ge0.0001336$ are used; these are
smaller than the two corresponding transcriptions in Li II (18).
Every terminating decimal in the present comparison is an exact
rational number.

\begin{proposition}[A positive classical-level four-prefix gain]
\label{prop:safe-double-sieve-u10}
Assume the existing switching bound with coefficient $K=681/100$
on the positive prefix regions of $\Upsilon_{10}$, and Wu's classical
double-sieve bound for its class $\mathcal U_k(N)$.
The upper constant for $\Upsilon_{10}$ can be reduced, and hence the
current eleven-term Chen functional increased, by at least
\begin{equation}\label{eq:safe-u10-gain}
 \Delta c_\Box
 =\frac{685170246954473}{491297300191433880000}
 >\frac1{720000}.
\end{equation}
This is an increment to the explicit functional, with the usual
asymptotic error; it is not a newly rounded eventual counting constant.
\end{proposition}

\begin{proof}
Use the ascending exponents from the quotient weight and restrict to
\[
 \mathcal R=\{a<t_1<t_2<t_3<t_4<b\},\qquad
 S=t_1+t_2+t_3+t_4.
\]
Write $q=p_1p_2p_3p_4$ for the specified prefix and
$Q=N^{1/2-\delta}$.  Every ordered prime vector in this region satisfies
Wu's square conditions with a fixed margin, since
\[
 t_4+\cdots+t_{j+1}+2t_j\le5b=\frac{12}{25}<\frac12.
\]
For sufficiently small fixed $\delta>0$ this also proves membership
in the class at total level $Q$.  All prime slots exceed its auxiliary
roughness threshold, upon choosing the fixed construction parameters
sufficiently small.

There is no moving-cutoff difficulty in this region.  Indeed,
\[
 \frac{\log (Q/q)^{1/2}}{\log N}
 \le\frac14-2a<a<t_2.
\]
Consequently the original quotient-sifted count at $p_2$ is at most
the count sifted at $(Q/q)^{1/2}$, with the specified primes removed
from the sieving set.  This last removal can only enlarge an upper
bound, including on nonsquarefree exceptional integers.
We may apply Wu's bound at the fixed parameter $s=2$, for which
$A(2)=2F(2)/(2e^\gamma)=1$.

Here is the precise common-prefix argument.  Expand the prime
indicator of $[N^a,N^b)$ into Wu's short intervals with multiplicative
ratio $1+O((\log N)^{-4})$, and take its fourth convolution power.
Every resulting positive convolution belongs to $\mathcal U_k(N)$
for $k\ge4$, after ordering the interval endpoints; all its square
conditions follow from $5b<1/2$.  Apply (3.6) at $s=2$ to each
convolution and divide by $4!$.  Distinct tuples are counted exactly
$4!$ times before this division.  Repeated slots or repeated interval
labels cause no factorial ambiguity: the convolution identity is
used before division, and the whole expression has nonnegative
coefficients.  Diagonal prime tuples have negligible density, since
their reciprocal prime-square sum is $O(N^{-a})$.
The source bound is uniform over its prefix class, so the positive
sum of these bounds is legitimate.  First take $N\to\infty$ with
all margins fixed, and then remove the margins.

Wu's normalization
\[
 \Theta(N,\sigma)=4\operatorname{li}(N)
   \sum_q\frac{\sigma(q)C_{qN}}{\varphi(q)\log(Q/q)}
\]
therefore supplies the local upper coefficient
\begin{equation}\label{eq:safe-u10-old-kernel}
 \frac{4(1-H_0)}{1/2-S}\,
       \frac{dt_1\,dt_2\,dt_3\,dt_4}{t_1t_2t_3t_4}
\end{equation}
in units of $\HN$.  For squarefree $q$ coprime to $N$, the source
local factor is $\prod_{p\mid q}(p-2)^{-1}$, which can be replaced
by $1/q$ here with relative error $o(1)$.  Omitting the primes
dividing $N$ also costs $o(1)$ uniformly on this fixed positive
exponent range.  These observations retain the full singular series.

The existing switching coefficient on the same region is
\[
 \frac{K}{t_2}\omega\left(\frac{1-S}{t_2}\right)
       \frac{dt_1\,dt_2\,dt_3\,dt_4}{t_1t_2t_3t_4}.
\]
One has $\omega(u)\ge1/2$ for $u\ge1$: this holds directly on
$[1,2]$ and follows on successive unit intervals from
$(u\omega(u))'=\omega(u-1)$.  Also $(1-4b)/b>2$.
Thus switching minus \eqref{eq:safe-u10-old-kernel}, before the common
measure, is at least
\[
 \frac{K}{2b}-\frac{4(1-H_0)}{1/2-4b}
   =\frac{1019753}{580000}>0.
\]
Use switching on the complement of $\mathcal R$ and the old-level
bound on $\mathcal R$.  This is a choice between bounds for the same
positive count.  The existing switching proof allows the additional
prime-interval restriction; alternatively split into its positive
boxes and discard the measure-zero ordering boundaries.
No estimate for an individual specified prime is asserted.

Finally,
\[
 \int_{\mathcal R}\frac{dt_1\,dt_2\,dt_3\,dt_4}{t_1t_2t_3t_4}
 \ge\frac{(b-a)^4}{24b^4}
 =\frac{1343796482}{423532155337443}.
\]
The coefficient of $-U_{10}$ in the Chen functional is $1/4$.
Multiplying the last two rational bounds and dividing by four
gives \eqref{eq:safe-u10-gain}.
\end{proof}

The ordinary classical upper sieve alone already beats switching
throughout $\mathcal R$, since
\[
 \frac{K}{2b}-\frac4{1/2-4b}=\frac{915}{928}>0.
\]
Its guaranteed functional increment is
$204928963505/262025226768764736$.
After making this ordinary minimum choice, Wu's $H_0$ gives a
further, separately positive increment of at least
\[
 \frac{H_0}{1/2-4a}\frac{(b-a)^4}{24b^4}
 =\frac{150464220191299}{321611667129607500000}
 >\frac1{2200000}.
\]
This is a genuine double-sieve contribution.  It is already included
in \eqref{eq:safe-u10-gain}; the two displayed gain bounds must not
be added to each other.  Outward overestimates used in the existing
global switching certificate do not obstruct subtraction of the
proved local gain.

\paragraph{Why the printed tables do not improve the other current
linear-prefix terms.}
Let $H_{\rm tab},h_{\rm tab}$ be the decreasing step lower bounds
formed from Wu's published tables.  For $h_{\rm tab}$ use zero
below $2$ or beyond the last tabulated point.  The claims below
concern these chosen lower bounds, not upper bounds for the true
functions $H,h$.  Set
\[
 \mathcal F(s)=\frac{F(s)}{2e^\gamma},\qquad
 \mathcal f(s)=\frac{f(s)}{2e^\gamma}.
\]
For $U_3,U_4$, Wu's one-prefix class requires $t<1/4$.
On $a\le t\le1/4$,
\[
 s_0=\frac{1/2-t}{a}\in[1149/400,949/200],\qquad
 s_1-s_0=\frac{v(t)-1/2}{a}\ge\frac{1149}{6400}.
\]
The exact rational interval calculation in
\texttt{certify\_safe\_double\_sieve\_gain.py} proves throughout this
interval that
\[
 \mathcal F(s_0)-\mathcal F(s_0+1149/6400)
       -\frac{H_{\rm tab}(s_0)}{s_0}>\frac1{4000}.
\]
It uses the exact formula $\mathcal F(s)=(1+J_0(s))/s$ for $s\le5$,
4,096 interval cells, and 60,000 outward Darboux cells for $J_0$.
Monotonicity of $F$ then excludes a positive safe residual saving.
The same comparison also covers the current conservative replacement
by $F(5)$ where the new parameter exceeds five.  For $t>1/4$
the old one-prefix class is unavailable, even where a classical
ordinary upper sieve is used.

For $L_5^{(2)},L_6$, wherever $s_0\ge2$ the largest old parameter is
$s_0\le749/200$, and the existing total level is at least $v(t)$.
Its argument gain is bounded below by
\[
 s_1-s_0\ge\frac{v(d_*)-1/2}{a}=\frac{2373}{6400}=:\xi.
\]
For a short hand-check put $g(s)=\log(s-1)/s$.  On the three old
parameter ranges $[2,3]$, $[3,7/2]$, $[7/2,749/200]$, respectively,
bound $s_0+\xi$ by $27/8,31/8,33/8$ and use
\[
 \log(19/8)<7/8,\qquad \log(23/8)<17/16,\qquad
 \log(25/8)<8/7.
\]
The formula
$g'(s)=\{s/(s-1)-\log(s-1)\}/s^2$, together with
$\mathcal f(s_0)=g(s_0)$ and $\mathcal f(s_0+\xi)\ge g(s_0+\xi)$,
gives lower bounds for
$\mathcal f(s_1)-\mathcal f(s_0)-h_{\rm tab}(s_0)/s_0$ of
\[
 \frac{113680691}{18468000000},\qquad
 \frac{92697877}{20721562500},\qquad
 \frac{38107351}{12705000000},
\]
all greater than $29/10000$.  The step caps used are
$h(2)$, $h(3)$ and $h(3.5)$ from the printed table, respectively.
For $s_0<2$ those entries cannot be extrapolated: the source says
$h$ is increasing there, so only $h_{\rm tab}=0$ is retained.

For $L_2$, the arguments are $309/100$ and $3399/1000$.
On this interval $g'(s)>8/245$ follows from $\log(5/2)<1$.
Subtracting $0.0066236/(309/100)$ still leaves a normalized gap
greater than $7/1000$.  For $L_1$, the old argument is $5.745$,
beyond the $h$ table.  Even the positive extension obtained by
integrating the printed $H$ steps from $s_0-1$ to $4.9$ leaves
the old lower bound more than $1/1000$ below the existing certified
$L_1\ge12.9014$.  This last comparison is enclosed in the same
certificate, using the exact lower-sieve formula at the old parameter.
It does not bound the unknown true value of $h(5.745)$ above.

\paragraph{The remaining prefix classes.}
The entire old-level classes needed for $U_7,U_8,U_9,U_{11}$ fail,
respectively, because a required square exponent is at least
\[
 3d_*>1/2,\qquad 2/3>1/2,\qquad
 d_*+2h>1/2,\qquad h+4a>1/2.
\]
For the last inequality, order the four specified primes decreasingly
and examine the final square condition.  Thus the published $H/h$
tables cannot simply be attached to these switched sequences.
The valid four-prefix subregion of $U_{10}$ proved above is the
positive exception found in this comparison.

The saved \texttt{safe-double-sieve-gain.json} records the exact table
entries, rational exclusions and gain, and hashes of its arithmetic
dependencies.  It certifies these numerical comparisons only.
It neither verifies Wu's table derivation nor resolves the separate
cross-level issue behind Li's quoted $1.9728$.

\section{Further analytic refinements and route tests}
The following constructions provide new integral options and clarify a
limitation of a proposed double-sieve shortcut. They are not included
in the rounded coefficients or imbalance exponents stated earlier.
The saved floating-point evaluations guide further work; they are not
outward numerical certificates. The limitation on a universal mixed
bound below does not decide the more restricted formulas in Li's paper.

% Derived analytic integration fragment, 2026-10-08.
% No numerical integration or new global Chen constant is certified here.
% Uses the existing prefix interface and the selected-subproduct repair.

\subsection{A safe ordinary upper option for four specified primes}
\label{sec:four-prefix-safe-minimum}

Fix $0<a<h<d<1/3$.  Write the four specified prime exponents in
decreasing order as $t=(x,y,z,w)$ and put
\[
 S=x+y+z+w,\quad
 W(t)=\max\{2x,x+2y,x+y+2z,x+y+z+2w\},
 \qquad d\mu(t)=\frac{dx\,dy\,dz\,dw}{xyzw}.
\]
The two tuple domains, with the full fourth-prime range, are
\[
 \Omega_{10}=\{a<w<z<y<x<h\},\qquad
 \Omega_{11}=\{a<w<z<y<h<x<d\}.
\]
The quotient in each original count is sifted at the second smallest
specified prime, whose exponent is $z$.  The four specified primes
need not be the four largest prime factors of that quotient count.
They will, however, be the four largest primes in every modulus atom
of the ordinary upper sieve constructed below.  These are different
statements, and only the latter is used in factorization.

Put
\[
 \mu=\frac{25}{128},\quad \Lambda=\frac{153}{256},\quad
 v(s)=\min\left\{\frac{1+s}{2},\frac{57}{64}-\frac32s\right\},
 \qquad B(z)=\Lambda-\frac38z.
\]
For each nonempty $I\subseteq\{1,2,3,4\}$ define
\begin{align*}
 s_I(t)&=\sum_{i\in I}t_i,\\
 G_I(t)&=\max_{j\notin I}
       \left(s_I(t)+\sum_{\substack{k<j\\k\notin I}}t_k+2t_j\right),
       \qquad \max\varnothing=0.
\end{align*}
Include the endpoint $v(s_I(t))$ in the finite set $\mathcal C(t)$
if and only if
\begin{equation}\label{eq:four-prefix-subset-gate}
 v(s_I(t))>\max\{1/2,W(t),G_I(t)\}.
\end{equation}
Also include $B(z)$ if
\begin{equation}\label{eq:four-prefix-W-gate}
 x<\mu,\qquad B(z)>\max\{1/2,W(t),v(x)\}.
\end{equation}
Define $\mathcal E=\{t:\mathcal C(t)\ne\varnothing\}$ and
$L_4(t)=\max\mathcal C(t)$ on $\mathcal E$.  This is a finite
piecewise affine prescription, with strict eligibility gates.  All
its endpoints are at most $\Lambda<3/5$.

Each actual sieve level $\nu$ used for a selected endpoint must
satisfy
\[
 \max\{1/2,W(t),G_I(t)\}<\nu<v(s_I(t)),
\]
or the analogous interval between the two sides of
\eqref{eq:four-prefix-W-gate}.  An endpoint cap does not authorize
levels below its gate.  On a compact subregion of these strict
inequalities there is a fixed positive margin on every side.
Above $x=\mu$ only the explicitly tested selected-subproduct branches
are used.  In particular the full raw eight-branch envelope is not
extended beyond its proved largest-prime range.

\begin{lemma}[Averaged ordinary upper sieve with four prime slots]
\label{lem:four-prefix-upper-interface}
On any compact positive prefix box contained in a strict branch of
$\mathcal E$, choose $\nu$ between its gate and cap, with fixed
margins.  The ordinary residual upper-sieve weights of level
$N^{\nu-S+o(1)}$, sifted below the smallest specified prime, have
the averaged distribution remainder of
Lemma~\ref{lem:prefix-interface}, now with four specified primes.
Consequently their limiting upper coefficient at the prefix point
$t$ is
\begin{equation}\label{eq:four-prefix-ordinary-kernel}
 \mathcal O_\nu(t)
 =\frac4{\nu-S}\,
       \mathcal A\left(\frac{\nu-S}{w}\right),
 \qquad
 \mathcal A(s)=\frac{sF(s)}{2e^\gamma}\quad(s\ge2),
 \quad \mathcal A(s)=1\quad(0<s<2).
\end{equation}
The coefficient is in units of $\HN$ and is integrated against
$d\mu(t)$.  The assertion is averaged over the prefix box and does
not assert an estimate for each particular quadruple of primes.
\end{lemma}

\begin{proof}
The required change to Lemma~\ref{lem:prefix-interface} is at the
modulus atom, including its factorization; it is not a substitution
of $r=4$ in a theorem stated only for $r=1,2$.
Choose four disjoint ordered intervals
$I_i=(N^{\ell_i},N^{u_i}]$, a fixed level $\nu$, and
$T_+=\sum_i u_i$.  Make the intervals short enough that the upper
endpoint vectors satisfy every required square test with margin
$20\delta$, and that the entire chosen cap exceeds $\nu$ with
margin $20\delta$.  Set
\[
 R=N^{\nu-T_+-20\delta},\qquad
 Z=N^{u_*},\qquad
 u_* =\min\left\{\ell_4-\xi,
                 \frac{\nu-T_+-20\delta}{2+\rho}\right\}>0,
\]
where $\xi,\rho>0$ are fixed and sufficiently small.  Thus $Z$ is
below every specified prime and below the original quotient cutoff.
Lowering the cutoff gives a pointwise upper bound for the original
positive quotient count.  Also $\log R/\log Z\ge2+\rho$; no
linear-sieve formula with an inadmissible small parameter is invoked.

Construct the ordinary modified residual upper weights $\rho_R^+$
independently of the particular primes inside the four intervals.
Use the atoms of
\href{https://arxiv.org/html/2109.02851v2}{Lichtman 2021,
(5.14)--(5.18) and Theorem 2.10}, at residual level $R$.
Convolve each residual atom with
$\chi_{I_1}*\chi_{I_2}*\chi_{I_3}*\chi_{I_4}$.
Choose the construction precision $\eta$ small enough that its
presieving threshold $z_0=N^{c(\eta)}$ is below $Z$ and every
specified interval.  Every residual rough prime is below $Z$.
The specified slots are therefore exactly the largest four slots
in the concatenated rough vector.  There is no residual-prime
interleaving to be justified.

The four initial $\mathcal D^{\rm well}$ inequalities hold because
$W(t)<\nu$ with the prescribed margin.  For every later rough prime
$q_j$, residual well support gives
\[
 (p_1p_2p_3p_4)q_1\cdots q_{j-1}q_j^2
 \le N^{T_+}R^{1+O(\eta)}\le N^{\nu-19\delta}
\]
after decreasing $\eta$.  Thus the complete rough modulus is in
the required buffered $\mathcal D^{\rm well}$ support.  This check
also covers the residual atom with no rough slots: that atom still
has four specified rough slots, and is not reclassified as a
rank-zero modulus of level below $1/2$.

For a selected branch put $b_0=\prod_{i\in I}p_i$.
The inverse critical-interval identity in
Lemma~\ref{lem:selected-subproduct-factorization} applies because
$\nu<v(s_I)$ with a margin.  For each unselected specified slot its
remaining-square inequality is exactly the corresponding term in
$G_I<\nu$.  Before any subsequent residual prime all selected and
unselected specified slots have occurred, so its remaining-square
inequality is exactly the complete-vector inequality just checked.
The selected-subproduct lemma supplies a factorization for every
allowed Type II parameter $A$, including its $C<1$ case by
Lemma~\ref{lem:two-factor-greedy-repair}.  For the extra $B(z)$
branch use the repaired $W$ construction: on $x<\mu$ the inner
$W(x,y,z)$ cap is $B(z)$ and its gate is $v(x)$, as verified in
Lemma~\ref{lem:variable-envelope-simplification}.  Its hypotheses
and the complete well support have fixed margins as well.

Here are the coefficient details needed for this transfer.  The
factorization is performed on the upper endpoint vector, which
gives a fixed partition of its prime slots valid for every tuple
in the rough box.  The selected lemma explicitly supplies this
uniform partition on both sides of $C=1$.  Retain each original
scalar coefficient of the residual sieve atom, its signs, and its
reciprocal factorials for repeated residual interval labels.
No label crosses from a specified slot to a residual slot, because
$Z<\inf I_4$.  If a factor group contains $m_i$ rough slots,
divide its tuple convolution by $m_i!$.  The resulting coefficient
is $1$-bounded, even before removing nonsquarefree products; multiply
the original scalar by the product of these group factorials.
This is an exact convolution identity, and does not introduce a
second copy of the original repeated-label normalization.
The finite number of atoms and the sum of the absolute values of
their new scalars depend only on the fixed construction precision.

Combine the original signed small-prime weight $\beta$ with the
third factor.  Its support is $e\le N^{O(\eta)}$ on primes below
$z_0$, whereas that factor's rough primes exceed $z_0$.
Unique small/rough factorization proves that its coefficient stays
$1$-bounded.  Attaching $e$ to the third factor multiplies the three
factorization monomials by $e^2,e^2,e^{2+2\sigma}$, respectively;
the fixed slack pays these losses when $\eta\ll\delta$.
As in the existing prefix interface, allowing unrestricted prime
convolutions and then removing rough prime-square collisions costs
\[
 \ll (N+N^{\nu+O(\eta)})z_0^{-1}(\log N)^{C}
   =O_A(N(\log N)^{-A}).
\]
Thus neither squarefreeness nor signed atoms are replaced by an
unproved positivity assertion.

The distribution step is now precisely the existing atom argument.
The three factorization inequalities feed
\href{https://arxiv.org/html/2309.08522v1}{Lichtman 2023,
Proposition 5.2}, with $A_0=UX^{-\zeta}$ in the Type II range
on a local prime scale $X$.
The smallest saving after replacing $A_0$ by $U$ is
$(1-2\sigma)\zeta=9\zeta/16$.
The rough Siegel--Walfisz and bounded-complement steps, and the
one-long and two-long cases, are exactly those of the existing
prefix interface, with the corrected exact coprime mean of
Section~\ref{sec:one-long-coprime-mean}.
They only require a fixed number of normalized bounded atom factors
and total support below $X^{3/5}$, both just checked; they do not
require the number of specified slots to be at most two.

For clarity the moving residue is handled on the actual prime
scale.  Discard primes $P\le N^{1-\kappa}$, at cost
$N^{1-\kappa+o(1)}$ with the fixed divisor weights, and use dyadic
$P\asymp X$.  Then $\alpha=\log N/\log X=1+O(\kappa)$.
The strict cap and gate margins absorb this change and the
simultaneous rescaling of all exponent coordinates.  For the
selected branches this is the verified
$V_{\alpha,7/32}$ interval, not an unperturbed largest-prime
envelope.  The $W$ branch also retains $x<\mu$ with a fixed
margin before this change of scale.  Choose $\kappa$ small enough
that $N<X^{1+\epsilon}$; Proposition 5.2 explicitly permits this
moving residue.  This proves the uniform averaged remainder
\[
 \sum_{\substack{p_i\in I_i\\(p_1p_2p_3p_4,N)=1}}
 \sum_{\substack{q\mid P(Z)\\(q,N)=1}}
 \rho_R^+(q)E_N(p_1p_2p_3p_4q)
 =O_A(N(\log N)^{-A}).
\]

For the main term keep the full presieve factor $V_N(z_0)$, as in
the final part of Lemma~\ref{lem:prefix-interface}.  The one-sided
dimension-one estimate below $z_0$ is uniform in even $N$; above
$z_0$ the omitted primes dividing $N$ have total reciprocal weight
$o(1)$.  Consequently the main term is
\[
 \pi(N)V_N(Z)\bigl(F(\log R/\log Z)+o_\eta(1)\bigr)
 \sum_{\substack{p_i\in I_i\\(p_1p_2p_3p_4,N)=1}}
       \frac1{\varphi(p_1p_2p_3p_4)}.
\]
Here $V_N(Z)=(2e^{-\gamma}+o(1))C(N)/(u_*\log N)$ uniformly,
and the reciprocal prime sums tend to the product measure $d\mu$.
First let $N\to\infty$ at fixed margins and precision, then refine
the boxes and precision within those margins, and finally send
$\delta,\xi,\rho$ to zero.  The limiting cutoff exponent is
$\min\{w,(\nu-S)/2\}$.  If $(\nu-S)/w\ge2$, this gives
$2e^{-\gamma}F((\nu-S)/w)/w$; otherwise lowering gives exactly
$4/(\nu-S)$.  These are \eqref{eq:four-prefix-ordinary-kernel}.
\end{proof}

Define $\mathcal O(t)=\mathcal O_{L_4(t)}(t)$ on $\mathcal E$,
and $\mathcal O(t)=+\infty$ elsewhere.  The endpoint here is a limit
from strictly admissible levels.  Since $F$ is nonincreasing and
$4/(\nu-S)$ decreases in $\nu$, the largest eligible endpoint
gives the smallest ordinary option.  Moreover $W\ge S+w$, so
$L_4-S>w>a$ on $\mathcal E$; no singularity from zero residual
level occurs at its boundary.

\begin{proposition}[Four-prefix safe minimum]
\label{prop:four-prefix-safe-minimum}
Assume the existing local switching upper bound with coefficient
$K$ on positive prefix boxes of $\Omega_j$, for $j=10,11$.
Assume also $d+4h<1$, so the following Buchstab arguments exceed
one throughout the two domains.  Then a valid upper constant is
\begin{equation}\label{eq:four-prefix-safe-minimum}
 U_j^{\rm safe}
 =\int_{\Omega_j}
   \min\left\{\frac Kz\omega\left(\frac{1-S}{z}\right),
                   \mathcal O(t)\right\}\,d\mu(t),
 \qquad \Upsilon_j\le (U_j^{\rm safe}+o(1))\HN.
\end{equation}
Here $\Upsilon_{11}$ means its full fourth-prime range.  If
$a>1/12$, the ordinary option has the simpler exact form
\begin{equation}\label{eq:four-prefix-rational-kernel}
 \mathcal O(t)=\frac4{L_4(t)-S}\quad(t\in\mathcal E).
\end{equation}
This proposition supplies an analytic integral to be bounded; it
does not certify a numerical value of that integral.
\end{proposition}

\begin{proof}
Take a common finite positive partition for both upper bounds.
Inside a box separated from the strict cap, gate, and ordering
boundaries, select a branch and an admissible level with fixed
margins.  Both estimates bound the same averaged original
quotient count.  Select the smaller of the two box-level upper
bounds, then sum.  This does not take the minimum of unproved
estimates for individual prime quadruples.
Use switching alone on the remaining boxes, including those
meeting a boundary.  The finite gates and endpoint comparisons
are piecewise affine; their nontrivial equality sets have measure
zero.  Remove thin strips about these sets before refining the
partition.  On these bounded domains $t_i\ge a>0$ and
$(1-S)/z>1$, so the switching kernel and its measure are bounded.
The strips therefore contribute $o(1)$ as their widths tend to
zero, by the existing local switched estimate.  Refine within the
fixed margins, approach each eligible endpoint, and then remove
the margins.  This proves \eqref{eq:four-prefix-safe-minimum}.
The parameter inequality displayed in the proposition is a simple
sufficient condition for this kernel domain: $S+z<d+4h<1$.
Other constraints needed for an entire eleven-term Chen weight
remain separate and are not inferred from this local proposition.

To prove the simplification, let $k=|I|$ and use
$S\ge s_I+(4-k)a$ and $w\ge a$.  Whenever the selected endpoint
is eligible its residual exponent is positive, so
\[
 \frac{v(s_I)-S}{w}\le
 \begin{cases}
  1/(2a)-7/2,&k=1,\\
  1/(2a)-3,&k=2,\\
  57/(64a)-17/2,&k=3,\\
  57/(64a)-10,&k=4.
 \end{cases}
\]
The first two inequalities use $v(s)\le(1+s)/2$ and
$s_I\ge ka$; the last two use $v(s)\le57/64-3s/2$.
For the $B$ endpoint,
\[
 \frac{B(z)-S}{w}\le\frac\Lambda a-\frac{35}{8}.
\]
If $a>1/12$, all five upper bounds are strictly below $3$.
At $a=1/12$ their values are $5/2,3,35/16,11/16,179/64$,
respectively.  The normalized function $\mathcal A(s)$ is exactly
$1$ for $0<s\le3$, including the cutoff-lowering convention.
This proves \eqref{eq:four-prefix-rational-kernel}.  For smaller
$a$ the general $F$ kernel remains the justified formula; no
extrapolation of $\mathcal A=1$ is made.
\end{proof}

\paragraph{Why the domain is filtered.}
The four square tests cannot be dropped even below $x=\mu$.
For the original parameters $a=100/1149$, $h=50/309$,
$d=1/2-3a$, the interior point
$(x,y,z,w)=(.155,.150,.145,.140)\in\Omega_{10}$ has
$W=.730>\Lambda$.
The point $(.200,.140,.130,.120)\in\Omega_{11}$ has
$W=.710>\Lambda$.  Neither admits an endpoint of this construction.
They are handled by switching in \eqref{eq:four-prefix-safe-minimum}.
Conversely $(.170,.110,.100,.090)\in\Omega_{11}$ admits the
subset $I=\{3,4\}$: $s_I=.190$, $v(s_I)=.595$,
$G_I=.580$, and $W=.560$.  This demonstrates why all fifteen
subsets can be useful, even though only three specified slots
appeared in the earlier envelope.  Above the largest-prime cutoff,
$(.202,.092,.090,.088)\in\Omega_{11}$ admits $I=\{2,3\}$:
its cap is $.591$, its selected gate is $.586$, and $W=.560$.
All these terminating decimals are exact rational numbers.

\paragraph{Accounting for the earlier classical double-sieve option.}
Bounds from different methods may be retained as a third entry in
the same pointwise minimum when their common positive prefix class
has been proved.  Improvements from overlapping regions may not
be added.  In fact the previous small region
$100/1149<w<z<y<x<12/125$ is already improved by the ordinary
option here beyond its stated classical $H(2)$ option.  On it
$W<5(12/125)<1/2<v(x)$, $x<\mu$, and
$v(x)-1/2>a/2>H_0\Lambda$, where $H_0=223939/10^7$.
Since $v(x)-S\le\Lambda$ and $1/2-S>0$, this proves
\[
 \frac4{v(x)-S}
 <\frac{4(1-H_0)}{1/2-S}.
\]
Thus the previously certified gain exceeding $1/720000$ is
subsumed on its region and must not be added to a computation of
\eqref{eq:four-prefix-safe-minimum}.

The distribution and switching inputs used here are the established
derived interfaces and their cited published theorems.  The saved
JSON records the exact numerical domain and dependencies.  A
floating quadrature of this integral is a diagnostic only; a new
rounded Chen coefficient requires its own outward certificate.

% Generalization of the existing eleven-term functional, 2026-10-08.
% Numerical diagnostics using this fragment are not certified bounds.
\subsection{Allowing the three weight thresholds to vary}
\label{sec:free-weight-thresholds}

The identity $d=1/2-3a$ is not used in the direct quotient-weight
proof. The following range records sufficient conditions for retaining
that weight and all of the ordinary-sieve inputs used in the current
separate calculation. Put
\[
 \mu=\frac{25}{128},\quad e=\frac{57}{224},\quad
 \nu_*=\frac{249}{448},\quad w_*=\frac{25}{224},
 \qquad \bar f=\frac{f}{2e^\gamma}.
\]
Take fixed parameters satisfying
\begin{equation}\label{eq:free-threshold-range}
 \frac1{12}<a<h<d<e,\quad w_*<h<\min(\mu,2a),\quad
 h+d>\frac13,\quad h+d<\frac12,\quad
 4a+d<1,\quad 4h+d<1.
\end{equation}
The conditions are sufficient, not claimed necessary or optimal.
All appearances of $Y_i$ and $\Upsilon_i$ in this subsection use these
new thresholds in their original quotient definitions, with the full
fourth-prime interval $h<t_4<d$.

\begin{proposition}[Parameterized ordinary-sieve functional]
\label{prop:free-weight-thresholds}
Under the already derived distribution and switching interfaces,
the eleven-term lower bound continues to hold under
\eqref{eq:free-threshold-range}. It gives
\[
 \liminf_{\substack{N\to\infty\\N\ {\rm even}}}
 \frac{D_{1,2}(N)}{\HN}\ge c(a,h,d),
\]
where
\[
 c(a,h,d)=\frac{3L_1+L_2+L_5+L_6-U_3-U_4
                 -2U_7-U_8-U_9-U_{10}-U_{11}}4.
\]
The upper terms have the integral definitions in the original functional,
with every switching coefficient $8$ replaced by the established cap
$K=681/100$. The last domain is always the full domain
$a<t_1<t_2<t_3<h<t_4<d$.
For the single-prefix upper terms use
\[
 U_3=\frac{2e^{-\gamma}}a
 \left\{\int_a^e F\left(\frac{v(t)-t}{a}\right)\frac{dt}{t}
       +\int_e^{1/3}F\left(\frac{1/2-t}{a}\right)\frac{dt}{t}\right\},
 \qquad
 U_4=\frac{2e^{-\gamma}}a
       \int_a^dF\left(\frac{v(t)-t}{a}\right)\frac{dt}{t}.
\]
If a displayed upper argument is in $(0,1)$, the notation $F(s)$
means the extension $2e^\gamma/s$. It is implemented by lowering the
sifting cutoff; no linear-sieve theorem at an argument below one is
being invoked.
The lower terms are
\[
 L_1=\frac4a\bar f(\nu_*/a),\qquad
 L_2=4\left\{(1/h+1)\bar f\left(\frac{1/h+1}{2}\right)
                         -\bar f(249/50)\right\}.
\]
For $a\le u\le\min(t,h)$, define
\[
 V_{2,a}(t,u)=
 \begin{cases}
 \max\{v(t),v(t+u),153/256-3a/8\},&t\le\mu,\\
 v(t),&t>\mu.
 \end{cases}
\]
Use $V_{2,a}$ in the existing two-prefix lower integrals for $L_5,L_6$.
These are integral statements; floating-point evaluations of $c$ do not
supply a rigorous lower bound.
\end{proposition}

\begin{proof}
The pointwise weight proof uses the counts $r,s,t$ of prime factors
below $N^h,N^d,N^{1/3}$, respectively. Its algebra and the identities
$Y_{10}=\binom{r-2}{2}$ and $Y_{11}=(r-2)(s-r)$ for $r\ge2$
do not require a relation between $d$ and $a$.
The quotient cutoff in $Y_9$ is at least $N^{(1-d)/4}>N^a$ by
$4a+d<1$, so every term has the required initial roughness.
In the case $r=0,s=1,t\le2,k\ge4$, the product of the first four
factors is greater than $N^{h+d+2/3}>N$; this uses $h+d>1/3$.
All other cases of the previous pointwise proof are unchanged.
The exceptional square factors and $n<N^{1-a}$ still cost
$O(N^{1-a})$. In particular, every prime-product sequence in the
switching argument has at most eleven factors, since $a>1/12$.

For the one-prefix upper bounds, $2t<v(t)$ holds for $t<e$.
Beyond $e$ the classical upper sieve supplies the stated option.
If its residual parameter is less than two, lower the sifting cutoff
to half the residual level exponent. Since $sF(s)=2e^\gamma$ on
$0<s\le3$, this gives the same displayed kernel. Boundary strips
are removed using the existing order of limits.

The uniform restricted-lower argument only requires that the actual
largest rough prime exponent be at most $h<\mu$. For three or more
rough slots the repaired full envelope is at least
\[
 \max\{(1+t_1)/2,153/256-3t_1/8\}\ge\nu_*.
\]
For at most two slots, $2h<1/2$ permits the same ordinary
Bombieri--Vinogradov treatment. Thus the proof of $L_1$ applies at
cutoff $a$. For $L_2$, start the exact least-prime Buchstab identity
at $w_*<h$ and use the one-prefix upper bound on $[w_*,h]$.
Its smallest limiting upper-sieve argument is
$(1-h)/(2h)>2$, and the initial square condition has a positive
margin since $h<\mu<1/3$. The antiderivative already derived in
Section~\ref{sec:buchstab-restricted-lower} yields the stated full
$f$ expression. Its later specialization to an elementary logarithm
at the old value of $h$ is not imposed on this parameter range.

In the two-prefix lower sieve, all residual rough primes are below
$N^a$, while the two specified primes exceed it. The proof of the
$V_{2,a}$ envelope applies because $1-2h>153/256$ and $a<h<\mu$.
The empty residual atom has total support below $N^{d+h}<N^{1/2}$.
On the part where the lower function is positive,
$V_{2,a}-t-u>2a>h\ge u$ proves the second prefix square condition.
The first condition follows from $2t<v(t)$ for $t<d<e$.
Where the function is zero, nonnegativity suffices.
These checks replace numerical inequalities special to the old
thresholds; the signed box construction and its density are unchanged.

Finally, $S+z<1$ throughout the two four-prefix switching domains,
where $S$ is the sum of the four exponents and $z$ is the second
smallest. For $U_{11}$ this follows from $S+z<d+4h<1$;
for $U_{10}$ it follows from $5h<1$, since $h<\mu<1/5$.
The quotient-unit boundary is therefore avoided. The already derived
finite-prime-product switching proof also permits the current $a$ in
place of the original fixed roughness lower bound. Its uses of that
bound are a fixed positive prime-factor exponent, at most eleven
factors, the choice $12\zeta<a/2$, and errors of sizes
$N^{1-a+o(1)}$ and $N^{1-a/2+o(1)}$. Choose $\zeta$ again for this
fixed $a$; those errors remain negligible. The coefficient integral
of the modulus sieve is independent of $a$, so the established
$K$ applies to all five terms. This is an adaptation of the earlier
proof, rather than an application of its old fixed-$a$ statement.
Retain fixed positive margins, let $N$ tend to
infinity, and only then remove the construction parameters and
boundary strips. Combining the eleven estimates proves the assertion.
\end{proof}

The rational tuple $(a,h,d)=(43/500,3/20,47/200)$ satisfies every
inequality in \eqref{eq:free-threshold-range}. It also has
$a>(153/256)/7$, so the four-prefix ordinary upper option from the
separate minimum construction uses its simpler kernel on this tuple.
No numerical value for its integral functional is certified here.

% Independent source-audit fragment; no documentclass and not inserted in main.
% 2026-10-08. This does not refute Li's theorem or certify its constant.
\subsection{A limitation on transporting double-sieve savings}
\label{sec:mixed-level-bridge-followup}

Write $c_\gamma=2e^\gamma$. On one positive averaged prime-prefix
box, let $M$ be the common dimension-one main factor at a fixed sieve
point $z$. If the prefix has exponent $T$ and $z=N^u$, put
\[
 s_j=(\nu_j-T)/u,\qquad \nu_1>\nu_0.
\]
All formulas below with a single $T$ denote the limit of shrinking
exponent boxes. Before that limit the corresponding expressions are
averaged with the same positive main-factor measure.

\paragraph{The unnormalized defect has the opposite elementary monotonicity.}
For the actual count $S$ on a fixed box, set
\[
 E_j^+=M F(s_j)-S,\qquad E_j^-=S-Mf(s_j).
\]
Then, identically,
\[
 E_1^+=E_0^+-M\{F(s_0)-F(s_1)\},\qquad
 E_1^-=E_0^--M\{f(s_1)-f(s_0)\}.
\]
Both defects decrease as the linear-sieve baseline improves. The same
identities hold for the positive averages before shrinking the box.
Consequently an old lower bound for either defect is not preserved by
a monotonicity argument. This does not preclude a stronger bound for
the actual count on a special family; such a bound needs an additional
argument.

Suppose both new-level linear bounds are available. An upper mixed
formula $S/M\le F(s_1)-c_\gamma H_0(s_0)/s_0$ would require
\begin{equation}\label{eq:mixed-width-necessary-upper}
 \frac{H_0(s_0)}{s_0}
 \le \frac{F(s_1)-f(s_1)}{c_\gamma}.
\end{equation}
Likewise the mixed lower formula requires the same inequality with
$h_0$ in place of $H_0$. These are necessary conditions, not sufficient
conditions for either mixed formula.

\paragraph{An exact width bound.}
The delay equations for the linear sieve imply
\begin{equation}\label{eq:linear-width-dickman}
 \frac{F(s)-f(s)}{c_\gamma}=\frac{\rho(s-1)}{s}\qquad(s\ge1),
\end{equation}
where $\rho$ is the Dickman function. Indeed, setting
$q(s)=s(F(s)-f(s))/c_\gamma$ gives $q(s)=1$ for $1\le s\le2$ and
$q'(s)=-q(s-1)/(s-1)$ for $s>2$. Thus $q(s)=\rho(s-1)$.
The identity
\[
 v\rho(v)=\int_{v-1}^v\rho(t)\,dt\qquad(v\ge1)
\]
and nonnegativity and monotonicity give
$\rho(n)\le\rho(n-1)/n$, hence $\rho(5)\le1/120$. In particular,
\begin{equation}\label{eq:mixed-width-six}
 s\ge6\quad\Longrightarrow\quad
 \frac{F(s)-f(s)}{c_\gamma}\le\frac1{720}.
\end{equation}
No floating-point evaluation of sieve functions is used here.

\paragraph{A buffered four-prefix obstruction to a universal bridge.}
Consider four disjoint prime intervals with limiting exponents
\[
 (t_1,t_2,t_3,t_4)=
 \left(\frac15,\frac{149}{1000},\frac{37}{500},\frac{37}{1000}\right).
\]
Their total is $T=23/50$ and their four square-prefix exponents are
\[
 \frac25,\quad\frac{249}{500},\quad
 \frac{497}{1000},\quad\frac{497}{1000}.
\]
They are all strictly below $1/2$, with minimum margin $1/500$.
Take
\[
 u=\frac2{105},\qquad \nu_1=\frac{117}{200}.
\]
At the unbuffered center the two arguments are
\[
 s_0=\frac{1/2-T}{u}=\frac{21}{10},\qquad
 s_1=\frac{\nu_1-T}{u}=\frac{105}{16}>6.
\]
The sieve point is below the smallest specified prime. Moreover,
the single-variable L23 level at the largest prime is
\[
 v(1/5)=\frac{189}{320},\qquad
 v(1/5)-\nu_1=\frac9{1600}>0.
\]
This uses the single-variable branch of L23, Proposition~6.1;
no assertion about its three-variable formula above $\mu$ is made.

Fix the small positive level margins and exponent-box widths so that
the old total level is $Q=N^{1/2-\delta_0}$, all four old square-prefix
conditions persist, every old argument lies in $(2,11/5)$, and every
new residual argument, including the linear-sieve support margins,
exceeds $6$. This is possible because all displayed margins are
strict. On dyadic scales $x\ge N^{1-\kappa}$, choose $\kappa$ still
smaller so that the perturbation of $\alpha=\log N/\log x$ and of all
exponent coordinates is absorbed by those margins. The initial
short range and the residue-uniform estimates are treated as in the
prime-prefix interface.

Wu's class $\mathcal U_k$ allows the convolution of four different
prime indicators on intervals $[V_i/\Delta,V_i)$, for every $k\ge4$,
with a common $1+(\log N)^{-4}\le\Delta<1+2(\log N)^{-4}$.
The displayed strict square-prefix inequalities verify its support
conditions. Fixed exponent boxes can be partitioned into such
intervals. Their positive sum preserves the relative main-term
estimate; interval endpoint and variable-cutoff errors tend to zero.
The condition $p\nmid N$ is retained in every indicator.

The same ordinary prime-prefix linear-sieve interface extends from its
stated one- and two-prefix cases to these four fixed disjoint slots:
convolve each residual rough atom with the four fixed prime indicators.
The initial four $\mathcal D^{\rm well}$ conditions are the displayed
ones, and the remaining conditions follow from the residual level
after multiplication by the prefix. The largest prime stays in the
first interval, so the single-variable factorization bound applies to
every atom. The fixed slot partition and its factorial normalization
are unchanged in form; only the fixed number of slots increases.
Thus the usual atom proof supplies both linear bounds with the common
main factor
\[
 M_N=\sum_{d\text{ in the positive prefix box}}
       \frac{\operatorname{li}(N)}{\varphi(d)}V_N(z)>0,
 \qquad
 V_N(z)=\prod_{\substack{\ell<z\\\ell\nmid N}}
                 \left(1-\frac1{\ell-1}\right).
\]
All prefix primes exceed $z$, so excluding primes dividing $d$ does
not change this product. Errors can be made $o(M_N)$ on fixed boxes.
This paragraph is the derived fixed-four-slot extension, not the
literal statement of L23, Proposition~6.6.

Wu's published Table~1 gives $H(11/5)\ge H_*=223939/10^7$.
As $H$ decreases, on every old local argument in $(2,11/5)$,
\[
 \frac{H(s_0)}{s_0}\ge\frac{5H_*}{11}
 =\frac{223939}{22000000}>\frac1{720}.
\]
Therefore a universal mixed upper bound on this family would give
\[
 S\le M_N\left(F(6)-c_\gamma\frac{223939}{22000000}+o(1)\right)
       <M_N(f(6)-c+o(1))
\]
for some fixed $c>0$, whereas the new-level lower linear sieve gives
$S\ge M_N(f(6)+o(1))$. This is a contradiction for large $N$.
Table~2 gives $h(11/5)\ge191556/10^7$; its monotonicity on $[2,10]$
gives the same contradiction for a universal mixed lower bound.
Fix all buffers and boxes first, let $N\to\infty$, and only then
shrink the buffers. The strict contradiction survives without taking
the old argument at the endpoint $s=2$.

This disproves the proposed bridge as a universal rule over the
common prime-prefix class. These four exponents are not claimed to
lie in Li's particular correction domains. The argument neither
disproves any of Li's individual mixed integral estimates nor
disproves the published Chen coefficient $1.9728$.

\paragraph{What positivity of the iteration does and does not give.}
Wu 2004, Lemmas~5.1--5.2, obtain the source terms $\Psi_i$ and the
positive $H,h$ integrals by evaluating the weighted inequalities with
a common total level. Their baseline cancellations use that same
level, notably the identities in (5.10). If replacing their individual
input bounds by better bounds lowers an old upper bound
$B_0-D_H$ by a total amount $I\ge0$, the desired expression
$B_1-D_H$ additionally requires
\[
 I\ge B_0-B_1.
\]
Positivity supplies only $I\ge0$. The four-prefix obstruction shows
that the stronger inequality cannot be assumed universally.

L21, Lemma~6.4, supplies a different valid interface on its stated
family: all extended and switched sets must be equidistributed at
the same new level, and the correction is
$c_\gamma H_{\nu_1}(s_1)/s_1$, at the new argument. Its monotonicity
is in the level at a fixed argument. It neither gives the mixed
old-argument saving nor states a Goldbach result uniform in residue
$a=N$; those extensions need their own distribution verification.

\paragraph{A conditional constructive lower correction at the new argument.}
There is nevertheless a simple useful consequence once suitable
new-level upper corrections have actually been proved. Suppose all
extended prefix boxes arising in Buchstab's identity have the upper
bound
\[
 S(\mathcal A_{dp},p)\le
 M_{dp}(p)\left(F(w)-\frac{c_\gamma}{w}J(w)+o(1)\right),
 \qquad w=\frac{\nu_1-T-t}{t},\quad p=N^t,
\]
where $J\ge0$ is supported in $[2,3]$, and suppose the ordinary
new-level lower bound is valid at the smaller sieve point with
argument $4$. Then for $2\le s\le4$, Buchstab's identity and positive
box summation give
\begin{equation}\label{eq:derived-new-level-h}
 \frac{S(\mathcal A_d,z)}{M_d(z)}\ge
 f(s)+\frac{c_\gamma}{s}
       \int_{\max(2,s-1)}^3\frac{J(w)}w\,dw+o(1).
\end{equation}
Indeed write $R=\nu_1-T$, $u=R/s$, $v=R/4$, and subtract the
upper bounds in
$S(\mathcal A_d,N^u)=S(\mathcal A_d,N^v)
-\sum_{N^v\le p<N^u}S(\mathcal A_{dp},p)$.
The main-factor ratios give the baseline
\[
 \frac4s f(4)-u\int_v^u\frac{F(R/t-1)}{t^2}\,dt=f(s),
\]
by $(sf(s))'=F(s-1)$. The same substitution $w=R/t-1$ turns the
positive upper saving into the integral in
\eqref{eq:derived-new-level-h}. No old-argument saving is inserted.
For example, if a justified same-new-level double-sieve input allows
$J(w)=H_{1/2}(w)$ on $[2,3]$, the printed bound $H(3)\ge0.0072943$
gives a positive correction at every $2\le s<4$.
This is a conditional new-argument construction, not a verification
of those extended/switched distribution hypotheses for all of Li's
domains and not a numerical recertification of $1.9728$.

% Research fragment, 2026-10-08.  This is not a new counting theorem.
% Uses the quotient Y_i and free-threshold conventions of the main manuscript.
\subsection{An exact slack term and obstructions to trimming the full weight}
\label{sec:richer-weight-route}

Throughout this subsection the free-threshold hypotheses hold, $n<N$
is squarefree, and $N^{1-a}\le n$.  Let $W$ be the full eleven-term
weight, and let $r,s,t$ count factors below $N^h,N^d,N^{1/3}$.
The analytic extension below uses the previously derived signed-box
distribution proof as an input; its new insertion and coefficient checks
are supplied explicitly.  No new numerical coefficient is certified.

\begin{proposition}[A positive three-prefix term]
Define
\[
 Z(n)=\sum_{N^a\le p_1<p_2<N^h\,,\ N^d\le p_3<N^{1/3}}
          Q_{p_1p_2p_3}(p_2;n).
\]
On the rough support of the weight,
\[
 Z(n)=\mathbf1_{r\ge2}(t-s),\qquad
 W(n)=s-t\quad(r\ge2).
\]
Consequently
\[
 W(n)+Z(n)\le4\mathbf1_{\Omega(n)\le2}.
\]
After the same square-factor and small-$n$ exceptions as for $W$,
the summed inequality gains the positive term $\sum_P Z(N-P)$.
\end{proposition}

\begin{proof}
The quotient cutoff makes $p_1,p_2$ the two smallest factors of $n$.
Thus the displayed tuple sum counts each factor in $[N^d,N^{1/3})$
once exactly when the two smallest factors are below $N^h$.
For $r\ge2$ the algebra of the first six terms and the two full
four-prefix terms is exactly $s-t$, rather than merely at most $s-t$.
Both $Y_7$ and $Y_8$ vanish because their second specified factor would
have to be the second smallest factor.  Every $Y_9$ quotient cutoff
exceeds $N^{(1-d)/4}>N^h$, by $4h+d<1$; a factor below $N^h$ remains
in its quotient, so $Y_9=0$ as well.  This proves the identities.
For $r<2$ the new term vanishes.  For $\Omega(n)\le2$ it also vanishes.
Every new tuple requires initial $N^a$-roughness, and all multiplicities
are bounded by a constant depending on $a$, so the earlier exceptional
integer argument applies unchanged.
\end{proof}

This lemma saturates the $r\ge2$ part of the elementary weight:
it sets its value to zero on all the unwanted configurations in that
part.  It does not itself lower-bound the new prime average.

\paragraph{What an ordinary lower estimate would need.}
Write the three specified exponents as $u<v<h$ and $d\le x<1/3$.
The quotient is sifted at $N^v$, and the prefix has exponent $u+v+x$.
A linear lower sieve at total distribution level $\nu$ is positive only if
\begin{equation}\label{eq:richer-Z-positivity}
 \nu>x+u+3v.
\end{equation}
In particular its lower corner needs $\nu>d+4a$.
With the three-prefix extension proved below, an admissible level
$\nu(u,v,x)$ supplies the lower main term
\[
 L_Z=4\int_a^h\int_u^h\int_d^{1/3}
  \frac{\bar f((\nu(u,v,x)-u-v-x)/v)}{u v^2 x}
  \,dx\,dv\,du,
 \qquad \bar f=f/(2e^\gamma),
\]
with zero assigned wherever the interface or its positive argument
is unavailable.
In particular $p_1<p_2$, whereas residual sieve primes can lie between
them; a proof must respect their actual order and cannot blindly call
a lemma requiring all named primes to precede every residual prime.
An exact possible starting point is the least-prime identity
\[
 Q_b(N^v;n)=Q_b(N^u;n)
 -\sum_{N^u\le q<N^v\,,\ q\nmid b}Q_{bq}(q;n),
 \qquad b=p_1p_2p_3,
\]
on squarefree integers.  Thus replacing $v$ by $u$ alone gives an
upper comparison, not a valid lower estimate.  A proof through this
identity would also have to upper-bound the four-prefix subtraction.

\begin{lemma}[Inserting a retained small prefix into residual well support]
Let $X>V>U>1$ be three distinct specified primes and put
$R=D/(XVU)$.  Suppose
\[
 X^2\le D,\qquad XV^2\le D,\qquad U\le R.
\]
Let $q_1>\cdots>q_m$ be distinct primes, none equal to $X,V,U$,
all less than $X$, with
$q_1\cdots q_{j-1}q_j^2\le R$ for every $j$.
Then the decreasingly sorted union of $X,V,U,q_1,\ldots,q_m$
belongs to $\mathcal D^{\mathrm{well}}(D)$.
\end{lemma}

\begin{proof}
The test at $X$ is the first displayed assumption.
For any residual prime, the specified primes preceding it form a
subproduct of $XVU$; its full-vector square product is therefore at
most $XVUq_1\cdots q_{j-1}q_j^2\le D$.
At $V$, if no residual prime precedes it, use $XV^2\le D$.
Otherwise write $Q_V=\prod_{q_j>V}q_j$ and take its final factor
$q_*>V$.  Then $Q_Vq_*\le R$, hence $Q_VV<R$, and the full square
product is $XQ_VV^2<XVR<D$.
Write $Q=\prod_{q_j>U}q_j$.  If this is empty, the test at $U$
is $XVU^2\le D$, by $U\le R$.
Otherwise let $q_*$ be the last factor in $Q$.  Residual well support
gives $Qq_*\le R$, hence $QU<R$.  Therefore the square product at
$U$ is $XVQU^2<D$ as well.
\end{proof}

This elementary insertion lemma resolves the ordering obstruction
for the single-largest-prime branch at the level of rough atoms.
With $D=N^\nu$, the conditions follow with fixed margins from
$2x<\nu$, $x+u+3v<\nu$, and $u<v<x$.
The corresponding signed box proof can retain every specified and
residual slot as a marked slot.  When their interval labels coincide,
keep the original residual scalar, including its original residual
reciprocal factorial; do not renormalize the joint sorted product to
an indicator, which would remove the genuine choices of the marked
prefix.  Partition and divide each resulting convolution group by
the factorial of its group size, restoring that factorial in the
scalar.  This is the same exact convolution identity used in the
existing prefix proof and produces bounded coefficients even for
coincident labels.  Finitely many interleavings and endpoint inflation
are absorbed by the fixed box margin.

Collisions between $U$ and a residual rough prime are removed by the
same rough-square error
$(N+N^{\nu+O(\eta)})z_0^{-1}(\log N)^C$
as in the existing atom proof.  On the squarefree sequence the sieve
may omit primes dividing the fixed prefix.  Only $U$ is below the
exact cutoff $V$; an upper box-endpoint cutoff may include $V$ as
well.  Removing either or both changes the residual Mertens factor
by $1+O(U^{-1})$, uniformly over the prefix box.  The density away from
those excluded primes remains $1/(\ell-1)$.  These observations show
how the existing signed-atom distribution proof extends to this
interleaving configuration, with largest exponent $x$ and endpoint
$v_0(x)$; they do not justify any higher selected-subproduct endpoint
without its additional gates.  A final box-interface writeup and
independent numerical comparison should retain these details.

\begin{proposition}[Derived interleaving three-prefix lower interface]
Assume the modified-sieve atom decomposition and the uniform
signed-box prime-distribution argument of
Section~\ref{sec:prefix-interface}, including its one-long-variable
coprime-mean correction.  On each fixed interior box with
$a<u<v<h<d<x<e$, choose a total level satisfying fixed margins in
\[
 \max\{1/2,2x,x+2v,x+v+2u,x+u+3v\}<\nu<v_0(x).
\]
Then the lower average of $Z$ on that box has the ordinary multiplier
$f((\nu-x-v-u)/v)$ and density
$2e^{-\gamma}/(xuv^2)$ relative to $N\mathcal C(N)/\log^2N$.
Consequently the displayed integral for $L_Z$ is valid with
the supremum $\nu\uparrow v_0(x)$ on these boxes and zero elsewhere.
\end{proposition}

\begin{proof}
Use the upper endpoints of the three fixed prime boxes to define the
residual level $R=N^{\nu-T_+-c\delta}$, and sieve at the upper endpoint
of the middle-prime box.  The latter cutoff is at least the actual
middle prime, so its sifted count is a valid lower comparison for
$Q_b(p_2;n)$.  Its residual rough primes can slightly exceed the actual
middle prime but remain below the largest fixed prime.  The insertion
lemma was stated in precisely this more general form.  The three
fixed-prime square tests and $U\le R$ have the required margins;
all residual rough vectors therefore concatenate into full well
support at level $N^{\nu-c'\delta}$.

The largest rough slot remains the fixed prime of exponent $x$.
Apply the single-largest-prime factorization at endpoint $v_0(x)$.
Keep all marked slots and their original scalar factors and normalize
the factor groups as explained above.  The small-prime factor is
attached by the same unique rough/small decomposition as in the
existing prefix proof.  Removing collisions costs its existing
rough-square error.  Thus its moving-residue Type II and one-long
and two-long estimates apply to every signed atom with the same
fixed positive margins.  They are insensitive to the number of
specified marked slots after their bounded-coefficient normalization.
Reassemble the residual lower sieve, preserving its signs.

On squarefree support remove every fixed-prefix prime below the
box sifting cutoff from the residual sifting set.  Their Mertens correction is
$1+O(N^{-u})$.  Apart from that negligible correction, the residual
density and presieve are those of the existing lower interface,
so the main multiplier is the usual lower function at
$(\nu-x-v-u)/v$.  Prime harmonic summation over the three boxes
gives $2e^{-\gamma}/(xuv^2)$.  Let $N$ tend to infinity with fixed
box widths and margins, then shrink the boxes and sieve precision,
and finally remove the fixed margins.  Boxes at boundaries or with
nonpositive lower function can be dropped by nonnegativity.  This
proves the asserted integral with $4\bar f=2e^{-\gamma}f$.
\end{proof}

\begin{lemma}[A zero region for the currently available caps]
Put
\[
 v_0(z)=\min\{(1+z)/2,57/64-3z/2\},\quad
 \mu=25/128,\quad e=57/224.
\]
Suppose $7/32\le d<e$ and
\begin{equation}\label{eq:richer-Z-zero}
 d+4a\ge v_0(d),\quad\hbox{equivalently}\quad
 d\ge\frac{57/64-4a}{5/2}.
\end{equation}
Then the ordinary lower-sieve candidate for $Z$ is zero throughout its
domain when using the classical cap $1/2$, the largest-prime cap
$v_0(x)$, and any of the currently proved selected-subproduct caps.
The complete small-largest-prime envelope is inapplicable since $x>\mu$.
\end{lemma}

\begin{proof}
First consider any selected subproduct whose exponent is $b$ and which
contains the largest exponent $x$.  Then $b\ge x>\mu$, hence
$v_0(b)\le v_0(x)$.  If the selected subproduct does not contain $x$,
its necessary gate includes $2x+b$.  Eligibility would imply
\[
 2x+b<v_0(b)\le(1+b)/2,\qquad b<1-4x,
 \qquad v_0(b)<1-2x.
\]
For $7/32\le x\le1/4$, one has
$1-2x\le57/64-3x/2=v_0(x)$.
For $x\ge1/4$, no positive $b$ can pass that gate.
Thus none of the selected-subproduct branches beats $v_0(x)$,
even if additional small residual factors are made available for
selection.  On $x\ge d$, the total cap furnished by these choices
and the classical choice is at most $v_0(d)>1/2$.
But $x+u+3v>d+4a\ge v_0(d)$ in the open integration domain.
The normalized lower function is therefore identically zero.
\end{proof}

For $(a,h,d)=(43/500,3/20,47/200)$, the two sides of the corner test are
\[
 d+4a=579/1000,\qquad v_0(d)=861/1600,
 \qquad d+4a-v_0(d)=327/8000>0.
\]
The sampled tuple $(0.08455645453148589,0.15049959287189102,
0.2352266768449549)$ likewise lies well inside the zero region.
For smaller $d$, the corner test can change sign.  For example, with
$a=43/500$ and $d=43/200$, one has $v_0(d)>d+4a$.
This gives an open positive block for the preceding derived
three-prefix interface, but does not establish an improvement over
the complete old functional after changing $d$.
For instance, the strict interior exponents
$(u,v,x)=(173/2000,87/1000,431/2000)$ have
\[
 x+u+3v=563/1000<113/200<4539/8000=v_0(x).
\]
This produces an open positive block at level $113/200$.

\begin{proposition}[Every interior full-$Y_{11}$ tuple needs compensation]
Assume in addition that $d>1/6$ and $5h+d<1$.
Take any strict exponent tuple
\[
 a<x_1<x_2<x_3<h<y<d,
 \qquad S=x_1+x_2+x_3+y.
\]
It can be completed to a sum-one exponent profile with five or six
distinct positive entries, having exactly three entries below $h$,
and having $t=s$ (no entry in $[d,1/3)$).
The original full weight has value zero on that profile, while its
$Y_{11}$ count includes the prescribed tuple.
Consequently no deletion of a nonempty open subdomain of the full
$Y_{11}$ tuple domain, with the other terms fixed and with no added
compensation, is a valid pointwise improvement.  The conclusion also
holds if the three small specified primes are required to be the
three smallest prime factors.
\end{proposition}

\begin{proof}
If $S<2/3$, append the exponent $1-S>1/3$.
There are then three small factors, one middle factor, and one factor
above $1/3$: $r=3,s=t=4$, so $W=0$.
If $S>2/3$, put $R=1-S$.  The hypotheses give
\[
 R>1-3h-d>2h,\qquad R<1/3<2d.
\]
Choose two distinct numbers in $(h,d)$ with sum $R$, avoiding $y$.
This is possible on an open interval around $R/2$.  These six entries
give $r=3,s=t=6$, again with $W=0$.  At $S=2/3$, either use the
single remaining exponent $1/3$ or use the same two-factor completion;
the latter has an open interval of choices by the strict hypotheses.
In every case $x_1,x_2,x_3$ are the three smallest entries.

For completeness the strict profiles yield actual arithmetic
counterfamilies, not just a formal exponent model.  For a large real
$X$, choose distinct primes $q_i=X^{\xi_i+o(1)}$ for the constructed
entries $\xi_i$, put $n=\prod_iq_i$ and $N=n+3$.
Prime existence in fixed relative intervals follows from the prime
number theorem.  Then $N$ is even, $n=N-3$, $(n,N)=1$, and
$\log q_i/\log N\to\xi_i$.  All strict domain conditions persist.
The quotient cutoff for the prescribed tuple holds since any remaining
factors exceed the third-smallest factor.  Thus deleting a tuple from
an open neighborhood gives a positive weight on an integer with
$\Omega(n)>2$.  This contradicts the proposed pointwise inequality.
The equality slice $S=2/3$ is not needed for the open-domain conclusion.
\end{proof}

Two rational examples at $(43/500,3/20,47/200)$ are
\[
 (0.09,0.10,0.14,0.23,0.44),\qquad
 (0.141,0.145,0.149,0.162,0.170,0.233).
\]
The first has the offending small-plus-middle sum $0.14+0.23=0.37$
above $1/2-2a=0.328$.  The second shows why merely retaining the
part $S\le2/3$ is insufficient: the tuple formed by its three smallest
entries and $0.233$ has $S=0.668>2/3$, yet $s=t=6$.
Both have $W=0$.  Their role is to refute uncorrected pointwise
domain deletion, not to prove that all more general weights fail.

\begin{lemma}[Exact $Y_9$ multiplicity near the full scale]
Fix
\[
 0<\epsilon<\min\{a,(1-3d)/2\},\qquad N^{1-\epsilon}\le n<N.
\]
For squarefree $n$, $Y_9(n)$ can be nonzero only for three prime factors
$q_1<q_2<q_3$ with $N^h\le q_1<N^d\le q_2$.
In that case
\[
 Y_9(n)=1+\mathbf1_{q_3/q_2<N/n}.
\]
In particular its excess over the conservative value can occur only
where $\log(q_3/q_2)/\log N<\epsilon$.
\end{lemma}

\begin{proof}
If the quotient in a $Y_9$ tuple has at least two prime factors,
both are at least $\sqrt{N/(p_1p_2)}$; this gives $n\ge N$, impossible.
If it is a unit, then
$n=p_1p_2<\sqrt{Np_1}<N^{(1+d)/2}<N^{1-\epsilon}$.
So there are exactly three prime factors.
If two of them are below $N^d$, the specified $p_2$ is the largest;
the remaining factor is below $N^d$, giving
$n<N^{(1+3d)/2}<N^{1-\epsilon}$, again impossible.
Thus $p_1=q_1$ is unique and lies in $[N^h,N^d)$.
Since $n\ge N^{1-\epsilon}$ and $q_2,q_3\ge N^d>N^\epsilon$,
the quotient cutoff holds for either choice of the second prefix.
The choice $p_2=q_2$ always passes its upper endpoint because
$q_1q_2^2<n<N$.  The other choice passes precisely when
$q_1q_3^2<N$, equivalent to $q_3/q_2<N/n$.
\end{proof}

The lemma already proves that $Y_9$ multiplicity cannot compensate
any of the $r\ge2$ omissions of $Y_{11}$.  Moreover, under the existing
finite-prime-product switching upper interface, its excess has zero
limiting main coefficient.  Indeed discard $n<N^{1-\epsilon}$,
which costs $O(N^{1-\epsilon})$, and switch the remaining three-prime
sequence.  Its limiting exponents satisfy
$t_1+t_2+t_3=1$ and $0<t_3-t_2<\epsilon$.
The corresponding mass is bounded by
\[
 \int_h^d\frac{dt_1}{t_1}
 \int_{(1-t_1-\epsilon)/2}^{(1-t_1)/2}
 \frac{\mathbf1_{t_2,t_3\ge d}}{t_2t_3}\,dt_2
 \le\frac{\epsilon}{2d^2}\log(d/h).
\]
For each fixed $\epsilon$ the usual buffered boxes justify this upper
estimate; multiply it by the established switching cap and then let
$\epsilon\downarrow0$.  Hence exploiting this multiplicity alone
does not improve the leading coefficient of the eleven-term LP.

\paragraph{Status of the route.}
The new proven information is a genuine positive three-prefix slack
term, a strict zero region for its current ordinary-sieve estimates,
and an all-interior obstruction to uncompensated full-$Y_{11}$ trimming.
The positive block at lower $d$ is supported by the derived
interleaving interface above; a useful numerical consequence would
still require a certified comparison against the losses in the other
terms after changing $d$.  The floating probe
\texttt{probe\_positive\_three\_prefix.json} found, at
$(a,h,d)=(0.086,0.15,0.215)$, a gain $L_Z/4$ of about
$3.73\cdot10^{-7}$, whereas the original coefficient was about
$1.8049$.  At $d=0.20$ the gain was about $0.000198$ and the original
coefficient about $1.69590$.  Its parameter search returned to a
point with $Z=0$ near the earlier diagnostic coefficient $1.87414$
after the separate four-prefix minima.  These are floating route
diagnostics, not integral certificates or a global optimizer proof.
Any more ambitious trimming
must supply negative compensation on both the five-factor and the
six-factor families above, or use an aggregate identity with an
independent analytic proof.  No new Chen coefficient, imbalance
exponent, global optimality statement, or priority claim is made here.

% Derived one-step double-sieve fragment, 2026-10-08.
% Not inserted in the main manuscript; no global Chen constant asserted.
% The full Wu H-table migration was rejected during independent review.

\subsection{A same-new-level one-step correction on a Buchstab block}
\label{sec:same-new-level-double-sieve-route}

Write $\mathcal A_N=\{N-P:P<N,\ P\text{ prime}\}$,
$K=681/100$, $\mu=25/128$, and $v(t)=(1+t)/2$ for $t<\mu$.
We use
\href{https://www.impan.pl/shop/en/publication/transaction/download/product/82961}
{Wu 2004, Lemma 4.1 and equations (5.1)--(5.10)},
\href{https://arxiv.org/html/2109.02851v2}
{Lichtman, arXiv:2109.02851v2, Lemma 6.4},
and the already derived moving-residue prefix interface based on
\href{https://arxiv.org/html/2309.08522v1}
{Lichtman, arXiv:2309.08522v1, Propositions 5.2 and 6.1}.
The $K$ input is the previously certified positive prime-product upper
bound from \texttt{extended-prefix-upper.tex}. The exact small
calculation below is saved in
\texttt{certify-same-new-level-double-sieve-route.py} and
\texttt{same-new-level-double-sieve-route.json}.

This construction does not infer the signed switched-distribution
hypothesis (5.22) of L21 from a bound for the number of primes. Instead
we use the positive switched count at its actual place in Wu's first
weighted inequality. We do not use the $H,h$ tables or their full
iteration. The resulting explicit function $J$ is a new-argument
correction to a linear upper bound at the same new total level.

\begin{proposition}[A one-step correction for the $L_2$ Buchstab terms]
\label{prop:same-new-level-one-step}
Let
\[
 \frac{31}{200}\le t\le\frac4{25},\qquad
 p=N^{t+o(1)},\qquad \nu=v(t)=\frac{1+t}{2},\qquad
 s=\frac{\nu-t}{t}=\frac{1-t}{2t}.
\]
On every positive averaged prefix box in this interval, after the
usual fixed level buffers and exponent-box limits, the permissible
upper multiplier for $S(\mathcal A_{N,p},p)$ is
\begin{equation}\label{eq:same-new-one-step-upper}
 F(s)-\frac{2e^\gamma}{s}J_0,\qquad J_0=\frac1{250}.
\end{equation}
It is a bound averaged over the prefix primes, with their usual
positive main-factor measure, not an estimate for each individual
prime. In the existing variable-level Buchstab lower bound for $L_2$,
this yields an additional lower coefficient
\begin{equation}\label{eq:same-new-one-step-L2}
 \Delta L_2\ge8J_0\log\frac{676}{651}.
\end{equation}
Consequently the complete Chen functional increases by at least
\begin{equation}\label{eq:same-new-one-step-Chen}
 \Delta c\ge2J_0\log\frac{676}{651}>
 \frac2{6635}.
\end{equation}
The last inequality is purely rational after using
$\log(676/651)=2\operatorname{atanh}(25/1327)>50/1327$.
\end{proposition}

\begin{proof}
Fix a compact exponent box $t\in[t_-,t_+]$ strictly inside
$(31/200,4/25)$.
Use a single total level $Q=N^\nu$ with
$\nu<v(t_-)$ and a fixed positive margin. Then choose all sieve,
factorization, local-scale and mesh precisions inside that margin.
The centre notation below denotes the subsequent shrinking-box limit.
The local exponents are
\[
 \nu\in[231/400,29/50],\qquad
 R=Q/p=N^{L+o(1)},\qquad L=\nu-t=\frac{1-t}{2},
\]
and
\[
 \frac{21}{8}\le s\le\frac{169}{62},\qquad s_0=\frac72,\qquad
 y=R^{1/s_0},\qquad z=p=R^{1/s}.
\]
All inequalities used below have strict surplus, so fixed buffers
preserve them.

For $\sigma(p)$ equal to the anchor prime-box indicator, or a fixed
finite nonnegative constant combination of such indicators, define
\[
 \Phi_s=\sum_p\sigma(p)S(\mathcal A_{N,p},p),\qquad
 \Theta_Q=4\operatorname{li}(N)
   \sum_p\frac{\sigma(p)\mathcal C(pN)}
                 {\varphi(p)\log(Q/p)}.
\]
The ordinary normalized main multipliers are
$A(s)=sF(s)/(2e^\gamma)$ and $a(s)=sf(s)/(2e^\gamma)$.
In particular $A(s)=1$ on this block.

Wu's first weighted inequality has the form
\begin{equation}\label{eq:same-new-one-step-weight}
 2\Phi_s\le
  2\sum_p\sigma(p)S(\mathcal A_{N,p},y)
 -\sum_p\sigma(p)\sum_{y\le q<z}S(\mathcal A_{N,pq},y)
 +\sum_p\sigma(p)\!\sum_{y\le q_1<q_2<q_3<z}
        S(\mathcal A_{N,pq_1q_2q_3},q_2)
 +o(\Theta_Q).
\end{equation}
The prime set excludes divisors of $pN$, with the same additional
prefix exclusions as in Wu's formula. Repeated-prime exceptions can
be removed with a power saving. The inequality follows from the
pointwise Buchstab identities in Wu (4.1)--(4.5); it does not require
the original short-box class after this fixed positive box has been
specified. Thus no error relative to a vanishing short-box mass is
being inferred from a fixed-mass asymptotic.

\paragraph{The two non-switched terms.}
The first term has the ordinary upper multiplier $2A(s_0)$.
Its residual primes are below $y<p$, so the largest modulus prime
remains the anchor. For the second term write
$q=N^{Lu+o(1)}$, $1/s_0\le u\le1/s$. Then
\[
 N^{3/25+o(1)}\le q\le p,\qquad
 t+2Lu\le3t\le\frac{12}{25}<\nu,\qquad
 2t\le\frac8{25}<\nu.
\]
The endpoint equality $q=p$ is trimmed before taking limits.
The lower residual argument is
\[
 s_0(1-u)\in[13/6,5/2].
\]
These are exactly the ordered two-prefix configurations covered by
the existing residual prime-prefix interface: the cutoff $y$ is at
most the smaller specified prime $q$, both initial square conditions
hold with fixed margins, and every later rough square condition
follows from its residual level. Empty residual rough atoms are
included. Both signs retain their original Iwaniec atom coefficients,
small-prime weights, and multiplicity factorials. Every concatenated
rough atom has largest exponent in $[t_-,t_+]$, and hence cap
$v(t_-)> \nu$; L23 Proposition 6.1 supplies its factorization.
The moving-residue argument uses Proposition 5.2 on the true local
scales, as already proved in \texttt{prefix-sieve-interface.tex}.

The negative term in \eqref{eq:same-new-one-step-weight} is therefore
bounded below, after division by $\Theta_Q$, by
\[
 B(s)=\int_{1/s_0}^{1/s}
  \frac{\log(s_0(1-u)-1)}{u(1-u)}\,du
 =\int_{1-1/s}^{1-1/s_0}
  \frac{\log(s_0x-1)}{x(1-x)}\,dx.
\]
No lower double-sieve correction is used here.

\paragraph{The exact positive switched mass.}
The last term counts primes $N-m$ in a labelled product multiset
\[
 m=pq_1q_2q_3n\le N,\qquad
 y\le q_1<q_2<q_3<z,\qquad P^-(n)\ge q_2,
\]
after the same coprimality and repeated-prime exceptions as in
Wu (5.3)--(5.4). Canonically factor $n$ into its ordered prime factors.
Every prime factor of $m$ is at least $N^{3/25+o(1)}$, which is
strictly larger than the $N^{100/1149}$ cutoff of the existing
$K$ theorem. The fixed level buffers are chosen small enough to
preserve this strict gap; the exact exponent $3/25$ is the unbuffered
limiting value. There are only finitely many prime-factor patterns.
All ordering, roughness, product, and prefix-box constraints become
finite monomial inequalities; their labels are retained. Thus the
existing positive-count theorem applies to this exact multiset.

The quotient $n=1$ contributes at most $N^{4t+o(1)}\le N^{16/25+o(1)}$
to the unsifted mass and to the desired count. It is harmless here
by this explicit power bound. Also, throughout the triple integration,
\[
 \alpha=\frac{\log(N/p)}{\log(Q/p)}\ge2+o(1),\qquad
 \frac{\alpha-u_1-u_2-u_3}{u_2}\ge2s-3\ge\frac94.
\]
The standard rough-number estimate is therefore used away from its
argument-one boundary. In particular no six-prime unit-quotient
boundary is present in this one-step construction.

Let $X$ be the unsifted labelled mass. The same partial summation as
in Wu (5.8) gives
\[
 X\le (I(s)+o(1))N
       \sum_p\frac{\sigma(p)}{\varphi(p)\log(Q/p)},
\]
where
\[
 I(s)=\sup_{\alpha\ge2}
 \int_{1/s_0\le u_1\le u_2\le u_3\le1/s}
 \omega\!\left(\frac{\alpha-u_1-u_2-u_3}{u_2}\right)
 \frac{du_1\,du_2\,du_3}{u_1u_2^2u_3}.
\]
For completeness $\omega(w)\le3/5$ for $w\ge2$ follows from
$\omega(w)=(1+\log(w-1))/w$ on $[2,3]$ and the delay equation.
The maximum of
$1+\log(w-1)-(3/5)w$ on $[2,3]$ is at $w=8/3$ and is negative:
$\log(5/3)\le((5/3)-(3/5))/2=8/15<3/5$.
The delay equation propagates this upper bound on successive unit
intervals. Hence
\begin{equation}\label{eq:same-new-one-step-Ibar}
 I(s)\le\overline I(s)=\frac35
 \left\{(s_0+s)\log\frac{s_0}{s}-2(s_0-s)\right\}.
\end{equation}
The elementary integration here first integrates $u_1,u_3$, leaving
$\int_{1/s_0}^{1/s}\log(s_0u)\log(1/(su))u^{-2}\,du$.

The $K$ upper bound, in its common-mass normalization, is
$(K+o(1))\mathcal C(N)X/\log N$. Since
$\mathcal C(N)\le\mathcal C(pN)$ and $\operatorname{li}(N)\sim N/\log N$,
the last term in \eqref{eq:same-new-one-step-weight}, divided by
$\Theta_Q$, is at most $(K/4)\overline I(s)+o(1)$.
The factor is $K/4$, not $K/(4/(\nu-t))$.
This is an upper bound for a positive prime count, not a signed
distribution assertion for a specified global sieve weight.
The $K$ proof's one common positive smooth outer approximation may be
used here on each fixed exponent box, followed by its smoothing limit;
there is no need to claim uniformity over Wu's shrinking short boxes.

\paragraph{An explicit new-argument saving.}
Since
\[
 A(7/2)=1+C,\qquad
 C=\int_2^{5/2}\frac{\log(x-1)}x\,dx,
\]
the three estimates in \eqref{eq:same-new-one-step-weight} give
\begin{equation}\label{eq:same-new-one-step-J}
 \Phi_s\le\{1-J(s)+o(1)\}\Theta_Q,\qquad
 J(s)=-C+\frac12B(s)-\frac{681}{800}\overline I(s).
\end{equation}
Every argument and every baseline in this formula uses the same $Q$.
No saving evaluated at $(1/2-t)/t$ is added.

Here is a short outward certificate for $J(s)>1/250$ on the entire
interval $[s_-,s_+]=[21/8,169/62]$. The function $B$ decreases because
its positive-integrand interval shrinks with $s$. Also
\[
 \overline I'(s)=\frac35\left(\log(s_0/s)+1-s_0/s\right)\le0.
\]
Consequently
\[
 J(s)\ge-C+\frac12B(s_+)-\frac{681}{800}\overline I(s_-).
\]
The integrands defining $C$ and $B(s_+)$ are increasing on their
integration intervals. For $C$ this follows from
$x/(x-1)>\log(x-1)$; for $B$, the positive numerator increases and
$x(1-x)$ decreases because $x>1/2$.
Use a right Riemann sum for $C$ and a left Riemann sum for $B(s_+)$,
each with $32768$ equal rational subintervals.
The existing $80$-bit integer interval logarithm routine, with its
explicit atanh remainder, gives the exact lower bound
\[
 -C_{\rm right}+\frac12B_{\rm left}
       -\frac{681}{800}\overline I(s_-)
 \ge
 \frac{4850531719234607697725}{1208925819614629174706176}
 >\frac1{250}.
\]
No floating point number decides this inequality. The saved certificate
also records a stronger optional check, using $16$ rational
$s$-subintervals and $8192$ cells in each one-dimensional Riemann sum;
its minimum lower bound exceeds $1/200$. The proposition only uses
the simpler bound $1/250$.

In terms of the ordinary main factor at cutoff $p$, the normalized
defect in \eqref{eq:same-new-one-step-J} is
$2e^\gamma J_0/s$, proving \eqref{eq:same-new-one-step-upper}.
The existing Buchstab lower bound for $L_2$ subtracts the ordinary
upper kernels for $S(\mathcal A_{N,p},p)$ over
$25/224<t<50/309$. Our entire block lies strictly inside this range.
Its ordinary kernel uses precisely the level $v(t)$, so the saving
can be inserted in that identity:
\[
 \Delta L_2\ge
 4J_0\int_{31/200}^{4/25}\frac{dt}{t(v(t)-t)}
 =8J_0\log\frac{(4/25)(169/200)}{(31/200)(21/25)}
 =8J_0\log\frac{676}{651}.
\]
The $L_2$ coefficient in the complete eleven-term weight is $1/4$.
This proves \eqref{eq:same-new-one-step-Chen}.
All parameters are fixed before $N\to\infty$; next remove the
smoothing and sieve-precision errors, shrink the boxes and their
level buffers within the open anchor interval, and only then exhaust
that interval by compact subintervals. On each such compact interval
the buffered arguments remain inside the certified $s$ interval.
The endpoints have zero main-factor mass. The strict numerical
surplus above $J_0$ absorbs the remaining errors.
\end{proof}

\paragraph{Why the full old $H$ table was not imported.}
The preceding proposition proves a positive correction directly.
It does not prove the stronger proposed statement that the entire
old table, such as $H(3)\ge72943/10^7$, transfers to all new-level
states. Independent review found a concrete interface gap in that
proposed extension.

In Wu's Lemma 4.2, terms $\Gamma_1$ and $\Gamma_5,\ldots,\Gamma_8$
are positive non-switched upper terms; $\Gamma_2,\Gamma_3,\Gamma_4$
are negative lower terms; and $\Gamma_9,\ldots,\Gamma_{21}$ are
positive switched terms. The latter consist of respectively three
additional primes for indices $9$--$15$, four for $16$--$19$, five
for $20$, and six for $21$. Replacing their final positive prime-count
upper bounds by $K$ would have the correct $K/4$ normalization.
That replacement alone does not check terms $5$--$8$.

Specifically, the published seed at $s=2.2$ uses $\phi_2=2.90$.
The $\Gamma_5$ domain allows two new exponents $u<v<1/\phi_2$,
measured relative to $R=Q/b$. Near its upper corner,
$v+2u>1$. The enlarged prefix then has square exponent
$T+L(v+2u)>\nu$, and is outside the full
$\mathcal D^{\rm well}(Q)$ class of the presently proved prefix
interface. Its small bare product does not solve this: residual
modulus atoms can raise the full modulus above $N^{1/2}$.
A new factorization theorem for that explicit enlarged-prefix domain,
or a replacement upper kernel and a recomputed recursion, is needed
before importing the old $\Psi_2$ seeds or their nine-by-nine matrix.

There is a second mass check to retain for such a future construction.
In the five- and six-extra-prime terms, the unit quotient may lie on
a product boundary at scale $N$. Keeping it in the exact multiset
does not by itself show that an integral involving only the ordinary
Buchstab function includes that boundary mass. Its negligibility,
or its separate mass term, must be checked before reusing those
$I_{2,i}$ envelopes. Neither issue occurs for the one-step triple
term proved above.

Thus the direct L21 black-box route still lacks its requested signed
switched-distribution statement at level $v(t)$. The proposition gives
a different, explicitly proved positive-count route on a concrete
$L_2$ block. It does not restore a universal mixed-level bridge and
does not independently certify a new global counting coefficient.

% Exact consequence of the independently checked single-step correction.
\subsection{A numerical consequence of the single-step correction}
\label{sec:same-level-one-step-corollary}

This consequence uses only the three-term inequality from
Section~\ref{sec:same-new-level-double-sieve-route}. It does not use a
transport of Wu's full correction table, the new positive $Z$ term,
or the floating free-threshold and four-prefix diagnostics.

\begin{theorem}[Separate consequence after one double-sieve step]
\label{thm:same-level-one-step-chen}
Under the derived distribution and positive prime-product interfaces
in this manuscript, for every sufficiently large even $N$,
\[
 D_{1,2}(N)\ge\left(\frac{909}{500}-o(1)\right)\HN,
 \qquad
 D_{7659/10000}(N)\ge\frac1{1000}\HN.
\]
In particular there is a representation $N=p+rq$, with $p,q$ prime,
$r=1$ or prime, and $r\le q^{0.7659}$. This is a consequence within
the research proof draft, not a certification of priority over all
existing results.
\end{theorem}

\begin{proof}
The existing outward-rounded table gives $c_0=9089/5000$.
The separately checked one-step correction has $J\ge1/250$ on
$31/200<t<4/25$, at the same total level $v(t)=(1+t)/2$ as the
existing Buchstab subtraction for $L_2$. Its contribution to the
Chen functional, including the coefficient $1/4$ of $L_2$, is
\[
 \Delta c\ge\frac{2}{250}\int_{31/200}^{4/25}
                  \frac{dt}{t(1-t)}
 =\frac{1}{125}\log\frac{676}{651}
 >\frac{2}{6635}.
\]
The strict inequality follows from
$\log(676/651)=2\operatorname{atanh}(25/1327)>50/1327$.
Thus
\[
 c_0+\Delta c>
 \frac{9089}{5000}+\frac{2}{6635}>
 \frac{909}{500}.
\]
The first statement follows with a fixed positive rounding surplus.
For the second set $b=7659/10000$ and
$u=(1-b)/(1+b)=2341/17659$. The elementary convergent series gives
\[
 \log(1/b)
 <2\left(u+\frac{u^3}{3}+\frac{u^5}{5(1-u^2)}\right).
\]
Exact rational arithmetic verifies
\[
 \frac{909}{500}
 -\frac{681}{100}\,2\left(u+\frac{u^3}{3}
                      +\frac{u^5}{5(1-u^2)}\right)
 >\frac1{1000}.
\]
The existing transfer with $K<681/100$ proves the result. The
strict surplus absorbs the asymptotic errors. The script
\texttt{certify\_same\_level\_one\_step\_transfer.py} records these
arithmetic comparisons and binds their analytic and numerical sources;
it does not certify those analytic sources by itself.
\end{proof}


\paragraph{A related roughness preprint and the scope of comparison.}
Ratliff's 2026 preprint \cite{Ratliff2026} uses the same starting
values $25/224$ and $249/448$ in a one-start Buchstab argument and
claims a roughness condition $P^-(N-p)>N^{0.2665-\epsilon}$ with
$\Omega(N-p)\le2$. This is relevant prior literature for that
construction, not an input to the present proof. A lower bound on
the smaller factor and the present upper bound $r\le q^b$ are
different restrictions: the roughness statement still permits
balanced semiprimes. The separate source audit records a failure of
the cited one-prefix distribution interface on
$(57/224,25/96)$ and a possible Bombieri--Vinogradov repair; it does
not disprove the final roughness assertion. No roughness theorem or
priority claim is deduced from the accompanying floating diagnostics.

% Derived route fragment, 2026-10-08. Not inserted in main/state.
% Companion JSON distinguishes exact algebra from floating diagnostics.

\subsection{A fourth-prefix refinement of the product upper sieve}
\label{sec:fourth-prefix-K-refinement-route}

This construction keeps every first- and third-prefix decision of
\texttt{extended-prefix-upper.tex}. It changes only the later odd
cube tests inside selected surviving third-prefix trees.
It therefore avoids making a rank-three decision depend on an
unseen fourth prime. The resulting coefficient improvement is
strictly positive on an explicit rational box.

Write
\[
 \rho=\frac{19}{96},\quad \mu=\frac{25}{128},\quad
 \Theta=\frac{299}{500},\quad \Lambda=\frac{153}{256},
 \qquad v(t)=\min\{(1+t)/2,57/64-3t/2\}.
\]
Let \(L_3(x,y,z)\) be exactly the limiting cap of the extended
construction, and put \(S_3=x+y+z\), \(S_4=S_3+w\).
Below \(x=\mu\), its existing simplified formula is
\[
 L_3=\max\{v(x),v(x+y),v(S_3),
                 \Lambda-3z/8,v(x+z),v(y+z)\};
\]
above \(\mu\) it is the original fallback \(v(x)\).
The minimum with \(\Theta\) is retained if needed.

For four ordered slots \(t=(x,y,z,w)\), define
\[
 W_4=\max\{2x,x+2y,x+y+2z,S_3+2w\}.
\]
For a nonempty subset \(I\subseteq\{1,2,3,4\}\), put
\[
 s_I=\sum_{i\in I}t_i,\qquad
 G_I=\max_{j\notin I}
       \left(s_I+\sum_{\substack{k<j\\k\notin I}}t_k+2t_j\right),
 \qquad \max\varnothing=0.
\]
The already derived selected-subproduct interface licenses
the endpoint \(v(s_I)\) when it strictly exceeds
\(\max\{1/2,W_4,G_I\}\), with the actual level taken strictly
between that gate and endpoint.

\paragraph{Admissible compact patches.}
Choose a finite union \(\mathcal C\) of compact positive exponent
boxes with strict ordering and fixed positive margins such that
\[
 0<w<z<y<x<\rho,\qquad S_3+2z<L_3,
 \qquad L_3<M<\Theta .
\]
Here \(M(t)\) is one fixed eligible selected-subproduct endpoint
on each branch of the patch, with
\[
 \max\{1/2,W_4,G_I\}<M=v(s_I).
\]
The inequality \(S_3+2z<L_3\) is the old third cube:
only third prefixes already retained by the old sieve are changed.
One may take inner approximations to a bounded piecewise-affine
region and then remove its gate and ordering strips.
The explicit patch below stays away from every such boundary.

Let
\[
 s_i=\frac{L_i-S_4}{w},\qquad (L_3,L_4)=(L_3,M),
 \qquad \overline F=F/(2e^\gamma).
\]
The old third cube gives \(L_3-S_4>2z-w>w\), so \(s_3>1\).
Define
\begin{equation}\label{eq:fourth-prefix-K-gain}
 \Delta_{\mathcal C}
 =4\int_{\mathcal C}
       \frac{\overline F(s_3)-\overline F(s_4)}
            {xyzw^2}\,dx\,dy\,dz\,dw.
\end{equation}
On a patch with \(s_4<3\), both arguments lie in \((1,3)\),
and this becomes the rational kernel
\begin{equation}\label{eq:fourth-prefix-K-rational-gain}
 \Delta_{\mathcal C}
 =4\int_{\mathcal C}
 \frac{M-L_3}{xyzw(L_3-S_4)(M-S_4)}\,dx\,dy\,dz\,dw.
\end{equation}

\begin{proposition}[Fourth-prefix descendant refinement]
\label{prop:fourth-prefix-descendant-refinement}
Assume the existing extended upper-weight, finite-prime-product
distribution, and selected-subproduct interfaces with their stated
strict margins. For every
\[
 K_*>K_{\rm ext}-\Delta_{\mathcal C},
\]
the same upper estimate with coefficient \(K_*\) holds for all the
original admissible labelled positive product families.
Their masses and multiplicities are unchanged.
The proposition provides a new weight construction and its analytic
coefficient; the large floating gain reported below is not a
certified numerical value.
\end{proposition}

\begin{proof}
Freeze the old coarse third labels and their caps \(\lambda_B\).
Inside a retained third tree, subdivide the fourth-prime range by
fixed coarse labels. On an active fourth box \(C\), select one
branch and a buffered cap \(\lambda_C\) satisfying
\[
 \lambda_B<\lambda_C<\nu<\Theta,\qquad
 \max_C\{1/2,W_4,G_I\}<\lambda_C<\inf_C v(s_I).
\]
Here the notation suppresses fixed positive margins, and
\(\nu\) is the old buffered global exponent. Refine the third and
fourth boxes enough that these intervals are nonempty.
Outside active fourth boxes, keep the old cap.

Let \(\mathcal R_3\) be the old raw upper support.
For ranks at most four, retain precisely its membership rule.
For rank at least five in an active box, retain the old first
and third cube tests and replace each later odd cube by
\[
 t_1+\cdots+t_{j-1}+3t_j<\lambda_C,
       \qquad j=5,7,\ldots.
\]
Do not retain the later old \(\lambda_B\) tests as extra conditions:
those are the conditions being relaxed.
Call the resulting support \(\mathcal R_4\).
A tuple whose third prefix failed the old test remains excluded.

This support is prefix closed. The fourth box makes no new
rank-four exclusion, and its cap is frozen for every descendant.
Every first exclusion is therefore still odd.
Equivalently, in the lower residual tree after the third prime,
the first possible changed exclusion has residual rank two,
which is full rank five. Thus
\(\mu(d)\mathbf1_{\mathcal R_4}(d)\) is an upper weight by the
same exact first-exit identity as the old construction.
This argument does not assume coefficientwise monotonicity of
M\"obius sums under support enlargement.

For factorization, all first four square conditions are below
\(\lambda_B\): the first three follow from the old upper support,
while
\[
 S_3+2w\le S_3+2z<\lambda_B.
\]
At every later odd slot the new cube implies its square test.
At a subsequent even slot, ordering gives
\[
 t_1+\cdots+t_{j-2}+t_{j-1}+2t_j
 \le t_1+\cdots+t_{j-2}+3t_{j-1}<\lambda_C .
\]
Hence every new complete rough vector belongs to the buffered
\(\mathcal D^{\rm well}(N^{\lambda_C})\).
The selected four-slot gate then supplies the same three-factor
partition as the previously derived four-prefix argument.
It remains fixed on the fine box. Appending the small factor,
perturbing the true local scale and the residue exponent, and
normalizing slot convolutions cost only the already budgeted
fixed margins. The cap remains below \(3/5\).

Apply the old parity boxing to both the inherited early cube tests
and the new later cube tests. The number of frozen levels is finite.
The coefficientwise upper comparison and the tagged-boundary estimate
therefore remain valid, with the same error type
\[
 O(C_{\rm grid}\tau\omega^{-5}
        +e^{-c/\eta}\omega^{-2}+o_N(1)).
\]
Attach the upper small-prime weight at even rough rank and the lower
one at odd rough rank. Unique small/rough decomposition, factorial
normalization, and the exact coprime-mass correction are unchanged.
The existing double Heath--Brown product distribution proof only
needs the verified factorization of these bounded atoms, not the
whole weight to have an unproved global factorability property.
All parameters are fixed before \(N\to\infty\).

It remains to justify the density subtraction.
For an active four-prime prefix \(P=p_1p_2p_3p_4\), its rank is even.
Its full subtree contribution is
\[
 g_N(P)V_N^+(N^\lambda/P,p_4).
\]
Indeed the subsequent full odd cubes are precisely residual odd cubes.
The active four-prefix trees are disjoint, so the exact identity is
\begin{equation}\label{eq:fourth-prefix-exact-density-difference}
 V_{\mathcal R_3}-V_{\mathcal R_4}
 =\sum_{\text{active }P}g_N(P)
 \left\{V_N^+(N^{\lambda_B}/P,p_4)
              -V_N^+(N^{\lambda_C}/P,p_4)\right\}.
\end{equation}
A one-sided upper estimate for each term is not enough for this
identity's limiting evaluation. The required ordinary upper-density
equality is proved in the next paragraph.
Together with prime-number-theorem integration it turns
\eqref{eq:fourth-prefix-exact-density-difference} into
\eqref{eq:fourth-prefix-K-gain}, since
\(2e^{-\gamma}F=4\overline F\).
Outside the compact patch, every coefficient and every tail estimate
is the old one. Thus there is no new improper-cutoff limit.
Taking the same ordered limits and using strict slack proves the
proposition.
\end{proof}

\paragraph{Uniform ordinary upper-density equality.}
For \((\beta,t)\) in a compact set with \(\beta/t>1\),
\[
 V_N^+(N^\beta,N^t)/V_N(N^t)=F(\beta/t)+o(1)
\]
uniformly in even \(N\). Here \(g_N(p)=
\mathbf1_{p>2,\ p\nmid N}/(p-1)\), and the support is the ordinary
cubic Rosser upper support.

The proof is the odd-exit counterpart of the already written
uniform residual lower-density lemma. The exact identity is
\[
 V_N^+/V_N=1+\sum_{\substack{j\ge1\\j\ {\rm odd}}}T_{j,N},
 \qquad T_{j,N}\ge0.
\]
For each fixed depth \(j\), the preceding square inequalities imply
that the last prime of a first exit exceeds \(N^{c_j}\), uniformly
on the compact parameter set. Prime-number-theorem integration gives
\(T_{j,0}=a_j(\beta/t)+o(1)\).
Deleting primes dividing \(N\) changes this by \(o(1)\), since
\(\sum_{p\mid N,\ p\ge N^{c_j}}1/(p-1)
 \ll N^{-c_j}\log N\).
The classical compact-uniform positive-exit series is
\[
 1+\sum_{\substack{j\ge1\\j\ {\rm odd}}}a_j(s)=F(s).
\]
It is the upper half of the same recurrence and convergence input
used in the lower-density lemma; see
\href{https://www.suzuyu1729.com/notes/NagoyaSeminar2022/RosserIwaniec.pdf}
{Suzuki, Sections 9--12, Propositions 9.4 and 11.8}.
A finite sum gives the required lower limit. The ordinary one-sided
linear upper-density bound gives the reverse limit.
Take \(N\to\infty\) first, then the fixed depth to infinity.
For \(1<s\le3\), one may also replace the cutoff by \(N^{\beta/3}\):
the raw support is unchanged because its first cube already forces
every supported prime below that threshold.
This proves the claimed equality without subtracting two one-sided
bounds or assuming a two-sided dimension estimate uniform in \(N\).

\paragraph{An exact positive box.}
Consider
\[
 \mathcal B=
 [17/100,89/500]\times[27/250,111/1000]
 \times[97/1000,99/1000]\times[17/200,19/200].
\]
Its ordering is strict. On this entire box the old simplified
envelope is \(L_3=(1+x)/2\).
For example \(y+z\ge.205>\mu\) gives
\(v(y+z)\le.583125<.585\le v(x)\); the other nonbaseline
branches are smaller as well.

Choose \(I=\{3,4\}\).
Then \(z+w\le.194<\mu\), and
\[
 M=v(z+w)=\frac{1+z+w}{2},\qquad M-L_3\ge\frac1{500}.
\]
The old third-cube margin and the nontrivial selected gate margin are
\[
 L_3-(x+y+3z)\ge\frac3{1000},\qquad
 M-(z+w+x+2y)\ge\frac3{1000}.
\]
The other selected gate \(z+w+2x\) has greater surplus;
the four square tests follow from the old third cube.
Moreover
\[
 1<\frac{106}{95}\le s_3<s_4\le\frac{131}{85}<3,
 \quad L_3-S_4\le\frac18,\quad M-S_4\le\frac{131}{1000}.
\]
Thus the rational kernel applies everywhere.
The box volume is \(3/6250000000\).
Bounding its positive numerator below and every denominator above gives
\begin{align}
 \Delta_{\mathcal B}
 &\ge
 \frac{4(1/500)(3/6250000000)}
 {(89/500)(111/1000)(99/1000)(19/200)(1/8)(131/1000)}
 \notag\\
 &=\frac{1024}{811431423}>\frac1{10^6}.
 \label{eq:fourth-prefix-exact-small-gain}
\end{align}
Every displayed comparison here is rational.
In particular this proves a strict improvement of \(K_{\rm ext}\)
within the inherited analytic inputs.
Combining it with the existing strict cap \(K_{\rm ext}<6.810\)
permits the deliberately small new cap \(6.809999\).
This is a local existence demonstration, not a useful optimized
numerical claim.

\paragraph{A coarse rigorous gain certificate.}
The separate script \texttt{certify\_fourth\_prefix\_K\_refinement.py}
now supplies a substantially larger exact lower bound. It uses grid
interiors \((k/800,(k+1)/800)\), with
\(32\le i_w<i_z<i_y<i_x<156\). Thus \(w>1/25\) and
\(x<39/200<\mu\). It tests just the three subsets
\(\{1,4\},\{2,4\},\{3,4\}\); it need not use all fifteen.
Every old-cap interval bound, selected gate, square condition and
\(s_4<3\) test uses integers with common affine denominator
\(D=25600\). On each retained box choose one certified branch.
The old cap is bounded above by an integer \(l\) divided by \(D\),
and the chosen new cap is bounded between \(m_-/D\) and \(m_+/D\).
Writing \(I=i_x+i_y+i_z+i_w\), that box contributes at least
\[
 \frac{4D(m_--l)}
 {(i_x+1)(i_y+1)(i_z+1)(i_w+1)
             (l-32I)(m_+-32I)}.
\]
The factor \(800^{-4}\) from the box volume has cancelled that
same factor in the product of coordinate upper endpoints.
The denominator is positive on every retained box.
Round each displayed rational number down onto the common
\(2^{80}\) lattice and sum. Of 9381251 disjoint candidate boxes,
483089 pass all tests. The resulting rigorous lower bound is
\[
 \Delta_{\rm grid}\ge
 \frac{10720642480494928338187}{604462909807314587353088}
 >\frac{3}{200}=0.015.
\]
All decisive comparisons and the branch selection are integer exact.
The companion JSON records the complete replay parameters and a
hash of the selected-box stream. No old certificate is rerun.
In particular, the previously established strict bound
\(K_{\rm ext}<681/100\) gives the safe new coefficient
\[
 K=\frac{1359}{200}=6.795.
\]
The old certificate's sharper exact endpoint may permit a smaller
rounded value; that comparison and all Chen-transfer arithmetic
are separate integration steps. The decimal value
\(0.01773581522795561\) is informational only.

\paragraph{A larger diagnostic region.}
For a fixed \(a_0>0\), take all ordered points \(w\ge a_0\),
\(x<\rho\), with the old third cube strictly present.
Test all fifteen subsets and put \(M\) equal to the largest
eligible endpoint above \(L_3\).
Keep only points with \((M-S_4)/w<3\).
The gain has exactly the rational kernel
\eqref{eq:fourth-prefix-K-rational-gain}; no interpolation of
sieve functions is involved.

Two independently scrambled Sobol samples of \(2^{19}\) points
gave the following floating diagnostics:
\[
\begin{array}{c|cc}
 a_0 & \text{seed }271 & \text{seed }902\\ \hline
 .03&.0279464644930&.0279862348447\\
 .04&.0278865355751&.0279691400924\\
 .05&.0276469125785&.0274150278547
\end{array}
\]
These are neither error bounds nor a certificate for a \(0.027\)
improvement. The domain includes sampled gate boundaries and the
samples cannot decide a theorem constant.
A finite collection of rational inner boxes with outward rational
kernel bounds would provide an appropriate next certificate.
The old \(K\) certificate need not be rerun to certify a subtraction
from its already established strict cap.

The route proves no improvement for old rejected third prefixes.
It imports no old double-sieve table and changes no Chen weights
or main/state file. Larger gains and their final consequences
require independent review and their own rigorous integration.

% Derived integration fragment, 2026-10-08. No documentclass.
% Numerical results in the companion JSON are floating diagnostics only.
\subsection{A parameterized one-step double sieve at a fixed total level}
\label{sec:generalized-one-step-double-sieve}

This extends the directly proved one-step construction in
\texttt{same-new-level-double-sieve-route.tex}; it does not transfer an
old $H,h$ table. The published combinatorial input is
\href{https://www.impan.pl/shop/en/publication/transaction/download/product/82961}
{Wu 2004, Lemma 4.1, (4.1)--(4.5), and (5.1)--(5.10)}.
The signed non-switched distribution input is the existing prefix
interface, using
\href{https://arxiv.org/html/2109.02851v2}{Lichtman 2021, (5.14)--(5.18)}
and
\href{https://arxiv.org/html/2309.08522v1}
{Lichtman 2023, Propositions 5.2 and 6.1}, with the already proved
selected-subproduct repair. The switched input is the separately
derived positive-count theorem with $K=681/100$. None of these uses
identifies that positive-count conclusion with the signed switched
distribution hypothesis in L21 (5.22).

Write $b=p_1\cdots p_r$, with decreasing specified exponents
$t_1>\cdots>t_r>0$, and set
\[
 T=\sum t_i,\quad Q=N^\nu,\quad R=Q/b=N^{L+o(1)},\quad
 L=\nu-T>0,\quad \alpha=\frac{1-T}{L},\qquad
 A(s)=\frac{sF(s)}{2e^\gamma},\quad a(s)=\frac{sf(s)}{2e^\gamma}.
\]
Here $a(s)=0$ for $s\le2$. Every formula below is the limit of
fixed positive exponent boxes, fixed level buffers, and the actual
residual weights at that one total level $Q$.

\paragraph{The precise non-switched family.}
We use either of the following two verified interfaces. In the enhanced
case $\nu>1/2$, require the original full-prefix square conditions
\begin{equation}\label{eq:generalized-one-step-squares}
 t_1+\cdots+t_{j-1}+2t_j<\nu\quad(1\le j\le r),
\end{equation}
and a factorization branch that remains valid after appending primes
below the smallest specified prime. The branches used below are the
largest-prime branch $v(t_1)$; a selected-subproduct branch involving
only the original specified slots, with all its remaining-square
gates; and a $W$ branch whose first three slots or upper bound on its
third slot persist. All require strict cap and gate margins.
The full level is below $153/256<3/5$.

If the current physical cutoff is $z=N^{L/s}$, we require
\begin{equation}\label{eq:generalized-one-step-cutoff}
 s\ge2,\qquad L/s\le t_r.
\end{equation}
For the added prime $q=N^{Lu}$ in Wu's negative term,
$1/s_0\le u\le1/s$, the new square condition is
$T+2Lu\le T+L=\nu$. The previous square conditions are unchanged.
Every subsequent residual rough square condition follows by
multiplying the residual condition by $bq$; its cutoff is
$y=R^{1/s_0}\le q$. Thus these are checks of the complete concatenated
vector, not just its bare specified product. For a selected branch,
the remaining-square tests for old unselected slots are unchanged;
before any new slot all original slots have already occurred, so
each new selected-branch test is the complete square test above.
The original signed atom coefficients, factorials, and small-prime
weight are retained by the box interface. This proves the signed
upper estimate for $S(\mathcal A_{N,b},y)$ and signed lower estimate
for $S(\mathcal A_{N,bq},y)$ at the same $Q$.
Where its lower argument is at most two, count nonnegativity supplies
the multiplier zero, without a distribution claim for that lower term.

In the classical case take $\nu<1/2$ with fixed margin. The complete
moduli of both residual sieve sums have support at most $N^\nu$.
Every specified prime is at least a fixed positive power $N^\epsilon$.
Consequently the number of ways any modulus can contain these fixed
labelled prefix slots is bounded by a constant depending only on
$r,\epsilon$. The residual modified weights have an absolute
coefficient bound depending only on their fixed construction
precision. Summing them at each complete modulus therefore leaves a
bounded coefficient, and the absolute Bombieri--Vinogradov remainder
applies, with residue $N$ on coprime moduli. The finitely many sieve
atoms are counted with their original signs; their number depends on
the fixed precision, not on $N$. This case does not require
\eqref{eq:generalized-one-step-squares}. In particular it covers the
classical part of $U_3$ even when $2t>1/2$.

\begin{proposition}[Explicit one-step upper defect]
\label{prop:generalized-one-step-defect}
On a compact positive prefix box satisfying one of the preceding
interfaces, suppose $2\le s<s_0\le5$, $\alpha\ge2$, and
\eqref{eq:generalized-one-step-cutoff}. Define
\begin{align}
 I(s,s_0,\alpha)
 &=\int_{1/s_0<u_1<u_2<u_3<1/s}
 \omega\!\left(\frac{\alpha-u_1-u_2-u_3}{u_2}\right)
                  \frac{du_1\,du_2\,du_3}{u_1u_2^2u_3},\label{eq:generalized-one-step-I}\\
 J(s,s_0,\alpha)
 &=A(s)-A(s_0)+\frac12\int_{1/s_0}^{1/s}
       \frac{a(s_0(1-u))}{u(1-u)}\,du
                   -\frac K8 I(s,s_0,\alpha).
                   \label{eq:generalized-one-step-J}
\end{align}
The ordinary normalized upper multiplier $A(s)$ can be replaced by
$A(s)-\max(0,J(s,s_0,\alpha))$. Any finite choice of $s_0$ on a common
positive prefix partition can be optimized box by box. In the limit
one can use the supremum over a fixed compact set of choices. This is
an averaged prefix estimate, not a bound for an individual prefix.
\end{proposition}

\begin{proof}
For a box indicator $\sigma(b)$, or a finite nonnegative constant
combination of such labelled indicators, put
\[
 \Phi_s=\sum_b\sigma(b)S(\mathcal A_{N,b},R^{1/s}),\qquad
 \Theta_Q=4\operatorname{li}(N)
       \sum_b\frac{\sigma(b)\mathcal C(bN)}{\varphi(b)\log(Q/b)}.
\]
Wu's pointwise identity and its prime-square error give
\[
 2\Phi_s\le2\sum_b\sigma(b)S(\mathcal A_{N,b},y)
 -\sum_b\sigma(b)\sum_{y\le q<z}S(\mathcal A_{N,bq},y)
 +\sum_b\sigma(b)\sum_{y\le q_1<q_2<q_3<z}
            S(\mathcal A_{N,bq_1q_2q_3},q_2)+o(\Theta_Q).
\]
The two non-switched inputs just proved yield, after division by
$\Theta_Q$, respectively $2A(s_0)$ and the negative of the integral
in \eqref{eq:generalized-one-step-J}. The lower argument is precisely
$s_0(1-u)$, not an argument computed at a different total level.

The positive switched multiset has the exact labelled form
\[
 m=bq_1q_2q_3n\le N,\quad y\le q_1<q_2<q_3<z,\quad P^-(n)\ge q_2.
\]
Canonically order the prime factors of $n$, retain all prefix labels,
and split into its finitely many possible factor counts. Every factor
is at least $N^{\epsilon_1}$ for some fixed $\epsilon_1>0$: the old
prefix is on a positive compact box, the new primes have exponent at
least $L/s_0$, and the quotient primes have the still larger threshold
$q_2$. All order, box, roughness, and total-product constraints are
finite monomial inequalities. Repeated factors have a power-saving
mass bound. The unit quotient has mass at most
$N^{T+3L/s+o(1)}=o(N)$ by a fixed power, since
$\alpha-3/s\ge1/2$. Thus it is not an omitted boundary atom.
The ordinary rough-number mass calculation gives the actual integral
\eqref{eq:generalized-one-step-I}, with the actual box value of
$\alpha$, rather than a supremum over $\alpha$. Its smallest Buchstab
argument is at least $\alpha s-3\ge1$. Any argument-one corner is
removed by a thin strip before taking limits.

For its exact unsifted mass $X$, the positive switching theorem gives
$(K+o(1))\mathcal C(N)X/\log N$. The partial summation in Wu (5.8),
with the same prefix measure, gives
\[
 X=(I(s,s_0,\alpha)+o(1))N
              \sum_b\frac{\sigma(b)}{\varphi(b)\log(Q/b)}.
\]
Using $\mathcal C(N)\le\mathcal C(bN)$ shows that the switched count
divided by $\Theta_Q$ is at most $KI/4+o(1)$. Dividing the weighted
inequality by two proves \eqref{eq:generalized-one-step-J}; the cost
is $KI/8$. No residual-level factor is inserted into $K$.
The required extension of the original fixed-cutoff statement of the
positive switching theorem is proved immediately below.
\end{proof}

\paragraph{Why the same $K$ permits every fixed positive factor cutoff.}
Here the original cutoff $100/1149$ must be reduced; it cannot simply
be reused. Fix any $\epsilon_1>0$ before $N\to\infty$. In the proof of
\texttt{uniform-prime-products-upper.tex}, replace its $a$ by
$\epsilon_1$ and its factor bound by $k_0=\lfloor1/\epsilon_1\rfloor$.
The modulus weights and the density integral in
\texttt{extended-prefix-upper.tex} are unchanged. The following checks
give the needed extension of that proof.

Apply the $k_{\rm HB}=3$ Heath--Brown identity to each original prime
slot at its own scale. There are at most $6k_0$ variables. Choose
$12\zeta<\epsilon_1/2$, in addition to all previous modulus margins,
and then choose the smaller distribution and local-scale precisions.
Before Mellin separation discard each dyadic rectangle incompatible
with its original prime-slot scale. The product of each group's at
most six lower endpoints is then at least $N^{\epsilon_1}/O(1)$,
including on the separated rectangular extension. If no individual
variable is in $[x^{2\zeta},x^{1/3+\zeta/2}]$, every original group
must contain a long variable: otherwise its product is at most
$x^{12\zeta+o(1)}$. There are at most two long variables globally.
Thus the no-middle-variable case is impossible for three or more
original groups. For one or two groups, the same one-long and
two-long estimates apply; the small complementary products have
the original bounds $x^{10\zeta+o(1)}$ and $x^{20\zeta+o(1)}$.
In the middle-variable case the individual coefficient is $1$ or
$\mu$, up to a logarithm and Mellin twist; its rough Siegel--Walfisz
input and the divisor-bounded complement remain the original ones.
Increasing the requested logarithmic saving absorbs the fixed larger
factor count, number of rectangles, and Mellin losses.

The exact coprime-mean replacement costs
$O(X_{\rm out}N^{-\epsilon_1}\log^C N)$ by the same union bound over
prime factors of a product. Repeated primes cost
$N^{1-\epsilon_1+o(1)}$, and prime powers introduced by the identities
cost $N^{1-\epsilon_1/2+o(1)}$; the divisor-moment treatment of the
progression errors preserves these savings. Products below
$N^{1-\kappa}$ are discarded as before. Smooth outer approximations
retain the lower cutoff $N^{\epsilon_1}/2$, the labels, and the exact
canonical quotient multiplicities. At fixed positive exponent boxes
their excess mass tends to zero. This proves the same positive-count
coefficient $K$ for the present finite product families. It asserts
neither uniformity as $\epsilon_1\downarrow0$ nor distribution for an
arbitrary arithmetic reweighting. In particular no shrinking Wu
short-box class is being substituted for these fixed boxes.

\paragraph{A single Buchstab transfer to the two-prime lower terms.}
Define a nonnegative chosen correction $j(s,\alpha)$ from the
proposition on $2\le s\le3$, and put it equal to zero outside that
interval. On a stable prefix family with cutoff $N^a$ and residual
argument $s=L/a$, use this transfer only for $2\le s<4$:
\begin{equation}\label{eq:generalized-one-step-h}
 h_1(s,\alpha)=\int_{\max(2,s-1)}^3
 j\left(w,\alpha+\frac{\alpha-1}{w}\right)\frac{dw}{w}.
\end{equation}
The normalized lower multiplier $a(s)$ can be replaced by
$a(s)+h_1(s,\alpha)$. To prove this, use the exact least-prime
Buchstab identity from cutoff $R^{1/4}$ to $N^a=R^{1/s}$.
The starting lower multiplier is $a(4)$. For a subtracted prime
$q=N^{L/(w+1)}$, the child residual argument at cutoff $q$ is $w$,
and its size ratio is
\[
 \alpha_{bq}=\frac{1-T-L/(w+1)}{L-L/(w+1)}
            =\alpha+\frac{\alpha-1}{w}.
\]
The normalization changes by $dw/w$. The ordinary part is exactly
$a(4)-\int_{s-1}^3 A(w)\,dw/w=a(s)$, with the usual lowered-cutoff
upper convention $A(w)=1$ for $w<2$. The corrected child part is
\eqref{eq:generalized-one-step-h}. Its primes satisfy
$q\le N^a\le p_r$, and its square condition is $2\log_Nq\le L$.
The second added prime used inside the child's one-step inequality
is at most $q$, with its own residual square condition. Thus both
generations retain the original stable cap. This is one finite tree;
no abstract $h,H$ iteration is being assumed. We choose zero extra
lower correction for $s<2$ and $s\ge4$.

\paragraph{The actual eleven-term domains.}
Take the open parameter range in \texttt{free-weight-thresholds.tex}:
$1/12<a<h<d<e=57/224$, $w_*=25/224<h<\min(25/128,2a)$,
$1/3<h+d<1/2$, and $4a+d,4h+d<1$.
The following are additive gains in the complete Chen coefficient,
already including its factor $1/4$.

For $L_2$ use $T=t$, $\nu=v(t)=(1+t)/2$, and cutoff $p=N^t$.
Here $2t<\nu$ and the largest-prime branch persists; $\alpha=2$.
The gain is
\[
 \Delta c_{L_2}=\int_{w_*}^h
 \frac{j((v(t)-t)/t,2)}{t(v(t)-t)}\,dt.
\]
Only $t\ge1/7$ contributes when $j$ is supported on $[2,3]$.

For $U_3,U_4$ the cutoff is lowered, when necessary, to exponent
$\min\{a,(\nu-t)/2\}$. Their $j$ argument is
$\max\{2,(\nu-t)/a\}$ and their ratio is
$(1-t)/(\nu-t)$. Use $\nu=v(t)$ for $a<t<e$, and for the
$U_3$ tail $e<t<1/3$ use the classical endpoint $\nu=1/2$ approached
from below. On the enhanced branch $2t<v(t)$; the physical cutoff
does not exceed $a\le t$. The gain integrand in either case is
\[
 \frac{j(\max\{2,(\nu-t)/a\},(1-t)/(\nu-t))}{t(\nu-t)}.
\]
Integrate over $a<t<1/3$ for $U_3$ and $a<t<d$ for $U_4$.

For $L_5,L_6$ use exactly the current two-prefix cap
\[
 \nu=V_{2,a}(t,u)=
 \begin{cases}
 \max\{v(t),v(t+u),153/256-3a/8\},&t\le25/128,\\
 v(t),&t>25/128.
 \end{cases}
\]
On the part $L=\nu-t-u\ge2a$, the first square test is
$2t<v(t)\le\nu$ and the second follows from
$L\ge2a>h\ge u$, namely $t+2u<\nu$. The cutoff in the
starting lower term is $R^{1/4}\le N^a$, and every added prime in
the exact Buchstab difference is also at most $N^a$. The original
two slots therefore remain the largest. The $v(t+u)$ branch keeps
its leading selected pair, and the $W$ branch keeps third exponent
at most $a$. It still has the gate supplied by $v(t)$, exactly as
in \texttt{double-prefix-level.tex}. Empty or short rough atoms use
the same classical fallback as that proof. Thus the gain is
\[
 \Delta c_{L_i}=\iint_{\Omega_i}
 \frac{h_1(L/a,(1-t-u)/L)}{tuL}\,du\,dt\quad(i=5,6),
\]
where $\Omega_5=\{a<u<t<h\}$ and
$\Omega_6=\{a<u<h<t<d\}$; the integrand is chosen zero outside
$2\le L/a<4$. Each cap is at most $(1+t+u)/2$, so $\alpha\ge2$.
For its constant branch this follows from
$153/256-3a/8\le1/2+a$, valid already for $a\ge25/352<1/12$.

For the two four-prefix upper terms retain all eligibility gates and
the cap $\nu=L_4(t)$ of \texttt{four-prefix-safe-minimum.tex}.
Choose cutoff exponent $\min\{w,(\nu-T)/2\}$ for the new upper
option. All four initial square tests hold by that construction;
all original selected slots, selected-branch gates, and the first
three slots for the $B(z)$ branch persist after either one-step
extension. Here $T\ge4a>1/3$, so $\nu\le153/256<(1+T)/2$.
Use the smaller corrected kernel inside the existing minimum:
\[
 \Delta c_{U_i}=\frac14\int_{\Omega_i}
 \left[\min\!\left\{\frac Kz\omega\left(\frac{1-T}{z}\right),
                         \frac4{\nu-T}\right\}
 -\min\!\left\{\frac Kz\omega\left(\frac{1-T}{z}\right),
 \frac{4(1-j(\max\{2,(\nu-T)/w\},(1-T)/(\nu-T)))}{\nu-T}\right\}
 \right]d\mu(t),\quad i=10,11.
\]
The integrand is zero off the eligible set. The original ordinary
normalized multiplier is one here because $a>1/12$, as already
proved for that minimum. This replacement avoids adding overlapping
gains from its switching and ordinary options.

Every positive contribution just used has a uniform fixed minimum
prime exponent on compact boxes. For the searched choices $s_0\le5$,
the four-prefix residual satisfies $L>w\ge a$; the corrected
two-prefix children have residual at least $4a/3$; and single-prefix
residuals are at least $1/6$ in the classical tail. For the displayed
numerical parameter tuple one can choose $\epsilon_1=a/6$ after
sufficiently small buffers. This is below the old fixed cutoff but
is exactly covered by the preceding Heath--Brown adaptation.

All boundaries are first trimmed. Take $N\to\infty$ with boxes,
branch choices, and buffers fixed; then remove the sieve and smooth
approximation errors, shrink the boxes and level buffers within the
strict domains, and finally exhaust those domains. A common finite
positive partition licenses every minimum and every choice of
$s_0$. The equality sets have zero integral mass. A numerical grid
interpolation alone is not a lower certificate for these limiting
integrals.

\paragraph{Floating feasibility, not a new certified constant.}
The companion \texttt{probe\_generalized\_one\_step.py} searches a finite
$s_0$ grid in $[3,5]$, tabulates $2\le s\le3$, $2\le\alpha\le8$,
and chooses zero correction outside that $s$ range. Its JSON records
coarse and fine Gauss quadratures, interpolation meshes, and the
separate gains. The baseline parameter tuple is
\[
 (a,h,d)=(.08455645453148589,.15049959287189102,
                           .2352266768449549).
\]
The current ordinary plus four-prefix-minimum diagnostic is
$c\approx1.87414321$. The coarse mesh gives an extra
$\Delta c\approx .01598683$, and the fine mesh gives $.01600361$.
In the fine calculation the contributions are
\[
\begin{array}{c|rrrrrrr}
 &L_2&U_3&U_4&L_5&L_6&U_{10}&U_{11}\\\hline
 \Delta c&.000136371&.013460255&0&.000094742&.002272783&
 .000039405&.000000051
\end{array}
\]
The resulting diagnostic is $1.89014682$, while $K\log(1/.749)$
is $1.96820097$; the remaining gap is about $.07805415$.
The zero-margin transfer diagnostic is $b\approx .75763420$.
The floating evaluation clamps a roundoff value just below two back
to the analytically proved bound $\alpha\ge2$, avoiding an erroneous
out-of-grid zero at the exact $\alpha=2$ boundary.
Thus this one-step family supplies genuine
analytic options but the sampled gain alone is far too small for
the requested frontier. The search is not an upper bound on all
possible one-step choices or on further finite recursions. Neither
the baseline nor these floating integrals certify a new rounded
counting coefficient; the previous certified constants are unchanged.

% Research fragment, 2026-10-08.  No change to the main manuscript/state.
% A finite second seed with explicit replacement kernels and unit masses.
\subsection{A repaired finite second double-sieve seed}
\label{sec:replacement-second-seed}

The purpose here is to replace the failed non-switched prefix extension
in Wu's second weighted inequality, not to import its old correction
table.  The argument uses the previously derived moving-residue
signed-box interface and positive finite-prime-product upper bound
$K=681/100$.  Its numerical certificate is separate from these inputs.
The source inequality is Wu 2004, Lemma 4.2; its 21 domains and all
their multiplicities are retained below.

Put
\[
 s=\frac{11}{5},\quad s_0=\frac{227}{50},\quad
 \phi_1=\frac{353}{100},\quad\phi_2=\frac{29}{10},\quad
 \phi_3=\frac{61}{25},\qquad
 a_i=1/\phi_i\ (1\le i\le3),\quad a_0=1/s_0,\quad b=1/s.
\]
Thus $a_0<a_1<a_2<a_3<b$.  Letters $a_i$ in this subsection are
relative endpoints, not the absolute initial sieve threshold.

\paragraph{An exact merge lemma.}
If a squarefree specified prefix $B$ belongs to
$\mathcal D^{\rm well}(D)$, and the distinct residual primes
$q_1>\cdots>q_m$ belong to $\mathcal D^{\rm well}(D/B)$ and do
not divide $B$, their sorted union belongs to
$\mathcal D^{\rm well}(D)$.
Indeed, at a residual prime its square product is at most the
residual square product times $B$.  At a fixed prefix prime $p$,
write $P$ for the product of larger fixed primes and $Q$ for the
product of larger residual primes.  If $Q=1$, use the old prefix
test $Pp^2\le D$.  Otherwise the last factor $q_*$ of $Q$ satisfies
$Qq_*\le D/B$ and $q_*>p$, hence
$PQp^2\le(Pp)(D/B)\le D$.
The lemma does not require the residual primes to lie below the
smallest specified prime.

\paragraph{The finite family on which the seed is proved.}
Retain a specified anchor prime of exponent
$t_0\in[29/200,3/20]$.  Fix a total level
$Q=N^\nu$, strictly below $(1+t_0)/2$ on its anchor box.
Every state used here has a nonnegative marked prime-box prefix
$B=N^{T+o(1)}$, containing that anchor, and satisfies
\begin{equation}\label{eq:second-seed-family}
 B\in\mathcal D^{\rm well}(Q),\quad
 T\ge t_0>2\nu-1,\quad T<\nu,\quad
 \text{every specified prime has exponent at most }\mu=25/128.
\end{equation}
Here the coefficient $\sigma$ is a finite nonnegative combination of
characteristic functions of marked prime boxes, with fixed constant
coefficients; no arbitrary prefix-dependent weight is asserted.
All boundary conditions carry fixed margins before passage to limits.
We only use finite collections of boxes, so their prime exponents
and $L=\nu-T$ have fixed positive lower bounds.  Write
\[
 R=Q/B=N^{L+o(1)},\qquad
 \alpha=\frac{1-T}{\nu-T}\ge2,\qquad
 L\le\frac{171}{400},\qquad
 \frac{L}{s}\le\frac{171}{880}<\mu.
\]
The last strict gap is $7/7040$.
Ordinary residual upper and lower weights are constructed at their
actual residual levels.  The merge lemma puts their marked rough
atoms into full well support.  The largest prime of every admitted
atom is at least the anchor and at most $\mu$, so its single-largest
cap $(1+x)/2$ is at least $(1+t_0)/2>\nu$.
Retain all marked slots, original residual factorials and signed
scalars; normalize each final convolution group by its group-size
factorial, restoring that factorial in the scalar.  Repeated-prime
collisions have the existing rough-square power saving.  Omitting
the finitely many specified primes below a residual cutoff changes
the Mertens factor by $1+O(N^{-\epsilon})$ on each fixed box.
Thus the existing signed atom proof, including its coprime mean and
local-scale moving-residue checks, applies after every indicated merge.

The positive product estimate is also needed for each fixed prime
cutoff $N^\epsilon$, not uniformly as $\epsilon$ tends to zero.
Its extension from the old displayed cutoff is obtained by the same
proof: allow at most $\lfloor1/\epsilon\rfloor$ original prime groups
and choose $12\zeta<\epsilon/2$.  In the double Heath--Brown
decomposition a group without a middle variable must contain a long
variable, because its six short variables could not produce a prime
at least $N^\epsilon$.  Three such groups would give three long
variables with product exceeding the total scale.  One and two groups
use the existing long-variable estimates.  A middle variable uses
the same Type II theorem; the complementary divisor order is still
fixed.  All other losses are powers of $N^{-c\epsilon}$ or fixed
logarithmic powers.  The modulus weights and their certified cap $K$
do not change.  This permits finite positive labelled monomial domains,
not arbitrary arithmetic reweightings.  For the single initial seed,
with $\nu$ sufficiently close to its stated cap,
the cutoff already satisfies $L/s_0>100/1149$; its finite descendants may need this fixed
$\epsilon$ extension.

\paragraph{The 21-term weight and its count check.}
Normalize by
\[
 \Theta_Q=4\operatorname{li}(N)
   \sum_B\frac{\sigma(B)\mathcal C(BN)}
                 {\varphi(B)\log(Q/B)},\qquad
 A(w)=\frac{wF(w)}{2e^\gamma},\quad
 \mathfrak a(w)=\frac{wf(w)}{2e^\gamma}.
\]
Write $\Phi_s$ for the prefix average sifted at $R^{1/s}$.
Wu's inequality is
\[
 5\Phi_s\le\Gamma_1-\Gamma_2-\Gamma_3-\Gamma_4
                       +\sum_{i=5}^{21}\Gamma_i+o(\Theta_Q).
\]
It also has a direct squarefree quotient verification.  On the
$R^{a_0}$-rough support let $A,B,C,D$ count prime factors in the
four successive intervals $[a_0,a_1),[a_1,a_2),[a_2,a_3),[a_3,b)$,
with exponents relative to $R$.  Then
\[
 \Gamma_1=4+\mathbf1_{A=0},\quad
 (\Gamma_2,\Gamma_3,\Gamma_4)=(A+B+C+D,A+B,A+B+C).
\]
The positive terms, pointwise and with $x_+=\max(x,0)$, are
\[
\begin{array}{c|l@{\qquad}c|l}
5&\binom{A+B}{2}&6&AC\\
7&(A-1)_+&8&\mathbf1_{A\ge1}B\\
9&\mathbf1_{A=0}(B+C-2)_+&10&\mathbf1_{A=0,B\ge2}(C+D)\\
11&\mathbf1_{A=0,B=1}(C-1)_+&12&\mathbf1_{A\ge2}D\\
13&\mathbf1_{A=1,B\ge1}(C+D)&14&\mathbf1_{A=1,B=0}(C+D-1)_+\\
15&\mathbf1_{A=0,B=1,C\ge1}D&16&\mathbf1_{A=B=0}(C-3)_+\\
17&\mathbf1_{A=B=0,C\ge3}D&18&\mathbf1_{A=B=0,C=2}(D-1)_+\\
19&\mathbf1_{A=0,B=1,C=0}(D-2)_+&20&\mathbf1_{A=B=0,C=1}(D-3)_+\\
21&\mathbf1_{A=B=C=0}(D-5)_+.&&
\end{array}
\]
These formulas follow because a quotient cutoff at the next-to-last
specified prime fixes all preceding specified primes as the smallest
factors.  Their signed total is nonnegative except at $A=B=C=D=0$,
where it is five.  More explicitly, for $A\ge2$ it is
$(A+B-2)(A+B-3)/2+(A-2)C$; for $A=1,B\ge1$ it is
$(B-1)(B-2)/2$; for $A=1,B=0$ it is $(1-C-D)_+$.
For $A=0,B\ge2$ it is $(B-2)(B-3)/2$.
For $A=0,B=1$, it is zero if $C\ge1$ and $(2-D)_+$ otherwise.
For $A=B=0$, it is zero if $C\ge3$, and respectively
$(1-D)_+,(3-D)_+,(5-D)_+$ when $C=2,1,0$.
This proves the pointwise inequality, independently of a table of
correction functions.  Off rough support all terms vanish; square
factors and intersections with the anchor are the usual power-saving
exceptions.

\paragraph{Non-switched kernels and the repair of $\Gamma_5$.}
The baseline upper term is $4A(s_0)+A(\phi_1)$.
The three negative terms have lower main terms
\[
 B(c)=\int_{a_0}^{c}
 \frac{\log(s_0(1-u)-1)}{u(1-u)}\,du,
 \qquad c=b,a_2,a_3.
\]
Their new prefix prime satisfies $u\le b<1/2$, so its singleton
belongs to residual well support; merge it with the old prefix.
Its lower argument $s_0(1-u)$ lies in $[681/275,177/50]\subset(2,4)$.
All needed ordinary lower atoms therefore belong to the family just
checked.  No unproved positive $h$ value is inserted.

For $\Gamma_5,\ldots,\Gamma_8$ the two additional exponents obey
$u<v$.  Their domains and ordinary upper kernels are
\[
\begin{array}{c|c|c}
i&\text{domain}&\text{ordinary kernel against }du\,dv/(uv)\\\hline
5&a_0<u<v<a_2&(1-u-v)^{-1}\\
6&a_0<u<a_1,\ a_2<v<a_3&(1-u-v)^{-1}\\
7&a_0<u<v<a_1&(1-u-v)^{-1}\\
8&a_0<u<a_1,\ a_1<v<a_2&(1-u-v)^{-1}.
\end{array}
\]
All effective upper parameters are at most $s_0-2=127/50<3$.
If an original parameter is below two, lower the sifting cutoff to
half the residual level exponent; this preserves an upper bound and
the displayed kernel because $wF(w)/(2e^\gamma)=1$ up to three.

The residual pair has full well support precisely when
$2v<1$ and $v+2u<1$.  The first inequality always holds.  The second
holds throughout domains $6,7,8$, since their respective suprema are
\[
 a_3+2a_1<1,\qquad3a_1<1,\qquad a_2+2a_1<1.
\]
In domain 5 keep the ordinary kernel on $v+2u<1$.
On its bad corner
\[
 \mathcal B=\{a_0<u<v<a_2:\ v+2u>1\}
\]
use the positive $K$ count.  Its quotient cutoff is $R^{a_0}$,
its unit quotient has a power-saving product range, and its rough
argument is at least $s_0(2-2a_2)>5/3$ for every $\alpha\ge2$.
Thus $\omega\le3/5$ gives the valid replacement kernel
\begin{equation}\label{eq:second-seed-kernel}
 k_5(u,v)=
 \begin{cases}
 (1-u-v)^{-1},&v+2u<1,\\
 3Ks_0/20,&v+2u>1.
 \end{cases}
\end{equation}
The common normalization is $K/4$ before multiplying by $\omega/a_0$.
The bad domain splits exactly at
$u=(1-a_2)/2$, $u=1/3$, and $u=a_2$: its lower $v$ endpoint is
$1-2u$ on the first interval and $u$ on the second.
This is an actual replacement upper bound, not a subtraction of a
missing old $H$ value.

\paragraph{All positive switch domains and their unit masses.}
Use the notation $[\ell,r]^k_\uparrow$ for $k$ ordered exponents
in that interval; juxtaposition denotes independent successive bands.
The unchanged domains $9,\ldots,21$ are
\[
\begin{array}{c|l@{\qquad}c|l}
9&[a_1,a_3]^3_\uparrow&10&[a_1,a_2]^2_\uparrow[a_2,b]\\
11&[a_1,a_2][a_2,a_3]^2_\uparrow&12&[a_0,a_1]^2_\uparrow[a_3,b]\\
13&[a_0,a_1][a_1,a_2][a_2,b]&14&[a_0,a_1][a_2,b]^2_\uparrow\\
15&[a_1,a_2][a_2,a_3][a_3,b]&16&[a_2,a_3]^4_\uparrow\\
17&[a_2,a_3]^3_\uparrow[a_3,b]&18&[a_2,a_3]^2_\uparrow[a_3,b]^2_\uparrow\\
19&[a_1,a_2][a_3,b]^3_\uparrow&20&[a_2,a_3][a_3,b]^4_\uparrow\\
21&[a_3,b]^6_\uparrow.&&
\end{array}
\]
For $r$ extra primes the quotient cutoff is the next-to-last prime,
of exponent $w=u_{r-1}$.  The normalized unsifted mass is the sum
\begin{equation}\label{eq:second-seed-full-mass}
 \int_{\substack{\boldsymbol u\in\mathcal D_i\,,\ \sum u_j\le\alpha-w}}
 \frac{\omega((\alpha-\sum u_j)/w)}{w\prod u_j}\,d\boldsymbol u
 +\int_{\substack{\boldsymbol u\in\mathcal D_i\,,\ \sum u_j=\alpha}}
             \frac{du_1\cdots du_{r-1}}{\prod u_j}.
\end{equation}
In the surface term eliminate a specified coordinate with derivative
one; it counts quotient $1$ and is not part of the ordinary Buchstab
function.  The formula follows by the prime number theorem and the
rough-number estimate on fixed positive boxes.  Boundary strips can
be removed before the limit; an upper bound for the surface term is
also obtained by counting the eliminated prime last.

For $9\le i\le15$, the rough argument is at least $5/3$ already at
$\alpha=2$, hence for every $\alpha\ge2$.
For example the two smallest corner bounds are $47/25$ and
$3136/1765>5/3$; the other five are larger.
Use $\omega\le3/5$.  For $16\le i\le21$ use $\omega\le1$ and
enlarge the nonunit domain if necessary.  There is no unit mass for
$9\le i\le19$, since every sum is at most $4b=20/11<2$.

Put $\ell_{ij}=\log(\phi_i/\phi_j)$,
$\ell_{01}=\log(s_0/\phi_1)$, and
$\ell_{3b}=\log(\phi_3/s)$.
The following explicit masses $M_i$ majorize
\eqref{eq:second-seed-full-mass} for every $\alpha\ge2$:
\begin{align*}
 M_9&=\tfrac35\{(\phi_1+\phi_3)\log(\phi_1/\phi_3)-2(\phi_1-\phi_3)\},\\
 M_{10}&=\tfrac35\{\phi_1-\phi_2(1+\ell_{12})\}\log(\phi_2/s),\\
 M_{11}&=\tfrac35\ell_{12}\{\phi_2(\ell_{23}-1)+\phi_3\},\\
 M_{12}&=\tfrac35\{s_0-\phi_1(1+\ell_{01})\}\ell_{3b},\\
 M_{13}&=\tfrac35\ell_{01}(\phi_1-\phi_2)\log(\phi_2/s),\\
 M_{14}&=\tfrac35\ell_{01}\{\phi_2(\log(\phi_2/s)-1)+s\},\\
 M_{15}&=\tfrac35\ell_{12}(\phi_2-\phi_3)\ell_{3b},\\
 M_{16}&=\phi_2\ell_{23}^4/24,\qquad
 M_{17}=\phi_2\ell_{23}^3\ell_{3b}/6,\\
 M_{18}&=\phi_3\ell_{23}^2\ell_{3b}^2/4,\qquad
 M_{19}=\phi_3\ell_{12}\ell_{3b}^3/6,\\
 M_{20}&=\phi_3\ell_{23}\ell_{3b}^4/24
                         +\phi_2\ell_{3b}^4/24,\\
 M_{21}&=\phi_3\ell_{3b}^6/720+\phi_3\ell_{3b}^5/120.
\end{align*}
The second summand in $M_{20}$ bounds the unit surface by eliminating
$u_1$, using $1/u_1\le\phi_2$, and dropping its interval restriction.
The other four variables are ordered in $[a_3,b]$.
For $M_{21}$ eliminate $u_6$, use $1/u_6\le\phi_3$, and integrate
the other five ordered variables.  Thus these are uniform mass bounds
even when a descendant has $\alpha>2$.

At the initial value $\alpha=2$ there is a sharper useful check:
domain 21 is empty, and domain 20 has no nonunit quotient because
$a_2+5a_3>2$.  Its unit surface has
\[
 \delta=2-a_2-4a_3=28/1769,\qquad
 M_{20}^{\rm unit}(2)\le
 \frac{\delta^4}{(4!)^2a_2a_3^4}
 =\frac{4802}{428712890625}.
\]
Indeed write its last four coordinates as $a_3+x_i$;
$x_i\ge0$, $\sum x_i\le\delta$, ordered, and the first coordinate
is $a_2+\delta-\sum x_i$.  Its other endpoint restrictions are
automatic here.  The ordered simplex volume is $\delta^4/(4!)^2$.
This sharper initial mass is not used when iterating: the uniform
$M_{20},M_{21}$ above are used instead.

\begin{proposition}[A rigorously bounded finite replacement seed]
\label{prop:replacement-second-seed}
On the finite family \eqref{eq:second-seed-family}, with the fixed
positive-cutoff version of the $K$ input just described, one has
\[
 \Phi_{11/5}\le\{1-J_*+o(1)\}\Theta_Q,
 \qquad J_*=9/500.
\]
This statement uses ordinary lower baselines only and invokes no
previous $H,h$ table or unproved child correction.
\end{proposition}

\begin{proof}
Let $U_i$ for $5\le i\le8$ be the integrals of the displayed ordinary
kernels against $du\,dv/(uv)$, and put
\[
 E_5=\iint_{\mathcal B}
 \left\{\frac{3Ks_0}{20}-\frac1{1-u-v}\right\}\frac{du\,dv}{uv}.
\]
It follows from the established termwise bounds and the common
$K/4$ normalization that
\[
 J_2=\frac15\left\{5-4A(s_0)-A(\phi_1)
       +B(b)+B(a_2)+B(a_3)-\sum_{i=5}^8U_i-E_5
                         -\frac K4\sum_{i=9}^{21}M_i\right\}
\]
is a valid saving.  Every quantity here is evaluated at the same
total level $Q$.

The file \texttt{certify\_replacement\_second\_seed.py} uses 80-bit
outward integer intervals, 32 positive atanh terms with explicit
remainders, and 4096 rational Darboux cells for each one-dimensional
integral.  The double ordinary kernels are integrated once explicitly:
\[
 \int_l^r\frac{dv}{uv(1-u-v)}
  =\frac1{u(1-u)}\log\frac{r(1-u-l)}{l(1-u-r)}.
\]
The split bad domain uses its two exact lower endpoints stated above.
The switch bounds are the explicit logarithmic expressions $M_i$;
there is no floating quadrature input.  The certificate gives
\[
 J_2\ge
 \frac{22417233603428441622629}{1208925819614629174706176}
 >\frac9{500}.
\]
The strict surplus absorbs all previously fixed asymptotic and box
errors.  This proves the assertion under the named analytic inputs.
\end{proof}

\paragraph{A finite propagation, with its lower-child domain retained.}
There is a small fully specified use at the original Buchstab cutoff.
Start with an anchor prefix and upper target argument three.  Buchstab
from arguments three to four produces lower children with arguments
$y\in[2,3]$ and remaining exponent at most $3L/4$.
The ordinary upper baseline at four is licensed by the same merge
argument.  On $2\le y\le5/2$, apply a second Buchstab identity from
$y$ to $16/5$ to that child.  Its ordinary lower baseline is licensed
because all residual primes are at most
$(3L/4)/(16/5)<\mu$.  Its upper grandchildren have arguments
$z\in[y-1,11/5]$.
For $z\le11/5$, monotonicity of the positive sifted count gives the
seed bound at argument $z$ from the proved bound at $11/5$;
the normalized $A$ equals one on this entire interval.
It only lowers the comparison cutoff.

Both newly appended singleton primes have square at most the current
residual level, so the merge lemma preserves full prefix well support.
The first lower-child cutoff is at most $3L/8<\mu$, and all second
appended primes have that same upper bound.  Hence no largest prime
leaves $[t_0,\mu]$.  Total prefix exponents increase, so
$\alpha=(1-T)/(\nu-T)$ increases and remains at least two;
the uniform masses, not the initial $\alpha=2$ masses, are used.
Only finitely many boxes and a fixed positive prime cutoff are needed.

Cancellation of the ordinary $A,\mathfrak a$ Buchstab main terms gives
the explicit defect bounds
\begin{equation}\label{eq:replacement-second-seed-h3}
 h_{\rm child}(y)\ge J_*\log\frac{11/5}{y-1}
 \quad(2\le y\le5/2),\qquad
 H_{\rm anchor}(3)\ge
 J_*\log\frac54\log\frac{22}{15}
 >\frac7{4625}.
\end{equation}
The last inequality uses only
$\log(5/4)>2/9$ and $\log(22/15)>14/37$.
Here $h_{\rm child}$ and $H_{\rm anchor}$ name the bounds on these
particular finite nodes, not an unrestricted old one-variable table.
At the actual anchor cutoff $p$, the argument tends to a value in
$[17/6,171/58]\subset(2,3)$ as the level approaches its cap;
count monotonicity therefore retains
the saving $7/4625$.

If the $L_2$ threshold exceeds $3/20$, the block
$29/200<t<3/20$ is available and yields
\begin{equation}\label{eq:replacement-second-seed-l2-gain}
 \Delta c\ge2\frac7{4625}\log\frac{513}{493}
 >\frac{56}{465275}>0.00012.
\end{equation}
This is a concrete small consequence.  It falls far short of the
approximately $0.094$ gap in the best sampled ordinary counting route.
The repaired finite seed itself is about $0.0185$ under the rigorous
coarse envelopes (the floating exact-$\omega$ probe is about $0.02081$).
It may be used as an input to a larger, explicitly checked finite tree;
it does not authorize restoring the old nine-by-nine matrix unchanged.
If another correction already applies on one of these boxes, take the
stronger valid bound there rather than adding overlapping defects.
No new global Chen coefficient, imbalance exponent, or priority claim
is asserted by this fragment.

% Integration fragment, 2026-10-08.  Does not modify the locked second seed.
\subsection{The second seed on stable specified-prefix families}
\label{sec:second-seed-stable-prefix}

The numerical seed used here is exactly the proved $J_*=9/500$ of
Proposition~\ref{prop:replacement-second-seed}, with
\[
 s_*=11/5,\quad s_0=227/50,\quad
 (\phi_1,\phi_2,\phi_3)=(353/100,29/10,61/25).
\]
Its proof is now extended to the stable factorization branches used
in Section~\ref{sec:generalized-one-step-double-sieve}.  No old
correction table is used and no new numerical integration is needed.
The extension keeps the repaired $\Gamma_5$ kernel, all other
ordinary kernels, and the separate unit masses in $\Gamma_{20},
\Gamma_{21}$ without change.

Write the specified squarefree prefix as $B=p_1\cdots p_r$, with
decreasing exponents $t_1>\cdots>t_r>0$, and put
\[
 T=\sum_i t_i,\qquad Q=N^\nu,\quad R=Q/B=N^{L+o(1)},\quad
 L=\nu-T>0,\qquad \alpha=(1-T)/L\ge2.
\]
All prefix coefficients are finite nonnegative constant combinations
of marked prime-box indicators.  The prime exponents and $L$ have
fixed positive lower bounds on every box.  As usual all box,
factorization, and level inequalities have strict margins before
taking their endpoint limits.

\paragraph{Precisely which factorization data are retained.}
In the enhanced case $1/2<\nu<153/256$, require the full specified
square conditions
\[
 t_1+\cdots+t_{j-1}+2t_j<\nu\qquad(1\le j\le r).
\]
Choose one of the already proved cap branches, below its cap and
above every gate: the largest-prime branch; a selected-subproduct
branch using specified slots and their remaining-square gates; or
the repaired $W$ branch with its first three slots or a retained
bound on its third slot.  The chosen branch must persist on adding
primes of exponent below a specified ceiling $\chi\le t_r$.
For branches determined by existing slots one can take $\chi=t_r$.
If there are fewer than three specified slots and the $W$ branch
uses only the assertion that a future third exponent is at most
$c$, one must also impose $\chi\le c$.  The condition
$\chi\le t_r$ alone does not preserve such a third-slot bound.

In the classical case $\nu<1/2$ with a fixed margin, use the absolute
Bombieri--Vinogradov remainder for complete moduli.  There is no
full-well or ordering requirement on the specified prefix in that
case.  Each complete modulus has only boundedly many ways to contain
the finitely many labelled prime slots, because every prime slot
exceeds a fixed positive power of $N$.  The residual construction
has bounded coefficients for each fixed precision, so the complete
coefficient remains bounded.  This is the classical interface of
Section~\ref{sec:generalized-one-step-double-sieve}, including its
coprimality and collision treatment.

\begin{proposition}[Second seed with stable prefix data]
\label{prop:second-seed-stable-prefix}
Under the named moving-residue and positive $K=681/100$ inputs of
the preceding seed, suppose $\alpha\ge2$ and one of these interfaces
holds.  In the enhanced case assume additionally
\begin{equation}\label{eq:second-seed-comparison-cutoff}
 L/s_*<\chi.
\end{equation}
Then, with the same normalization $\Theta_Q$, one has
\[
 \Phi_{s_*}\le(1-J_*+o(1))\Theta_Q.
\]
For any $0<w\le s_*$, positivity and cutoff monotonicity imply the
same upper bound for $\Phi_w$.  In particular the normalized ordinary
upper multiplier $A(w)=1$ for $1\le w\le s_*$ may be replaced by
$1-J_*$.
\end{proposition}

\begin{proof}
All specified primes added by the 21-term inequality are below
$R^{1/s_*}$.  For every non-switched term the subsequently sieved
prime slots are smaller still.  Under
\eqref{eq:second-seed-comparison-cutoff} all these slots are below
the old specified prefix and below its stable ceiling $\chi$.
Thus the original largest prime, original selected subset, its old
remaining-square gates, and the applicable $W$ data remain unchanged.

The residual singleton in each negative term satisfies $2u<1$,
where $u\le1/s_*<1/2$ is relative to $R$.  The residual pair in
$\Gamma_6,\Gamma_7,\Gamma_8$ satisfies both $2v<1$ and $v+2u<1$.
These same tests hold on the good portion of $\Gamma_5$.
The exact merge lemma proved with the seed combines each such
residual vector with the old full-well prefix.  It then combines
their ordinary residual sieve atoms at the correct residual level.
Equivalently, since all old slots precede all new slots, the old
complete and selected-remaining tests are unchanged, and each new
test is the corresponding residual square test multiplied by $B$.
This explicitly checks the complete atom, not just its bare prefix.

Every factorization is performed on buffered endpoint vectors.
Keep all marked slots, the original scalar signs, and the original
residual factorials; use the existing bounded-group normalization
and moving-residue argument.  The small-prime weight and the
exclusion of primes dividing the specified prefix are treated by
the same signed interface.  Thus all non-switched ordinary terms
retain exactly their previous main kernels.

On the bad portion of $\Gamma_5$ no signed distribution statement
is used.  The positive product estimate gives the existing
replacement $3Ks_0/20$.  The proof of every positive switched mass
uses only $\alpha\ge2$, a fixed positive exponent cutoff, and the
retained labelled domains.  It is independent of the former anchor
interval and of the largest exponent of $B$.  In particular the
uniform $M_{20},M_{21}$ still include their separate unit surfaces;
one does not reuse their special $\alpha=2$ simplifications.

In the classical case, all complete ordinary moduli have support
at most $Q$, including on the bad portion of $\Gamma_5$.  Absolute
Bombieri--Vinogradov therefore licenses even that ordinary kernel.
For a common conservative formula we may instead keep the positive
replacement there; either choice is a valid upper bound.  No
selected-subproduct or full-well assertion is needed for this case.

All main kernels in the certified functional are consequently
unchanged.  Its already saved outward rational certificate gives
$J_2>9/500$.  Finally, if $w\le s_*$, then
$R^{1/w}\ge R^{1/s_*}$ and
$S(\mathcal A_{N,B},R^{1/w})\le
S(\mathcal A_{N,B},R^{1/s_*})$ pointwise.  The comparison cutoff
is lowered; the normalization uses the same total level $Q$.
\end{proof}

\paragraph{The correction to use and the exact lower transfer.}
Write
\[
 j_2(w)=J_*\mathbf1_{1\le w\le s_*}.
\]
At an upper node with $1\le w\le s_*$ and physical cutoff
$R^{1/w}\le N^\chi$, condition
\eqref{eq:second-seed-comparison-cutoff} follows automatically,
after the usual trimming of equality sets.  More generally it is
enough to check the comparison cutoff directly.  If $j_1$ is any
proved first-seed correction on the same physical prefix box, use
$\max\{j_1,j_2\}$ there.  Their defects are alternative upper
bounds for the same count and are not added.

For a stable lower prefix with cutoff $N^a\le N^\chi$, write
$s=L/a\ge2$.  Buchstab from $R^{1/4}$ to $R^{1/s}$, when $s<4$,
gives children indexed by $w\in[s-1,3]$.  The appended prime and
the child ratio are
\[
 q=N^{L/(w+1)},\quad L'=\frac{Lw}{w+1},\qquad
 \alpha'=\alpha+\frac{\alpha-1}{w}\ge2.
\]
Every child used by the second seed has $1\le w\le s_*$.
Its new square condition follows from $w\ge1$, namely
$2\log_Nq\le L$.  Also $q\le N^a\le N^\chi$, so the original
branch persists, and its seed comparison cutoff satisfies
\begin{equation}\label{eq:second-seed-child-cutoff}
 \frac{L'}{s_*}=\frac{L}{w+1}\frac{w}{s_*}
 \le\log_Nq.
\end{equation}
In particular the seed's new prime slots are below this child's
new smallest specified prime.  This remains true for $1\le w<2$.
It would lose a valid part of the estimate to truncate the integral
at $w=2$.

The ordinary identity
$a(4)-\int_{s-1}^3 A(w)\,dw/w=a(s)$ holds for $2\le s<4$,
including $A(w)=1$ on $1\le w<2$.  After inserting the uniform seed
defect and retaining its exact prefix measure $dw/w$, the new lower
correction is
\begin{equation}\label{eq:second-seed-stable-lower}
 h_2(s)=
 \begin{cases}
 \displaystyle\frac9{500}\log\frac{11/5}{s-1},&2\le s<16/5,\\
 0,&s\ge16/5.
 \end{cases}
\end{equation}
Thus the normalized lower multiplier is at least $a(s)+h_2(s)$.
This is a lower bound from one finite tree; it is not an abstract
old $h$ value.  At $s=2$, its limiting value is
$(9/500)\log(11/5)$, approximately $0.01419$.
When combining seeds, integrate their pointwise maximum on the
child boxes; do not add $h_2$ to a first-seed correction that already
uses any of the same children.

\paragraph{Actual domains in the complete eleven-term count.}
Use the parameters and level prescriptions of
Section~\ref{sec:generalized-one-step-double-sieve}.  All the
following assertions concern the limiting cap approached from
strictly admissible levels.
\begin{itemize}
\item For $L_2$, the largest-prime branch has
$\nu=(1+t)/2$ and target $w=(1-t)/(2t)$, on the relevant
$t<h<25/128$ range.  The direct second seed applies exactly when
$5/27\le t<h$.  This interval is empty for the current sampled
$h\simeq0.1505$; the separately proved finite $H(3)$ propagation
on $0.145<t<0.15$ is a different available option.
\item For $U_3,U_4$, set $L=\nu-t$ and use the current physical
upper argument $w=\max\{2,L/a\}$.  The direct seed is applicable
on $0<L\le(11/5)a$, since its comparison cutoff is at most $a<t$.
Here $\nu=v(t)$ on the enhanced part.  On the classical $U_3$
tail $57/224<t<1/3$, it applies on
$t\ge1/2-(11/5)a$.  Keep the corrected coefficient
$4(1-\max\{j_1,j_2\})/L$ in place of the ordinary option.
\item For $L_5,L_6$, retain their existing level $V_{2,a}(t,u)$
and all its gates.  Use \eqref{eq:second-seed-stable-lower} where
$2\le(\nu-t-u)/a<16/5$.  In this application the first appended
prime $q$ is at most $N^a$, so the two original primes remain the
largest.  If the chosen $W$ cap used a third-slot bound $a$, this
child now has an actual third slot $q\le N^a$.  Every seed prime
is below $q$ by \eqref{eq:second-seed-child-cutoff}, and the bound
is therefore preserved.  The selected leading pair is preserved
in the same way.  The gain integrand is
$h_2((\nu-t-u)/a)/(tu(\nu-t-u))$ before combining with any
overlapping first-seed correction.
\item For $U_{10},U_{11}$, retain the full eligibility set and
all branch gates of the four-prefix safe minimum.  Put
$T=x+y+z+w_0$, $L=\nu-T$ and let $w_0$ be the smallest specified
exponent.  The physical upper argument is
$\max\{2,L/w_0\}$.  The direct second seed applies when
$L\le(11/5)w_0$, and its comparison cutoff is then below the
smallest old slot.  All original selected slots and the first
three $W$ slots persist.  Replace the ordinary option inside the
existing switching/ordinary minimum by
$4(1-\max\{j_1,j_2\})/L$; do not add this saving to an already
overlapping switched improvement.
\end{itemize}
For a common finite positive partition these are genuine averaged
prefix bounds.  The endpoint comparisons and cutoff equality sets
are approached through buffered boxes.  There is no assertion for
an individual fixed tuple of primes.

\paragraph{Finite recursive closure, not a restored old table.}
The state conditions also permit further finite Buchstab trees.
At an upper or lower parent with physical argument $s\ge2$ and
cutoff below its stable ceiling, an appended prime indexed by
$w\ge s-1\ge1$ has square at most the current residual level and
is below that ceiling.  Retaining the original branch data and
using this new prime as the child's cutoff preserves the same
typed family.  Its actual ratio changes to
$\alpha+(\alpha-1)/w$; it is not reset to two.

For example, if uniform defect bounds $H,h$ have already been proved
on all nodes of a finite such tree, and $S_{\max}+1\ge s$, ordinary
Buchstab cancellation gives the further valid choices
\[
 h_{\rm new}(s)=\int_{s-1}^{S_{\max}}H(w)\frac{dw}{w},\qquad
 H_{\rm new}(s)=\int_{s-1}^{S_{\max}}h(w)\frac{dw}{w},
 \quad s\ge2,
\]
where unproved corrections are set to zero; in particular $h(w)=0$
for $w<2$, whereas the proved seed gives $H(w)=J_*$ for
$1\le w\le s_*$.  These are alternative lower bounds for the
defect, so one may take their maximum with a previous bound rather
than adding overlapping estimates.  Every generation is finite,
and its positive prime cutoff remains fixed before $N\to\infty$.
The first- and second-seed identities can similarly use previously
proved child defects at their actual states, with their original
coefficients.  This paragraph licenses an explicitly checked finite
construction; it does not supply an unevaluated numerical matrix
or an infinite-iteration convergence claim.

One explicit consequence needs no new integration certificate.
If the upper comparison cutoff $R^{1/3}$ is below the stable ceiling,
the same two Buchstab steps used for
\eqref{eq:replacement-second-seed-h3}, now in this enlarged family,
give
\begin{equation}\label{eq:second-seed-stable-h3}
 H(3)\ge J_*\log(5/4)\log(22/15)>7/4625.
\end{equation}
The first new prime is at most $R^{1/3}$; the second is at most
the first, and the final seed comparison cutoff is at most the
second.  Thus every old branch datum persists.  The proof uses
only the uniform $\alpha\ge2$ seed, so this consequence has no
restriction to the original anchor interval.  By count monotonicity
it also bounds upper nodes with physical argument $1\le w\le3$,
provided the comparison cutoff $R^{1/3}$ is licensed.  It is another
option to maximize with the direct seed and other proved defects.

The saved companion JSON binds the already certified scalar seed
and records these analytic domains.  No new total counting
coefficient, transfer exponent, global optimality claim, or
resolution of the requested frontier follows from this extension
alone.  Any numerical improvement after taking the common-box
maxima requires its own outward integral certificate.

% Derived integration fragment, 2026-10-08. Floating diagnostics separate.
\subsection{Monotone cutoff envelopes for the finite double-sieve defects}
\label{sec:monotone-one-step-double-sieve}

Use the same fixed-level prefix states and notation as
\texttt{generalized-one-step-double-sieve.tex}. In particular the
non-switched family is either a hereditary enhanced-level branch
with complete square tests, or a classical complete-modulus family
of level below $1/2$. Keep the prescribed original positive prefix
boxes and the original total level $Q$. This supplement only enlarges
the explicitly chosen defects; it imports no old $H,h$ table.

Write $R=Q/B=N^{L+o(1)}$, $\alpha=(1-T)/L\ge2$, and let the
physical target cutoff be $z=R^{1/s}$. If a normalized upper bound
$\Phi_v\le(A(v)-j_v+o(1))\Theta_Q$ is proved at $v\ge s$, then
\[
 R^{1/v}\le R^{1/s},\qquad \Phi_s\le\Phi_v.
\]
For $s\le3$ and $2\le v\le3$ both ordinary normalized multipliers
are one, including the lowered-cutoff convention at $s<2$.
Thus a legal explicit upper defect is
\begin{equation}\label{eq:monotone-one-step-envelope}
 j_{\rm mono}(s,\alpha)
 =\sup_{\max(2,s)\le v\le3}\ \sup_{v<s_0\le5}
                  \max(0,J(v,s_0,\alpha)).
\end{equation}
Set it equal to zero for $s>3$. In the numerical probe the inner
supremum is restricted to a finite grid in $[3,5]$. In a proof one
chooses only finitely many comparison parameters on a common positive
box partition, then takes finer partitions. The interpolation in a
floating table is a diagnostic, not a numerical certificate.

When an enhanced branch uses a bound $c$ on a future third slot,
its stable ceiling is $\chi\le\min\{t_r,c\}$, not just $t_r$.
All physical cutoffs used below are at most this ceiling. For the
two-prefix $W$ option we take $c=a$; for an existing three-slot
branch its original slots persist below the smallest specified prime.

This comparison does not change $T,L,\alpha$ or $\Theta_Q$. The
comparison cutoff is smaller than the physical target cutoff, so
it remains below the smallest specified prime in every application
of the enhanced hereditary interface. The inner one-step construction
then appends only smaller primes, and its two signed non-switched
estimates have the same already checked square conditions. In the
classical branch the complete moduli still have support below $Q$.
In particular the constant value at $s=2$ need not be zero: a
comparison argument near $2.3$ is allowed even though the unmodified
one-step expression evaluated exactly at $s=2$ gives zero.

\paragraph{Taking the maximum with the replacement second seed.}
The direct finite replacement seed has defect $J_*=9/500$ at
$v_*=11/5$, uniformly in $\alpha\ge2$.
The extension to the present hereditary and classical families is
recorded separately in
\texttt{second-seed-stable-prefix-extension.tex}. Its condition is
that the comparison cutoff $R^{1/v_*}$ does not exceed the smallest
specified prime. On an enhanced branch every new prime in the seed
is at most this comparison cutoff; the old selected slots and their
gates persist. The good residual pairs have their explicitly checked
well support, and the bad part of term 5 uses the positive $K$ count.
The classical family can use the same conservative replacement,
with ordinary Bombieri--Vinogradov for the non-switched terms.

For the concrete domains below, whenever $s\le11/5$ the physical
cutoff is at most the smallest specified prime, so this comparison
condition holds automatically. Therefore we may use
\begin{equation}\label{eq:monotone-and-second-seed}
 j_*(s,\alpha)=\max\left\{j_{\rm mono}(s,\alpha),
                   \frac9{500}\mathbf1_{s\le11/5}\right\}
 \quad(s\le3),
\end{equation}
and zero for $s>3$. Maxima are taken on common positive boxes.
The two savings are alternatives for one count and are never added.

\paragraph{The lower transfer includes physical child arguments below two.}
For a two-prefix lower state with physical cutoff $N^a=R^{1/s}$,
use the following improvement only when $2\le s<4$:
\begin{equation}\label{eq:monotone-lower-transfer}
 h_*(s,\alpha)=\int_{s-1}^3
    j_*\left(w,\alpha+\frac{\alpha-1}{w}\right)\frac{dw}{w}.
\end{equation}
The earlier conservative formula used only $w\ge2$. There is no
need to discard $1\le w<2$ once the cutoff comparison has been
proved. The exact Buchstab identity starts at $R^{1/4}$; the child
prime is $q=R^{1/(w+1)}\le N^a$, its residual level is $R/q$, and
its physical cutoff is $q$. On $1\le w<2$, replace that physical
cutoff by $(R/q)^{1/v}$ with $v\ge2>w$. It is smaller than $q$.
The comparison's added primes are therefore below the new specified
slot $q$ and below the original two slots. The first new square test
is $q^2\le R$, and all subsequent tests are those of the comparison
construction at $v\ge2$. No residual sieve at an inadmissible small
parameter is being invoked.

The exact ratio is still
$\alpha_{Bq}=\alpha+(\alpha-1)/w\ge2$. The ordinary identity
\[
 a(4)-\int_{s-1}^3\frac{A(w)}w\,dw=a(s)
\]
includes $A(w)=1$ on $1\le w<2$, so the new positive lower defect
is precisely \eqref{eq:monotone-lower-transfer}. On the two-prefix
cap $V_{2,a}$, the original largest two slots persist because
$q\le N^a$. For the $W$ branch the third slot remains at most $a$;
the largest-prime gate and selected leading pair remain unchanged.
One can choose a fixed positive minimum prime exponent on every
compact subdomain. In the displayed parameter search $s_0\le5$,
the child residual exponent is at least half its parent's residual
exponent, hence at least $a$; the previous choice $\epsilon_1=a/6$
continues to suffice after fixed buffers.

\paragraph{Accounting in the actual functional.}
For $L_2,U_3,U_4$ and the four-prefix upper alternatives substitute
$j_*$ for $j$ in the formulas of the generalized fragment.
For $L_5,L_6$ substitute \eqref{eq:monotone-lower-transfer} for the
earlier truncated lower defect, and continue to choose zero extra
lower correction for $s<2$ or $s\ge4$.
In a four-prefix state the physical comparison cutoff is
$\min\{N^w,R^{1/2}\}$. Every monotone comparison and seed cutoff
is at most that value. Thus all four original slots, the selected
subproduct gates, and the top three slots of the $B(z)$ branch
persist. The upper kernel becomes $4(1-j_*)/L$ only inside the
already established minimum with the original switched kernel.
Its Chen gain still has the factor $1/4$; no separate switched gain
is added on the same four-prefix box.

There is also a small independent lower option for $L_2$. The audited
finite propagation of the replacement seed proves defect $7/4625$
on anchor exponents $29/200\le t\le3/20$, when the $L_2$ threshold
contains this interval. On that interval take the maximum of this
defect and $j_*((v(t)-t)/t,2)$. This is the only propagated-seed
contribution included by the probe.

First keep all margins and comparison choices fixed while
$N\to\infty$. Then remove the smoothing and sieve errors, refine
the common positive box partition, and approach admissible endpoints
from the interior. Finally exhaust the strict domains. The jump in
the chosen lower envelope at $s=11/5$ is harmless: it is a choice
between valid bounds, and its equality set has zero integration mass.
The former narrow certified correction and all global certificates
remain unchanged by this floating exploration.

\paragraph{Reproducible feasibility calculation.}
The probe \texttt{probe\_monotone\_one\_step.py} records separate
results for the monotone first seed alone and for its maximum with
the replacement seed, including the narrow propagated $L_2$ option.
It uses the same diagnostic baseline $1.8741432097004047$ and tuple
$(a,h,d)=(.08455645453148589,.15049959287189102,.2352266768449549)$.
The companion JSON records the numerical meshes and all termwise
gains. These are floating feasibility results; outward numerical
certification and any new global theorem constant remain separate.
After correcting the interpolation boundary at the analytically
exact value $\alpha=2$, the fine mesh gives
\[
 \Delta c_{\rm mono}=.0213357740\ldots,\qquad
 \Delta c_{\rm mono+seed}=.0219069552\ldots .
\]
The latter has combined diagnostic $1.8960501649\ldots$ and
zero-margin transfer diagnostic $b=.7569777160\ldots$, with remaining
gap $.0721508072\ldots$ to $K\log(1/.749)$. Its coarse-mesh extra
gain is $.0219011482\ldots$. These figures use $K=6.81$ and the
existing sampled baseline, so they do not include the separately
certified fourth-prefix improvement to $K$.

\subsection{Certified consequences of the fourth-prefix cap and second seed}
\label{sec:fourth-prefix-second-seed-corollary}

The following consequence uses only the old outward table, the new
compact fourth-prefix certificate, and two disjoint certified
one-prefix correction blocks. It includes none of the free-threshold
or monotone-envelope floating gains described elsewhere.

The previously certified extended-sieve integral satisfies
\[
 K_{\rm ext}\le
 \frac{1230679674488580033396846301}
      {180734410032387061618573312}.
\]
The integer inner-box certificate of
Proposition~\ref{prop:fourth-prefix-descendant-refinement} gives
\[
 \Delta_{\mathcal C}\ge
 \frac{10720642480494928338187}{604462909807314587353088}.
\]
Their difference is strictly below
\begin{equation}\label{eq:fourth-prefix-rounded-K}
 K_4=\frac{849}{125}=6.792.
\end{equation}
The strict gap is greater than $0.0004$, as verified by exact rational
comparison. Thus the original five labelled switched masses and the
bad-semiprime mass admit this new coefficient. Their labels, quotient
roughness restrictions and multiplicities are unchanged.

Keep the earlier threshold tuple
\[
 a=\frac{100}{1149},\qquad h=\frac{50}{309},\qquad
 d=\frac12-3a.
\]
The earlier outward table at $K_0=681/100$ gives $c_0=9089/5000$.
Its five switched upper entries, with their Chen multiplicities, sum to
\[
 2U_7+U_8+U_9+U_{10}+U_{11}^{\rm full}
 \le \frac{101203}{10000}.
\]
Each new upper entry can be taken to be its old outward entry times
$K_4/K_0$: the underlying unsifted mass is the same and only its
positive-count coefficient has changed. Recomputing the same table
therefore raises its lower coefficient by
\[
 \frac{101203}{40000}\left(1-\frac{K_4}{K_0}\right).
\]
This is a calculation with valid upper entries; it does not infer a
lower bound on any of the five actual masses.

The first double-sieve block has
$31/200<t<4/25$ and contributes more than $2/6635$.
The repaired second-seed block has $29/200<t<3/20$ and contributes
more than $56/465275$ by
\eqref{eq:replacement-second-seed-l2-gain}.
Both intervals lie below $h$, and they are disjoint.
Their proofs may continue to use the older positive cap $K_0$;
there is no need to replace its fixed-positive-cutoff extension by
an unproved stronger input. Hence the two lower gains add.
Altogether the coefficient is strictly greater than
\begin{equation}\label{eq:fourth-prefix-second-seed-c}
 \frac{9089}{5000}
 +\frac{101203}{40000}\left(1-\frac{K_4}{K_0}\right)
 +\frac2{6635}+\frac{56}{465275}
 =\frac{2046151322624013}{1121234583800000}
 >\frac{18249}{10000}.
\end{equation}

\begin{theorem}[Updated separate research-draft consequence]
\label{thm:fourth-prefix-second-seed-chen}
Under the published distribution and linear-sieve inputs and the
derived interfaces explicitly established in this manuscript,
\[
 D_{1,2}(N)\ge\left(\frac{18249}{10000}-o(1)\right)\HN.
\]
For every sufficiently large even $N$, there are at least $\HN/1000$
representations
\[
 N=p+rq,\qquad p,q\in\PP,\quad r=1\text{ or }r\in\PP,
 \qquad r\le q^{1529/2000}.
\]
No effective lower threshold for $N$ is asserted.
\end{theorem}

\begin{proof}
Equation~\eqref{eq:fourth-prefix-second-seed-c} proves the first
claim using the same complete eleven-term weight.
Apply the bad-semiprime removal estimate with $K_4$.
For $b=1529/2000$ put $u=(1-b)/(1+b)$.
The positive atanh series gives the strict upper bound
\[
 \log(1/b)<2\left(u+\frac{u^3}{3}
                         +\frac{u^5}{5(1-u^2)}\right).
\]
Exact rational arithmetic verifies
\[
 \frac{18249}{10000}-\frac{849}{125}\,
 2\left(u+\frac{u^3}{3}+\frac{u^5}{5(1-u^2)}\right)
 >\frac1{1000}.
\]
The strict surplus absorbs the asymptotic losses after the fixed
construction parameters are chosen. This proves the second claim.
\end{proof}

If the separately audited published Chen coefficient $1233/625$ is
accepted as an external input, the same new removal constant also
gives $b=74793/100000$ and at least $\HN/50000$ representations.
The earlier qualifications concerning that external proof audit
remain in force. This is a project deduction, not an imbalance
theorem located verbatim in that source. Neither consequence
certifies priority over every existing result.

The file \texttt{certify\_fourth\_prefix\_second\_seed\_transfer.py}
checks the rational subtractions, table rescaling and final margins.
The local seed also has a separate numerical implementation,
\texttt{replay\_replacement\_second\_seed.py}; the fourth-prefix
integer certificate was replayed exactly and its domain/kernel bounds
were reviewed separately. These are numerical and targeted analytic
checks, not formal verification of all cited mathematics.

% Targeted literature-review fragment, 2026-10-08.
% Companion record: frontier-primary-gap-review-2026-10-08.json.
% This fragment neither changes the manuscript nor certifies a global record.

\subsection{Targeted primary-source comparison and two remaining coverage gaps}
\label{sec:frontier-primary-gap-review-20261008}

The target is the assertion that every sufficiently large even integer
has a representation
\[
 N=p+rq,\qquad p,q\ \hbox{prime},\quad
 r=1\ \hbox{or prime},\qquad r\le q^b.
\]
The bounded check below found no new stronger theorem with these exact
quantifiers. It does not establish that none exists: the full texts of
two specifically identified older papers remain unavailable in this audit.

\paragraph{The current directly stated benchmark.}
The current \href{https://arxiv.org/abs/2606.05224}{arXiv record}
for Jiamin Li and Jianya Liu still lists version~2, submitted on
8 August 2026, with no version~3. Its
\href{https://arxiv.org/html/2606.05224v2}{complete HTML text}
was saved separately and checked at Proposition~$(1+a)$, (1.6),
Theorem~1.1, its following remark, and Theorem~1.3.
The numbered unconditional theorem states \(b=0.9\), with a lower
count \(0.0004\,C(N)N/\log^2N\).
The following remark asserts \(b=0.894\).
That number occurs only once in the complete text, and no corresponding
parameter computation was located. Thus \(0.9\) is the strongest
numbered same-scope theorem in the checked text, while \(0.894\) is
the strongest explicit author-asserted refinement there.
The conditional \(b=0.4\) requires \(\mathrm{WEH}(0.999)\).
The twin-prime \(b=0.75\) has a different quantifier.
These classifications check what the source states; they are not
independent validations of all its analytic proofs.

\paragraph{Cai's 2008 paper: metadata obtained, full text not obtained.}
The paper in question is Yingchun Cai,
\emph{On Chen's theorem (II)}, Journal of Number Theory
128(5) (2008), 1336--1357,
\href{https://doi.org/10.1016/j.jnt.2007.04.010}
{doi:10.1016/j.jnt.2007.04.010}.
The ordinary publisher page returned HTTP~403.
A standard unauthenticated request to the publisher's article API
returned HTTP~200 and an XML record of 1895 bytes. Inspection showed
only a \texttt{coredata} element: there was no article body.
The metadata marks the article open access/open archive and gives
an Elsevier user licence, but these flags do not constitute full-text
coverage. The API's \texttt{view=FULL} request returned HTTP~401,
requiring an API key. Both responses are saved in the source-audit
directory.

Exact-title, identifier, author/institutional and reference-chain
searches yielded no usable original full text in this turn.
Consequently the paper's full theorem list and internal factor
restrictions remain unchecked. A quantitative Chen abstract cannot
exclude an additional structural statement elsewhere in the paper.
This article must also be distinguished from Cai's
\emph{A remark on Chen's theorem (II)} in Chinese Annals B
29(6) (2008), and the distinct 2011 small-additive-prime paper.

\paragraph{Fouvry--Grupp 1989: the original coverage gap remains.}
The second paper is
\'Etienne Fouvry and F.~Grupp,
\emph{Weighted sieves and twin prime type equations},
Duke Mathematical Journal 58(3) (1989), 731--748,
\href{https://doi.org/10.1215/S0012-7094-89-05834-1}
{doi:10.1215/S0012-7094-89-05834-1}.
The official Project Euclid \texttt{.full} and \texttt{.pdf} endpoints
returned verification/iframe pages, rather than the article.
The author's institutional page and targeted institutional/HAL
searches did not supply a copy.

Li--Liu's Lemma~3.6 transcribes a mean-value estimate attributed to
Theorem~2 of this 1989 paper. Its error constant explicitly depends
on the residue \(a\), so this transcription alone does not give
uniformity for \(a=N\). It also does not inspect the original paper's
other theorems. The separate 1986 Fouvry--Grupp paper is a different
reference and must not be silently substituted.

\paragraph{Exact comparison rules.}
Put \(m=rq\). For \(0\le b<1\), the target condition has the exact
equivalent forms
\[
 r\le q^b
 \quad\Longleftrightarrow\quad
 r\le m^{b/(1+b)}
 \quad\Longleftrightarrow\quad
 m\le q^{1+b}.
\]
Thus a theorem stating \(r\le m^\delta\), \(0\le\delta<1/2\),
has factor-imbalance exponent \(b=\delta/(1-\delta)\).
By contrast, \(r\le N^\delta\) alone is insufficient because
\(m=N-p\) may be much smaller than \(N\).
A hypothesis \(m\ge\varepsilon N\), together with a fixed strict
exponent margin, permits the usual limiting conversion.

A roughness assertion \(P^-(rq)\ge N^\alpha\) is a lower bound on
the smaller factor. Balanced factors satisfy such bounds for suitable
\(\alpha<1/2\), while their logarithmic ratio tends to \(1\);
they need not satisfy any fixed \(b<1\).
Accordingly the separately reviewed Ratliff roughness claim is not
a stronger same-scope imbalance result.
Likewise, a bound on the additive prime \(p\), or on
\(\lvert p-N/2\rvert\), does not bound the ratio of the factors of
\(N-p\).

The relevant almost-prime definition is \(\Omega(m)\le2\), with
multiplicity. The weaker \(\omega(m)\le2\) allows unrestricted prime
powers and does not match the target. An almost-all-even-integers
result also leaves a possible infinite exceptional set.
Since the target allows \(r=1\), almost-all Goldbach already gives
almost-all representations for every \(b\ge0\); it does not settle
the all-sufficiently-large-even-integers assertion.

\paragraph{Published statements and project deductions.}
The existing project comparison obtains safe \(b=0.782\) or
\(b=0.749\) from an accepted published Chen counting coefficient
and separately justified removal estimates.
These are deductions of the project from specified inputs, not
imbalance theorems printed by the source author.
The project's external-input \(b=0.7485\) has the same distinction
and retains its recorded input-audit limitations.
The separate research-draft \(b=0.7659\) is numerically below the
checked Li--Liu remark exponent, but that comparison alone does
not establish a record over all prior methods.

This review therefore supplies a precise current direct-statement
benchmark and preserves the two original full-text gaps.
It makes no exhaustive literature, global priority, or completed-goal
claim. The companion JSON records URLs, acquisition results,
source hashes, and the limits of the check.


\paragraph{Later numerical update to this comparison.}
The subsequent certified consequence in
Theorem~\ref{thm:fourth-prefix-second-seed-chen} replaces the separate
draft exponent $0.7659$ by $0.7645$, and the external-input deduction
$0.7485$ by $0.74793$. The literature review's distinctions between
printed results, project deductions and unverified full-text coverage
remain unchanged.

% Local analytic supplement; numerical dependencies remain separately audited.
\subsection{A disjoint fourth-prefix supplement beyond the elementary range}
\label{sec:fourth-prefix-tail-refinement}

Use the construction and exact density identity of
Proposition~\ref{prop:fourth-prefix-descendant-refinement}.
Its density formula is valid on compact patches with $s_3>1$;
only the simple rational evaluation had required $s_4<3$.
We now extend the certified inner patches to $s_4<4$.
The first and third decisions, and the fourth-prefix choice of a
frozen descendant cap, are unchanged.

Write $A(s)=sF(s)/(2e^\gamma)$ and
$\overline F(s)=A(s)/s$. For $3\le s\le4$, the linear-sieve
delay equation gives
\[
 A'(s)=\frac{\log(s-2)}{s-1},\qquad
 -\overline F'(s)=
 \frac{A(s)-\dfrac{s}{s-1}\log(s-2)}{s^2}.
\]
The factor $s/(s-1)$ in this expression is essential.
The earlier unintegrated candidate certificate omitted this factor;
its claimed additional gain $0.005860760569$ and cap $6.786$
are withdrawn. The corrected certificate below does not use them.

Let $S<4$ be an upper bound for $s_4$ throughout a grid cell.
For $S\le3$ set $c(S)=1$; for $3<S<4$ let
\[
 z=\frac{S-3}{S-1},\quad
 U(S-2)=2\left(z+\frac{z^3}{3(1-z^2)}\right),\quad
 c(S)=1-\frac{S}{S-1}U(S-2).
\]
The positive atanh expansion shows $\log(S-2)\le U(S-2)$.
The function $s\log(s-2)/(s-1)$ is increasing on $[3,4]$:
its derivative is
\[
 \frac{s}{(s-1)(s-2)}-\frac{\log(s-2)}{(s-1)^2}>0.
\]
The corresponding rational upper approximation is also increasing:
each term of $sU(s-2)/(s-1)$, after putting
$z=(s-3)/(s-1)$, is a positive multiple of
$(3-z)z^{2k+1}$, whose derivative is positive for
$0<z\le1/3$ and nonnegative at zero. Consequently
\[
 c(S)>1-\frac43\frac{25}{36}=\frac2{27}>0.
\]
Alternatively that last positivity can be checked directly using
the rational numerator recorded for every selected cell.
Since $A(s)\ge1$, and $-\overline F'(s)=1/s^2$ below three,
integration over $s_3<s_4\le S$ gives
\begin{equation}\label{eq:fourth-prefix-tail-kernel}
 \overline F(s_3)-\overline F(s_4)
 \ge c(S)\left(\frac1{s_3}-\frac1{s_4}\right).
\end{equation}
This gives a valid lower multiple of the old rational gain kernel.

The new certificate uses the same mesh $1/800$ as the previous
fourth-prefix certificate, extending its lower fourth-coordinate
cutoff from $1/25$ to $1/50$. Before choosing a branch, it excludes
every cell selected by any of the three old branches. There are
exactly $483089$ such excluded cells. This full-cell exclusion is
necessary: excluding only points with the new argument below three
would not prevent overlap with a different old branch.

For each remaining cell, use the old affine cap bounds and the
same three selected subsets. In the notation of the old certificate,
put $I=i_x+i_y+i_z+i_w$, $b_3=l-32I$, $b_4=m_+-32I$,
and $S=b_4/(32i_w)<4$. If $c(S)=c_n/c_d$, its integral
contribution is bounded below by
\[
 \frac{4D(m_--l)c_n}
 {(i_x+1)(i_y+1)(i_z+1)(i_w+1)b_3b_4c_d},
 \qquad D=25600.
\]
All gates, endpoints and logarithm bounds are rational. Each
contribution is rounded down to the $2^{-80}$ lattice.
There are $295979$ newly selected cells, giving
\begin{equation}\label{eq:fourth-prefix-tail-gain}
 \Delta_{\rm tail}\ge
 \frac{5155254219572575727579}{1208925819614629174706176}
 >0.004264326343.
\end{equation}
The interiors of all cells are disjoint. To use compact patches,
first trim the finitely many cell faces, then let the trimming width
tend to zero; the positive final rounding slack absorbs this limit.

Adding only the previous disjoint certified gain and this supplement,
and subtracting from the previously certified extended endpoint,
gives
\[
 K_{\rm new}\le
 \frac{2453406983762171899504910655}
 {361468820064774123237146624}
 <\frac{1697}{250}=6.788.
\]
Here the displayed endpoint is an upper bound for the limiting
density coefficient; the rounded strict cap is used for the
finite buffered construction. No floating estimate is included.
The independent replay reconstructs the old winner stream and gain,
and the corrected new stream, without importing the author module.
It verifies the same full-cell disjointness and rational sum.
These checks establish this local functional bound, conditional on
the inherited analytic interfaces; they do not verify the complete
number-theoretic proof or establish priority.

% Derived research fragment, 2026-10-08. No main/state changes.
\subsection{Legal fifth- and sixth-prefix refinements of the product sieve}
\label{sec:higher-prefix-K-refinement}

Keep the first and third decisions of the extended upper tree and
the already constructed fourth-prefix tree. This note changes a
fifth-prefix decision at its own odd rank, or freezes a sixth-prefix
cap for use only from rank seven onward. Neither construction allows
a rank-three decision to depend on a later prime.

Write $\rho=19/96$, $\mu=25/128$, $\Lambda=153/256$, and
$v(t)=\min\{(1+t)/2,57/64-3t/2\}$. For ordered exponents
$t_1>\cdots>t_r>0$, write $S_j=\sum_{i\le j}t_i$ and
\[
 W_r=\max_{j\le r}(S_{j-1}+2t_j),\qquad
 s_I=\sum_{i\in I}t_i,\qquad
 G_I=\max_{j\notin I}
 \left(s_I+\sum_{\substack{k<j\\k\notin I}}t_k+2t_j\right).
\]
The maximum over an empty set is zero. Every new selected-subproduct
branch used below must satisfy the entire gate
\begin{equation}\label{eq:higher-prefix-complete-gate}
 \max\{1/2,W_r,G_I\}<M_r=v(s_I)\le\Lambda.
\end{equation}
The actual frozen level is strictly between gate and endpoint.
All selected subsets are permitted, but no remaining-square test
is discarded. On later slots the remaining-square inequality is
exactly the complete residual square inequality, since all selected
and unselected specified slots have already occurred.

Let $L_3$ be the exact cap of the existing extended tree, and let
$L_4(t_1,t_2,t_3,t_4)\ge L_3$ be its chosen fourth-prefix cap.
It can be the already certified compact-patch construction.
For diagnostics one can instead use the maximum of $L_3$ and all
eligible four-slot endpoints. These are different numerical bases.
In every case only old retained third prefixes
\begin{equation}\label{eq:higher-prefix-old-third}
 S_2+3t_3<L_3
\end{equation}
are changed. Fourth-prefix membership is unchanged.

\paragraph{The odd fifth-prefix construction.}
On a finite union of compact positive fifth-prefix boxes select
one branch $M_5>L_4$ satisfying
\eqref{eq:higher-prefix-complete-gate}. At rank five use the new
test
\[
 S_4+3t_5<M_5.
\]
If it passes, freeze this cap and use $M_5$ in all subsequent odd
cube tests $7,9,\ldots$. Outside active boxes retain the old rule.
The old rank-five test is replaced, not retained as an additional
condition. An old rank-five exit can therefore be restored.
Rank-five information is available at this odd decision itself.
Every first exclusion is still odd, and the support is prefix
closed; the exact first-exit identity therefore gives an upper
weight. This does not use coefficientwise monotonicity of Mobius
sums under support enlargement.

Every admitted fifth prefix has complete well support at $M_5$:
the first four squares follow from the old third cube, and its own
cube implies the fifth square. Later odd cubes imply their own
squares, and ordering gives each subsequent even square. Together
with all tests $G_I<M_5$, this licenses the selected-subproduct
factorization for the complete modulus atom. The small factor,
signed boxing, factorial normalization, moving residue, and local
scale changes use the existing fixed-margin argument; the cap is
still below $3/5$.

There is an exact density identity. For $P=p_1\cdots p_5$ let
$\epsilon_i(P)$ indicate that its own cube passes at cap
$\lambda_i$, where $(\lambda_0,\lambda_1)$ are the old and new
frozen caps. Then
\begin{equation}\label{eq:rank5-exact-density}
 V_{\rm old}-V_{\rm new}
 =\sum_{\text{active }P}g_N(P)
 \left\{\epsilon_1(P)V_N^-(N^{\lambda_1}/P,p_5)
       -\epsilon_0(P)V_N^-(N^{\lambda_0}/P,p_5)\right\}.
\end{equation}
The sign is important: a full rank-five subtree contributes
$-g_N(P)V_N^-$, since its following full odd cubes are residual
even cubes. The active fifth-prefix trees are disjoint.

Put $\overline f=f/(2e^\gamma)$, extended by zero on $s\le2$.
For a retained root the ordinary lower-density equality applies
at $s>2$; for a rejected root the displayed indicator is zero.
Thus no lower-density assertion at an inadmissible parameter is
needed. Applying the already established uniform lower-density
equality on compact sets and then prime-number integration yields
\begin{equation}\label{eq:rank5-density-gain}
 \Delta_5=4\int_{\mathcal C_5}
 \left\{\overline f\left(\frac{M_5-S_5}{t_5}\right)
       -\overline f\left(\frac{L_4-S_5}{t_5}\right)\right\}
 \frac{d\mu_5}{t_5},\qquad
 d\mu_r=\prod_{i=1}^r\frac{dt_i}{t_i}.
\end{equation}
The surfaces $s=2$ are trimmed before taking limits. For
$2<s\le4$ the kernel is explicit:
$\overline f(s)=\log(s-1)/s$.

\paragraph{The even sixth-prefix construction.}
Start either after the old fourth tree or after the preceding
fifth refinement, and denote the inherited frozen cap at rank
five by $L_5$. Only use prefixes whose fifth cube has passed:
$S_4+3t_5<L_5$. Select an eligible $M_6>L_5$ on compact
sixth-prefix boxes. Preserve all membership through rank six,
and replace later odd cubes $7,9,\ldots$ by this frozen cap.
There is no new even-rank exclusion. In particular
\[
 S_5+2t_6\le S_4+3t_5<L_5,
\]
so all six initial squares are already licensed. The new gate
\eqref{eq:higher-prefix-complete-gate} supplies factorization of
every complete descendant, exactly as above.

The even subtree now has the opposite sign and ordinary upper
parity. Its exact difference is
\begin{equation}\label{eq:rank6-exact-density}
 V_{\rm old}-V_{\rm new}
 =\sum_{\text{active }P_6}g_N(P_6)
 \{V_N^+(N^{\lambda_0}/P_6,p_6)
             -V_N^+(N^{\lambda_1}/P_6,p_6)\}.
\end{equation}
The old fifth cube gives $(L_5-S_6)/t_6>1$. The uniform ordinary
upper-density equality proved in the fourth-prefix fragment gives
\begin{equation}\label{eq:rank6-density-gain}
 \Delta_6=4\int_{\mathcal C_6}
 \left\{\overline F\left(\frac{L_5-S_6}{t_6}\right)
       -\overline F\left(\frac{M_6-S_6}{t_6}\right)\right\}
       \frac{d\mu_6}{t_6},\qquad \overline F=F/(2e^\gamma).
\end{equation}
On $(M_6-S_6)/t_6<3$ it is again a rational kernel. A fifth gain
and a sixth gain computed in this order telescope between three
explicit trees. They are not independent savings against the same
old tree, and a sixth gain using a different base cannot be added.

\paragraph{Useful pruning and exact inner-box bounds.}
If the inherited $(r-1)$-slot cap is the maximum of all eligible
selected endpoints, an improving $r$-slot subset must contain the
new slot $r$. Otherwise deleting slot $r$ preserves all old gate
tests, and that endpoint was already available. Thus only $16$
fifth-slot and $32$ sixth-slot subsets need be tested in that
version of the construction. The probe retains all subsets as
an additional check of this implication.

On $2<s<4$ an entirely rational lower derivative is available:
\[
 \overline f'(s)
 =\frac{s/(s-1)-\log(s-1)}{s^2}
 \ge\frac{4-s}{2s(s-1)}>0.
\]
Here $\log x\le(x^2-1)/(2x)$ for $x\ge1$, as follows either by
differentiation or the positive atanh series. The final rational
bound decreases on $[2,4]$.
For a box with upper coordinate endpoints $u_i$, cap bounds
$M_-\le M_5\le M_+$, old-cap upper bound $L_+$, and
$\sigma_+=(M_+-S_-)/\ell_5<4$, define
\[
 d_-=M_--\max\{L_+,S_++2u_5\}>0.
\]
Then its contribution to \eqref{eq:rank5-density-gain} is at least
\begin{equation}\label{eq:rank5-rational-box}
 \operatorname{vol}(\mathcal B)\,
 \frac{4d_-}{(\prod_i u_i)u_5^2}
 \frac{4-\sigma_+}{2\sigma_+(\sigma_+-1)}.
\end{equation}
All ordering, old-third, full-square, and selected gates must hold
with margins on the whole box. This bound also covers restored
old rank-five exits, because the old lower multiplier is zero
below two. An upper envelope over all potentially eligible old
four-slot branches gives a safe $L_+$ even when the already
certified fourth tree uses only some of those branches. Thus a
new fifth-box certificate can be added to that existing compact
fourth-tree certificate without claiming its entire floating gain.

Here is the combination argument without any support-monotonicity
assumption. Subdivide each new five-dimensional box by the old
four-prefix labels. On each resulting cell the existing fourth
tree has some frozen cap $\lambda_4$, with $\lambda_4\le L_+$.
Its exact odd-fifth difference is still
\eqref{eq:rank5-exact-density}. Both its old and new normalized
lower arguments are at most $\sigma_+<4$. The clipped function
$\overline f$ is nondecreasing on $(-\infty,4]$, directly by its
zero value below two and the derivative just proved. Therefore
replacing $\lambda_4$ by the larger $L_+$ can only decrease the
lower bound for the new-minus-old lower density. Bound
\eqref{eq:rank5-rational-box} remains valid on every subcell, and
their volumes sum to the original box volume. The earlier
four-prefix compact patches and any disjoint four-prefix tail
patches may consequently be retained exactly as constructed.

\paragraph{A replayable coarse fifth-prefix certificate.}
The script \texttt{certify\_fifth\_prefix\_K\_refinement.py} uses
ordered grid interiors of width $1/400$, with
$1/50<t_5<\cdots<t_1<39/200<\mu$. It tests all old four-slot
branches when bounding $L_+$, including branches that are only
potentially eligible somewhere on the box. For a new fifth branch
it requires the complete upper-endpoint gates, the old third
cube, and $2<\sigma_+<4$ throughout the box.

All these comparisons use integers with common affine denominator
$D=12800$. For lower grid indices $(i_1,\ldots,i_5)$, put
$A=D(M_+-S_-)$ and $B=32i_5$, so $\sigma_+=A/B$.
If $d=Dd_-$, the rational lower bound for the entire box is
\[
 \frac{400\,d\,B(4B-A)}
 {16A(A-B)\{\prod_{j=1}^5(i_j+1)\}(i_5+1)^2}.
\]
The grid volume has cancelled the five factors of $400$ in the
prefix measure. Floating scores choose one valid branch only;
the selected branch's contribution is recomputed with arbitrary
precision integers and rounded down to a $2^{80}$ lattice.
No floating comparison licenses geometry or bounds a contribution.

Among $12103014$ candidate boxes, $149251$ pass. The exact result is
\begin{equation}\label{eq:rank5-coarse-certificate}
 \Delta_5\ge
 \frac{2167182368869852589729}{1208925819614629174706176}
 >\frac{179}{100000}=.00179.
\end{equation}
The decimal $.0017926512393958863$ is informational. This is a
conservative part of the much larger floating fifth-prefix gain;
it does not certify that whole diagnostic. The companion JSON
contains the selected-box digest, representative boxes, and the
minimum positive slacks. The gain can be subtracted from the
previously certified fourth-prefix-plus-tail coefficient by the
combination argument above; final transfer arithmetic is separate.

\paragraph{The unused fourth-prefix tail.}
The previous exact certificate selected $s_4<3$. Its analytic
upper-density formula also covers $s_4\ge3$; this tail is not a
new tree rule. For $3\le s\le4$,
\[
 -\overline F'(s)=\frac{A(s)-\frac{s}{s-1}a(s-1)}{s^2}
 \ge\frac{1-\frac{s}{s-1}\log(s-2)}{s^2},
 \qquad A(s)=s\overline F(s),\quad a(s-1)=\log(s-2).
\]
The factor $s/(s-1)$ comes from
$A'(s)=\overline f(s-1)=a(s-1)/(s-1)$ and cannot be omitted.
An upper rational bound for the logarithm, and an upper bound
for $s/(s-1)$ on the box, give a positive box lower bound when
the resulting right side is positive.
Restricting to old $s_3\ge3$ and new $s_4\le4$ makes these boxes
disjoint from every old certified box with new $s_4<3$.
If instead the complete fourth-prefix integral is certified afresh,
it replaces the old partial subtraction; the two are not added.

\paragraph{Floating feasibility and scope.}
The saved probe samples ordered logarithmic coordinates. With
cutoff $a_0$ its measure is
$\log(\rho/a_0)^r/r!$ times a Sobol average, which is exactly
the change of variables for $d\mu_r$. It uses the complete
eligible four-slot cap, then the complete fifth cap, then the
complete sixth cap on retained fifth prefixes. The current
$2^{18}$ samples at $a_0=.01,.02$ suggest
\[
 \Delta_4\simeq .041\text{--}.042,
 \qquad\Delta_5\simeq .022,
 \qquad\Delta_6\simeq .010\text{--}.011.
\]
The fourth total includes about $.013$--$.014$ from $s_4\ge3$;
about $.020$ of the fifth gain has new argument below four.
The combined diagnostic subtraction is about $.074$, suggesting
$K$ near $6.735$ from the old $6.809326869$ endpoint. This is
a feasibility observation only. The sample spread, unresolved
small-cutoff tail, sieve-function ODE, and interpolation have no
outward error bounds. No new theorem constant is asserted here.
The companion JSON separates exact tree algebra, any subsequent
inner-box certificate, and these floating calculations.

% Integration fragment, 2026-10-08.  Old certificates remain unchanged.
\subsection{A certified second seed of size one fiftieth}
\label{sec:second-seed-omega-certified}

The ordinary-only repaired second seed can be increased from
$9/500$ to $1/50$ by using a sharper uniform Buchstab envelope.
This assertion uses no numerical $H,h$ feedback and keeps the
conservative coefficient $K=681/100$.

\begin{proposition}[Certified uniform ordinary second seed]
\label{prop:second-seed-omega-certified}
On the retained-anchor family of
Proposition~\ref{prop:replacement-second-seed}, or the stable
specified-prefix family of
Proposition~\ref{prop:second-seed-stable-prefix} with its comparison
cutoff and branch gates, one has
\[
 \Phi_{11/5}\le\left(1-\frac1{50}+o(1)\right)\Theta_Q.
\]
The result is uniform over all admitted states with $\alpha\ge2$.
It remains valid if a smaller positive switching coefficient is
subsequently justified by a separate argument.
\end{proposition}

\begin{proof}
Keep exactly
\[
 s=11/5,\quad s_0=227/50,\quad
 (\phi_1,\phi_2,\phi_3)=(353/100,29/10,61/25),
 \qquad c_\omega=70893/125000.
\]
The bound
\[
 \omega(x)\le\max\{1/x,c_\omega\}\qquad(x\ge1)
\]
has a short exact proof.  On $[1,2]$, use $\omega(x)=1/x$.
For $2\le x\le3$, put $v=x-1$.  The maximum of
$1+\log v-c_\omega(v+1)$ is
$-\log c_\omega-c_\omega<0$; the strict sign follows from
\[
 c_\omega\sum_{k=0}^{8}\frac{c_\omega^k}{k!}-1
 =\frac{65111700737847638377192824351410078164698277}
 {33378601074218750000000000000000000000000000000000}>0.
\]
Thus $\omega\le c_\omega$ on $[2,3]$.  The differential equation
$(x\omega(x))'=\omega(x-1)$ then propagates that same bound over
each successive unit interval.

For the unchanged switched domains $9,\ldots,15$, the exact
argument lower bounds, valid for every $\alpha\ge2$, are
\[
\begin{array}{c|rrrrrrr}
i&9&10&11&12&13&14&15\\\hline
m_i&47/25&273/110&1478/725&3801/1100&103299/38830&3136/1765&15388/7975.
\end{array}
\]
Each satisfies $m_i c_\omega>1$.  Therefore the factor $3/5$ in
each of their previously proved explicit logarithmic mass bounds
may be replaced by $c_\omega$.  This uses a uniform upper envelope,
not monotonicity of the Buchstab function.
On the bad part of $\Gamma_5$ the argument is at least
$4313/725>2$, so its positive replacement kernel is now
$Kc_\omega s_0/4$.
All non-switched ordinary kernels are unchanged.  In particular
the bad pair is still switched, and no missing complete-well
condition has been waived.
For $\Gamma_{16},\ldots,\Gamma_{21}$ keep the former conservative
$\omega\le1$ bounds; both uniform unit-surface summands for
$\Gamma_{20},\Gamma_{21}$ are retained in full.

Let $J_2^{\omega}$ be the saved repaired functional with precisely
those replacements and $H=h=0$.  The new separate certificate
\texttt{certify\_second\_seed\_omega\_envelope.py} uses the audited
80-bit outward integer interval implementation, 32 positive atanh
terms with explicit logarithm remainders, and 8192 rational
Darboux cells for each one-dimensional integral.  The double
ordinary kernels are integrated in one variable exactly, as in
the previous certificate.  No floating quadrature or floating
comparison contributes to the bound.
The output gives
\begin{equation}\label{eq:second-seed-omega-exact-bound}
 J_2^{\omega}\ge
 \frac{24479824718463837928127}
      {1208925819614629174706176}
 >\frac1{50}.
\end{equation}
The displayed decimal lower endpoint, only for orientation, is
$0.0202492364058096657\ldots$.
The strict surplus permits the existing fixed box and asymptotic
errors to be removed in their stated order.  This proves the
proposition under the same named analytic inputs as the original
seed.  Decreasing a valid $K$ decreases only positive switched
costs, so preserves the conclusion.
\end{proof}

The cutoff comparisons and finite Buchstab transfers already proved
for the stable family may now use $1/50$ instead of $9/500$.
For example the direct upper option is $H(w)\ge1/50$ for
$1\le w\le11/5$ whenever its seed comparison cutoff is admitted.
Its one-step lower transfer becomes
\[
 h(s)\ge\frac1{50}\log\frac{11/5}{s-1}
 \quad(2\le s<16/5),
\]
with the $[1,2)$ child arguments still retained.  The same two-step
comparison at argument three gives
\[
 H(3)\ge\frac1{50}\log(5/4)\log(22/15)
 >\frac{14}{8325},
\]
on the same stable comparison domain.  These are rescalings of
previously proved finite transfers, not numerical feedback inputs
used to establish \eqref{eq:second-seed-omega-exact-bound}.
All competing defects must still be combined by a common-box
maximum, and four-prefix alternatives remain inside their existing
minimum.  No integrated global Chen coefficient, imbalance
exponent, or completion of a quasi-Goldbach goal is asserted here.

% Research integration fragment, 2026-10-08.  Old sources are unchanged.
\subsection{Finite same-level feedback in the repaired second seed}
\label{sec:second-seed-finite-feedback}

This gives the full positive feedback operator for the non-switched
terms in Wu 2004, Lemma 4.2, after the repair of $\Gamma_5$.
It uses already proved child estimates at their actual states.
The numerical probe below is only a floating feasibility calculation.

\paragraph{States and parameter restrictions.}
Use the stable enhanced or classical family of
Proposition~\ref{prop:second-seed-stable-prefix}.  At a parent state
$\mathfrak s$ write
\[
 B=N^{T+o(1)},\quad Q=N^\nu,\quad R=Q/B=N^{L+o(1)},\quad
 L=\nu-T>0,\quad \alpha=(1-T)/L\ge2.
\]
In the enhanced case retain full specified well support and one
proved factorization branch, with its stable ceiling $\chi$ and
every old gate.  In particular any third-slot bound $c$ used before
three actual specified slots exist requires $\chi\le c$.
The classical case uses the absolute remainder below level $1/2$.
All statements are averaged on finite positive marked boxes.

Choose parameters
\[
 2<s<\phi_3<\phi_2<\phi_1<s_0\le5,
 \quad a_0=1/s_0,\quad a_i=1/\phi_i,\quad b=1/s,
\]
with
\begin{equation}\label{eq:feedback-parameter-gates}
 a_3+2a_1<1,\qquad 3a_1<1,\qquad a_2+2a_1<1.
\end{equation}
The practical table below uses $2.1\le s\le2.5$.
Require the parent's comparison cutoff $L/s<\chi$ on an enhanced
box; all its new specified primes are then below the old prefix.
The state at a child always retains the same total level $Q$.

Let $H_{\mathfrak t}(w)$ and $h_{\mathfrak t}(w)$ mean previously
proved defects, at a state $\mathfrak t$, in the bounds
\[
 \Phi_{\mathfrak t,w}\le
 (A(w)-H_{\mathfrak t}(w)+o(1))\Theta_{\mathfrak t},\qquad
 \Phi_{\mathfrak t,w}\ge
 (a(w)+h_{\mathfrak t}(w)+o(1))\Theta_{\mathfrak t}.
\]
Here $A(w)=wF(w)/(2e^\gamma)$, $a(w)=wf(w)/(2e^\gamma)$,
and unproved defects are zero.  On $1\le w<2$ a positive upper
defect must come from the already justified cutoff comparison;
it is not an invocation of an ordinary upper-sieve theorem at a
forbidden parameter.  The normalized comparison multiplier remains
$A(w)=1$ on $1\le w\le3$.  We use $h=0$ below two.

\paragraph{The exact child states and cutoffs.}
The two terms in $\Gamma_1$ keep the parent prefix and ratio
$\alpha$, with respective arguments $s_0,\phi_1$.
For each of $\Gamma_2,\Gamma_3,\Gamma_4$, append a prime
$q=N^{Lu}$, where respectively
\[
 a_0<u<b,\qquad a_0<u<a_2,\qquad a_0<u<a_3.
\]
Call its state $\mathfrak s_u$.  Its data are
\begin{equation}\label{eq:feedback-single-child}
 L_u=L(1-u),\qquad
 \alpha_u=\frac{\alpha-u}{1-u},\qquad
 \rho_u=s_0(1-u),\qquad
 \log_N(\hbox{physical cutoff})=La_0.
\end{equation}
The singleton square condition is $2u<1$, and the physical cutoff
is at most the new smallest specified prime $N^{Lu}$.

For the pair terms use $q_1=N^{Lu}<q_2=N^{Lv}$ and put
\[
 D_5=\{a_0<u<v<a_2\},\quad
 D_6=(a_0,a_1)\times(a_2,a_3),\quad
 D_7=\{a_0<u<v<a_1\},\quad
 D_8=(a_0,a_1)\times(a_1,a_2).
\]
The child state $\mathfrak s_{u,v}$ has
\begin{equation}\label{eq:feedback-pair-child}
 L_{u,v}=L(1-u-v),\qquad
 \alpha_{u,v}=\frac{\alpha-u-v}{1-u-v}.
\end{equation}
Its actual physical cutoff and argument are
\[
\begin{array}{c|c|c}
 i&\log_N(\hbox{physical cutoff})&\rho_i(u,v)\\\hline
 5,6&La_0&s_0(1-u-v)\\
 7,8&Lu&(1-u-v)/u.
\end{array}
\]
In particular $\Gamma_7,\Gamma_8$ do not have cutoff $R^{a_0}$.
Changing that cutoff would give the wrong child argument and kernel.

The residual pair is well supported exactly under
\begin{equation}\label{eq:feedback-pair-squares}
 2v<1,\qquad v+2u<1.
\end{equation}
Both hold everywhere in $D_6,D_7,D_8$ by
\eqref{eq:feedback-parameter-gates}.  On $D_5$ retain only
$D_5^+=D_5\cap\{v+2u<1\}$ for ordinary feedback; write
$D_5^-=D_5\setminus D_5^+$.  The first test $2v<1$ holds throughout.
There is no replacement of these asymmetric pair tests by a
condition such as three times one exponent being below the level.

On every retained pair domain, $1<\rho_i\le s_0-2\le3$.
Indeed $1-u-v>u\ge a_0$ proves the lower endpoint; maximizing
the displayed arguments at the lower endpoints gives the upper
endpoint.  For $\rho_i<2$, compare the positive child count at
any previously licensed larger argument, hence a smaller cutoff.
This keeps every compared prime below its actual cutoff and below
the child's smallest specified prime.  The argument $\rho_i$ in
the operator is the actual argument, with that proved comparison
already included in $H_{\mathfrak s_{u,v}}$.

The complete prefix tests follow from the merge lemma.  More
explicitly, all old slots precede the one or two new slots, so
old complete and selected-remaining tests are unchanged.  The
new tests are precisely $2u<1$ or
\eqref{eq:feedback-pair-squares}, multiplied by $B$.
Ordinary residual atoms are merged at the indicated child level.
Original selected slots and the relevant first three $W$ slots,
or their stable upper bound, are retained.  All comparisons lower
cutoffs, and no new slot can cross back above an old slot.
Finally
\[
 \alpha_u=1+\frac{\alpha-1}{1-u}\ge\alpha,\qquad
 \alpha_{u,v}=1+\frac{\alpha-1}{1-u-v}\ge\alpha.
\]
Thus ratios remain at least two, but are not reset to two.
Signed atom normalizations, repeated-label factorials, coprime
means and moving-residue margins are those of the existing stable
interface.  If a more specialized child defect has further gates,
it is used only where those gates also hold; elsewhere it is zero.

\paragraph{A uniform rough-number envelope with an exact check.}
Put $c_\omega=70893/125000=0.567144$.  Then
\begin{equation}\label{eq:feedback-omega-envelope}
 \omega(x)\le\max\{1/x,c_\omega\}\qquad(x\ge1).
\end{equation}
For $1\le x\le2$ this follows from $\omega(x)=1/x$.
For $2\le x\le3$, write $v=x-1$ and maximize
$1+\log v-c_\omega(v+1)$ at $v=1/c_\omega$.
The maximum is $-\log c_\omega-c_\omega<0$, because the exact
positive rational calculation
\[
 c_\omega\sum_{k=0}^{8}\frac{c_\omega^k}{k!}-1
 =\frac{65111700737847638377192824351410078164698277}
 {33378601074218750000000000000000000000000000000000}>0
\]
implies $c_\omega e^{c_\omega}>1$.
For $x\ge3$, integrate $(x\omega(x))'=\omega(x-1)$ and induct
over unit intervals, starting from $\omega(3)<c_\omega$.
This proves the envelope without a floating maximization.

On $D_5^-$ the rough argument is at least
$s_0(2-2a_2)>2$, so the retained positive replacement kernel may
be sharpened from $3Ks_0/20$ to $Kc_\omega s_0/4$.
There is still no child feedback on this domain.
For the switched triple domains $9,\ldots,15$, the following
numbers bound their argument from below for all $\alpha\ge2$:
\[
\begin{array}{c|l}
9&2\phi_3-3\\
10&2\phi_2-2-\phi_2/s\\
11&2\phi_3-\phi_3/\phi_2-2\\
12&2\phi_1-2-\phi_1/s\\
13&2\phi_2-\phi_2/\phi_1-1-\phi_2/s\\
14&2s-s/\phi_1-2\\
15&2\phi_3-\phi_3/\phi_2-1-\phi_3/s.
\end{array}
\]
If the entry is $m_i$, replace the former factor $3/5$ in its
explicit logarithmic mass by
$C_i=\max\{c_\omega,1/\max(1,m_i)\}$.
This also covers an enlarged parameter set where $m_i<5/3$;
nonunit terms only exist at argument at least one.
Use the previous coarse $\omega\le1$ bounds for $16,\ldots,21$,
including the separate uniform unit terms in $M_{20},M_{21}$.
There are no unit terms for $9,\ldots,19$, since $s>2$ gives
every specified sum at most $4/s<2\le\alpha$.

\paragraph{The positive feedback operator.}
Let $J_2^0(s,s_0,\phi_1,\phi_2,\phi_3)$ denote the ordinary-only
repaired functional, with the preceding uniform switched envelopes.
It is the functional of the saved seed with $A(s)$ replacing one
if needed.  Its ordinary pair kernels are $1/(1-u-v)$ on
$D_5^+,D_6,D_7,D_8$, and its bad-$\Gamma_5$ kernel is
$Kc_\omega s_0/4$.  No switched term contains an imported defect.

Define the same-level positive linear operator
\begin{align}
 \mathcal P_{2,\mathfrak s}[H,h]
 =\frac15\bigg\{&4H_{\mathfrak s}(s_0)+H_{\mathfrak s}(\phi_1)
 +\sum_{c\in\{b,a_2,a_3\}}\int_{a_0}^{c}
 \frac{h_{\mathfrak s_u}(s_0(1-u))}{u(1-u)}\,du\notag\\
 &+\iint_{D_5^+}
 \frac{H_{\mathfrak s_{u,v}}(s_0(1-u-v))}
      {uv(1-u-v)}\,du\,dv\notag\\
 &+\sum_{i=6}^8\iint_{D_i}
 \frac{H_{\mathfrak s_{u,v}}(\rho_i(u,v))}
      {uv(1-u-v)}\,du\,dv\bigg\}.
 \label{eq:second-seed-feedback-operator}
\end{align}
The valid new upper defect at the parent is
\begin{equation}\label{eq:second-seed-feedback-bound}
 H_{\mathfrak s}(s)\ge
 \max\{0,J_2^0+\mathcal P_{2,\mathfrak s}[H,h]\}.
\end{equation}
Here the defects on the right are previously proved inputs, so the
formula describes a finite additional proof step and is not circular.
To prove it, insert each child bound with its actual residual level
into the unchanged 21-term pointwise inequality.  An improvement
to $\Gamma_1$ subtracts $4H(s_0)+H(\phi_1)$ from its upper bound;
an improvement to a negative term adds its lower defect; and an
ordinary pair upper improvement subtracts its upper defect.
All three effects improve the final upper bound.  The change from
child normalization to parent normalization contributes
$1/(1-u)$ or $1/(1-u-v)$, exactly as displayed.  Division by five
then proves the formula.  Every coefficient of the operator is
nonnegative.

The implementation shared with the first-seed iteration uses the
special case of defects $H(w),h(w)$ uniform over all these stable
and classical states and every $\alpha\ge2$.  This removes the
state subscript in \eqref{eq:second-seed-feedback-operator}; it does
not change the actual ratios in
\eqref{eq:feedback-single-child}--\eqref{eq:feedback-pair-child}.
For a finite parameter bank, take the maximum of its candidate
bounds and already proved alternatives on a common positive
partition.  Different upper candidates are not added.

\paragraph{Finite parameter probe and its limits.}
The new module \texttt{probe\_second\_seed\_feedback.py} exports
\texttt{SecondSeedOperator(parameters,K,n).candidate(H,h)} and
\texttt{.value(H,h)} for numpy-compatible uniform input functions.
It treats the bad triangle by its exact two-piece boundary, with
no masked-grid replacement for that boundary.  It checks the
three gates in \eqref{eq:feedback-parameter-gates} before admitting
a parameter choice.

The initial upper inputs are the proved values $H=9/500$ on
$[1,11/5]$, $H=7/4625$ on $(11/5,3]$, and zero elsewhere.
The lower input is one Buchstab integral of their pointwise maximum,
with the lower endpoint $s-1$, including $[1,2)$ children.
There are 881 admitted parameter choices in the saved finite search.
It refines selected choices with quadrature orders 24, 48 and 80.
For $K=6.81$, the order-80 outcomes at the selected candidates are
\[
\begin{array}{c|c|c|c}
s&\hbox{ordinary-only seed}&\hbox{feedback}&\hbox{candidate}\\\hline
2.1&.02047094&.00253880&.02300973\\
2.2&.02057652&.00243968&.02301620\\
2.3&.01879876&.00245380&.02125256\\
2.4&.01571500&.00244468&.01815968\\
2.5&.01164637&.00229828&.01394466.
\end{array}
\]
These are sampled choices, not certified optima.
At the original parameters $(2.2,4.54,3.53,2.90,2.44)$ the
diagnostic is $.02033936+.00258581=.02292517$.
The feedback splits approximately as $.00093624$ from the three
negative terms and $.00164957$ from the four ordinary pair terms;
the first-generation $\Gamma_1$ contribution is zero because these
initial $H$ inputs vanish above three.  The same original parameters
give $.02299950$ if the separately supplied $K=6.792$ is used.
No validity or certification of that alternative $K$ is inferred
from this probe.

The discontinuities in chosen input defects make these quadratures
especially unsuitable as outward bounds without further partitioning.
Their agreement is a diagnostic, not a certificate.  The operator
shows a nonzero additional feedback of order $0.0025$ in a seed,
but that is not itself a gain of $0.0025$ in the complete Chen count.
The contribution must still pass through the full eleven-term
integrals and their existing minima.  These computations alone do
not fill the remaining frontier gap, establish a new global count,
or justify importing any old double-sieve table.  A subsequent
finite typed iteration may use this operator, but its resulting
integrals require their own certificates.

\subsection{Removing the incidental one-twelfth threshold restriction}
\label{sec:free-threshold-fixed-cutoff-extension}

The lower bound $a>1/12$ in
Proposition~\ref{prop:free-weight-thresholds} only fixed the number
of original prime groups at eleven. The later fixed-positive-cutoff
product extension permits that number to depend on the fixed $a$.
Consequently the same functional and pointwise weight proof apply
in the enlarged sufficient range
\begin{equation}\label{eq:extended-free-threshold-range}
 0<a<h<d<57/224,\quad
 25/224<h<\min(25/128,2a),\quad
 1/3<h+d<1/2,\quad 4a+d<1,\quad4h+d<1.
\end{equation}
All three parameters remain fixed before $N$ tends to infinity.

To check this extension, the direct polynomial weight inequalities
are valid for every number of rough prime factors; their case
analysis uses $h+d>1/3$ and $4a+d<1$, not an upper bound of eleven.
The one- and two-prefix factorization gates, restricted lower
argument and four-prefix quotient conditions are exactly those
checked in the preceding proof and are retained here.
The exceptional square and small-size errors remain
$O(N^{1-a+o(1)})$ for each fixed positive $a$.

In the switched product estimate there are now at most
$\lfloor1/a\rfloor$ original prime groups. Apply the already
proved fixed-positive-cutoff extension with cutoff exponent $a$,
choosing the auxiliary parameter anew so that $12\zeta<a/2$.
Each group has at most six Heath--Brown factors. A group without
a middle factor must have a long factor; otherwise its short
factors cannot produce a prime above $N^a$. There can still be
at most two groups with long factors, by their total product bound.
Thus the same one- and two-long alternatives and the same modulus
sieve coefficient apply. The number of groups is fixed and finite;
only harmless constants and the allowed logarithmic saving change.
No new distribution level is asserted by this extension.

Every sieve function in the integral functional must retain its
full definition. In particular the ordinary four-prefix option is
\[
 \frac{4A(\max\{2,(\nu-S)/w\})}{\nu-S},
 \qquad A(s)=\frac{sF(s)}{2e^\gamma},
\]
with the same physical cutoff comparison when its argument is
below two. It cannot be replaced by $4/(\nu-S)$ when the argument
exceeds three. The lower terms likewise retain the full $f$.

For applying the separately proved uniform double-sieve options,
all old stable-prefix gates remain necessary. The new range implies
$a>25/448$. Hence at a four-prefix state
\[
 T\ge4a>25/112>2(153/256)-1,
\]
which gives $(1-T)/(\nu-T)>2$ whenever $\nu\le153/256$ and
$\nu>T$. This observation does not supply that gate for a one-
or two-prefix state: those states must still be checked individually,
and an unlicensed defect is zero. Similarly, each comparison cutoff
must stay below its licensed stable ceiling.

This extension is an analytic option for a wider parameter search.
It provides no evaluated new global coefficient or exponent.
Any resulting numerical optimization remains a diagnostic until
all signed integrals and common-box maxima and minima have
outward bounds.

% New finite recursion, 2026-10-08. No imported old H/h tables.
\subsection{A finite recursion on typed stable-prefix states}
\label{sec:new-finite-double-sieve-recursion}

This construction uses only the explicitly proved first and replacement
second seeds, their stable-prefix extensions, and finite Buchstab
identities. It does not transplant any old correction table, claim an
optimal correction function, or take an infinite iteration limit.
Its numerical implementation is
\texttt{probe\_new\_finite\_double\_sieve.py}; the saved JSON is a
floating feasibility calculation, not an outward certificate.

At a state retain the marked specified prefix $B=N^{T+o(1)}$, one
total level $Q=N^\nu$, residual $R=Q/B=N^{L+o(1)}$, and
$\alpha=(1-T)/L\ge2$. The state also retains its enhanced
factorization branch, all its gates and square tests, and a ceiling
$\chi$ below which that branch persists. For a $W$ branch with a
future third-slot bound $c$, require $\chi\le c$ as well as
$\chi$ below the smallest specified slot. Alternatively the state
uses the classical absolute Bombieri--Vinogradov interface at
$\nu<1/2$. All physical cutoffs used below are below the retained
ceiling. Enhanced non-switched nodes retain full well support;
positive switched nodes use the separate finite-product theorem.

Put $A(s)=sF(s)/(2e^\gamma)$ and $a(s)=sf(s)/(2e^\gamma)$.
For fixed generation $n$, the claims represented by $H_n,h_n$ are
the simultaneous bounds
\[
 \Phi_s\le\{A(s)-H_n(s)+o(1)\}\Theta_Q,\qquad
 \Phi_s\ge\{a(s)+h_n(s)+o(1)\}\Theta_Q,
\]
uniformly over the indicated typed states and every $\alpha\ge2$.
Only upper nodes at $s\ge1$ and lower nodes at $s\ge2$ are used;
set $h_n(s)=0$ below two. The initial seed is
\[
 H_0(s)=\frac1{50}\mathbf1_{1\le s\le11/5},\qquad h_0(s)=0.
\]
Its strict uniform inequality is separately certified in
\texttt{second-seed-omega-certified.tex}. Zero is chosen outside
the displayed support. A finite positive prefix partition licenses
all choices of upper or lower estimates.

\paragraph{A common upper envelope for the Buchstab function.}
Let $C=70893/125000=.567144$. The exact inequality
$C\sum_{j=0}^{8}C^j/j!>1$ implies $Ce^C>1$, hence
$C+\log C>0$. On $2\le u\le3$ the maximum of
$1+\log(u-1)-Cu$ is $-\log C-C<0$; the delay equation propagates
$\omega(u)\le C$ for $u\ge3$. Since $\omega(u)=1/u$ on $[1,2]$,
\[
 \omega(v)\le\Omega(u):=\max\{1/u,C\}\qquad(v\ge u\ge1).
\]
Consequently the first-seed switched integral is bounded uniformly
for all $\alpha\ge2$ by
\[
 \overline I(s,s_0)=
 \int_{1/s_0<u_1<u_2<u_3<1/s}
 \Omega\!\left(\frac{2-u_1-u_2-u_3}{u_2}\right)
              \frac{du_1du_2du_3}{u_1u_2^2u_3},\qquad s\ge2.
\]
Its argument is at least $2s-3\ge1$. This is a pointwise tail
majorant, not an assumption that the switched integral attains its
supremum at $\alpha=2$. Unit quotients in this triple term are
power-saving; the second seed retains its separate five- and
six-prime unit masses.

\paragraph{The positive finite operators.}
Fix $M=10$, and choose zero defects at any unproved argument beyond
the finite table. For $s\ge2$, the following are valid new candidates:
\begin{align*}
 \mathcal B_H[h_n](s)&=\int_{\max(2,s-1)}^M h_n(w)\frac{dw}{w},\\
 \mathcal B_h[H_n](s)&=\int_{s-1}^M H_n(w)\frac{dw}{w},
\end{align*}
where a reversed interval contributes zero. These identities start
the ordinary comparison at argument $M+1$ and discard any correction
there. Ordinary baseline cancellation is exact at that common level.
In particular upper children with physical argument $1\le w<2$
are included in the lower operator. Their defect comes from a
smaller comparison cutoff, while the lower function below two is
still zero.

The first weighted inequality gives the further upper candidate
\begin{align*}
 \mathcal W_1[H_n,h_n](s;s_0)
 ={}& A(s)-A(s_0)+H_n(s_0)-\frac K8\overline I(s,s_0)\\
 &+\frac12\int_{1/s_0}^{1/s}
       \frac{a(s_0(1-u))+h_n(s_0(1-u))}{u(1-u)}\,du,
 \qquad 2\le s<s_0\le10.
\end{align*}
The feedback is evaluated at the same residual levels as the
ordinary terms. It is not added at an old argument.
The earlier first-seed search used $s_0\le5$. Wu's pointwise
twofold inequality also holds for every fixed $2\le s<s_0\le10$.
In that larger range use the full $A,a$, not their initial-interval
logarithmic simplifications. The negative-term prime still satisfies
$q^2\le R$, the triple unit-quotient gap remains
$\alpha-3/s\ge1/2$, and every switched prime has exponent at least
$L/s_0>0$. Thus the same signed-prefix and fixed-positive-cutoff
arguments prove exactly the extension used here.

For each admissible finite parameter tuple of the second seed, the
operator in \texttt{second-seed-finite-feedback.tex} contributes
\[
 \mathcal W_2[H_n,h_n](s)=J_2(s)+\frac15
 \left\{4H_n(s_0)+H_n(\phi_1)
       +\sum_{i=2}^4\int h_n(\rho_i)\,d\lambda_i
       +\sum_{i=5}^8\int_{\mathcal G_i}H_n(\rho_i)\,d\lambda_i
 \right\}.
\]
Here $\mathcal G_5$ is only the good square-supported part of term 5;
the bad part remains a positive switched cost. For terms 6--8 it is
the full verified domain. The coefficient measures and child
arguments are precisely those derived in that separate fragment.
The finite parameter bank is recorded by the probe; the chosen tuple
for every generation and target argument is saved in its history.

Take the maximum of the previous defect, zero, and these separately
valid candidates. Then apply positive-count monotonicity to the
upper bound, in its full normalization:
\[
 H_{n+1}(s)\ge
 \sup_{\max(2,s)\le v\le M}
       \{A(s)-A(v)+\widetilde H_{n+1}(v)\}.
\]
For $s\ge2$ retain the candidate at $v=s$ as well. The shift
$A(s)-A(v)$ is essential when an argument exceeds three. No lower
monotonicity with an enlarged physical cutoff is invoked. Maxima
select alternative estimates for a count; overlapping defects are
never added.

\paragraph{Why every generation stays in the typed family.}
A Buchstab child has $q=R^{1/(w+1)}$, $w\ge s-1\ge1$.
Thus $q^2\le R$, $q$ is below the parent's physical cutoff, and
its updated ratio is
\[
 \alpha' =\alpha+\frac{\alpha-1}{w}\ge2.
\]
The complete square condition and selected-remaining conditions
follow by multiplying the residual test by the old prefix.
The new cutoff is at most $q$, which is also below the old stable
ceiling. If $w<2$, its upper comparison is made at $v\ge2$ and
has cutoff $(R/q)^{1/v}\le q$. The subsequent square test is
$q(q')^2\le q(R/q)=R$; it does not require $3\log_Nq\le L$.

The same tests were checked separately for each first-seed and
second-seed non-switched term. Their cap data persist below the
current ceiling, and their ratio after adding exponents remains
at least two. Consequently the functions may be taken uniform in
$\alpha$, even though the actual ratio increases along the tree.
At any fixed finite depth the residual exponents and all newly
specified prime exponents have a fixed positive lower bound on
compact boxes. It can depend on the depth. The finite-product $K$
proof rechooses its Heath--Brown precision for that bound; no
uniformity in infinite depth or in a cutoff tending to zero is
claimed. The modulus coefficient $K$ and its density integral are
independent of this fixed factor-count bound.
For the newer fourth-prefix and tail modulus weights this extension
uses the same original-group argument as
\texttt{generalized-one-step-double-sieve.tex}: retain the original
Heath--Brown scales, choose $12\zeta<\epsilon/2$, and keep the number
of groups fixed before $N\to\infty$. Those new modulus weights
already have the normalized factorization and local-scale margins
needed by the same Type II and long-variable inputs. The coprime
mean and prime-power errors still save a fixed power. Changing this
fixed prime-factor cutoff therefore leaves the modulus-weight density
and its numerical $K$ unchanged. This is an extension of the positive
count theorem, not a signed switched-distribution assertion.

\paragraph{Insertion in the full Chen functional.}
Use the free-threshold eleven-term identity with the complete
$Y_{11}$ range. The probe uses $K=1697/250=6.788$, from the separately
certified fourth-prefix tail refinement, in the five original
switching terms and the positive costs of the two seeds.
The three-parameter range is the positive-$a$ extension of the
earlier range, retaining all its other strict inequalities.
The unrestricted $L_1$ term keeps its established ordinary baseline;
no unproved defect for its different prefix family is inserted.

For $L_2$ add the upper defect in its exact least-prime identity.
For $U_3,U_4$ replace the normalized ordinary upper multiplier by
$A(s)-H_n(s)$ at the actual selected level. For $L_5,L_6$ replace
$a(s)$ by $a(s)+h_n(s)$ only where their full square tests and
$\alpha\ge2$ hold. The latter gate matters for small $a$ on the
constant branch of $V_{2,a}$; the probe retains ordinary lower
bounds and chooses zero extra defect if that gate fails.

On each eligible four-prefix box use
\[
 \min\left\{\frac Kz\omega\left(\frac{1-T}{z}\right),
             \frac4{\nu-T}
               \left[A\left(\max\{2,(\nu-T)/w\}\right)
                     -H_n\left(\max\{2,(\nu-T)/w\}\right)\right]
       \right\}.
\]
All old cap and gate tests are retained. The full $A$ is required
when $a<1/12$; it is not replaced by one. This minimum is evaluated
directly in the functional, so neither its ordinary improvement nor
its switched improvement is counted twice. Each choice is made on
a common finite positive prefix partition before taking limits.

The implementation uses finite arrays, interpolation, Gauss
quadrature and local parameter searches to diagnose the resulting
explicit operators. Those approximations are not a proof that the
computed array entries are outward lower bounds, even though the
underlying exact finite operators preserve the stated inequalities.
An observed small generation change is not an infinite-iteration
error bound or an optimality certificate.

\paragraph{Numerical decision and a hypothetical coefficient test.}
With 32 generations, $\Delta s=.01$, $\Delta s_0=.05$, and ten
Gauss nodes in each triple coordinate, the optimized diagnostic tuple
is approximately
\[
 (a,h,d)=(.08315180,.15198811,.23634369).
\]
The complete functional at quadrature orders 16, 24 and 36 is
respectively $1.93228169$, $1.93234623$, and $1.93235164$.
At $K=6.788$ its zero-margin imbalance diagnostic is $.75226117$;
the remaining coefficient gap to the internal comparison $b=.749$
is approximately $.02949097$. At the original sampled tuple the
generation-24 to generation-32 increase is only $.00002118$.
This is evidence about this finite construction, not an upper bound
on all possible double sieves.

For a decision about future prefix work, keep this exact same
floating $H,h$ array and threshold tuple, and hypothetically lower
only the coefficient of the five original switched counts and the
bad-semiprime removal. Re-evaluating the full four-prefix minimum
gives
\[
\begin{array}{c|c|c}
 \text{assumed }K&c\text{ diagnostic}&c-K\log(1/.749)\\\hline
 6.786&1.93313658&-.02812800\\
 6.750&1.94726640&-.00359360\\
 6.745&1.94922902&-.00018589\\
 6.740&1.95119168& .00322185\\
 6.735&1.95315437& .00662962
\end{array}
\]
The crossing is near $K=6.74473$. In particular a future valid
$K=6.735$ would make this sampled finite construction a plausible
route past the internal comparison. That $K$ is hypothetical here;
the table is not a proof of it, a certified functional value, or a
new imbalance theorem. No reoptimization or further recursive gain
from lowering $K$ is included in this sensitivity test.

\subsection{Consequences of the fifth-prefix gain and sharper second seed}
\label{sec:fifth-prefix-omega-corollary}

Use the original fourth-prefix compact patches, the disjoint
corrected supplement of Section~\ref{sec:fourth-prefix-tail-refinement},
and the fifth-prefix boxes of
Section~\ref{sec:higher-prefix-K-refinement} in that order.
The fifth-prefix bound uses an upper envelope of every eligible
four-prefix cap, so it applies to this actual partially refined
tree. It does not require its full floating four-prefix gain.
Subtracting the three rigorously bounded gains from the old exact
extended endpoint gives
\[
 K_{\rm limiting}\le
 \frac{1230679674488580033396846301}{180734410032387061618573312}
 -\frac{10720642480494928338187}{604462909807314587353088}
 -\frac{5155254219572575727579}{1208925819614629174706176}
 -\frac{2167182368869852589729}{1208925819614629174706176}
 <\frac{3393}{500}=6.786.
\]
The strict slack permits compact trimming, fixed buffered levels
and the inherited ordered limits. Denote this final rounded cap
by $K_5$. The corrected fourth-prefix supplement alone does not
give this rounded number; the fifth-prefix gain is required here.

Keep the old threshold tuple and its outward eleven-term table
from Section~\ref{sec:fourth-prefix-second-seed-corollary}.
The five switched upper entries scale by $K_5/(681/100)$.
The first narrow correction block still gives $2/6635$.
Proposition~\ref{prop:second-seed-omega-certified} increases the
second seed from $9/500$ to $1/50$, with no new feedback or
floating integration input. The already proved finite propagation
is linear in this seed, so the disjoint second block now gives
\[
 \frac{10}{9}\frac{56}{465275}=\frac{112}{837495}.
\]
It replaces the old second-block contribution; the two versions
of that same block are not added. The new total coefficient is
strictly greater than
\begin{equation}\label{eq:fifth-prefix-omega-c}
 \frac{9089}{5000}
 +\frac{101203}{40000}\left(1-\frac{K_5}{681/100}\right)
 +\frac2{6635}+\frac{112}{837495}
 =\frac{4609497840932309}{2522777813550000}
 >\frac{18271}{10000}.
\end{equation}
The improved second seed continues to use the old positive cap
$6.81$ and its proved fixed-positive-cutoff extension. Thus its
proof does not assume an unproved extension of the new $K_5$.

\begin{theorem}[Further separate research-draft consequence]
\label{thm:fifth-prefix-omega-chen}
Under the named published inputs and the derived interfaces of
this manuscript, the complete eleven-term count satisfies
\[
 D_{1,2}(N)\ge(18271/10000-o(1))\HN.
\]
For every sufficiently large even $N$, there are at least
$\HN/1000$ representations
\[
 N=p+rq,\qquad p,q\in\PP,\quad r=1\ \hbox{or}\ r\in\PP,
 \qquad r\le q^{76407/100000}.
\]
No effective threshold for $N$ is asserted.
\end{theorem}

\begin{proof}
The first conclusion is the exact table calculation
\eqref{eq:fifth-prefix-omega-c}. For the second use the same
bad-semiprime removal inequality with coefficient $K_5$.
For $b=76407/100000$ and $u=(1-b)/(1+b)$, the positive
atanh expansion bounds
\[
 \log(1/b)<2\left(u+\frac{u^3}{3}+\frac{u^5}{5(1-u^2)}\right).
\]
The separate rational certificate verifies that $18271/10000$
minus $K_5$ times this upper bound is strictly greater than
$1/1000$. The strict surplus absorbs the asymptotic losses.
\end{proof}

Accepting the published coefficient $1233/625$ as an external
input gives instead $b=74774/100000$ and at least $\HN/50000$
representations by the same rational calculation. The recorded
limitations of that external proof audit remain in force.
This is a project deduction and is not a quoted imbalance theorem
from that source.

The finite-feedback iteration and enlarged-threshold searches
are separate floating diagnostics and are not used in either
conclusion. These are research-draft consequences under the
specified analytic inputs, supported by exact functional arithmetic
and targeted independent review. They are not formal verification
of the entire argument, a determination of bibliographic priority,
or completion of the requested stronger-than-frontier objective.

% New replacement certificate fragment, 2026-10-08. Main/state untouched.
\subsection{A stronger complete fourth-prefix certificate}
\label{sec:fourth-prefix-stronger-refinement}

This construction replaces the entire earlier partial fourth-prefix
plus fourth-tail subtraction.  Neither of those two old numerical
gains is added to the new gain.  It keeps the old extended tree through
rank four, and changes only subsequent odd cube tests on the new finite
collection of fourth-prefix boxes.  All first and third decisions,
selected-subproduct interfaces, signed atoms, and product-family
normalizations are those of Proposition
\ref{prop:fourth-prefix-descendant-refinement}.

Write $t=(x,y,z,w)$, $S=x+y+z+w$, and
\[
 v(s)=\min\{(1+s)/2,57/64-3s/2\},\quad
 L=\max\{v(x),v(x+y),v(x+y+z),v(x+z),v(y+z),153/256-3z/8\}.
\]
Every box has $3/200<w<z<y<x<39/200<25/128$, the old third
cube $x+y+3z<L$, a chosen endpoint $M=v(s_I)>L$, and the full
square and all selected-subproduct gates below $M$.  No gate is
replaced by a top-slot-only surrogate.  The new cap is frozen at rank
four and used for ranks $5,7,\ldots$; no fourth-rank exclusion is added.
The exact even-subtree identity gives the gain
\begin{equation}\label{eq:fourth-stronger-target}
 4\int_{\mathcal C}\frac{
   \overline F((L-S)/w)-\overline F((M-S)/w)}{xyzw^2}\,dt,
 \qquad \overline F=F/(2e^\gamma).
\end{equation}
The old third cube makes the first argument greater than one.
Only boxes whose certified upper argument is below five are used.

All fifteen subsets have been accounted for.  Subsets contained in
the first three slots already occur in $L$, except for the singleton
second and third slots, which are bounded by $v(x)$ since $v$ is
increasing below $25/128$.  The singleton fourth slot is likewise
bounded by $v(x)$.  Thus only masks $9,\ldots,15$ can improve $L$
in this region.  This pruning removes no potentially improving branch.

\paragraph{A positive derivative numerator through argument five.}
Define $\mathcal N(s)=s^2(-\overline F'(s))$.
For $1<s\le3$, $\mathcal N(s)=1$.  For $3\le s\le5$, the
ordinary sieve delay equation and the explicit lower function on
$2\le s-1\le4$ give
\begin{align*}
 A(s)&=1+\int_3^s\frac{\log(t-2)}{t-1}\,dt,
 &\overline F(s)&=A(s)/s,\\
 \mathcal N(s)&=A(s)-\frac{s}{s-1}\log(s-2),
 &\mathcal N'(s)&=\frac{s}{(s-1)^2}
       \left(\log(s-2)-\frac{s-1}{s-2}\right)<0.
\end{align*}
Indeed $1+1/u-\log u$ decreases on $1\le u\le3$, and is positive
at $u=3$ since $\log3<10/9<4/3$.  The same inequality proves
that $\log(t-2)/(t-1)$ increases on $[3,5]$.
The exact table below proves $\mathcal N(5)>0$, so the numerator
is positive and nonincreasing throughout the range used here.
In particular the factor $s/(s-1)$ has been retained.

Take $m=4096$, $Q=2^{80}$, and $s_j=3+j/m$ for $0\le j\le2m$.
For $z=j/(2m+j)$ use the first $24$ terms of
$2\sum_{k\ge0}z^{2k+1}/(2k+1)$, with each term rounded down to
the $Q^{-1}$ lattice.  An upper bound is obtained by adding
$24/Q$ and the geometric tail bound
\[
 \frac{2z^{49}}{49(1-z^2)}.
\]
Denote the resulting logarithm intervals by
$[\ell_j/Q,u_j/Q]$.  Left Darboux sums and integer rounding give
\[
 n_j=Q+\sum_{i<j}\left\lfloor\frac{\ell_i}{2m+i}\right\rfloor
       -\left\lceil\frac{(3m+j)u_j}{2m+j}\right\rceil
 \le Q\mathcal N(s_j).
\]
For arbitrary $s\in(3,5]$ use $n_{\lceil m(s-3)\rceil}/Q$.
For $s\le3$ use one.  This direction follows from the proved
monotonicity, independently of any floating interpolation.
The table's endpoint is
\[
 \frac{n_{8192}}Q
 =\frac{19821997013242695217835}{604462909807314587353088}
 >0.03279.
\]

\paragraph{Correlated affine bounds.}
Every old component is a minimum of one or two affine functions.
Let $\mathcal A$ be the two affine pieces of $M$, and let
$\mathcal B_j$ be the pieces of old component $L_j$.
For a box $B$, write $\inf_B a$ and $\sup_B a$ for exact affine
endpoint evaluation.  A valid whole-box improvement margin is
\[
 d_B=\min_j\min_{a\in\mathcal A}
                 \max_{b\in\mathcal B_j}\inf_B(a-b)
 \le\inf_B(M-L).
\]
This is valid because $a-L_j=\max_{b\in\mathcal B_j}(a-b)$.
It retains correlations between shared coordinates.  Similarly
$\max_j\min_{b\in\mathcal B_j}\inf_B(b-x-y-3z)$ is a lower
old-cube margin, and $\min_{a,g}\inf_B(a-g)$ checks every full
and selected gate $g$.  All these margins must be strictly positive.

These numerical rectangles certify pointwise correlated inequalities;
they need not satisfy $\inf_B M>\sup_B L$.  The actual sieve labels
are allowed a further finite common refinement.  The uniform positive
gap and gate margins, together with the piecewise-affine continuity of
$L$ and $M$, then give on each smaller label a constant buffered cap
above its inherited old cap and all gates, and below its new endpoint.
For every prescribed positive error, this refinement and buffering
approximate the displayed endpoint gain integral within that error.
Thus the exact integral certificate and a strict final numerical
margin yield a single finite legal tree; a frozen constant cap over
each original numerical rectangle is not being assumed.

For each old component choose one of its affine pieces $b_j$.
The implementation chooses the piece smallest at the box center,
but validity does not depend on that rule.  Then
$U=\max_j b_j\ge L$ is convex, and $M-U$ is concave.  Multilinear
corner interpolation therefore gives
\[
 \frac1{|B|}\int_B(M-L)\,dt
 \ge\frac1{16}\sum_{v\text{ vertex of }B}(M(v)-U(v)).
\]
Take $\widehat d_B$ to be the larger of $d_B$ and this average
bound rounded down to the common affine lattice.  Then
$\int_B(M-L)\,dt\ge|B|\widehat d_B$.
The separate condition $d_B>0$ continues to license the entire box;
a positive average never licenses a point with a failed gate.

\paragraph{Four positive subintervals of the cap difference.}
Put $r_0=L-S$, $r_1=M-S$.  Use the correlated residual bounds
\[
 R_0=\max_j\min_{b\in\mathcal B_j}\sup_B(b-S),\qquad
 R_1=\max\{R_0,\min_{a\in\mathcal A}\sup_B(a-S)\}.
\]
Thus $0<r_0\le R_0\le R_1$ and $r_1\le R_1$.
Write $R_\theta=(1-\theta)R_0+\theta R_1$ and let
$\ell_w$ be the lower fourth-coordinate endpoint.  Require
$R_1/\ell_w<5$.  If $c_j$ is the certified lower value for
$\mathcal N(R_{j/4}/\ell_w)$, then the box contributes at least
\begin{equation}\label{eq:fourth-stronger-box}
 \frac{4|B|\widehat d_B}{\prod_{i=1}^4 u_i}
 \sum_{j=1}^4\frac{c_j}{4R_{(j-1)/4}R_{j/4}}.
\end{equation}
To check the normalization, the cap difference equals
\[
 \overline F(r_0/w)-\overline F(r_1/w)
 =w(M-L)\int_0^1
 \frac{\mathcal N(((1-\theta)r_0+\theta r_1)/w)}
 {((1-\theta)r_0+\theta r_1)^2}\,d\theta.
\]
The displayed factor $w$ cancels one of the two fourth-coordinate
factors in \eqref{eq:fourth-stronger-target}.
The table and $R_\theta$ bound the remaining positive integrand
uniformly on the box.  Finally
$\int_{(j-1)/4}^{j/4}R_\theta^{-2}d\theta
=1/(4R_{(j-1)/4}R_{j/4})$, also when $R_0=R_1$.

\paragraph{Integer implementation and independent finite geometry.}
The coordinate lattice is $n=204800$ and the affine lattice is
$D=32n$.  Let a stored rectangle have integer lower and upper
coordinates $a_i,b_i$, and put $h_i=b_i-a_i$.
Write $d=D\widehat d_B$, $r=D R_0$, and $q=D R_1$.
For subinterval $j$ define
\[
 A_j=(5-j)r+(j-1)q,\qquad B_j=(4-j)r+jq.
\]
If $n_{k_j}$ is the table integer at the upward-rounded argument
$B_j/(128a_4)$, the subinterval contributes at least the following
number of $Q^{-1}$ units:
\[
 \left\lfloor
 \frac{16D d\{\prod_i h_i\}n_{k_j}}
      {\{\prod_i b_i\}A_jB_j}
 \right\rfloor.
\]
There is no floating arithmetic in this formula or in any gate.

The adaptive probe supplies only a finite witness of integer boxes
and branch labels.  The exact replay checks their dyadic alignment,
membership in individual $1/200$ root cubes, and pairwise disjointness
by recursive midpoint separation.  It does not trust the generation
history, the refinement priorities, the diagnostic gains, or the
floating sieve-function values.  All geometry and all contributions
are recomputed.  The NPZ witness, outward derivative table, script,
and final output are hash-bound by the companion files.

\begin{proposition}[Certified replacement fourth-prefix gain]
\label{prop:fourth-prefix-stronger-certificate}
Under the inherited fourth-prefix analytic interfaces, the stored
finite witness has limiting gain at least
\begin{equation}\label{eq:fourth-stronger-exact-gain}
 \Delta_{4,\mathrm{new}}\ge
 \frac{10601625247768616959055}{302231454903657293676544}
 >\frac7{200}=0.035.
\end{equation}
Consequently every coefficient
\[
 K_*>
 \frac{1224339902590414400455331411}
      {180734410032387061618573312}
 =6.774249034099353\ldots
\]
is permitted by the previous exact extended-tree bound and this
replacement gain.  In particular $K_*=67743/10000=6.7743$ is a
safe standalone fourth-prefix coefficient.  No earlier partial
fourth-prefix or fourth-tail gain is added to this subtraction.
\end{proposition}

The exact output is recorded in
\texttt{certify\_fourth\_prefix\_stronger.json}.
There are $3699127$ certified boxes, inside $9963$ original root
cubes.  The independent disjointness routine visits $7607322$ nodes.
The branch counts for masks $9,10,11,12,13,14,15$ are respectively
\[
 1235537,\quad1931050,\quad36,\quad430050,\quad152,\quad102302,\quad0.
\]
The minimum positive affine slack in the combined checks is $128/D$.
All geometry is exact; the final sum rounds each of the four positive
subinterval contributions on every box downward.  The complete
selected-stream SHA256 is
\texttt{6b068f041379c73d01624e0014980ca71bfef6baa0c178398b0298c847b4b936}.

Compactly trim
coordinate boundaries and buffer the finitely many frozen levels
before applying the inherited density theorem.  The finite limits
and a strict final numerical margin give one legal fourth-prefix
tree.  Its cap is everywhere bounded above by the maximum of all
eligible four-slot endpoints.  Hence a later fifth-prefix certificate
using an upper envelope over every potentially eligible four-slot
cap continues to apply to this replacement tree.

% Integration fragment: finite feedback, not Lean formalization.
\subsection{Four certified finite feedback steps}
\label{sec:second-seed-certified-finite-feedback}

The positive operator of Section~\ref{sec:second-seed-finite-feedback}
can be used a fixed number of times with an exactly certified input.
The construction below uses the original coefficient $K=681/100$
throughout. It does not depend on a numerical infinite iteration or
on the applicability of the newer modulus weights at a smaller cutoff.

Write $S=11/5$, $s_0=227/50$, and
$(\phi_1,\phi_2,\phi_3)=(353/100,29/10,61/25)$.
The family consists of the typed stable-prefix and classical states
of Section~\ref{sec:new-finite-double-sieve-recursion}; in particular
the seed comparison cutoff must satisfy $L/S<\chi$ on enhanced
states. Every descendant retains the full and selected square tests,
the actual residual level, its stable ceiling, and $\alpha\ge2$.
No assertion below discards this comparison-cutoff condition.

\begin{proposition}[A larger seed after four finite rounds]
\label{prop:second-seed-certified-finite-feedback}
Under the analytic inputs of
Proposition~\ref{prop:second-seed-stable-prefix}, on the admitted
states just specified, one has
\[
 \Phi_{11/5}\le
 \left(1-\frac{2327}{100000}+o(1)\right)\Theta_Q.
\]
The estimate is uniform in the admitted $\alpha\ge2$ and uses only
four finite feedback rounds starting from the seed $1/50$.
\end{proposition}

\begin{proof}
For a previously proved positive constant $J$, choose only the
following upper defect and its one-step Buchstab lower consequence:
\[
 H_J(w)=J\mathbf1_{[1,S]}(w),\qquad
 h_J(w)=J\log\frac{S}{w-1}\,\mathbf1_{[2,1+S]}(w).
\]
The upper option on $w<S$ follows by comparing positive counts
with the admitted smaller cutoff at $S$; the ordinary normalized
upper multiplier is one on this range. The lower option is
$\int_{w-1}^S J\,dt/t$, including the already justified upper
children with $1\le t<2$. It is used only on states satisfying
their inherited comparison and square conditions.

Here $h_J(w)$ needs admission of its own physical cutoff, not a
new assumption that this lower parent itself admits argument $S$.
Indeed, for a Buchstab child with argument $v\in[w-1,S]$, its new
prime has exponent $L/(v+1)$ and its residual exponent is
$Lv/(v+1)$. Its seed comparison exponent satisfies
\[
 \frac{Lv}{S(v+1)}\le\frac{L}{v+1},
\]
so the compared cutoff is at most that new prime, and remains
below the old stable ceiling. Since $v\ge1$, its new square test
also holds. This distinction is needed for the single-prime lower
children of the weighted inequality: they need not themselves
admit the larger cutoff associated with argument $S$.

Insert these previously proved options in
\eqref{eq:second-seed-feedback-operator}. The $\Gamma_1$ feedback
vanishes, since $s_0,\phi_1>S$. The remaining feedback is a positive
linear function of $J$; write it as $\lambda J$.
The three single-prime integrals use
$\rho=s_0(1-u)$ and hence start at
\[
 u_* = 1-\frac{1+S}{s_0}=\frac{67}{227}.
\]
Their upper endpoints are $1/S,1/\phi_2,1/\phi_3$ respectively.
All their actual arguments exceed two, so no lower defect below
two is used.

For the four pair terms, set $a_i=1/\phi_i$ and $a_0=1/s_0$.
The exact support restrictions for the constant upper defect are
\[
 u+v\ge\frac{117}{227}\quad(i=5,6),\qquad
 v\ge1-\frac{16}{5}u\quad(i=7,8).
\]
They are intersected with the original $D_5^+,D_6,D_7,D_8$;
there is no feedback on the bad part $D_5^-$.
For fixed $u$, if the resulting lower and upper endpoints are
$l(u),r(u)$, the exact pair contribution is
\[
 J\int \frac{1}{u(1-u)}
 \log\frac{r(u)(1-u-l(u))}{l(u)(1-u-r(u))}\,du.
\]
All endpoints are affine with rational coefficients. For $D_5^+$
the three $u$ intervals end at $117/454,19/58,1/3$, with initial
endpoint $50/227$. The rectangle $D_6$ needs no support truncation.
For $D_7$ the nonempty support starts at $1265/5648$ and its lower
boundary changes at $5/21$. For $D_8$ the lower boundary changes
at $1265/5648$. These cuts also follow directly by intersecting
the displayed lines; no floating boundary decision is required.

The certificate \texttt{certify\_second\_seed\_first\_feedback.py}
uses 4096 rational Darboux intervals on each resulting piece and
the previously audited 80-bit outward logarithm arithmetic. At
$J=1/50$, after division by five, it gives
\[
 \lambda\frac1{50}\ge
 \frac{3152244426576931077198}{1208925819614629174706176}
 >\frac{13}{5000},
 \qquad\text{hence}\qquad \lambda>\frac{13}{100}.
\]
Separately, the ordinary seed certificate
\eqref{eq:second-seed-omega-exact-bound} gives
\[
 J_2^\omega\ge
 \frac{24479824718463837928127}{1208925819614629174706176}
 >\frac{20249}{10^6}.
\]
Consequently an admitted uniform short-support seed $J_n$ proves
the next uniform seed
\[
 J_{n+1}=\frac{20249}{10^6}+\frac{13}{100}J_n.
\]
At each round first derive $h_{J_n}$ from the already proved
$H_{J_n}$, and then apply the second-seed inequality. Thus no
unknown estimate from the new round is used to prove itself.
Positive-count comparison then supplies $H_{J_{n+1}}$ on the
short interval, with the same comparison-cutoff qualification.

Starting with $J_0=1/50$, the four exact rational values are
\[
 \frac{22849}{10^6},\quad
 \frac{2321937}{10^8},\quad
 \frac{232675181}{10^{10}},\quad
 \frac{23273777353}{10^{12}}
 >\frac{2327}{100000}.
\]
They are checked in
\texttt{certify\_second\_seed\_finite\_iteration.json}.
There are only finitely many Buchstab and weighted steps. On
compact pieces their new prime exponents remain bounded below by
a fixed positive constant depending on this finite depth. The
already proved fixed-positive-cutoff product theorem with $K=6.81$
therefore applies after its precision parameters are chosen for
that depth. Finitely many errors remain $o(1)$, and the strict
final slack absorbs them in the inherited order of limits.
\end{proof}

For insertion in the earlier narrow lower-integral block, first
use the existing stable descendant propagation at argument three.
Retain the original threshold condition $h>3/20$, so the entire
anchor interval $29/200<t<3/20$ remains inside that lower integral.
It is linear in the short-support seed and gives
\[
 H(3)>\frac{28}{333}\frac{2327}{100000}.
\]
The second narrow-block contribution thus becomes
\[
 \frac{112}{837495}\,
 \frac{2327/100000}{1/50}.
\]
This replaces the older value for the same block. It is not added
to it. In particular, the seed at argument $11/5$ is not applied
directly to an anchor whose comparison cutoff exceeds its smallest
specified prime; the separate argument-three propagation is
essential. No entire floating $H,h$ table or global Chen coefficient
is asserted by this finite seed certificate alone.

\subsection{Combining the replacement fourth tree with finite feedback}
\label{sec:stronger-fourth-finite-seed-corollary}

Use the complete fourth-prefix refinement of
Proposition~\ref{prop:fourth-prefix-stronger-certificate} in place
of both the original partial fourth-prefix refinement and its
later tail supplement. Its gain is not added to either of those
two older gains. Then retain exactly the earlier fifth-prefix
boxes of Section~\ref{sec:higher-prefix-K-refinement}.
That fifth-prefix argument bounds the old cap by an envelope of
every potentially eligible four-slot endpoint. It therefore
continues to bound the actual cap of this replacement fourth tree.
A common finite refinement of their boxes preserves the old first
and third decisions, the full and selected square tests, and the
signed subtree identities at ranks four and five.

The sequential limiting gain is consequently bounded by
\begin{align*}
 K_{\rm limiting}
 &\le\frac{1230679674488580033396846301}
              {180734410032387061618573312}
 -\frac{10601625247768616959055}{302231454903657293676544}
 -\frac{2167182368869852589729}{1208925819614629174706176}\\
 &=\frac{2448031817652536714986333851}
          {361468820064774123237146624}
 <\frac{2709}{400}=6.7725.
\end{align*}
Compact trimming, finite frozen levels and their buffers consume
less than this strict slack. This proves a single admissible
product-sieve coefficient $K_*=2709/400$. The bound includes no
additional fifth-prefix diagnostic or sixth-prefix numerical gain.

Keep the old threshold tuple $h=50/309$ and the full eleven-term
table from Section~\ref{sec:fourth-prefix-second-seed-corollary}.
Its baseline lower coefficient is $9089/5000$. The five switched
upper entries, with the double multiplicity of $U_7$ included,
sum to $101203/10000$ and scale by $K_* /(681/100)$.
The first narrow correction block still gives $2/6635$.
For the second block use Proposition
\ref{prop:second-seed-certified-finite-feedback} and its existing
stable descendant propagation at argument three; its contribution
is $16289/104686875$. The two anchor intervals
$(29/200,3/20)$ and $(31/200,4/25)$ are disjoint and below $h$.
The new value replaces the older value on $(29/200,3/20)$.
Thus the resulting Chen coefficient is strictly greater than
\begin{align*}
 c_*^{\rm raw}
 &=\frac{9089}{5000}
 +\frac{101203}{40000}\left(1-\frac{2709/400}{681/100}\right)
 +\frac2{6635}+\frac{16289}{104686875}\\
 &=\frac{14791059723495739}{8072889003360000}
 >\frac{18321}{10000}=1.8321.
\end{align*}
Every feedback round itself used the earlier $K=6.81$ and its
fixed-positive-cutoff theorem. This calculation therefore requires
no unproved extension of the newly refined $K_*$ to those seed
states. It rescales only the original five positive switched terms
and uses the already licensed two disjoint lower corrections.

\begin{theorem}[Replacement-tree and finite-feedback consequence]
\label{thm:stronger-fourth-finite-seed-chen}
Under the named published inputs and the derived interfaces of
this manuscript, the complete eleven-term count satisfies
\[
 D_{1,2}(N)\ge(18321/10000-o(1))\HN.
\]
For every sufficiently large even $N$, there are at least
$\HN/1000$ representations
\[
 N=p+rq,\qquad p,q\in\PP,\quad r=1\ \hbox{or}\ r\in\PP,
 \qquad r\le q^{7631/10000}.
\]
No effective threshold for $N$ is asserted.
\end{theorem}

\begin{proof}
The Chen count follows from the unchanged pointwise eleven-term
inequality and the exact table computation above. For
$b=7631/10000$ set $u=(1-b)/(1+b)$. The positive atanh expansion
with its geometric tail gives
\[
 \log(1/b)<2\left(u+\frac{u^3}{3}
                          +\frac{u^5}{5(1-u^2)}\right).
\]
Subtracting $K_*$ times this upper bound from $18321/10000$ gives
exactly
\[
 \frac{582348998049910355353}{557636286566132280000000}
 >\frac1{1000}.
\]
The positive surplus absorbs the asymptotic errors in the same
bad-semiprime removal argument. The separate exact arithmetic is
saved in \texttt{certify\_stronger\_fourth\_finite\_seed\_transfer.json}.
\end{proof}

With the published coefficient $1233/625$ accepted as an external
input, the same $K_*$ gives $b=7473/10000$ and at least
$\HN/50000$ representations. This remains an external-input
deduction subject to the earlier audit qualifications, rather than
a quoted imbalance result from that paper.

The complete new fourth-prefix witness and the finite seed update
have independent numerical replays. Those checks certify the
specified functionals and finite arithmetic; the analytic
distribution and sieve interfaces remain mathematical dependencies.
The global floating double-sieve table and enlarged-threshold
searches are still excluded. This is an additional research-draft
consequence, not formal verification, established priority, or a
claim that the full requested research objective has been completed.

\begin{thebibliography}{9}
\bibitem{Li2026}
R. Li, \emph{On Chen's theorem, Goldbach's conjecture and almost prime
twins II}, Mathematical Reports 28(78) (2026), nos. 1--2, 39--61.
\href{https://doi.org/10.59277/mrar.2026.28.78.1.2.39}{doi:10.59277/mrar.2026.28.78.1.2.39}.
\href{https://imar.ro/journals/Mathematical_Reports/Pdfs/2026/1_2/3.pdf}{Journal PDF}.

\bibitem{Polymath2014}
D. H. J. Polymath, \emph{Variants of the Selberg sieve, and bounded
intervals containing many primes}, Research in the Mathematical Sciences
1 (2014), article 12. Claim 12 and Theorem 14.
\href{https://doi.org/10.1186/s40687-014-0012-7}{doi:10.1186/s40687-014-0012-7}.

\bibitem{Wu2008}
J. Wu, \emph{Chen's double sieve, Goldbach's conjecture and the twin
prime problem, 2}, Acta Arithmetica 131 (2008), 367--387.
\href{https://arxiv.org/abs/0709.3764}{arXiv:0709.3764}.

\bibitem{Lichtman2025}
J. D. Lichtman, \emph{A modification of the linear sieve, and the count
of twin primes}, Algebra \& Number Theory 19 (2025), 1--38.
\href{https://doi.org/10.2140/ant.2025.19.1}{doi:10.2140/ant.2025.19.1}.
Theorem 2.10 and Remark 2.9 in the
\href{https://arxiv.org/abs/2109.02851v2}{author version, arXiv:2109.02851v2}.

\bibitem{LiLiu2026}
J. Li and J. Liu, \emph{Theorem $(1+1.9)$ on the Goldbach Conjecture},
preprint, arXiv version 2 (2026).
\href{https://arxiv.org/abs/2606.05224v2}{arXiv:2606.05224v2}.

\bibitem{Lichtman2023}
J. D. Lichtman, with an appendix with S. Drappeau,
\emph{Primes in arithmetic progressions to large moduli, and Goldbach
beyond the square-root barrier}, preprint (2023).
\href{https://arxiv.org/abs/2309.08522}{arXiv:2309.08522}.
\bibitem{Maynard2020}
J. Maynard, \emph{Primes in arithmetic progressions to large moduli II:
Well-factorable estimates}, arXiv:2006.07088v1 (2020).
\href{https://arxiv.org/pdf/2006.07088v1}{Author version};
Lemma 7.1 and the proof of Proposition 5.2.

\bibitem{Pascadi2025}
A. Pascadi, \emph{On the exponents of distribution of primes and smooth
numbers}, arXiv:2505.00653v2 (2025).
\href{https://arxiv.org/html/2505.00653v2}{Author version};
Theorems 1.3 and 5.4, Proposition 4.4, and Lemmas 5.6--5.8.


\bibitem{Ratliff2026}
M. Ratliff, \emph{A Roughness-Explicit Chen Theorem for Goldbach
Representations}, preprint, version 2, 25 May 2026.
\href{https://zenodo.org/records/20377091}{Zenodo record and source}.
The stated result is not used as a validated input here.

\end{thebibliography}
\end{document}
